\documentclass[11pt,oneside,openany]{book}

\usepackage[T1]{fontenc}
\usepackage[utf8]{inputenc}
\usepackage[english]{babel}
\usepackage{lmodern}
\usepackage{microtype}
\usepackage[margin=1in,headheight=15pt]{geometry}
\usepackage{amsmath,amssymb,amsthm,mathtools}
\usepackage{stmaryrd}
\SetSymbolFont{stmry}{bold}{U}{stmry}{m}{n}
\usepackage{mathrsfs}
\usepackage{bm}
\usepackage{array,booktabs,longtable,tabularx,multirow}
\usepackage{enumitem}
\usepackage[dvipsnames]{xcolor}
\usepackage{tikz}
\usetikzlibrary{arrows.meta,positioning,calc,trees}
\usepackage{listings}
\usepackage{fancyhdr}
\usepackage{makeidx}
\usepackage{xurl}
\usepackage[hidelinks]{hyperref}
\usepackage[nameinlink,noabbrev]{cleveref}

\makeindex
\setcounter{tocdepth}{2}
\setcounter{secnumdepth}{3}
\setlist{nosep}
\setlist[enumerate,1]{label=\arabic*.,leftmargin=*}
\setlist[enumerate,2]{label=(\alph*),leftmargin=*}
\setlist[itemize]{leftmargin=*}
\allowdisplaybreaks

\definecolor{LogicBlue}{HTML}{17365D}
\definecolor{LogicGray}{HTML}{F2F4F7}
\definecolor{LogicRed}{HTML}{7A1F1F}

\hypersetup{
  pdftitle={Formal Logic: Language, Semantics, Deduction, and Metatheory},
  pdfsubject={A graduate-level reference in formal logic},
  pdfkeywords={formal logic, proof theory, model theory, computability, metatheory},
  colorlinks=true,
  linkcolor=LogicBlue,
  citecolor=LogicBlue,
  urlcolor=LogicBlue
}

\pagestyle{fancy}
\fancyhf{}
\fancyhead[L]{\small\nouppercase{\leftmark}}
\fancyhead[R]{\small Formal Logic}
\fancyfoot[C]{\thepage}
\renewcommand{\headrulewidth}{0.4pt}

\theoremstyle{definition}
\newtheorem{definition}{Definition}[chapter]
\newtheorem{example}[definition]{Example}
\newtheorem{exercise}[definition]{Exercise}
\newtheorem{method}[definition]{Method}
\newtheorem{convention}[definition]{Convention}

\theoremstyle{plain}
\newtheorem{theorem}[definition]{Theorem}
\newtheorem{lemma}[definition]{Lemma}
\newtheorem{corollary}[definition]{Corollary}

\theoremstyle{remark}
\newtheorem{remark}[definition]{Remark}
\newtheorem{warning}[definition]{Warning}

% Index the exact names assigned by definition and theorem environments.  The
% wrappers preserve the underlying amsthm environments and omit empty titles.
\let\LogicDefinition\definition
\let\endLogicDefinition\enddefinition
\renewenvironment{definition}[1][]{%
  \LogicDefinition[#1]%
  \if\relax\detokenize{#1}\relax\else\index{#1}\fi
}{\endLogicDefinition}
\let\LogicTheorem\theorem
\let\endLogicTheorem\endtheorem
\renewenvironment{theorem}[1][]{%
  \LogicTheorem[#1]%
  \if\relax\detokenize{#1}\relax\else\index{#1}\fi
}{\endLogicTheorem}

\newenvironment{solution}{\par\noindent\textit{Solution.}\ }{\hfill$\square$\par}
\newenvironment{proofplan}{\par\noindent\textit{Proof sketch.}\ }{\par}

\DeclarePairedDelimiter{\set}{\lbrace}{\rbrace}
\DeclarePairedDelimiter{\paren}{(}{)}
\DeclarePairedDelimiter{\bracket}{[}{]}
\DeclarePairedDelimiter{\angles}{\langle}{\rangle}
\DeclarePairedDelimiter{\abs}{\lvert}{\rvert}
\DeclarePairedDelimiter{\sem}{\llbracket}{\rrbracket}

\newcommand{\N}{\mathbb{N}}
\newcommand{\Z}{\mathbb{Z}}
\newcommand{\Q}{\mathbb{Q}}
\newcommand{\R}{\mathbb{R}}
\newcommand{\B}{\mathbb{B}}
\newcommand{\Pow}{\mathcal{P}}
\newcommand{\Lang}{\mathcal{L}}
\newcommand{\Meta}{\mathcal{M}}
\newcommand{\Theory}{\mathsf{T}}
\newcommand{\Prop}{\mathsf{Prop}}
\newcommand{\Var}{\mathsf{Var}}
\newcommand{\Term}{\mathsf{Term}}
\newcommand{\Fm}{\mathsf{Fm}}
\newcommand{\Sent}{\mathsf{Sent}}
\newcommand{\FV}{\operatorname{FV}}
\newcommand{\BV}{\operatorname{BV}}
\newcommand{\Sub}{\operatorname{Sub}}
\newcommand{\Mod}{\operatorname{Mod}}
\newcommand{\Th}{\operatorname{Th}}
\newcommand{\Cn}{\operatorname{Cn}}
\newcommand{\dom}{\operatorname{dom}}
\newcommand{\rng}{\operatorname{rng}}
\newcommand{\id}{\operatorname{id}}
\newcommand{\arity}{\operatorname{ar}}
\newcommand{\rank}{\operatorname{rk}}
\newcommand{\length}{\operatorname{len}}
\newcommand{\restr}{\mathbin{\upharpoonright}}
\newcommand{\entails}{\models}
\newcommand{\derives}{\vdash}
\newcommand{\notentails}{\not\models}
\newcommand{\notderives}{\not\vdash}
\newcommand{\True}{\mathbf{T}}
\newcommand{\False}{\mathbf{F}}
\newcommand{\bottom}{\bot}
\newcommand{\topv}{\top}
\newcommand{\forces}{\Vdash}
\newcommand{\ND}{\mathsf{ND}}
\newcommand{\LK}{\mathsf{LK}}
\newcommand{\LJ}{\mathsf{LJ}}
\newcommand{\PA}{\mathsf{PA}}
\newcommand{\ZF}{\mathsf{ZF}}
\newcommand{\ZFC}{\mathsf{ZFC}}
\newcommand{\K}{\mathsf{K}}
\newcommand{\Sfour}{\mathsf{S4}}
\newcommand{\Sfive}{\mathsf{S5}}
\newcommand{\IPL}{\mathsf{IPL}}
\newcommand{\CPL}{\mathsf{CPL}}
\newcommand{\FOL}{\mathsf{FOL}}
\newcommand{\HOL}{\mathsf{HOL}}
\newcommand{\SAT}{\mathsf{SAT}}
\newcommand{\UNSAT}{\mathsf{UNSAT}}
\newcommand{\CNF}{\mathsf{CNF}}
\newcommand{\DNF}{\mathsf{DNF}}
\newcommand{\NP}{\mathsf{NP}}
\newcommand{\coNP}{\mathsf{coNP}}
\newcommand{\Rule}[3][]{\frac{#2}{#3}\;\textsc{#1}}
\newcommand{\side}[1]{\quad\text{#1}}
\newcommand{\tuple}[1]{\angles{#1}}
\newcommand{\defeq}{\mathrel{:=}}
\newcommand{\eqdef}{\mathrel{=:}}

\lstdefinestyle{logiccode}{
  basicstyle=\ttfamily\small,
  columns=fullflexible,
  keepspaces=true,
  frame=single,
  rulecolor=\color{black!25},
  backgroundcolor=\color{LogicGray},
  breaklines=true,
  showstringspaces=false
}
\lstset{style=logiccode}

\begin{document}

\hypersetup{pageanchor=false}
\begin{titlepage}
  \centering
  \vspace*{2.2cm}
  {\Huge\bfseries Formal Logic\par}
  \vspace{0.6cm}
  {\LARGE Language, Semantics, Deduction, and Metatheory\par}
  \vspace{1.2cm}
  {\Large A Graduate-Level Reference\par}
  \vfill
  {\large Version 1.0\par}
  \vspace{0.2cm}
  {\large 9 August 2026\par}
  \vfill
  \begin{minipage}{0.82\textwidth}
  \small
  This document fixes one technical term for each concept and records the scope
  of every notation before use.  Definitions precede use.
  Every calculus is identified before the derivability symbol is used, and
  every semantic claim states the semantic data over which it is made.
  \end{minipage}
\end{titlepage}

\frontmatter
\hypersetup{pageanchor=true}

\chapter*{Preface}
\addcontentsline{toc}{chapter}{Preface}

Formal logic studies formal languages, interpretations of those languages,
and rule-governed derivations.  Correct work in formal logic therefore requires
more than symbol manipulation.  It requires exact control of the level of
language, the grammar of expressions, variable binding, interpretation,
assumptions, inference rules, and the metatheoretic hypotheses under which a
result is stated.

This reference is organized at graduate level rather than as a catalogue of
disconnected systems.  The dependency order is deliberate.  Syntax is
defined before semantics; semantics is defined before soundness; a calculus is
defined before derivability; classical propositional logic is established
before first-order logic; classical first-order logic is established before
model theory, computability, incompleteness, proof theory, and nonclassical
logic.  The organization follows the core treatments in
\cite{Enderton2001,EbbinghausFlumThomas1994,OpenLogic2026,Marker2024}, while the
historical and advanced results are tied to their original sources where
practical.

The document has four purposes:
\begin{enumerate}
  \item to teach the formal-logical content of the subject;
  \item to teach the construction and checking of derivations, proofs,
        countervaluations, structures, and countermodels of sentence
        arguments;
  \item to teach the writing and typesetting of formal logic; and
  \item to provide a dependency-ordered reference for graduate study, teaching,
        and preparation for qualifying examinations.
\end{enumerate}

The reference is self-contained at the level of notation and foundational
definitions.  Advanced model theory and proof theory are presented to the
point required for a broad graduate-level reference, but a research specialization
still requires the monographs cited in the relevant chapters.

\section*{Terminology and notation policy}

Literature in formal logic does not impose one universal notation.  Pretending
otherwise would be inaccurate.  This document instead fixes a registered
global convention.  A subsystem-specific departure is permitted only when its
scope and standing are declared before use, and no registered notation is used
for a second role within that scope.

Whenever this document introduces a new technical name, five conditions are
met at the same location:
\begin{enumerate}
  \item \emph{Necessity}: an existing term cannot express the required
        distinction without ambiguity;
  \item \emph{Definition}: an exact eliminable definition or an explicit axiom
        is supplied;
  \item \emph{Standing}: the grammatical and formal standing of the
        name is stated;
  \item \emph{Nonredundancy}: the distinction from previously defined terms is
        explained; and
  \item \emph{Use}: later definitions or theorems require the distinction.
\end{enumerate}
Until those conditions are met, the text describes the required object or
relation without assigning it a technical name.

Throughout this document, the word \emph{iff} abbreviates the
metalinguistic biconditional ``if and only if''; it is the registered
connective of definitions and metatheoretic equivalences.

The following distinctions are permanent:
\begin{itemize}
  \item A \emph{formula} is a grammatical object-language expression.  A
        \emph{sentence} is a formula with no free variable occurrence.
  \item A \emph{structure} supplies a domain and interpretations for the
        nonlogical symbols of a signature.  A \emph{model} is a structure that
        satisfies a specified set of sentences.
  \item A \emph{valuation} assigns truth values to propositional variables.  A
        \emph{variable assignment} assigns domain elements to first-order
        variables.
  \item \emph{Satisfaction} is the relation between a formula and the complete
        semantic data permitted by the active semantics.  \emph{Truth} is
        satisfaction relative to specified such data.  \emph{Validity} is
        satisfaction for every permitted tuple of semantic data.
        \emph{Semantic consequence} is preservation of satisfaction from a set
        of assumptions to a conclusion.  For a first-order sentence, assignment
        independence later permits the specialized phrase ``truth in a
        structure.''
  \item A \emph{derivation} is a finite rule-governed object in a specified
        calculus.  A \emph{proof} is a derivation with no undischarged
        assumptions other than the stated premises.
  \item \emph{Soundness} proceeds from derivability to semantic consequence.
        \emph{Completeness} proceeds from semantic consequence to derivability.
  \item \emph{Consistency} excludes derivability of a contradiction.
        A theory is \emph{deductively complete} iff every sentence or its
        negation is derivable from the theory.  A theory is
        \emph{semantically complete} iff it semantically entails every sentence
        or its negation.  For the fixed sound and complete first-order calculus,
        these two properties are equivalent; neither is consistency.
\end{itemize}

\section*{Evidence and verification}

The status of a formal result must be reported at the level actually checked.
These levels are not interchangeable:
\begin{enumerate}
  \item a displayed calculation checks only its displayed instances;
  \item a truth table checks all valuations of the finitely many variables it
        contains;
  \item a derivation checks derivability in the stated calculus;
  \item a proof establishes the stated result under its stated
        hypotheses;
  \item a proof-assistant kernel checks a formal proof relative to the kernel,
        imported declarations, and admitted axioms;
  \item successful \LaTeX{} compilation checks typesetting syntax, not logical
        correctness; and
  \item editorial or formal-logical review is independent evidence and must not be
        inferred from any preceding level.
\end{enumerate}

When this document states a result status explicitly, it uses exactly one of
three terms.  \emph{Complete proof} means that all substantive formal-logical
steps are supplied in this document.  \emph{Proof sketch} identifies the
construction and states which substantive verification is omitted.
\emph{Cited theorem} states exact hypotheses and a conclusion but refers to a
named source for the proof.  A dependency used within a complete proof is
stated separately as a cited theorem; it is not hidden in a qualified status.

\chapter*{Graduate-Level Scope}
\addcontentsline{toc}{chapter}{Graduate-Level Scope}

\section*{Prerequisites}

The reader should be able to work with sets, functions, relations, finite
strings, countable sets, induction on natural numbers, and ordinary rigorous
proofs.  Prior exposure to elementary logic is useful but not
required.  Chapters that use cardinal arithmetic, algebra, or topology state
the additional facts they require.

\section*{Learning outcomes}

A reader completing the reference should be able to:
\begin{enumerate}
  \item define a formal language by an inductive grammar and prove results by
        structural induction;
  \item define capture-avoiding substitution and prove a substitution lemma;
  \item prove soundness for a natural-deduction or sequent calculus;
  \item prove completeness for classical propositional logic and explain a
        Henkin proof of completeness for classical first-order logic;
  \item derive compactness and the downward and upward L\"owenheim--Skolem
        theorems with their cardinal hypotheses;
  \item use back-and-forth arguments, elementary chains, and ultraproducts;
  \item prove undecidability by a many-one reduction;
  \item state and apply the diagonal lemma and derive the first incompleteness
        theorem, the second incompleteness theorem, and undefinability of truth
        with exact hypotheses;
  \item explain cut elimination, normalization, and the subformula property;
  \item give relational semantics for intuitionistic and normal modal logic and
        prove representative soundness results;
  \item distinguish full semantics from Henkin semantics for second-order logic;
        and
  \item audit a written or mechanized proof for hidden assumptions, variable
        capture, rule misuse, and shifts between object language and
        metalanguage.
\end{enumerate}

\tableofcontents

\mainmatter

\part{Foundations of Formal Work}

\chapter{Sets, Relations, Functions, and Cardinality}
\label{ch:set-prerequisites}

This chapter fixes the metatheoretic set theory used before the dedicated
study of axiomatic set theory.  The treatment assumes ordinary classical set
theory in the metalanguage.  It does not treat an informal collection as a set
unless the ambient set theory licenses its existence.

\section{Sets and classes}
\label{sec:set-prerequisites-01}

\begin{definition}[Subset]
For sets \(A\) and \(B\), write \(A\subseteq B\) iff
\[
  \forall x\,(x\in A\to x\in B).
\]
The sets are equal iff each is a subset of the other.
\end{definition}

\begin{definition}[Power set]
For a set \(A\), its \emph{power set} is
\[
  \Pow(A)\defeq\set{X:X\subseteq A}.
\]
\end{definition}

\begin{definition}[Class]
A \emph{class} is a metalinguistically specified collection of sets.  A class
need not be a set.  A \emph{proper class} is a class that is not a set.
\end{definition}

An unrestricted collection described by a condition may be a proper class.
Quantification over such a collection is metalinguistic quantification in the
ambient foundation, not object-language quantification over a set.

\begin{warning}
Set-builder notation does not itself prove that a set exists.  The expression
\(\set{x:P(x)}\) may denote a proper class.  When existence matters, identify a
previously available bounding set or cite the relevant set-existence axiom.
\end{warning}

\section{Ordered pairs and products}
\label{sec:set-prerequisites-02}

\begin{definition}[Ordered pair]
The \emph{ordered pair} of sets \(a\) and \(b\) is
\[
  \tuple{a,b}\defeq\set{\set{a},\set{a,b}}.
\]
\end{definition}

\begin{lemma}[Ordered-pair identity]
For all sets \(a,b,c,d\),
\[
  \tuple{a,b}=\tuple{c,d}
  \quad\Longleftrightarrow\quad
  a=c\text{ and }b=d.
\]
\end{lemma}
\begin{proof}
The reverse implication follows by substitution.  For the forward implication,
if \(a=b\), then \(\tuple{a,b}=\set{\set a}\), so equality of the ordered
pairs forces \(c=d=a\).  If \(a\neq b\), then the ordered pair has the two
distinct members \(\set a\) and \(\set{a,b}\).  Equality forces
\(\set c=\set a\), hence \(c=a\), and then comparison of the two-element
members gives \(d=b\).  The remaining apparent case \(c=d\) would make the
right ordered pair a singleton and is excluded by \(a\neq b\).
\end{proof}

\begin{definition}[Cartesian product]
For sets \(A\) and \(B\),
\[
  A\times B\defeq\set{\tuple{a,b}:a\in A\text{ and }b\in B}.
\]
For every \(n\in\N\), define \(A^n\) as the set of functions from
\(n=\set{0,\ldots,n-1}\) to \(A\).  When \(n=0\), the domain is empty, so
\(A^0\) contains exactly the empty function.
\end{definition}

\section{Relations and functions}
\label{sec:set-prerequisites-03}

\begin{definition}[Binary relation]
A \emph{binary relation from \(A\) to \(B\)} is a subset of \(A\times B\).  A
\emph{binary relation on \(A\)} is a subset of \(A\times A\).
\end{definition}

\begin{definition}[Properties of a binary relation]
Let \(R\) be a binary relation on \(A\).  It is \emph{reflexive} iff
\(aRa\) for every \(a\in A\); \emph{irreflexive} iff \(aRa\) for no
\(a\in A\); \emph{symmetric} iff \(aRb\) implies \(bRa\) for all
\(a,b\in A\); and \emph{transitive} iff \(aRb\) and \(bRc\) imply \(aRc\)
for all \(a,b,c\in A\).

The relation \(R\) is a \emph{strict linear order} on \(A\) iff it is
irreflexive and transitive and, for every distinct \(a,b\in A\), exactly one
of \(aRb\) and \(bRa\) holds.  It \emph{strictly well-orders} \(A\) iff it is
a strict linear order and every nonempty \(B\subseteq A\) has an
\(R\)-least member: some \(b\in B\) such that no \(a\in B\) satisfies
\(aRb\).  These clauses define the relation properties used below; none is a
second name for another clause, and strict well-ordering is required in the
definition of ordinal.
\end{definition}

\begin{definition}[Function]
A \emph{function from \(A\) to \(B\)}, written \(f:A\to B\), is a binary
relation \(f\subseteq A\times B\) such that for every \(a\in A\) there exists
exactly one \(b\in B\) with \(\tuple{a,b}\in f\).  The set \(A\) is the
domain, and \(B\) is the codomain.
\end{definition}

The value at \(a\) is written \(f(a)\).  The range is
\(\rng(f)=\set{f(a):a\in A}\).  The codomain is part of the typed declaration
and need not equal the range.

\begin{definition}[Partial function]
A \emph{partial function from \(A\) to \(B\)}, written
\(f:A\rightharpoonup B\), is a function \(f:D\to B\) for some
\(D\subseteq A\).  Its domain of definition is \(D\).
\end{definition}

\begin{definition}[Injection, surjection, and bijection]
Let \(f:A\to B\).
\begin{enumerate}
  \item The function is \emph{injective} iff
        \(f(a)=f(a')\) implies \(a=a'\) for all \(a,a'\in A\).
  \item The function is \emph{surjective} iff for every \(b\in B\) there is
        an \(a\in A\) such that \(f(a)=b\).
  \item The function is \emph{bijective} iff it is injective and surjective.
\end{enumerate}
\end{definition}

\begin{definition}[Image and inverse image]
For \(X\subseteq A\) and \(Y\subseteq B\), define
\[
  f[X]\defeq\set{f(x):x\in X},
  \qquad
  f^{-1}[Y]\defeq\set{x\in A:f(x)\in Y}.
\]
The notation \(f^{-1}[Y]\) does not assert that \(f\) has an inverse function.
\end{definition}

\section{Equivalence relations and quotients}
\label{sec:set-prerequisites-04}

\begin{definition}[Equivalence relation]
A binary relation \(E\) on \(A\) is an \emph{equivalence relation} iff it is
reflexive, symmetric, and transitive.
\end{definition}

\begin{definition}[Equivalence class and quotient]
If \(E\) is an equivalence relation on \(A\), define
\[
  [a]_E\defeq\set{x\in A:xEa},
  \qquad
  A/E\defeq\set{[a]_E:a\in A}.
\]
\end{definition}

\begin{method}[Definition on a quotient]
To define \(\bar f:A/E\to B\) by \(\bar f([a]_E)=f(a)\), prove
\[
  aEa'\quad\Longrightarrow\quad f(a)=f(a').
\]
This proof establishes that the displayed value is independent of the chosen
representative.  Without it, the formula does not define a function.
\end{method}

\section{Finite, countable, and uncountable sets}
\label{sec:set-prerequisites-05}

\begin{definition}[Finite set]
A set \(A\) is \emph{finite} iff there exists \(n\in\N\) and a bijection from
\(n=\set{0,\ldots,n-1}\) to \(A\).
\end{definition}

\begin{definition}[Countable set]
A set \(A\) is \emph{countable} iff there exists an injection \(A\to\N\).
It is \emph{countably infinite} iff there exists a bijection \(\N\to A\).
\end{definition}

This document includes finite sets among countable sets.  The convention is
used in the countable-language versions of the completeness and omitting-types
theorems.

\begin{definition}[Cardinal comparison]
For sets \(A\) and \(B\), write \(\abs A\leq\abs B\) iff there exists an
injection \(A\to B\), and write \(\abs A=\abs B\) iff there exists a bijection
\(A\to B\).  Write \(\abs A<\abs B\) iff
\(\abs A\leq\abs B\) and \(\abs A\ne\abs B\).
\end{definition}

\begin{theorem}[Cantor--Schr\"oder--Bernstein]
If there is an injection \(f:A\to B\) and an injection \(g:B\to A\), then
there is a bijection \(h:A\to B\).
\end{theorem}
\begin{proof}
Let \(A_0=A\setminus g[B]\), and recursively let
\(A_{n+1}=g[f[A_n]]\).  Put \(C=\bigcup_{n\in\N}A_n\).  Define
\[
  h(a)=
  \begin{cases}
    f(a),&a\in C,\\
    g^{-1}(a),&a\notin C.
  \end{cases}
\]
The second case is defined because \(A_0\subseteq C\), so
\(a\notin C\) implies \(a\in g[B]\), and injectivity of \(g\) makes the
preimage unique.  The two ranges are disjoint: if \(a\in C\), then
\(g(f(a))\in C\); if \(a'\notin C\), then \(a'=g(b)\) with
\(g(b)\notin C\), so \(b\notin f[C]\).  Each branch is injective.  For
surjectivity, let \(b\in B\).  If \(b\in f[C]\), it is in the first range.
Otherwise \(g(b)\notin C\): since \(g(b)\in g[B]\), the element
\(g(b)\) is not in \(A_0=A\setminus g[B]\); and if
\(g(b)\in A_{n+1}=g[f[A_n]]\), injectivity of \(g\) would give
\(b\in f[C]\), contrary to the case assumption.  Thus
\(h(g(b))=b\).  Hence \(h\) is bijective.
\end{proof}

\begin{theorem}[Cantor]
For every set \(A\), there is no surjection \(f:A\to\Pow(A)\).  Consequently,
\(\abs A<\abs{\Pow(A)}\).
\end{theorem}
\begin{proof}
Suppose \(f:A\to\Pow(A)\).  Define
\[
  D\defeq\set{a\in A:a\notin f(a)}.
\]
If \(f\) were surjective, there would be \(d\in A\) with \(f(d)=D\).  Then
\[
  d\in D\quad\Longleftrightarrow\quad d\notin f(d)
  \quad\Longleftrightarrow\quad d\notin D,
\]
a contradiction.  Thus no such surjection exists.  The singleton map
\(a\mapsto\set a\) is an injection \(A\to\Pow(A)\), so the cardinal inequality
is strict.
\end{proof}

\begin{definition}[Transitive set]
\label{def:transitive-set}
A set \(A\) is \emph{transitive} iff
\[
  \forall x\in A\,\forall y\in x\,(y\in A),
\]
equivalently, iff \(\bigcup A\subseteq A\).
\end{definition}

\begin{definition}[Ordinal and cardinal]
\label{def:ordinal-cardinal}
An \emph{ordinal} is a transitive set \(\alpha\) such that \(\in\) strictly
well-orders \(\alpha\).  For ordinals \(\alpha,\beta\), write
\(\alpha<\beta\) iff \(\alpha\in\beta\).  A \emph{cardinal} is an ordinal
\(\kappa\) for which no smaller ordinal is bijective with \(\kappa\).  Thus a
cardinal is already the least ordinal of its bijection type; the phrase
``initial ordinal'' will not be used as a second term.  A cardinal is
\emph{infinite} iff it is not bijective with any \(n\in\N\).
\end{definition}

For an ordinal \(\alpha\), its \emph{successor ordinal} is
\(\alpha+1\defeq\alpha\cup\{\alpha\}\).  A \emph{limit ordinal} is a nonzero
ordinal that is not a successor ordinal.  These terms classify ordinals by
their immediate-predecessor behavior; they do not introduce a second order or
a second notion of cardinality.  Limit ordinals index the elementary chains
used later, and successor and limit stages distinguish the two clauses of the
Henkin recursion.

\begin{theorem}[Transfinite recursion]
\label{thm:transfinite-recursion}
Let \(\delta\) be an ordinal.  Suppose a functional rule \(G\) assigns a set
\(G(g)\) to every function \(g\) whose domain is an ordinal smaller than
\(\delta\).  Then there is a unique function \(F\) with domain \(\delta\) such
that
\[
 F(\alpha)=G(F\mathbin{\upharpoonright}\alpha)
 \qquad\text{for every }\alpha<\delta.
\]
\textit{Status: cited theorem.}
\end{theorem}

See \cite{Jech2003} for a complete proof.  The theorem is the exact
metalinguistic recursion principle used below; it is not a new inference rule
of any object-language calculus.

\section{Choice principles used later}
\label{sec:set-prerequisites-06}

\begin{definition}[Axiom of Choice]
The \emph{Axiom of Choice} states that for every set \(X\) of nonempty sets
there exists a function \(c:X\to\bigcup X\) such that \(c(A)\in A\) for every
\(A\in X\).
\end{definition}

\begin{theorem}[Well-ordering theorem]
Assuming the Axiom of Choice, every set can be well-ordered.
\textit{Status: cited theorem.}
\end{theorem}

\begin{corollary}[Cardinal assignment]
\label{cor:cardinal-assignment}
Assuming the Axiom of Choice, every set \(A\) is bijective with a unique
cardinal.
\textit{Status: cited theorem.}
\end{corollary}

See \cite{Jech2003} for complete proofs.  The cardinal bijective with
\(A\) is denoted by \(\abs A\), and its order agrees with the
injection-based comparison defined above.  The cardinal of \(\N\) is
denoted by \(\aleph_0\).  For a nonempty finite list of
cardinals, \(\max(\kappa_1,\ldots,\kappa_n)\) denotes its greatest member in
the cardinal order.  The well-ordering theorem makes that order total under
the stated Axiom of Choice, so this maximum exists.  These conventions are
the exact cardinal conventions used by the Henkin and
L\"owenheim--Skolem constructions.

\begin{theorem}[Finite coding over an infinite cardinal]
\label{thm:infinite-cardinal-finite-coding}
Assume the Axiom of Choice, and let \(\kappa\) be an infinite cardinal.
Then:
\begin{enumerate}
  \item \(\abs{\kappa^n}=\kappa\) for every positive \(n\in\N\);
  \item the set of all finite sequences from \(\kappa\) has cardinality
        \(\kappa\);
  \item if \(\alpha<\kappa\) and \(A_\beta\) is finite for every
        \(\beta<\alpha\), then
        \(\abs{\bigcup_{\beta<\alpha}A_\beta}<\kappa\); and
  \item if \(\abs{A_n}\leq\kappa\) for every \(n\in\N\), then
        \(\abs{\bigcup_{n\in\N}A_n}\leq\kappa\).
\end{enumerate}
\textit{Status: cited theorem.}
\end{theorem}

\begin{corollary}[Syntax counting]
\label{cor:syntax-counting}
Assume the Axiom of Choice and let \(\kappa\) be an infinite cardinal.
\begin{enumerate}
  \item Every set \(A\) with \(\abs{A}\leq\kappa\) satisfies
        \(\abs{A^{<\omega}}\leq\kappa\), where \(A^{<\omega}\) is the set of
        finite sequences from \(A\).
  \item If \(\Sigma\) is an alphabet with \(\abs{\Sigma}\leq\kappa\), then
        the finite strings over \(\Sigma\), the finite tuples of such
        strings, and the formulas generated by any fixed finitary grammar
        whose alphabet extends \(\Sigma\) by at most countably many further
        symbols each form a set of cardinality at most \(\kappa\).
\end{enumerate}
\textit{Status: complete proof.}
\end{corollary}

\begin{proof}
For item~1, fix an injection \(f:A\to\kappa\); such an injection is the
meaning of \(\abs{A}\leq\kappa\).  The map
\((a_1,\ldots,a_n)\mapsto(f(a_1),\ldots,f(a_n))\) is a length-preserving
injection of \(A^{<\omega}\) into \(\kappa^{<\omega}\), and item~2 of
\cref{thm:infinite-cardinal-finite-coding} gives
\(\abs{\kappa^{<\omega}}=\kappa\).

For item~2, the finite strings over \(\Sigma\) lie in \(\Sigma^{<\omega}\):
apply item~1 with \(A=\Sigma\).  The finite tuples of such strings lie in
\((\Sigma^{<\omega})^{<\omega}\): apply item~1 with \(A=\Sigma^{<\omega}\),
a set of cardinality at most \(\kappa\) by the previous step.  Finally,
every expression generated by a finitary grammar is a finite string over the
grammar's alphabet; the string chapters below record this fact for formulas.
Write \(X\) for the at most countable set of adjoined symbols hypothesized
in the statement---logical symbols, punctuation, and variables.  Item~4 of
\cref{thm:infinite-cardinal-finite-coding}, applied to the sequence
\(\Sigma,X,\varnothing,\varnothing,\ldots\), gives
\(\abs{\Sigma\cup X}\leq\kappa\), so the formula set injects into
\((\Sigma\cup X)^{<\omega}\) and item~1 applies.
\end{proof}

See \cite[Theorem~3.5 and Chapter~5]{Jech2003} for complete proofs of
\cref{thm:infinite-cardinal-finite-coding}: Theorem~3.5 gives
\(\aleph_\alpha\cdot\aleph_\alpha=\aleph_\alpha\), and with the
well-ordering theorem cited above every infinite cardinal is an aleph;
Chapter~5 gives the cardinal sums and the union bound
\(\abs{\bigcup_{i\in I}A_i}\leq\abs{I}\cdot\sup_{i\in I}\abs{A_i}\) under
the Axiom of Choice.  \Cref{thm:infinite-cardinal-finite-coding}
and \cref{cor:syntax-counting} supply exactly the cardinal arithmetic
invoked by the Henkin and L\"owenheim--Skolem proofs; no claim about
arbitrary infinite products is being made.

\begin{definition}[Ultrafilter Lemma]
\label{def:ultrafilter-lemma}
The \emph{Ultrafilter Lemma} states the following.  Let \(I\) be a set and let
\(\mathcal F\subseteq\Pow(I)\) satisfy
\[
 I\in\mathcal F,\quad\varnothing\notin\mathcal F,\quad
 X,Y\in\mathcal F\Rightarrow X\cap Y\in\mathcal F,
\]
and
\[
 X\in\mathcal F\text{ and }X\subseteq Y\subseteq I
 \Rightarrow Y\in\mathcal F.
\]
Then there is \(\mathcal U\subseteq\Pow(I)\) containing \(\mathcal F\) that
satisfies the same conditions and, for every \(X\subseteq I\), contains
exactly one of \(X\) and \(I\setminus X\).  This is a metatheoretic choice
principle about families of subsets, not a formula of a first-order theory
under study.  It is strictly weaker over \(\ZF\) than the Axiom of Choice and
therefore is not a second name for that axiom.  The later filter chapter gives
the short technical terms for the two displayed families before applying this
principle.
\end{definition}

Some model-theoretic constructions use the Ultrafilter Lemma, and some cardinal
arguments use the Axiom of Choice.  Later theorems state the required principle
rather than silently treating either principle as part of pure first-order
logic.

\begin{exercise}
\label{ex:set-prerequisites-01}
Prove that the image of a countable set under a function is countable.
\end{exercise}

\begin{exercise}
\label{ex:set-prerequisites-02}
Let \(E\) be congruence modulo \(5\) on \(\Z\).  Define addition on
\(\Z/E\) using representatives and prove that it is well-defined.
\end{exercise}

\begin{exercise}
\label{ex:set-prerequisites-03}
Find the exact point at which the proof that a countable union of countable sets
is countable requires a simultaneous choice of enumerations.
\end{exercise}

\section*{Dependency summary}

The later syntax chapters use finite strings, functions, induction, and
countability.  Semantics uses relations, products, and power sets.  Quotient
constructions use equivalence classes and well-definedness.  Model theory and
set theory use cardinal comparison, Cantor's theorem, and stated choice
principles.

\chapter{Levels, Objects, and Consequence}
\label{ch:levels}

\section{Formal logic and a logical system}
\label{sec:levels-01}

\begin{definition}[Formal language]
A \emph{formal language} is a set of finite expressions determined from a
specified alphabet by a specified grammar.  A formal language is an abstract
syntactic object.  The grammar determines which finite strings count as
expressions of each grammatical category.
\end{definition}

\begin{definition}[Calculus]
A \emph{calculus} is a specified collection of initial derivations and
inference rules that inductively determines a derivability relation.  The
standing of a calculus is syntactic: it concerns expressions and rules, not
their interpretations.
\end{definition}

\begin{definition}[Semantics]
A \emph{semantics} for a formal language specifies the complete semantic data
permitted for that language and a satisfaction relation between those data and
formulas.  The data may be a propositional valuation, a first-order structure
with a variable assignment, a modal structure with a world, or another tuple
declared explicitly by the semantics.  The satisfaction relation supplies the
truth conditions of the language.
\end{definition}

\begin{definition}[Logical system]
A \emph{logical system} is an ordered triple
\[
  \tuple{\Lang,\mathcal C,\mathcal S},
\]
where \(\Lang\) is a formal language, \(\mathcal C\) is a calculus for
\(\Lang\), and \(\mathcal S\) is a semantics for \(\Lang\).  The term is used
only when all three components have been identified.
\end{definition}

The ordered triple prevents a common error.  A symbol such as \(\derives\)
does not name one universal relation: it abbreviates derivability in a
particular calculus.  Likewise, \(\entails\) depends on a language and its
specified semantic data.  When more than one calculus or semantics is active,
this document writes subscripts explicitly.

\begin{example}
For classical propositional logic, \(\Lang\) contains propositional formulas,
\(\mathcal C\) may be a natural-deduction calculus, and \(\mathcal S\) uses
two-valued valuations.  Replacing natural deduction by a sequent calculus
changes \(\mathcal C\) but not necessarily \(\Lang\) or \(\mathcal S\).
Replacing two-valued valuations by three-valued matrices changes
\(\mathcal S\), and the resulting semantic consequence relation must be
defined again.
\end{example}

\section{Object language and metalanguage}
\label{sec:levels-02}

\begin{definition}[Object language]
The \emph{object language} is the formal language currently being studied.
Its symbols and formulas are objects of formal-logical investigation.
\end{definition}

\begin{definition}[Metalanguage]
The \emph{metalanguage} is the language in which the syntax, semantics, and
derivations of the object language are defined and studied.  In this document
the metalanguage is precise academic English supplemented by ordinary set
theory.
\end{definition}

The distinction is functional, not typographical.  A Greek letter such as
\(\varphi\) is normally a metavariable ranging over object-language formulas;
it is not itself a symbol of the object language.  Quotation marks or corner
quotes may be used when an object-language inscription must be mentioned
literally.  Semantic predicates, including satisfaction, belong to the
metalanguage unless the object language has been explicitly extended by a
predicate with specified axioms.

Tarski's analysis of truth makes the separation of object language and
metalanguage metalogically substantive, especially when a language can encode
its own syntax \cite{Tarski1983}.  The distinction will be required in the
definition of satisfaction, in arithmetization, and in the undefinability
theorem.

\begin{example}[Use and mention]
The display
\[
  P \land Q
\]
uses an object-language formula in a derivation.  The sentence
``\((P\land Q)\) contains five symbol occurrences, exactly one of which is a
logical-connective occurrence'' mentions that formula and makes a
metalinguistic statement about its syntax.  The mentioned formula carries the
parentheses required by the grammar of \cref{ch:prop-syntax-semantics};
displayed uses may omit the outermost pair under the typographical
abbreviation declared there.
\end{example}

\begin{warning}
Do not write \(\varphi\entails\psi\) when \(\varphi\) and \(\psi\) are being
used as object-language inscriptions unless the semantic consequence relation
has been lifted to metavariables and the intended universal quantification is
clear.  A precise metalinguistic statement is:
\[
  \text{for all formulas }\varphi,\psi\text{ of the language},
  \quad \set{\varphi}\entails\psi\ \text{iff}\ \cdots.
\]
\end{warning}

\section{Syntax, semantics, and proof theory}
\label{sec:levels-03}

\begin{definition}[Syntax]
The \emph{syntax} of a formal language consists of its alphabet, grammar, and
operations on expressions defined without reference to interpretation.
\end{definition}

\begin{definition}[Proof theory]
The \emph{proof theory} of a calculus studies derivations and transformations
of derivations by syntactic methods.
\end{definition}

Syntax and proof theory are related but distinct.  Syntax determines which
formulas exist.  A calculus determines which derivations exist.  Semantics
determines when formulas are satisfied.  Metatheorems connect these domains.

\begin{definition}[Semantic consequence]
Let \(\Gamma\) be a set of formulas and let \(\varphi\) be a formula in a
language with a specified semantics.  Then \(\varphi\) is a \emph{semantic
consequence} of \(\Gamma\), written
\[
  \Gamma\entails\varphi,
\]
iff, for every complete tuple of semantic data permitted by that semantics,
satisfaction of every member of \(\Gamma\) implies satisfaction of
\(\varphi\).  The permitted data may be a propositional valuation, a
first-order structure with a variable assignment, a modal structure with a
world, or another tuple explicitly declared by the active semantics.
\end{definition}

\begin{definition}[Derivability]
Let \(\mathcal C\) be a specified calculus.  A formula \(\varphi\) is
\emph{derivable} from a set of formulas \(\Gamma\), written
\[
  \Gamma\derives_{\mathcal C}\varphi,
\]
iff there exists a derivation in \(\mathcal C\) whose conclusion is
\(\varphi\) and whose undischarged assumptions are members of \(\Gamma\).
\end{definition}

\begin{definition}[Validity]
A formula \(\varphi\) is \emph{valid} for a specified semantics iff
\(\varnothing\entails\varphi\).  We then write \(\entails\varphi\).
\end{definition}

\begin{definition}[Theorem of a calculus]
A formula \(\varphi\) is a \emph{theorem} of a specified calculus
\(\mathcal C\) iff \(\varnothing\derives_{\mathcal C}\varphi\).  We then
write \(\derives_{\mathcal C}\varphi\).
\end{definition}

\begin{definition}[Soundness]
A calculus \(\mathcal C\) is \emph{sound} for a specified semantics iff, for
every set of formulas \(\Gamma\) and formula \(\varphi\),
\[
  \Gamma\derives_{\mathcal C}\varphi
  \quad\Longrightarrow\quad
  \Gamma\entails\varphi.
\]
\end{definition}

\begin{definition}[Completeness]
A calculus \(\mathcal C\) is \emph{complete} for a specified semantics iff,
for every set of formulas \(\Gamma\) and formula \(\varphi\),
\[
  \Gamma\entails\varphi
  \quad\Longrightarrow\quad
  \Gamma\derives_{\mathcal C}\varphi.
\]
\end{definition}

\begin{warning}[Direction check]
Soundness begins with a derivation.  Completeness begins with semantic
consequence.  Writing the two displayed implications before beginning a proof
prevents reversal of their roles.
\end{warning}

\section{Truth, satisfaction, validity, and models}
\label{sec:levels-04}

The notions of structure, variable assignment, sentence, and first-order
formula used in this section receive their exact inductive definitions in
\cref{ch:fol-syntax-semantics} and \cref{ch:fol-semantics}.  The
definitions below are stated here schematically, to fix the registered
senses of the semantic vocabulary; nothing before those chapters depends on
more than the registered senses.

\begin{definition}[Satisfaction]
\emph{Satisfaction} is the relation defined by the semantic clauses of the
language.  In first-order logic it has the form
\(\mathcal M,s\entails\varphi\), where \(\mathcal M\) is a structure,
\(s\) is a variable assignment, and \(\varphi\) is a formula.
\end{definition}

\begin{definition}[Truth in a structure]
A sentence \(\sigma\) is \emph{true in} a structure \(\mathcal M\) iff
\(\mathcal M,s\entails\sigma\) for every variable assignment \(s\).
Because a sentence has no free variable occurrence, the choice of \(s\) will
later be proved irrelevant.  We then write \(\mathcal M\entails\sigma\).
\end{definition}

\begin{definition}[Model]
Let \(T\) be a set of sentences.  A structure \(\mathcal M\) is a
\emph{model of} \(T\) iff \(\mathcal M\entails\sigma\) for every
\(\sigma\in T\).  The class of models of \(T\) is
\[
  \Mod(T)\defeq\set{\mathcal M: \mathcal M\entails T},
\]
where \(\mathcal M\entails T\) abbreviates: \(\mathcal M\entails\sigma\)
for every \(\sigma\in T\).
\end{definition}

Thus a structure is not called a model without an accompanying theory or
sentence.  The phrase ``the model \(\mathcal M\)'' is permitted only after the
relevant theory has been fixed.

\section{Arguments and formalization}
\label{sec:levels-05}

\begin{definition}[Argument]
An \emph{argument} is an ordered pair \(\tuple{\Gamma,\varphi}\), where
\(\Gamma\) is a set of formulas called the premises and \(\varphi\) is a
formula called the conclusion.
\end{definition}

An argument \(\tuple{\Gamma,\varphi}\) is \emph{semantically valid} iff
\(\Gamma\entails\varphi\); for \(\Gamma=\varnothing\) this agrees with
validity of the formula \(\varphi\), and the phrase names no relation
beyond semantic consequence.  A proof of
the conclusion from the premises is a derivation witnessing
\(\Gamma\derives_{\mathcal C}\varphi\).  If every premise and the conclusion
are sentences, an invalid argument is refuted by a countermodel.  For formulas
with free variable occurrences, the refuting data are stated explicitly as a
structure and a variable assignment.

\begin{definition}[Countermodel]
Let every member of \(\Gamma\cup\set{\varphi}\) be a sentence.  A
\emph{countermodel} to \(\tuple{\Gamma,\varphi}\) is a model of \(\Gamma\)
that does not satisfy \(\varphi\).
The restriction to sentences preserves the fixed standing of a model as a
structure satisfying a specified set of sentences.
\end{definition}

\begin{method}[Formalizing an argument]
Perform the following steps in order.
\begin{enumerate}
  \item State the intended domain and the conditions under which an expression
        counts as a premise or conclusion.
  \item Declare every nonlogical symbol, including its arity and intended
        interpretation.
  \item Resolve pronouns, tense, modality, plurality, restricted quantification,
        and uniqueness before translating.
  \item Translate each premise separately.  Preserve quantifier scope with
        parentheses.
  \item Translate the conclusion separately.
  \item Test the translation against intended positive and negative cases.
  \item Only then determine derivability.  For a sentence argument, search for
        a countermodel when semantic consequence is in doubt; for formulas
        with free variable occurrences, search for an explicit structure and
        variable assignment that refute semantic consequence.
\end{enumerate}
\end{method}

\begin{example}[Quantifier order]
Over a domain of people, let \(A(x,y)\) mean that \(x\) admires \(y\).
The sentence
\[
  \forall x\,\exists y\,A(x,y)
\]
allows the admired person to depend on \(x\).  The sentence
\[
  \exists y\,\forall x\,A(x,y)
\]
requires one person admired by everyone.  Neither sentence may be substituted
for the other without changing the claim.
\end{example}

\begin{exercise}
\label{ex:levels-01}
Classify each expression as object-language or metalanguage after fixing the
object language to be classical propositional logic: \(P\), \(\varphi\),
\(P\land Q\), ``is valid,'' \(\entails\), and the left parenthesis occurring in
\(P\land Q\).
\end{exercise}

\begin{exercise}
\label{ex:levels-02}
Give a valuation that satisfies both premises and falsifies the conclusion
of the argument with premises \(P\to Q\) and \(Q\) and conclusion
\(P\).  State explicitly which formulas
are true and which formula is false under the valuation.
\end{exercise}

\chapter{Finite Strings, Induction, and Recursion}
\label{ch:strings}

Formal languages are usually built from finite strings.  This chapter states
the set-theoretic machinery used throughout the reference.

\section{Finite sequences and strings}
\label{sec:strings-01}

\begin{definition}[Finite sequence]
Let \(A\) be a set and let \(n\in\N\).  A \emph{finite sequence in \(A\) of
length \(n\)} is a function \(w:\set{0,\ldots,n-1}\to A\).  Its length is
\(\length(w)=n\).  The unique sequence of length \(0\) is denoted by
\(\epsilon\).
\end{definition}

\begin{definition}[String]
Let \(A\) be an alphabet.  A \emph{string over \(A\)} is a finite sequence in
\(A\).  The set of all strings over \(A\) is
\[
  A^*\defeq\bigcup_{n\in\N}A^n.
\]
\end{definition}

The proof below requires one injection that records two natural numbers in one
natural number.  Fix the metalinguistic abbreviation
\[
  \operatorname{pair}_{\N}(a,b)\defeq
  \frac{(a+b)(a+b+1)}2+b
  \qquad(a,b\in\N).
\]
This map is a bijection \(\N\times\N\to\N\).  Indeed, put
\(t_s=s(s+1)/2\).  For each \(z\in\N\), there is a unique \(s\) with
\(t_s\leq z<t_{s+1}\), because the consecutive finite intervals
\([t_s,t_{s+1})\) partition \(\N\).  Set
\[
  b=z-t_s,
  \qquad
  a=s-b.
\]
The upper inequality gives \(0\leq b\leq s\), so \(a,b\in\N\), and direct
substitution gives \(\operatorname{pair}_{\N}(a,b)=z\).  Conversely, if
\(z=\operatorname{pair}_{\N}(a,b)\), then \(s=a+b\) is the unique interval index and the displayed
inverse recovers \(b\) and then \(a\).  Thus both surjectivity and injectivity
hold.  The abbreviation is used only to encode pairs; it does not add a new
kind of sequence.

\begin{lemma}[Countability of finite strings]
If \(A\) is countable, then \(A^*\) is countable.
\end{lemma}
\begin{proof}
Fix an injection \(e:A\to\N\) and use the numerical pairing bijection just
defined.  Define injections
\(e_n:A^n\to\N\) by induction: \(e_0(\epsilon)=0\), and for
\(w\in A^{n+1}\), let
\[
  e_{n+1}(w)=\operatorname{pair}_{\N}(e_n(w\restr n),e(w(n))).
\]
Induction on \(n\) proves that every \(e_n\) is injective.  The map sending
\(w\in A^n\) to \(\operatorname{pair}_{\N}(n,e_n(w))\) is therefore an injection
\(A^*\to\N\).
\end{proof}

\begin{corollary}
The set of formulas of a language with a countable alphabet is countable.
\end{corollary}
\begin{proof}
Every formula is a finite string over the alphabet, so the formula set is a
subset of the countable set \(A^*\).
\end{proof}

If \(u\in A^m\) and \(v\in A^n\), their concatenation \(uv\in A^{m+n}\) is
the sequence defined by
\[
  (uv)(i)=
  \begin{cases}
    u(i),&i<m,\\
    v(i-m),&m\leq i<m+n.
  \end{cases}
\]
Concatenation is associative and \(\epsilon\) is its identity.

\begin{lemma}
For all strings \(u,v,w\in A^*\),
\((uv)w=u(vw)\), \(\epsilon u=u\), and \(u\epsilon=u\).
\end{lemma}
\begin{proof}
Two finite sequences are equal iff they have the same length and the same
value at every argument.  Write \(m=\length(u)\), \(n=\length(v)\), and
\(p=\length(w)\).

For associativity, first compare lengths.  By the definition of
concatenation, \(\length(uv)=m+n\), so \(\length((uv)w)=(m+n)+p\); likewise
\(\length(vw)=n+p\), so \(\length(u(vw))=m+(n+p)\).  Addition on \(\N\) is
associative, so both strings are functions on \(\set{0,\ldots,m+n+p-1}\).

Now fix an argument \(i<m+n+p\) and distinguish the three exhaustive and
pairwise disjoint intervals \(i<m\), \(m\leq i<m+n\), and
\(m+n\leq i<m+n+p\).

If \(i<m\), then also \(i<m+n\), so
\(((uv)w)(i)=(uv)(i)=u(i)\); and \((u(vw))(i)=u(i)\) directly, since
\(i<m\).

If \(m\leq i<m+n\), then \(((uv)w)(i)=(uv)(i)=v(i-m)\), by the second
clause for \(uv\) applied at \(i\).  On the other side, \(i\geq m\) gives
\((u(vw))(i)=(vw)(i-m)\), and \(0\leq i-m<n\) puts \(i-m\) in the first
interval for \(vw\), so \((vw)(i-m)=v(i-m)\).

If \(m+n\leq i<m+n+p\), then \(i\geq m+n\) gives
\(((uv)w)(i)=w(i-(m+n))\).  On the other side, \(i\geq m\) gives
\((u(vw))(i)=(vw)(i-m)\), and \(i-m\geq n\) puts \(i-m\) in the second
interval for \(vw\), so \((vw)(i-m)=w((i-m)-n)=w(i-(m+n))\).

In each case the two strings agree at \(i\), so \((uv)w=u(vw)\).

For the identity equations, \(\length(\epsilon u)=0+m=m=\length(u)\), and
every argument \(i<m\) satisfies \(0\leq i<0+m\), so the second clause gives
\((\epsilon u)(i)=u(i-0)=u(i)\); hence \(\epsilon u=u\).  Similarly
\(\length(u\epsilon)=m+0=m\), and every argument \(i<m\) falls in the first
interval, so \((u\epsilon)(i)=u(i)\); hence \(u\epsilon=u\).
\end{proof}

\section{Inductive definitions}
\label{sec:strings-02}

\begin{definition}[Inductively generated set]
Let \(U\) be a set.  Let \(B\subseteq U\) be a set of base elements, and let
\(F\) be a set of finitary operations on \(U\).  The \emph{set inductively
generated by \(B\) under \(F\)} is the intersection of all subsets \(X\) of
\(U\) such that \(B\subseteq X\) and \(X\) is closed under every operation in
\(F\).
\end{definition}

This definition is minimal: it includes the base elements and required
closures, and nothing else.  A grammar for formulas is an inductive definition
of this form.

\begin{theorem}[Structural induction]
Let \(G\) be generated from base elements \(B\) by finitary constructors
\(F\).  Let \(P(x)\) be a metalinguistic property of members of \(G\).  Suppose
that:
\begin{enumerate}
  \item \(P(b)\) holds for every \(b\in B\); and
  \item for every \(k\)-ary constructor \(f\in F\) and all
        \(x_1,\ldots,x_k\in G\), if \(P(x_i)\) holds for every
        \(1\leq i\leq k\), then \(P(f(x_1,\ldots,x_k))\) holds.
\end{enumerate}
Then \(P(x)\) holds for every \(x\in G\).
\end{theorem}
\begin{proof}
Let \(X\defeq\set{x\in G:P(x)}\).  Condition 1 gives \(B\subseteq X\), and
Condition 2 says that \(X\) is closed under every constructor in \(F\).  By the
minimality clause in the definition of \(G\), \(G\subseteq X\).  Hence
\(P(x)\) holds for every \(x\in G\).
\end{proof}

\begin{warning}
A proof by structural induction must contain one case for every base clause and
one case for every constructor clause in the grammar actually in force.  Cases
from a familiar but different grammar do not establish the result.
\end{warning}

\section{Structural recursion}
\label{sec:strings-03}

\begin{theorem}[Structural recursion]
Assume that every member of an inductively generated set \(G\) has a unique
construction from the stated constructors.  To define a function
\(h:G\to Y\), it is enough to specify:
\begin{enumerate}
  \item a value \(h(b)\in Y\) for every base element \(b\); and
  \item for every constructor \(f\) of arity \(k\), a function
        \(H_f:Y^k\to Y\) such that
        \[
          h(f(x_1,\ldots,x_k))
          =H_f(h(x_1),\ldots,h(x_k)).
        \]
\end{enumerate}
Then there exists a unique function \(h\) satisfying those equations.
\end{theorem}

\begin{proof}
Define \(G_0\) as the set of base elements and \(G_{m+1}\) as \(G_m\)
together with every value of a constructor applied to members of \(G_m\);
since \(G\) is generated by the constructors from the base elements,
\(G=\bigcup_{m\in\N}G_m\).  Assign to each member \(x\) of \(G\) the least
\(m\) with \(x\in G_m\), called its \emph{stage}.  Define \(h\) by
recursion on the stage.  At stage \(0\), use the specified
value for the base element.  If
\(x=f(x_1,\ldots,x_k)\) has positive stage, each \(x_i\) has smaller stage,
so its value has already been defined; set
\[
  h(x)=H_f(h(x_1),\ldots,h(x_k)).
\]
Unique construction makes this value independent of any alternative clause,
so the recursion defines a function on all of \(G\).  The defining equations
hold by construction.  If \(h'\) satisfies the same equations, induction on
the stage gives \(h'(b)=h(b)\) for every base element and then
\(h'(f(x_1,\ldots,x_k))=h(f(x_1,\ldots,x_k))\); hence \(h'=h\).
\end{proof}

The unique-construction hypothesis matters.  Without it, two construction
histories for the same expression could impose incompatible values.  Prefix
notation or fully parenthesized infix notation provides an immediate route to
unique readability; conventional precedence rules provide another only after
they have been formally fixed.

\begin{example}[Formula complexity]
For a propositional grammar with variables, negation, and conjunction, define
the complexity \(c:\Fm\to\N\) by
\[
  c(p)=0\text{ for every propositional variable }p,\qquad
  c(\neg\varphi)=c(\varphi)+1,\qquad
  c((\varphi\land\psi))=c(\varphi)+c(\psi)+1.
\]
This is a structural recursion.  It is not an induction proof.
\end{example}

\section{Induction on derivations}
\label{sec:strings-04}

A derivation is itself inductively generated: initial derivations form the
base cases, and inference rules form the constructors.  Consequently, a proof
that every derivation has a property must inspect its final inference.

\begin{method}[Proof by induction on a derivation]
To prove that every derivation \(D\) has property \(P(D)\):
\begin{enumerate}
  \item state the induction measure or invoke structural induction on the
        inductive definition of derivations;
  \item prove the property for every initial derivation;
  \item for each inference rule, assume the property for the immediate
        subderivations and prove it for the derivation formed by that rule; and
  \item verify every side condition, including freshness and discharged
        assumptions.
\end{enumerate}
\end{method}

Soundness proofs will repeatedly use this method.  The local semantic
preservation of each rule becomes the constructor case.

\section{Relations and closures}
\label{sec:strings-05}

\begin{definition}[Reflexive transitive closure]
Let \(R\subseteq X\times X\).  The \emph{reflexive transitive closure} of
\(R\), written \(R^*\), is the least reflexive and transitive relation on
\(X\) that contains \(R\).
\end{definition}

\begin{lemma}[Path characterization of the reflexive transitive closure]
\label{lem:rtc-path-characterization}
Let \(R\subseteq X\times X\) and let \(x,y\in X\).  Then \(xR^*y\) iff there
are \(n\in\N\) and \(x_0,\ldots,x_n\in X\) such that \(x_0=x\), \(x_n=y\),
and \(x_iRx_{i+1}\) for every \(i<n\).
\textit{Status: complete proof.}
\end{lemma}
\begin{proof}
Call a finite sequence \(x_0,\ldots,x_n\) with the displayed three properties
a path from \(x\) to \(y\) along \(R\), and let \(S\subseteq X\times X\) be
the relation holding between \(x\) and \(y\) iff such a path exists.  The
claim is \(R^*=S\), and each inclusion is proved separately.

For \(R^*\subseteq S\), it suffices to show that \(S\) is reflexive and
transitive and contains \(R\), since \(R^*\) is the least such relation.
Reflexivity: for any \(x\), take \(n=0\) and \(x_0=x\); the condition
\(x_iRx_{i+1}\) for every \(i<0\) holds vacuously.  Containment of \(R\): if
\(xRy\), take \(n=1\), \(x_0=x\), \(x_1=y\).  Transitivity: let
\(x_0,\ldots,x_m\) be a path from \(x\) to \(y\) and \(y_0,\ldots,y_n\) a
path from \(y\) to \(z\).  Define \(u_0,\ldots,u_{m+n}\) by \(u_i=x_i\) for
\(i\leq m\) and \(u_{m+j}=y_j\) for \(j\leq n\); the two clauses agree at
the one shared index \(m\), since \(x_m=y=y_0\).  Then \(u_0=x\),
\(u_{m+n}=z\), and for every \(i<m+n\): if \(i<m\), then
\(u_iRu_{i+1}\) by the first path; if \(m\leq i<m+n\), write \(i=m+j\) with
\(j<n\), and then \(u_i=y_jRy_{j+1}=u_{i+1}\) by the second path.  So
\(u_0,\ldots,u_{m+n}\) is a path from \(x\) to \(z\).

For \(S\subseteq R^*\), induction on \(n\) proves that every path
\(x_0,\ldots,x_n\) from \(x\) to \(y\) along \(R\) gives \(xR^*y\).  Base
case \(n=0\): then \(y=x_0=x\), and \(xR^*x\) because \(R^*\) is reflexive.
Inductive step: assume the claim for \(n\), and let \(x_0,\ldots,x_{n+1}\)
be a path from \(x\) to \(y\).  The initial segment \(x_0,\ldots,x_n\) is a
path from \(x\) to \(x_n\), so \(xR^*x_n\) by the induction hypothesis.
Also \(x_nRx_{n+1}\) and \(R\subseteq R^*\), so \(x_nR^*y\).  Transitivity
of \(R^*\) gives \(xR^*y\).
\end{proof}

\begin{definition}[Equivalence closure]
The \emph{equivalence closure} of a relation \(R\subseteq X\times X\) is the
least equivalence relation on \(X\) containing \(R\).
\end{definition}

These closures will be used for reduction sequences, proof normalization, and
quotient constructions.  They must not be confused: a reflexive transitive
closure need not be symmetric.

\begin{exercise}
\label{ex:strings-01}
Give an inductive definition of balanced-parenthesis strings and prove by
structural induction that every generated string contains equally many left
and right parentheses.
\end{exercise}

\begin{exercise}
\label{ex:strings-02}
Define the height of a derivation by structural recursion.  Then prove that
every proper subderivation has strictly smaller height.
\end{exercise}

\chapter{Writing Formal Definitions and Proofs}
\label{ch:writing-foundations}

\section{The dependency order of a formal development}
\label{sec:writing-foundations-01}

A formal development should normally introduce material in the following
order:
\begin{enumerate}
  \item ambient sets and metalinguistic conventions;
  \item signature and arities;
  \item grammar of terms and formulas;
  \item free variable occurrences and substitution;
  \item structures and variable assignments;
  \item term evaluation and satisfaction;
  \item calculus and derivations;
  \item semantic and proof-theoretic properties; and
  \item connections between the calculus and semantics.
\end{enumerate}
An author may reorder material for a stated pedagogical reason, but no
definition may depend on an object that has not been defined.

\section{Definitions}
\label{sec:writing-foundations-02}

\begin{method}[Definition audit]
For every definition, answer the following questions in the definition or its
immediate context.
\begin{enumerate}
  \item What is the type or standing of the defined item?
  \item Over which previously defined objects does it range?
  \item Are all quantifiers explicit?
  \item Are all functions typed by domain and codomain?
  \item Are all arities and binding positions specified?
  \item Is the definition total, or is its domain of definition stated?
  \item If the definition is recursive, what is the well-founded measure or
        inductive grammar?
  \item If the definition uses a representative of an equivalence class, why is
        the result independent of that representative?
  \item If a new term is used, what previously unavailable distinction requires
        it?
\end{enumerate}
\end{method}

\begin{example}[Typed definition]
Let \(X\) be a set and let \(R\subseteq X\times X\).  Define the successor set
function \(R[-]:X\to\Pow(X)\) by
\[
  R[x]\defeq\set{y\in X:xRy}.
\]
The set \(X\), relation type, function type, argument, and value are all
explicit.
\end{example}

\begin{warning}[Circular definition]
A clause of the form ``a formula is an expression that is a formula according
to the usual rules'' does not define a set.  A reference to prior literature
may identify a convention, but a self-contained formal development must state
the grammar it uses.
\end{warning}

\section{Theorem statements}
\label{sec:writing-foundations-03}

A theorem statement should be readable independently of its proof.  It must
identify:
\begin{itemize}
  \item every object quantified over;
  \item every hypothesis, including finiteness, countability, consistency, and
        effectiveness assumptions;
  \item the exact language, calculus, and semantic class;
  \item the direction of each implication; and
  \item whether the result concerns formulas, sentences, theories, structures,
        or derivations.
\end{itemize}

The words ``clearly,'' ``obviously,'' and ``routine'' do not discharge a proof
obligation.  A short proof may cite a previous theorem and verify its
hypotheses; a missing proof may be marked as omitted, but it must not be
presented as established.

\begin{method}[Proof plan before proof]
Before writing details, record:
\begin{enumerate}
  \item the target statement with all quantifiers;
  \item the governing definition or theorem;
  \item the proof method;
  \item the induction parameter or minimal counterexample when applicable;
  \item the object to construct when the proof is existential; and
  \item the invariant maintained at each step.
\end{enumerate}
Delete the plan from a polished proof only if the proof remains easy to audit.
\end{method}

\section{Proof methods}
\label{sec:writing-foundations-04}

\subsection{Direct proof}

To prove \(A\to B\), assume \(A\) and derive \(B\).  In a metatheorem, the
assumption and conclusion must be stated at the same level.  Do not replace a
semantic hypothesis by a syntactic one unless an already proved theorem
licenses the replacement.

\subsection{Proof by contrapositive}

In classical metalogic, proving \(A\to B\) by proving
\(\neg B\to\neg A\) is legitimate.  In intuitionistic object logic, the
formula \((\neg B\to\neg A)\to(A\to B)\) is not generally derivable.  The
ambient logic must therefore be stated when this distinction matters.

\subsection{Proof by contradiction}

To prove \(A\) by contradiction in classical logic, assume \(\neg A\) and
derive \(\bottom\).  In intuitionistic logic this establishes
\(\neg\neg A\), and an additional principle is required to conclude \(A\).

\subsection{Induction}

State the induction domain, induction predicate, base cases, and induction
step.  Structural induction follows the grammar; induction on derivations
follows the final rule; induction on formula complexity follows the defined
complexity measure.  These methods may coincide extensionally in a particular
case, but the proof must identify the one used.

\subsection{Construction}

For an existence theorem, define the proposed witness and verify every required
property.  For a model-existence theorem, this includes a domain, an
interpretation for each symbol, and a satisfaction proof.  Naming a desired
object without constructing it or invoking an existence theorem is not a
construction.

\subsection{Counterexample}

To refute a universal statement, provide an object in the quantified domain and
verify that the predicate fails.  To refute semantic consequence for a
sentence argument, provide a countermodel.  For formulas with free variable
occurrences, provide a structure and variable assignment that satisfy the
premises and falsify the conclusion.  Either form of semantic refutation is
sufficient to refute derivability only after soundness of the stated calculus
has been proved.

\section{Displayed formulas and prose}
\label{sec:writing-foundations-05}

Displayed formulas are parts of sentences and receive punctuation.  Each
symbol must be introduced before the display in which it first appears.
Long transformations should use an aligned environment with a justification at
each nontrivial step:
\[
\begin{aligned}
  \neg(P\to Q)
    &\equiv \neg(\neg P\lor Q) && \text{equivalence of }P\to Q
    \text{ with }\neg P\lor Q\\
    &\equiv P\land\neg Q       && \text{De Morgan and double-negation
    equivalences}.
\end{aligned}
\]
The symbol \(\equiv\) in this example is a metalinguistic equivalence sign; it
is not an object-language connective.  The object-language biconditional will
be written \(\leftrightarrow\).

Halmos, Knuth, Larrabee, and Roberts emphasize that rigorous prose must
remain grammatical and that notation should reduce, not replace, explanation
\cite{Halmos1970,KnuthLarrabeeRoberts1989}.  The AMS author guidance provides
the corresponding publication and \LaTeX{} conventions
\cite{AMSAuthorHandbook}.

\section{Inference-rule typography}
\label{sec:writing-foundations-06}

An inference rule has premises above a horizontal line and a conclusion below
it.  A rule name appears to the right, and side conditions appear explicitly:
\[
  \Rule[\(\forall I\)]{\Gamma\derives\varphi}{\Gamma\derives\forall x\,\varphi}
  \side{provided \(x\) is not free in any member of \(\Gamma\)}.
\]
The displayed rule is only illustrative here.  It does not become part of the
active calculus until that calculus is defined in \cref{ch:fol-deduction}.

\section{Error taxonomy}
\label{sec:writing-foundations-07}

\begin{longtable}{@{}p{0.23\textwidth}p{0.34\textwidth}p{0.35\textwidth}@{}}
\toprule
Error & Diagnostic question & Repair \\
\midrule
\endhead
Level error & Is an object-language formula being confused with a
metalinguistic statement? & Mention the formula by quotation, or quantify
over it with a metavariable, and state the metalinguistic relation. \\
\addlinespace
Grammar error & Does the expression belong to the inductively defined set of
formulas? & Parse it from the grammar or correct it. \\
\addlinespace
Binding error & Which quantifier binds each variable occurrence? & Rename bound
variables and restore parentheses. \\
\addlinespace
Capture error & Does substitution change a free occurrence into a bound
occurrence? & Rename the conflicting bound variable before substitution. \\
\addlinespace
Semantic error & Is truth being asserted without a structure or required
assignment? & Supply the semantic parameters. \\
\addlinespace
Rule error & Is the cited inference rule part of the active calculus, and are
its side conditions met? & State the rule instance and verify every side
condition. \\
\addlinespace
Scope error & Does a connective or quantifier govern the intended subformula? &
Insert parentheses or rewrite with explicit restricted quantification. \\
\addlinespace
Metatheorem error & Have soundness, completeness, consistency, or decidability
been used in the wrong direction? & Write the theorem as a quantified
implication and match its hypotheses. \\
\addlinespace
Strength error & Does the proof establish a weaker statement than the theorem
states? & Strengthen the proof or weaken and relabel the theorem statement. \\
\bottomrule
\end{longtable}

\begin{exercise}
\label{ex:writing-foundations-01}
Rewrite the following statement so that its quantifiers, calculus, and semantic
class are explicit: ``If it is true, it is provable.''  Give one corrected
statement that is a completeness theorem and one corrected statement that is
false.
\end{exercise}

\begin{exercise}
\label{ex:writing-foundations-02}
Audit the definition ``A theory is complete when it proves everything that is
true.''  Identify at least three undefined or ambiguous components and replace
the sentence with the definition of deductive completeness used in
\cref{ch:fol-theories}; then state the separate semantic property and the
hypothesis under which the two properties are equivalent.
\end{exercise}

\part{Classical Propositional Logic}
\label{part:classical-propositional}

Classical propositional logic is the first complete instance of the general
pattern fixed in \cref{ch:levels}: an inductively generated language, a
two-valued semantics, and calculi connected to that semantics by soundness and
completeness.  The treatment follows the standard separation of syntax,
semantics, and derivation in \cite{EbbinghausFlumThomas1994,OpenLogic2026}.

\chapter{Propositional Syntax and Semantics}
\label{ch:prop-syntax-semantics}

\section{Alphabet and grammar}
\label{sec:prop-syntax-semantics-01}

\begin{definition}[Propositional alphabet]
The \emph{propositional alphabet} consists of:
\begin{enumerate}
  \item propositional variables
        \(\Prop\defeq\set{p_n:n\in\N}\);
  \item the nullary connective \(\bottom\);
  \item the unary connective \(\neg\);
  \item the binary connectives \(\land,\lor,\to\); and
  \item the punctuation symbols \((\) and \()\).
\end{enumerate}
All listed symbols are distinct.  A propositional variable is an
object-language symbol.  The letters \(\varphi,\psi,\chi\) used below are
metavariables ranging over formulas.  For readable examples, \(P,Q,R,S\)
denote specified members of \(\Prop\); the indexed names \(p_n\) are used when
the identity of a variable matters.
\end{definition}

\begin{definition}[Propositional formula]
The set \(\Fm_{\CPL}\) of \emph{propositional formulas} is the least set of
strings satisfying the following formation clauses:
\begin{enumerate}
  \item \(p_n\in\Fm_{\CPL}\) for every \(n\in\N\);
  \item \(\bottom\in\Fm_{\CPL}\);
  \item if \(\varphi\in\Fm_{\CPL}\), then
        \(\neg\varphi\in\Fm_{\CPL}\); and
  \item if \(\varphi,\psi\in\Fm_{\CPL}\) and
        \(\circ\in\set{\land,\lor,\to}\), then
        \((\varphi\circ\psi)\in\Fm_{\CPL}\).
\end{enumerate}
No other string is a propositional formula.
\end{definition}

The displayed grammar is
\[
  \varphi ::= p_n\mid\bottom\mid\neg\varphi
  \mid(\varphi\land\varphi)\mid(\varphi\lor\varphi)
  \mid(\varphi\to\varphi).
\]
Parentheses required by the grammar are part of the object language.  In prose
and displays, outermost parentheses may be omitted, and the precedence order
is \(\neg\), \(\land\), \(\lor\), \(\to\).  A chain of repeated
\(\land\), repeated \(\lor\), or repeated \(\to\) associates to the right.
These are typographical abbreviations, not additional formation clauses.

\begin{definition}[Defined connectives]
The symbols \(\topv\) and \(\leftrightarrow\) are metalinguistic
abbreviations:
\[
  \topv\defeq\neg\bottom,
  \qquad
  (\varphi\leftrightarrow\psi)
  \defeq((\varphi\to\psi)\land(\psi\to\varphi)).
\]
They do not add constructors to \(\Fm_{\CPL}\).
\end{definition}

For an ordered finite sequence \(\varphi_1,\ldots,\varphi_n\), define its
finite conjunction and disjunction as right-associated formulas.  When
\(n>0\), put \(C_n=\varphi_n\), \(D_n=\varphi_n\), and recursively put
\[
 C_i=(\varphi_i\land C_{i+1}),
 \qquad
 D_i=(\varphi_i\lor D_{i+1})
 \quad(1\leq i<n);
\]
then
\[
 \bigwedge_{i=1}^{n}\varphi_i\defeq C_1,
 \qquad
 \bigvee_{i=1}^{n}\varphi_i\defeq D_1.
\]
For \(n=0\), define the finite conjunction as \(\topv\) and the finite
disjunction as \(\bottom\).  An indexed finite conjunction or disjunction is
used only after its index order has been stated.  This convention applies to
both propositional and first-order formulas and gives every later finite
connective display one formula-string value.

\begin{lemma}[Parenthesis counting]
\label{lem:paren-count}
Every member of \(\Fm_{\CPL}\) contains equally many left and right
parenthesis occurrences.  Moreover, for every formula and every
binary-connective occurrence in it, the number of left parenthesis
occurrences strictly preceding that occurrence exceeds the number of right
parenthesis occurrences strictly preceding it.
\end{lemma}
\begin{proof}
Both claims proceed by structural induction on the formula.  A propositional
variable and \(\bottom\) contain no parentheses and no binary-connective
occurrences.  The negation clause prefixes one token and adds no
parentheses, so both claims pass through it unchanged.  For
\((\varphi\circ\psi)\), balance holds because the clause adds one left and
one right parenthesis to balanced parts.  A binary-connective occurrence in
\((\varphi\circ\psi)\) is the displayed occurrence of \(\circ\), an
occurrence in \(\varphi\), or an occurrence in \(\psi\).  The tokens
preceding the displayed \(\circ\) are the initial left parenthesis followed
by the balanced string \(\varphi\), so the excess of left over right
parentheses there is one.  For an occurrence in \(\varphi\), the initial
left parenthesis raises the positive excess supplied by the induction
hypothesis by one.  For an occurrence in \(\psi\), the preceding tokens are
the initial left parenthesis, the balanced string \(\varphi\), the
connective, and a prefix of \(\psi\), so the excess is one more than the
positive excess supplied by the induction hypothesis within \(\psi\).
\end{proof}

\begin{theorem}[Unique readability]
\label{thm:prop-unique-readability}
Every member of \(\Fm_{\CPL}\) satisfies exactly one formation clause.  In
the negation case its immediate subformula is unique.  In a binary case its
main connective and its two immediate subformulas are unique.
\end{theorem}
\begin{proof}
A propositional variable begins with a variable token, \(\bottom\) is a
distinct one-token string, a negation begins with \(\neg\), and every binary
formula begins with a left parenthesis.  Hence the four kinds of formation
clause have disjoint ranges.

For a binary formula \((\varphi\circ\psi)\), scan from immediately after
its first left parenthesis while maintaining the number of left minus right
parenthesis occurrences already scanned.  By the balance clause of
\cref{lem:paren-count}, this number is zero at the displayed occurrence of
\(\circ\), because the tokens scanned before it form exactly the balanced
string \(\varphi\).  By the excess clause of \cref{lem:paren-count},
applied to \(\varphi\) and to \(\psi\), every binary connective inside
either immediate subformula occurs while the number is positive.  Hence the
count-zero binary-connective position is unique.  The first formula ends
exactly before that unique position, and the second begins immediately after
it and ends immediately before the final right parenthesis.  Thus the split, the
main connective, and the two immediate subformulas are unique.  Removing an
initial \(\neg\) uniquely determines the immediate subformula in the unary
case.  Therefore every formula has one and only one outermost construction.
\end{proof}

\begin{definition}[Complexity and subformula]
The \emph{complexity} \(c(\varphi)\in\N\) is defined by
\[
\begin{aligned}
 c(p_n)&=0,&c(\bottom)&=0,\\
 c(\neg\varphi)&=c(\varphi)+1,
 &c((\varphi\circ\psi))&=c(\varphi)+c(\psi)+1.
\end{aligned}
\]
The set \(\Sub(\varphi)\) of \emph{subformulas} is defined by
\[
\begin{aligned}
 \Sub(p_n)&=\set{p_n},&\Sub(\bottom)&=\set{\bottom},\\
 \Sub(\neg\varphi)&=\set{\neg\varphi}\cup\Sub(\varphi),\\
 \Sub(\varphi\circ\psi)&=
 \set{\varphi\circ\psi}\cup\Sub(\varphi)\cup\Sub(\psi).
\end{aligned}
\]
Here and below \(\circ\in\set{\land,\lor,\to}\).
\end{definition}

\section{Valuations and recursive truth clauses}
\label{sec:prop-syntax-semantics-02}

\begin{definition}[Truth values and valuation]
Let \(\B\defeq\set{\False,\True}\).  A \emph{valuation} is a function
\(v:\Prop\to\B\).  Its unique extension
\(\widehat v:\Fm_{\CPL}\to\B\) is defined by structural recursion:
\[
\begin{array}{rcl}
\widehat v(p_n)&=&v(p_n),\\
\widehat v(\bottom)&=&\False,\\
\widehat v(\neg\varphi)&=&\True
  \quad\text{iff}\quad\widehat v(\varphi)=\False,\\
\widehat v(\varphi\land\psi)&=&\True
  \quad\text{iff}\quad
  \widehat v(\varphi)=\True\text{ and }\widehat v(\psi)=\True,\\
\widehat v(\varphi\lor\psi)&=&\True
  \quad\text{iff}\quad
  \widehat v(\varphi)=\True\text{ or }\widehat v(\psi)=\True,\\
\widehat v(\varphi\to\psi)&=&\False
  \quad\text{iff}\quad
  \widehat v(\varphi)=\True\text{ and }\widehat v(\psi)=\False.
\end{array}
\]
We write \(v\entails\varphi\) iff
\(\widehat v(\varphi)=\True\).
\end{definition}

\begin{lemma}[Coincidence for propositional formulas]
\label{lem:prop-coincidence}
If valuations \(v\) and \(w\) agree on every propositional variable occurring
in \(\varphi\), then
\(\widehat v(\varphi)=\widehat w(\varphi)\).
\end{lemma}
\begin{proof}
Proceed by structural induction on \(\varphi\).  The variable case is the
hypothesis, and the \(\bottom\) case is immediate.  Each connective case
follows by applying the induction hypotheses to the immediate subformulas and
then applying the corresponding recursive truth clause.
\end{proof}

\begin{definition}[Semantic notions]
Let \(\Gamma\subseteq\Fm_{\CPL}\) and
\(\varphi,\psi\in\Fm_{\CPL}\).
\begin{enumerate}
  \item A valuation \(v\) \emph{satisfies} \(\Gamma\), written
        \(v\entails\Gamma\), iff \(v\entails\gamma\) for every
        \(\gamma\in\Gamma\).
  \item The set \(\Gamma\) is \emph{satisfiable} iff some valuation satisfies
        it.
  \item The formula \(\varphi\) is \emph{valid}, written
        \(\entails\varphi\), iff every valuation satisfies it.
  \item The formula \(\varphi\) is a \emph{semantic consequence} of
        \(\Gamma\), written \(\Gamma\entails\varphi\), iff every valuation
        satisfying \(\Gamma\) satisfies \(\varphi\).
  \item The formulas \(\varphi\) and \(\psi\) are \emph{logically
        equivalent}, written \(\varphi\equiv\psi\), iff every valuation gives
        them the same truth value.
\end{enumerate}
\end{definition}

The empty set is satisfied by every valuation.  Consequently,
\(\varnothing\entails\varphi\) iff \(\entails\varphi\).  A countervaluation to
\(\Gamma\entails\varphi\) is a valuation satisfying \(\Gamma\) and falsifying
\(\varphi\).

\begin{example}[Truth table]
The formula \((P\to Q)\leftrightarrow(\neg P\lor Q)\) is valid:
\[
\begin{array}{cc|c|c}
P&Q&P\to Q&\neg P\lor Q\\ \hline
\True&\True&\True&\True\\
\True&\False&\False&\False\\
\False&\True&\True&\True\\
\False&\False&\True&\True
\end{array}
\]
The table checks all valuations of the two variables occurring in the formula;
\cref{lem:prop-coincidence} shows that values assigned to other variables are
irrelevant.
\end{example}

\section{Translation from controlled prose}
\label{sec:prop-syntax-semantics-03}

A translation requires a key assigning one propositional variable to each
atomic statement.  The key is part of the translation; a formula without its
key does not preserve the content of the prose.

\begin{example}[Necessary and sufficient conditions]
Use the key
\[
  P:\text{the archive is readable},
  \qquad Q:\text{the checksum is valid}.
\]
Then:
\begin{center}
\begin{tabular}{@{}ll@{}}
\toprule
Controlled prose & Formula \\ \midrule
If the archive is readable, then the checksum is valid. & \(P\to Q\) \\
The archive is readable only if the checksum is valid. & \(P\to Q\) \\
The archive is readable if the checksum is valid. & \(Q\to P\) \\
The archive is readable iff the checksum is valid. & \(P\leftrightarrow Q\) \\
The archive is readable unless the checksum is invalid. & \(\neg\neg Q\to P\) \\
\bottomrule
\end{tabular}
\end{center}
The word ``only'' marks a necessary condition, and it reverses the effect
of a bare ``if'': the ``if'' of the third row introduces the antecedent,
while the ``only if'' of the second row introduces the consequent, so the
main clause of the second row supplies the antecedent and the first two
rows share one formula.
For the final row, ``\(P\) unless \(R\)'' is \(\neg R\to P\).  Here
\(R\) is ``the checksum is invalid,'' namely \(\neg Q\), so the translation
is \(\neg\neg Q\to P\), and that formula is what the table records.
Simplifying it to \(Q\to P\) is a further step, justified by an equivalence
of classical propositional logic and not part of the translation.
\end{example}

\begin{example}[Exclusive disjunction]
``Exactly one of \(P\) and \(Q\)'' is translated by
\[
  (P\lor Q)\land\neg(P\land Q).
\]
The formula \(P\lor Q\) alone is inclusive: it is true when both disjuncts are
true.
\end{example}

\section{Normal forms}
\label{sec:prop-syntax-semantics-04}

\begin{definition}[Propositional literal, clause, and formula interpretation]
A \emph{propositional literal} is a propositional variable or the negation of
a propositional variable.  A \emph{clause} is a finite set of propositional
literals.  Thus clause equality is set equality: order and repeated
occurrences are not part of a clause.

To map this set object back to the formula language, put
\[
 \operatorname{code}(p_n)=2n,
 \qquad
 \operatorname{code}(\neg p_n)=2n+1.
\]
For a clause \(C\), let \(k=\abs{C}\); when \(k>0\), list its members as
\(L_1,\ldots,L_k\) in strictly increasing code order and define, by downward
recursion,
\[
 D_k\defeq L_k,
 \qquad
 D_i\defeq(L_i\lor D_{i+1})\quad(1\leq i<k).
\]
Then define
\[
 \operatorname{form}_{\CPL}(C)\defeq
 \begin{cases}
  \bottom,&C=\varnothing,\\
  D_1,&C\ne\varnothing.
 \end{cases}
\]
This is an exact metalinguistic function from clauses to
\(\Fm_{\CPL}\), not a new formula constructor.  A valuation \(v\)
\emph{satisfies} \(C\) iff
\(v\entails\operatorname{form}_{\CPL}(C)\).  In particular, the empty
clause is false under every valuation.
\end{definition}

\begin{definition}[Negation, conjunctive, and disjunctive normal form]
A formula is in \emph{negation normal form} iff it is built from literals and
\(\bottom,\topv\) using only \(\land\) and \(\lor\).  It is in
\emph{conjunctive normal form} iff it is \(\topv\), the empty conjunction,
or there are clauses \(C_1,\ldots,C_m\), with \(m>0\), such that it is
\(N_1\), where the formula sequence \(N_m,\ldots,N_1\) is defined by
\[
 N_m\defeq\operatorname{form}_{\CPL}(C_m),
 \qquad
 N_i\defeq
 \bigl(\operatorname{form}_{\CPL}(C_i)\land N_{i+1}\bigr)
 \quad(1\leq i<m).
\]
Thus a conjunctive-normal-form formula is a right-associated finite
conjunction of the formula interpretations of clauses.  In every finite
conjunction and finite disjunction used in this chapter, the empty conjunction
is \(\topv\) and the empty disjunction is \(\bottom\).  A formula is in
\emph{disjunctive normal form} iff it is a finite disjunction of finite
conjunctions of literals under these conventions.
\end{definition}

\begin{lemma}[Satisfaction of clauses and of clause conjunctions]
\label{lem:prop-clause-satisfaction}
Let \(v\) be a valuation.
\begin{enumerate}
  \item For every clause \(C\):
        \(v\entails\operatorname{form}_{\CPL}(C)\) iff
        \(v\entails L\) for at least one literal \(L\in C\).
  \item For all clauses \(C_1,\ldots,C_m\) with \(m>0\), and the formulas
        \(N_m,\ldots,N_1\) built from them as in the definition of
        conjunctive normal form: \(v\entails N_1\) iff
        \(v\entails\operatorname{form}_{\CPL}(C_j)\) for every \(j\) with
        \(1\leq j\leq m\).
\end{enumerate}
\end{lemma}
\begin{proof}
For 1, if \(C=\varnothing\), then
\(\operatorname{form}_{\CPL}(C)=\bottom\), which no valuation satisfies,
and \(C\) has no member, so both sides fail.  Otherwise let
\(k=\abs{C}>0\), list \(C\) in strictly increasing code order as
\(L_1,\ldots,L_k\), and let \(D_k,\ldots,D_1\) be the formulas of the
clause-interpretation recursion.  By downward induction along that
recursion, \(v\entails D_i\) iff \(v\entails L_j\) for some \(j\) with
\(i\leq j\leq k\): the base case \(D_k=L_k\) is immediate, and the case
\(D_i=(L_i\lor D_{i+1})\) follows from the recursive truth clause for
\(\lor\) together with the induction hypothesis.  Taking \(i=1\) and
\(\operatorname{form}_{\CPL}(C)=D_1\) proves 1.

For 2, by downward induction along the recursion defining
\(N_m,\ldots,N_1\), \(v\entails N_i\) iff
\(v\entails\operatorname{form}_{\CPL}(C_j)\) for every \(j\) with
\(i\leq j\leq m\): the base case \(N_m=\operatorname{form}_{\CPL}(C_m)\)
is immediate, and the case
\(N_i=(\operatorname{form}_{\CPL}(C_i)\land N_{i+1})\) follows from the
recursive truth clause for \(\land\) together with the induction
hypothesis.  Taking \(i=1\) proves 2.
\end{proof}

\begin{theorem}[Negation normal form]
Every propositional formula is logically equivalent to a formula in negation
normal form.
\end{theorem}
\begin{proof}
First eliminate every occurrence of \(\to\) using
\(\varphi\to\psi\equiv\neg\varphi\lor\psi\).  Proceed by structural
induction through the resulting formula.  Double negations are removed, and a
negation immediately above a binary connective is moved inward by
\[
 \neg(\varphi\land\psi)\equiv\neg\varphi\lor\neg\psi,
 \qquad
 \neg(\varphi\lor\psi)\equiv\neg\varphi\land\neg\psi.
\]
The base cases are variables, their negations, \(\bottom\), and \(\topv\).
Each displayed replacement decreases the number of connectives lying below a
negation, so the process terminates.  Every replacement is checked by the
recursive truth clauses and therefore preserves logical equivalence.
\end{proof}

\begin{theorem}[Conjunctive and disjunctive normal forms]
\label{thm:prop-normal-forms}
Every propositional formula is logically equivalent to a formula in
disjunctive normal form and to a formula in conjunctive normal form.
\end{theorem}
\begin{proof}
Let \(p_1,\ldots,p_n\) be the variables occurring in \(\varphi\), listed in
increasing index order.  In this
proof only, a \emph{valuation on} a finite variable set is a function from
that set to \(\B\); the recursive truth clauses apply verbatim to formulas
whose variables lie in the set, and \cref{lem:prop-coincidence} identifies
satisfaction under such a function with satisfaction under any total
valuation extending it.  For each
valuation \(v\) on these variables define
\[
  \delta_v\defeq\bigwedge_{i=1}^{n}\ell_i^v,
  \qquad
  \ell_i^v\defeq
  \begin{cases}
    p_i,&v(p_i)=\True,\\
    \neg p_i,&v(p_i)=\False.
  \end{cases}
\]
A valuation on the empty variable set is the empty function, and when
\(n=0\) the displayed conjunction is \(\topv\) by the preceding convention.
A valuation \(w\) satisfies \(\delta_v\) iff \(w(p_i)=v(p_i)\) for every
\(i\).  Let \(V_{\varphi}\) be the finite set of valuations on
\(p_1,\ldots,p_n\) that satisfy \(\varphi\).  Order these valuations
lexicographically by
\((v(p_1),\ldots,v(p_n))\), with \(\False<\True\), enumerate the resulting
order as \(v_1,\ldots,v_m\), and put
\[
  D_{\varphi}\defeq\bigvee_{j=1}^{m}\delta_{v_j}.
\]
If \(V_{\varphi}=\varnothing\), then \(m=0\) and this disjunction is
\(\bottom\).  By the
coincidence lemma, an arbitrary valuation satisfies \(D_{\varphi}\) iff
it satisfies \(\varphi\).  Hence \(D_{\varphi}\) is the required
disjunctive normal form.

For the conjunctive normal form, construct dually.  For each valuation
\(v\) on \(p_1,\ldots,p_n\) put
\[
  \bar\ell^{\,v}_i\defeq
  \begin{cases}
    \neg p_i,&v(p_i)=\True,\\
    p_i,&v(p_i)=\False,
  \end{cases}
  \qquad
  K_v\defeq\set{\bar\ell^{\,v}_1,\ldots,\bar\ell^{\,v}_n}.
\]
The coding in the clause-interpretation definition gives the variable of
index \(i\) the code \(2i\) and its negation the code \(2i+1\), both
strictly increasing in \(i\).  The enumeration \(p_1,\ldots,p_n\) is in
increasing index order---the indices need not form an initial
segment---so any literal on \(p_j\) has smaller code than any literal on
\(p_k\) for \(j<k\): if their indices are \(i_j<i_k\), then
\(2i_j+1<2i_k\).  The listing of \(K_v\) in strictly increasing code order
is therefore exactly \(\bar\ell^{\,v}_1,\ldots,\bar\ell^{\,v}_n\), and
\(\gamma_v\defeq\operatorname{form}_{\CPL}(K_v)\) requires no
re-ordering.  A valuation \(w\) falsifies \(\bar\ell^{\,v}_i\) iff
\(w(p_i)=v(p_i)\), so \cref{lem:prop-clause-satisfaction}(1) gives:
\(w\) falsifies \(\gamma_v\) iff \(w(p_i)=v(p_i)\) for every
\(i\leq n\).

Let \(W_\varphi\) be the set of valuations on \(p_1,\ldots,p_n\) that
falsify \(\varphi\), enumerated in the lexicographic order fixed above as
\(v_1,\ldots,v_m\).  Put \(C_\varphi\defeq\topv\) if \(m=0\); otherwise
put \(C_\varphi\defeq N_1\), where \(N_m,\ldots,N_1\) are built from the
clauses \(C_j\defeq K_{v_j}\) by the recursion in the definition of
conjunctive normal form.  In every case \(C_\varphi\) satisfies that
definition verbatim, the degenerate cases included: for valid
\(\varphi\), \(m=0\) and \(C_\varphi=\topv\); for variable-free true
\(\varphi\), \(n=m=0\) and again \(C_\varphi=\topv\); for variable-free
false \(\varphi\), \(n=0\) and \(m=1\) with \(v_1\) the empty valuation,
so \(K_{v_1}=\varnothing\) and
\(C_\varphi=\operatorname{form}_{\CPL}(\varnothing)=\bottom\), the
conjunctive normal form whose single clause is empty.

Let \(w\) be an arbitrary valuation and \(w_0\) its restriction to
\(\set{p_1,\ldots,p_n}\).  Suppose \(m>0\).  By
\cref{lem:prop-clause-satisfaction}(2) and (1), \(w\) satisfies
\(C_\varphi\) iff for every \(j\leq m\) it satisfies some literal of
\(K_{v_j}\), iff \(w_0\ne v_j\) for every \(j\leq m\), iff
\(w_0\notin W_\varphi\), iff \(w_0\) satisfies \(\varphi\), iff, by the
coincidence lemma, \(w\) satisfies \(\varphi\).  If \(m=0\), every
valuation on the variable set satisfies \(\varphi\), so by the
coincidence lemma every valuation satisfies both \(\varphi\) and
\(C_\varphi=\topv\).  Hence \(C_\varphi\) is the required conjunctive
normal form, obtained from the definitional recursions,
\cref{lem:prop-clause-satisfaction}, and the coincidence lemma alone.
\end{proof}

An \(n\)-ary \emph{truth-functional connective} is a connective symbol
interpreted by a fixed function \(\B^n\to\B\).  For a set \(C\) of
truth-functional connectives, the \emph{formulas over \(C\)} are generated
by the propositional variables together with one formation clause per member
of \(C\), and the recursive truth clause of each member applies its
interpreting function to the values of the immediate subformulas.  For a
formula \(\varphi\) whose variables are among \(p_1,\ldots,p_n\), the
\emph{truth function of \(\varphi\) in those variables} is the function
\(f_\varphi:\B^n\to\B\) with
\(f_\varphi(a_1,\ldots,a_n)\defeq\widehat v(\varphi)\) for any valuation
\(v\) with \(v(p_i)=a_i\) for each \(i\); this is well defined by
\cref{lem:prop-coincidence}.

\begin{definition}[Functional completeness]
A set \(C\) of truth-functional connectives is \emph{functionally complete}
iff, for every \(n\in\N\) and every function \(f:\B^n\to\B\), some formula
built from variables \(p_1,\ldots,p_n\) and connectives in \(C\) has truth
function \(f\) in those variables.
\end{definition}

\begin{theorem}[A functionally complete basis]
The set \(\set{\neg,\land,\bottom}\) is functionally complete.
\end{theorem}
\begin{proof}
Given \(f:\B^n\to\B\), form the disjunctive normal form whose conjunctions of
literals correspond exactly to the inputs \(a\in\B^n\) for which
\(f(a)=\True\), as in \cref{thm:prop-normal-forms}.  Replace each occurrence
of \(\alpha\lor\beta\) by
\(\neg(\neg\alpha\land\neg\beta)\).  The empty disjunction is \(\bottom\).
The resulting formula uses only \(\neg,\land,\bottom\) and computes \(f\).
\end{proof}

\chapter{Natural Deduction and Sequents}
\label{ch:prop-deduction}

\section{Classical natural deduction}
\label{sec:prop-deduction-01}

\begin{definition}[Natural-deduction derivation]
A \emph{natural-deduction derivation} is a finite rooted tree generated by the
rules displayed below.  Leaves may be marked assumptions.  An application of
an introduction or elimination rule may discharge the marked assumptions
indicated in brackets.  An assumption is \emph{open} iff no rule on the path
from that assumption to the root discharges it.

Discharge labels are metalinguistic markers drawn from \(\N\).  A consistent
injective renaming of the labels on assumptions and of the matching labels on
rules preserves a derivation and which assumptions it discharges.  Since a
derivation is finite, copies of derivations can always be given pairwise
disjoint label sets fresh from every label in a fixed receiving derivation.

An application of \(\to I\), \(\neg I\), or \(\mathsf{RAA}\) carries a
label \(i\) and discharges exactly the open occurrences marked \(i\) of
the bracketed assumption above its premise; an application of \(\lor E\)
carries two labels \(i,j\), as displayed, and discharges exactly the open
occurrences of the first disjunct marked \(i\) above its second premise
and the open occurrences of the second disjunct marked \(j\) above its
third premise.  Every label carried by such an application discharges at
least one occurrence.  Marking any set of open occurrences of one formula
with a common label fresh from every label in the derivation preserves
the derivation and the occurrences discharged by each of its rule
applications, since no rule application carries the fresh label; this
merging relabeling is distinct from the injective renamings above, and is
the operation used when a newly added root inference is to discharge
several open occurrences at once.  There is no repetition step: where an
informal presentation repeats a formula, the tree reuses the node already
present.
\end{definition}

The active calculus \(\ND_{\CPL}\) has the following rules.
\[
\begin{array}{c@{\qquad}c@{\qquad}c}
\Rule[\(\land I\)]{\varphi\quad\psi}{\varphi\land\psi}
&
\Rule[\(\land E_1\)]{\varphi\land\psi}{\varphi}
&
\Rule[\(\land E_2\)]{\varphi\land\psi}{\psi}
\\[2.2ex]
\Rule[\(\lor I_1\)]{\varphi}{\varphi\lor\psi}
&
\Rule[\(\lor I_2\)]{\psi}{\varphi\lor\psi}
&
\Rule[\(\to E\)]{\varphi\quad\varphi\to\psi}{\psi}
\end{array}
\]
\[
\Rule[\(\lor E\)]{
  \varphi\lor\psi\qquad
  \begin{array}{c}[\varphi]^i\\ \vdots\\ \chi\end{array}
  \qquad
  \begin{array}{c}[\psi]^j\\ \vdots\\ \chi\end{array}}
 {\chi}
\side{discharging all assumptions marked \(i\) and \(j\)}.
\]
\[
\begin{array}{c@{\qquad}c@{\qquad}c}
\Rule[\(\to I\)]{
  \begin{array}{c}[\varphi]^i\\ \vdots\\ \psi\end{array}}
 {\varphi\to\psi}
&
\Rule[\(\neg I\)]{
  \begin{array}{c}[\varphi]^i\\ \vdots\\ \bottom\end{array}}
 {\neg\varphi}
&
\Rule[\(\neg E\)]{\varphi\quad\neg\varphi}{\bottom}
\\[4ex]
\Rule[\(\bottom E\)]{\bottom}{\varphi}
&
\Rule[\(\mathsf{RAA}\)]{
  \begin{array}{c}[\neg\varphi]^i\\ \vdots\\ \bottom\end{array}}
 {\varphi}
&
\end{array}
\]
The rules \(\to I\), \(\neg I\), and \(\mathsf{RAA}\) discharge the
indicated assumptions.  The last rule is the explicitly classical rule.

\begin{definition}[Propositional derivability]
For \(\Gamma\subseteq\Fm_{\CPL}\), write
\(\Gamma\derives_{\ND_{\CPL}}\varphi\) iff there is a
\(\ND_{\CPL}\)-derivation of \(\varphi\) whose open assumptions all belong to
\(\Gamma\).  A \emph{proof} is such a derivation together with its stated set
of permitted open assumptions.
\end{definition}

\begin{example}[Commutativity of conjunction]
The following derivation proves
\(\derives_{\ND_{\CPL}}(P\land Q)\to(Q\land P)\):
\[
\begin{array}{rll}
1.&[P\land Q]^1&\text{assumption}\cr
2.&Q&\land E_2,1\cr
3.&P&\land E_1,1\cr
4.&Q\land P&\land I,2,3\cr
5.&(P\land Q)\to(Q\land P)&\to I,1\text{--}4.
\end{array}
\]
The superscript marks the assumption discharged at line 5.
\end{example}

\begin{example}[Excluded middle]
The formula \(P\lor\neg P\) has a classical proof:
\[
\begin{array}{rll}
1.&[\neg(P\lor\neg P)]^1&\text{assumption}\cr
2.&\quad[P]^2&\text{assumption}\cr
3.&\quad P\lor\neg P&\lor I_1,2\cr
4.&\quad\bottom&\neg E,1,3\cr
5.&\neg P&\neg I,2\text{--}4\cr
6.&P\lor\neg P&\lor I_2,5\cr
7.&\bottom&\neg E,1,6\cr
8.&P\lor\neg P&\mathsf{RAA},1\text{--}7.
\end{array}
\]
The proof uses \(\mathsf{RAA}\).  The calculus obtained by omitting
\(\mathsf{RAA}\) is intuitionistic natural deduction, made exact in
\cref{ch:intuitionistic-logic}; no derivation in that calculus concludes
\(P\lor\neg P\) from no open assumptions.
\end{example}

\begin{lemma}[Substitution for open assumptions]
\label{lem:nd-substitution}
If \(D\) derives \(\psi\) with an open assumption \(\varphi\), and \(E\)
derives \(\varphi\), then replacing each designated open occurrence of
\(\varphi\) in \(D\) by a copy of \(E\), after consistently renaming every
label in every copy to be globally fresh from \(D\) and from all other copies,
produces a derivation of \(\psi\).  Its open assumptions are those of \(D\)
other than the replaced occurrences, together with the correspondingly
renamed open assumptions of the inserted copies of \(E\).
\end{lemma}
\begin{proof}
First reserve for each designated occurrence a set of labels disjoint from
the labels of \(D\) and from the sets reserved for the other occurrences, and
consistently rename its copy of \(E\) into that set.  Strengthen the induction
invariant by requiring every inserted copy to use its reserved globally fresh
labels.  Induct on the height of \(D\).  If \(D\) is the designated
assumption, use its renamed copy of \(E\).  If it is another assumption, leave
it unchanged.  For an inference, apply the induction hypothesis to each
immediate subderivation and then reapply the same rule with the original
discharge labels of \(D\).  Consistent renaming preserves every discharge
internal to a copy of \(E\), while global freshness makes every label of that
copy different from every label on a rule of \(D\).  Hence no rule above a
graft point can accidentally discharge an open assumption of the inserted
copy, and the stated open-assumption accounting follows.
\end{proof}

\section{Classical sequent calculus}
\label{sec:prop-deduction-02}

\begin{definition}[Sequent]
A \emph{finite multiset of formulas} is a function
\(\Gamma:\Fm_{\CPL}\to\N\) whose support
\[
 \operatorname{supp}(\Gamma)
 \defeq\{\varphi\in\Fm_{\CPL}:\Gamma(\varphi)>0\}
\]
is finite.  The empty multiset is the constant-zero function.  If
\(\Gamma\) and \(\Pi\) are finite multisets, then the comma
\(\Gamma,\Pi\) denotes their pointwise sum; a displayed formula \(\varphi\)
in a comma list denotes the multiset having value \(1\) at \(\varphi\) and
\(0\) elsewhere.  Thus \(\Gamma,\varphi,\varphi\) increases the multiplicity
of \(\varphi\) by two.  In the first-order sequent chapter the same definition
is used with \(\Fm(\Lang)\) in place of \(\Fm_{\CPL}\).

A \emph{sequent} is an expression \(\Gamma\Rightarrow\Delta\), where
\(\Gamma\) and \(\Delta\) are finite multisets of formulas.  It is satisfied
by a valuation iff either some formula in \(\operatorname{supp}(\Gamma)\) is
false or some formula in \(\operatorname{supp}(\Delta)\) is true.  With
respect to satisfaction, an empty antecedent therefore behaves as
\(\topv\) and an empty succedent as \(\bottom\).

When a finite multiset occurs to the left of \(\entails\) or \(\derives\),
this document defines
\[
 \Gamma\entails\varphi
 \quad\Longleftrightarrow\quad
 \operatorname{supp}(\Gamma)\entails\varphi,
 \qquad
 \Gamma\derives_{\mathcal C}\varphi
 \quad\Longleftrightarrow\quad
 \operatorname{supp}(\Gamma)\derives_{\mathcal C}\varphi.
\]
Multiplicity is therefore proof-structural data in a sequent, while semantic
consequence and natural-deduction assumptions depend only on the support.
\end{definition}

The calculus \(\LK_{\CPL}\) has two forms of \emph{initial sequent}:
identity sequents \(\Gamma,\varphi\Rightarrow\varphi,\Delta\) and
\(\bottom\)-sequents \(\Gamma,\bottom\Rightarrow\Delta\).  It also has
the following rules.  Exchange is
implicit because contexts are multisets.
\[
\begin{array}{c@{\qquad}c}
\Rule[\(W_L\)]{\Gamma\Rightarrow\Delta}{\Gamma,\varphi\Rightarrow\Delta}
&
\Rule[\(W_R\)]{\Gamma\Rightarrow\Delta}{\Gamma\Rightarrow\Delta,\varphi}
\\[2.4ex]
\Rule[\(C_L\)]{\Gamma,\varphi,\varphi\Rightarrow\Delta}
 {\Gamma,\varphi\Rightarrow\Delta}
&
\Rule[\(C_R\)]{\Gamma\Rightarrow\Delta,\varphi,\varphi}
 {\Gamma\Rightarrow\Delta,\varphi}
\end{array}
\]
\[
\begin{array}{c@{\qquad}c}
\Rule[\(\neg L\)]{\Gamma\Rightarrow\Delta,\varphi}
 {\Gamma,\neg\varphi\Rightarrow\Delta}
&
\Rule[\(\neg R\)]{\Gamma,\varphi\Rightarrow\Delta}
 {\Gamma\Rightarrow\Delta,\neg\varphi}
\\[2.6ex]
\Rule[\(\land L\)]{\Gamma,\varphi,\psi\Rightarrow\Delta}
 {\Gamma,\varphi\land\psi\Rightarrow\Delta}
&
\Rule[\(\land R\)]{
 \Gamma\Rightarrow\Delta,\varphi\qquad
 \Gamma\Rightarrow\Delta,\psi}
 {\Gamma\Rightarrow\Delta,\varphi\land\psi}
\\[2.6ex]
\Rule[\(\lor L\)]{
 \Gamma,\varphi\Rightarrow\Delta\qquad
 \Gamma,\psi\Rightarrow\Delta}
 {\Gamma,\varphi\lor\psi\Rightarrow\Delta}
&
\Rule[\(\lor R\)]{\Gamma\Rightarrow\Delta,\varphi,\psi}
 {\Gamma\Rightarrow\Delta,\varphi\lor\psi}
\\[2.6ex]
\Rule[\(\to L\)]{
 \Gamma\Rightarrow\Delta,\varphi\qquad
 \Gamma,\psi\Rightarrow\Delta}
 {\Gamma,\varphi\to\psi\Rightarrow\Delta}
&
\Rule[\(\to R\)]{\Gamma,\varphi\Rightarrow\Delta,\psi}
 {\Gamma\Rightarrow\Delta,\varphi\to\psi}.
\end{array}
\]
The cut rule is
\[
 \Rule[Cut]{\Gamma\Rightarrow\Delta,\varphi\qquad
             \Sigma,\varphi\Rightarrow\Pi}
            {\Gamma,\Sigma\Rightarrow\Delta,\Pi}.
\]
A \(\LK_{\CPL}\)-derivation is a finite tree whose leaves are initial
sequents and whose internal nodes are licensed rule instances.

\begin{example}[A sequent derivation]
The sequent \(\Rightarrow P\to(P\lor Q)\) has the cut-free derivation
\[
 \Rule[\(\to R\)]{
   \Rule[\(\lor R\)]{P\Rightarrow P,Q}{P\Rightarrow P\lor Q}}
  {\Rightarrow P\to(P\lor Q)}.
\]
The upper sequent is an identity sequent.
\end{example}

\chapter{Soundness, Completeness, Compactness, and Resolution}
\label{ch:prop-metatheory}

\section{Soundness}
\label{sec:prop-metatheory-01}

\begin{theorem}[Soundness of classical propositional natural deduction]
\label{thm:prop-nd-soundness}
For all \(\Gamma\subseteq\Fm_{\CPL}\) and
\(\varphi\in\Fm_{\CPL}\),
\[
 \Gamma\derives_{\ND_{\CPL}}\varphi
 \quad\Longrightarrow\quad
 \Gamma\entails\varphi.
\]
\end{theorem}
\begin{proof}
Induct on the construction of the derivation.  The induction statement is
stronger than the theorem: every valuation satisfying all open assumptions of
a subderivation satisfies its conclusion.

An assumption case is immediate.  The conjunction introduction and
elimination cases, disjunction introduction cases, implication elimination
case, negation elimination case, and bottom elimination case follow directly
from their recursive truth clauses.  For \(\lor E\), let \(v\) satisfy the
open assumptions of the conclusion.  It satisfies \(\varphi\lor\psi\), so it
satisfies \(\varphi\) or \(\psi\).  In the first case the induction hypothesis
for the left subsidiary derivation gives \(v\entails\chi\); in the second case
the right subsidiary derivation gives the same conclusion.

For \(\to I\), let \(v\) satisfy the open assumptions.  If
\(v\notentails\varphi\), then \(v\entails\varphi\to\psi\).  If
\(v\entails\varphi\), it satisfies every open assumption of the subsidiary
derivation, so the induction hypothesis gives \(v\entails\psi\).  Thus again
\(v\entails\varphi\to\psi\).  The \(\neg I\) case is analogous: if
\(v\entails\varphi\), the induction hypothesis would give
\(v\entails\bottom\), which is impossible; hence
\(v\entails\neg\varphi\).

For \(\mathsf{RAA}\), suppose \(v\) satisfies the open assumptions
and falsifies \(\varphi\).  Then it satisfies \(\neg\varphi\), so the
induction hypothesis for the subsidiary derivation gives
\(v\entails\bottom\), impossible.  Therefore \(v\entails\varphi\).  Every
rule preserves truth from open assumptions to conclusion, completing the
induction.
\end{proof}

\begin{theorem}[Soundness of the classical propositional sequent calculus]
\label{thm:prop-lk-soundness}
Every sequent derivable in \(\LK_{\CPL}\) is satisfied by every valuation.
\end{theorem}
\begin{proof}
Induct on the sequent derivation, showing that every derivable sequent is
satisfied by every valuation in the sense of the sequent-satisfaction
definition of \cref{sec:prop-deduction-02}.  Both forms of initial sequent
are satisfied by every valuation \(v\).  In an identity sequent
\(\Gamma,\varphi\Rightarrow\varphi,\Delta\), the formula \(\varphi\)
belongs to the support of the antecedent and to the support of the
succedent, so whichever truth value \(v\) gives \(\varphi\), either some
antecedent formula is false under \(v\) or some succedent formula is true
under \(v\).  In a \(\bottom\)-sequent
\(\Gamma,\bottom\Rightarrow\Delta\), the antecedent support contains
\(\bottom\), which is false under \(v\).

For the induction step, suppose toward a contradiction that some valuation
\(v\) does not satisfy the conclusion of a rule each of whose premises is
satisfied by every valuation.  By the sequent-satisfaction definition,
\(v\) makes every formula in the support of the conclusion's antecedent
true and every formula in the support of its succedent false.  For
weakening, the supports of the premise are included in those of the
conclusion, so \(v\) does not satisfy the premise.  For contraction,
premise and conclusion have the same supports, so exactly the same
valuations satisfy both.  For cut, \(v\) makes all formulas in
\(\Gamma,\Sigma\) true and all formulas in \(\Delta,\Pi\) false; because
\(v\) satisfies the left premise, \(\varphi\) is true under \(v\), and
because \(v\) satisfies the right premise, \(\varphi\) is false under
\(v\), which is impossible.

For each logical rule, use the truth clause for the principal connective to
produce a premise that \(v\) does not satisfy.  For example, if \(v\)
does not satisfy the conclusion of \(\land R\), then \(v\) makes every
formula in \(\Gamma\) true, every formula in \(\Delta\) false, and
\(\varphi\land\psi\) false; hence \(\varphi\) or \(\psi\) is false under
\(v\), so \(v\) does not satisfy the corresponding premise.  If \(v\)
does not satisfy the conclusion of \(\to L\), then \(v\) makes
\(\varphi\to\psi\) true.  If \(\varphi\) is false under \(v\), then
\(v\) does not satisfy the left premise; if \(\psi\) is true under
\(v\), then \(v\) does not satisfy the right premise.  The cases for
\(\neg L,\neg R,\land L,\lor L,\lor R,\to R\) follow by the same direct
application of their displayed truth clauses.  Thus every rule whose
premises are satisfied by every valuation has a conclusion satisfied by
every valuation, completing the induction.
\end{proof}

\section{Maximal consistent sets}
\label{sec:prop-metatheory-02}

\begin{definition}[Consistency]
A set \(\Gamma\subseteq\Fm_{\CPL}\) is \emph{consistent} iff
\(\Gamma\notderives_{\ND_{\CPL}}\bottom\).
\end{definition}

\begin{definition}[Maximal consistent set]
A \emph{maximal consistent set} is a consistent set
\(\Delta\subseteq\Fm_{\CPL}\) such that
\(\Delta\cup\set{\varphi}\) is inconsistent for every
\(\varphi\notin\Delta\).
\end{definition}

\begin{lemma}[Closure properties]
\label{lem:prop-mcs}
If \(\Delta\) is maximal consistent, then:
\begin{enumerate}
  \item if \(\Delta\derives_{\ND_{\CPL}}\varphi\), then
        \(\varphi\in\Delta\);
  \item exactly one of \(\varphi\) and \(\neg\varphi\) belongs to
        \(\Delta\);
  \item \(\varphi\land\psi\in\Delta\) iff
        \(\varphi\in\Delta\) and \(\psi\in\Delta\);
  \item \(\varphi\lor\psi\in\Delta\) iff
        \(\varphi\in\Delta\) or \(\psi\in\Delta\); and
  \item \(\varphi\to\psi\in\Delta\) iff
        \(\varphi\notin\Delta\) or \(\psi\in\Delta\).
\end{enumerate}
\end{lemma}
\begin{proof}
For 1, suppose \(\Delta\derives\varphi\) but \(\varphi\notin\Delta\).
Maximality makes \(\Delta\cup\set{\varphi}\) inconsistent.  Substitute the
derivation of \(\varphi\) for every open use of the added assumption by
\cref{lem:nd-substitution}; this yields a derivation of \(\bottom\) from
\(\Delta\), contrary to consistency.

The two formulas in 2 cannot both belong to \(\Delta\), by \(\neg E\).  If
\(\varphi\notin\Delta\), maximality yields a derivation \(D\) of
\(\bottom\) whose open assumptions lie in \(\Delta\cup\set{\varphi}\).
Because \(\Delta\) is consistent, at least one open occurrence of
\(\varphi\) remains in \(D\).  Mark all open occurrences of \(\varphi\) in
\(D\) with a common fresh label \(i\), by the merging relabeling of the
derivation definition, and apply \(\neg I\) with label \(i\) at the root,
discharging every such occurrence, to obtain
\(\Delta\derives\neg\varphi\); then 1 gives \(\neg\varphi\in\Delta\).
This proves 2.

For 3, conjunction elimination proves the forward implication and conjunction
introduction proves the reverse implication, using 1.  For 4, disjunction
introduction proves the reverse implication.  If
\(\varphi\lor\psi\in\Delta\) while neither disjunct belongs to \(\Delta\),
then 2 gives \(\neg\varphi,\neg\psi\in\Delta\).  Disjunction elimination,
using each negation to derive \(\bottom\) in its corresponding branch, derives
\(\bottom\) from \(\Delta\), a contradiction.

For 5, if \(\varphi\to\psi\in\Delta\) and \(\varphi\in\Delta\), implication
elimination and 1 give \(\psi\in\Delta\).  Conversely, if
\(\psi\in\Delta\), assume \(\varphi\) with label \(i\) and take \(\psi\)
as an open assumption; \(\land I\) gives \(\varphi\land\psi\), and
\(\land E_2\) gives \(\psi\); then \(\to I\) discharges the single open
occurrence marked \(i\), concluding \(\varphi\to\psi\) with open
assumptions in \(\Delta\).  The discharged occurrence is genuinely
present, and only displayed rules occur.  If \(\varphi\notin\Delta\), then
\(\neg\varphi\in\Delta\) by 2; from an assumed \(\varphi\), derive
\(\bottom\), then \(\psi\), and discharge the assumption, which occurs
open in that derivation.  In either case 1
places \(\varphi\to\psi\) in \(\Delta\).
\end{proof}

\begin{lemma}[Lindenbaum lemma]
\label{lem:prop-lindenbaum}
Every consistent subset of \(\Fm_{\CPL}\) is contained in a maximal
consistent set.
\end{lemma}
\begin{proof}
The set \(\Fm_{\CPL}\) is countable; enumerate it as
\(\varphi_0,\varphi_1,\ldots\).  Starting with a consistent set
\(\Gamma_0=\Gamma\), define
\[
 \Gamma_{n+1}=
 \begin{cases}
   \Gamma_n\cup\set{\varphi_n},&
     \text{if this set is consistent},\\
   \Gamma_n,&\text{otherwise}.
 \end{cases}
\]
Let \(\Delta=\bigcup_{n\in\N}\Gamma_n\).  If \(\Delta\) were inconsistent,
a finite derivation of \(\bottom\) would use only finitely many assumptions,
all lying in some \(\Gamma_n\), contradicting its consistency.  Thus
\(\Delta\) is consistent.  If \(\varphi_n\notin\Delta\), it was rejected at
stage \(n+1\), so \(\Gamma_n\cup\set{\varphi_n}\), and hence
\(\Delta\cup\set{\varphi_n}\), is inconsistent.  Therefore \(\Delta\) is
maximal consistent.
\end{proof}

\begin{lemma}[Truth lemma]
\label{lem:prop-truth}
Let \(\Delta\) be maximal consistent and define a valuation \(v_\Delta\) by
\[
  v_\Delta(p_n)=\True\quad\text{iff}\quad p_n\in\Delta.
\]
Then, for every propositional formula \(\varphi\),
\[
  v_\Delta\entails\varphi
  \quad\text{iff}\quad
  \varphi\in\Delta.
\]
\end{lemma}
\begin{proof}
Proceed by structural induction.  The variable case is the definition of
\(v_\Delta\), and \(\bottom\notin\Delta\) by consistency.  The negation case
uses \cref{lem:prop-mcs}(2) and the truth clause for \(\neg\).  The
conjunction, disjunction, and implication cases use respectively clauses 3,
4, and 5 of that lemma together with the corresponding recursive truth
clauses.  These are all constructors in the grammar, so the induction is
complete.
\end{proof}

\section{Completeness and compactness}
\label{sec:prop-metatheory-03}

\begin{theorem}[Strong completeness of classical propositional natural deduction]
\label{thm:prop-completeness}
For every \(\Gamma\subseteq\Fm_{\CPL}\) and every
\(\varphi\in\Fm_{\CPL}\),
\[
  \Gamma\entails\varphi
  \quad\Longrightarrow\quad
  \Gamma\derives_{\ND_{\CPL}}\varphi.
\]
\end{theorem}
\begin{proof}
We prove the contrapositive.  Suppose
\(\Gamma\notderives_{\ND_{\CPL}}\varphi\).  Then
\(\Gamma\cup\set{\neg\varphi}\) is consistent.  For suppose a derivation
\(D\) of \(\bottom\) has open assumptions in
\(\Gamma\cup\set{\neg\varphi}\).  If no open occurrence of
\(\neg\varphi\) appeared in \(D\), then \(D\) would witness
\(\Gamma\derives_{\ND_{\CPL}}\bottom\), and \(\bottom E\) would give
\(\Gamma\derives_{\ND_{\CPL}}\varphi\), contrary to the supposition.  So
at least one open occurrence remains; mark all open occurrences of
\(\neg\varphi\) in \(D\) with a common fresh label \(i\), by the merging
relabeling of the derivation definition, and apply \(\mathsf{RAA}\) with
label \(i\) at the root, discharging every such occurrence, to obtain
\(\Gamma\derives_{\ND_{\CPL}}\varphi\)---again contrary to the
supposition.  Hence \(\Gamma\cup\set{\neg\varphi}\) is consistent.  By
\cref{lem:prop-lindenbaum}, extend this set to a maximal
consistent set \(\Delta\).  The truth lemma gives
\(v_\Delta\entails\Gamma\) and
\(v_\Delta\entails\neg\varphi\).  Hence
\(v_\Delta\notentails\varphi\), so
\(\Gamma\notentails\varphi\).
\end{proof}

Combining \cref{thm:prop-nd-soundness,thm:prop-completeness} yields
\[
  \Gamma\derives_{\ND_{\CPL}}\varphi
  \quad\text{iff}\quad
  \Gamma\entails\varphi.
\]

\begin{theorem}[Propositional compactness]
\label{thm:prop-compactness}
For every \(\Gamma\subseteq\Fm_{\CPL}\) and every
\(\varphi\in\Fm_{\CPL}\):
\begin{enumerate}
  \item if \(\Gamma\entails\varphi\), then some finite
        \(\Gamma_0\subseteq\Gamma\) satisfies
        \(\Gamma_0\entails\varphi\); and
  \item \(\Gamma\) is satisfiable iff every finite subset of \(\Gamma\) is
        satisfiable.
\end{enumerate}
\end{theorem}
\begin{proof}
If \(\Gamma\entails\varphi\), completeness gives a finite derivation from
\(\Gamma\).  Only finitely many members of \(\Gamma\) occur as open
assumptions; call their set \(\Gamma_0\).  Soundness gives
\(\Gamma_0\entails\varphi\), proving 1.

The forward implication in 2 is immediate.  For the converse, if \(\Gamma\)
were not satisfiable, then \(\Gamma\entails\bottom\).  Clause 1 would give a
finite \(\Gamma_0\subseteq\Gamma\) with
\(\Gamma_0\entails\bottom\), so \(\Gamma_0\) would not be satisfiable.
\end{proof}

\begin{theorem}[Completeness of the sequent calculus]
\label{thm:prop-lk-completeness}
For all finite multisets \(\Gamma,\Delta\) of propositional formulas, the
sequent \(\Gamma\Rightarrow\Delta\) is derivable in \(\LK_{\CPL}\) iff
every valuation making every formula of \(\operatorname{supp}(\Gamma)\)
true makes some formula of \(\operatorname{supp}(\Delta)\) true.
\textit{Status: proof sketch.}
\end{theorem}
\begin{proof}[Proof sketch]
The exact scope is the displayed cut-admitting calculus with finite multiset
contexts.  Left to right is \cref{thm:prop-lk-soundness}, whose semantic
condition on a sequent is exactly the one displayed here.

For right to left, first record the translation subroutine.  Induction on
a \(\ND_{\CPL}\)-derivation translates every introduction rule to the
corresponding right rule, every elimination rule to the corresponding
left rule followed when necessary by cut, and assumption discharge to a
right rule.  Weakening and contraction align the finite contexts.  The
induction establishes that if
\(\Gamma_0\derives_{\ND_{\CPL}}\varphi\) for a finite set \(\Gamma_0\),
then the sequent \(\Gamma_0\Rightarrow\varphi\) is derivable in the
displayed calculus.  The individual rule-by-rule translation trees are
omitted, and the argument does not establish cut elimination.

Now suppose every valuation making every formula of
\(\operatorname{supp}(\Gamma)\) true makes some formula of
\(\operatorname{supp}(\Delta)\) true, and list
\(\operatorname{supp}(\Delta)=\set{\delta_1,\ldots,\delta_m}\) with
\(m\geq0\).
\begin{enumerate}
  \item Let \(S\) be the set of members of
        \(\operatorname{supp}(\Gamma)\) together with
        \(\neg\delta_1,\ldots,\neg\delta_m\).  A valuation satisfying
        all of \(S\) would make every member of
        \(\operatorname{supp}(\Gamma)\) true, hence some
        \(\delta_j\) true together with \(\neg\delta_j\).  So no
        valuation satisfies \(S\), and \(S\entails\bottom\).
  \item \Cref{thm:prop-completeness} gives
        \(S\derives_{\ND_{\CPL}}\bottom\).  The open assumptions of
        that derivation form a finite subset of \(S\), say
        \(\gamma_1,\ldots,\gamma_k\in\operatorname{supp}(\Gamma)\) and
        \(\neg\delta_{j_1},\ldots,\neg\delta_{j_r}\).  The translation
        subroutine turns it into an \(\LK_{\CPL}\)-derivation of
        \(\gamma_1,\ldots,\gamma_k,
        \neg\delta_{j_1},\ldots,\neg\delta_{j_r}\Rightarrow\bottom\).
  \item Cut with the initial \(\bottom\)-sequent
        \(\bottom\Rightarrow\), the instance of
        \(\Gamma,\bottom\Rightarrow\Delta\) with both contexts empty,
        empties the succedent.
  \item For each \(j\), the sequent
        \(\Rightarrow\delta_j,\neg\delta_j\) is derivable: start from
        the identity sequent \(\delta_j\Rightarrow\delta_j\) and apply
        \(\neg R\).  Cutting it against the current sequent, with cut
        formula \(\neg\delta_j\),
        \[
         \Rule[Cut]{\Rightarrow\delta_j,\neg\delta_j\qquad
          \Gamma',\neg\delta_j\Rightarrow\Pi}
          {\Gamma'\Rightarrow\delta_j,\Pi},
        \]
        moves \(\delta_j\) into the succedent; after \(r\) such cuts
        the derived sequent is
        \(\gamma_1,\ldots,\gamma_k\Rightarrow
        \delta_{j_1},\ldots,\delta_{j_r}\).
  \item Weakening (\(W_L\), \(W_R\)) adds the unused members of
        \(\Gamma\) and \(\Delta\), and contraction (\(C_L\),
        \(C_R\)) merges duplicates, giving exactly the multiset
        sequent \(\Gamma\Rightarrow\Delta\).
\end{enumerate}
When \(m=0\), the hypothesis asserts that no valuation makes every
member of \(\operatorname{supp}(\Gamma)\) true, step 4 performs no cuts,
and the same route derives \(\Gamma\Rightarrow\Delta\) with empty
succedent; an empty \(\Gamma\) is likewise covered.  The route uses cut,
which is primitive in the displayed calculus, and makes no cut-free
claim; no choice principle enters, because the propositional Lindenbaum
lemma runs on a countable enumeration.  The status is a sketch because
the individual translation trees of the subroutine remain omitted.
\end{proof}

\section{Propositional resolution}
\label{sec:prop-metatheory-04}

The clause object in this section is exactly the finite set defined in
\cref{sec:prop-syntax-semantics-04}.  A valuation satisfies a clause set \(S\) iff it
satisfies every \(C\in S\), so the set is interpreted conjunctively.  The
complement of \(p\) is \(\neg p\), and the complement of \(\neg p\) is
\(p\).  A clause containing both a literal and its complement is
\emph{tautological}; its formula interpretation is valid, so deleting it from
a clause set preserves the satisfying valuations of that set.

\begin{definition}[Resolvent]
Let \(E\) and \(F\) be clauses, let \(L\in E\), and let its complement
\(\overline L\) belong to \(F\).  Their \emph{resolvent on \(L\)} is the
clause
\[
  (E\setminus\set{L})\cup(F\setminus\set{\overline L}).
\]
A \emph{resolution derivation} from a clause set \(S\) is a finite sequence
whose members are clauses in \(S\) or resolvents of earlier members.
\end{definition}

\begin{theorem}[Soundness of propositional resolution]
Every valuation satisfying both parent clauses satisfies their resolvent.
Consequently, deriving the empty clause proves that the initial clause set is
not satisfiable.
\end{theorem}
\begin{proof}
Let \(E,F\) be the parents and let
\(R=(E\setminus\set{L})\cup(F\setminus\set{\overline L})\).  Suppose a
valuation falsifies \(R\).  It therefore falsifies every nonpivot literal in
both parents.  To satisfy \(E\), it must satisfy \(L\); to satisfy \(F\), it
must satisfy \(\overline L\).  No valuation satisfies complementary
literals.  Therefore no valuation can satisfy both parents while falsifying
the resolvent.  Induction on a resolution derivation preserves satisfaction.
The empty clause is false under every valuation.
\end{proof}

\begin{theorem}[Refutational completeness of propositional resolution]
\label{thm:prop-resolution-complete}
Let \(S\) be a finite set of clauses.  Then \(S\) is not satisfiable
iff the empty clause has a resolution derivation from \(S\).
\end{theorem}
\begin{proof}
Soundness proves the reverse implication.  For the forward implication,
delete tautological clauses.  Induct on the number of variables occurring in
\(S\).  If there are no variables, an unsatisfiable clause set contains the
empty clause.

Choose a variable \(p\).  Form \(S_p\) from all clauses not containing
\(p\) or \(\neg p\), together with every non-tautological resolvent on
\(p\) of a clause containing \(p\) and a clause containing \(\neg p\).  We claim that
\(S\) is satisfiable iff \(S_p\) is satisfiable.  Restricting a valuation
satisfying \(S\) to the other variables satisfies all retained clauses and all
resolvents by soundness.  Conversely, let \(w\) satisfy \(S_p\).  A clause
\(\set{p}\cup A\), where \(p\notin A\), whose remainder \(A\) is false
under \(w\) requires \(p=\True\); a clause
\(\set{\neg p}\cup B\), where \(\neg p\notin B\), whose remainder \(B\)
is false requires \(p=\False\).  Both requirements cannot occur:
otherwise the resolvent \(A\cup B\) on \(p\) would be false under
\(w\); being false under \(w\), it contains no complementary pair of
literals, so it is not tautological, belongs to \(S_p\), and contradicts
the assumption that \(w\) satisfies \(S_p\).  Choose either
permitted value for \(p\).  The resulting extension satisfies \(S\), proving
the claim.

If \(S\) is not satisfiable, neither is \(S_p\).  It contains fewer variables,
so the induction hypothesis gives a resolution derivation of the empty clause
from \(S_p\).  Every noninitial clause of \(S_p\) is itself a resolvent of
clauses in \(S\).  Splicing those one-step derivations into the derivation from
\(S_p\) gives a resolution derivation from \(S\).
\end{proof}

\begin{remark}[Attribution and scope of the finiteness hypothesis]
\label{rem:resolution-finiteness}
The variable-elimination step in the proof of
\cref{thm:prop-resolution-complete}---passing from \(S\) to \(S_p\) and
discarding the variable \(p\)---is the elimination rule of Davis and
Putnam \cite{DavisPutnam1960}.  The finiteness hypothesis on \(S\) is
consumed by the induction, not by the result: an infinite clause set may
contain infinitely many variables, and the induction on the number of
variables then has no measure to decrease.  The restriction is removed by
propositional compactness.  Suppose an arbitrary clause set \(S\) is not
satisfiable.  A valuation satisfies a clause \(C\) iff it satisfies the
formula \(\operatorname{form}_{\CPL}(C)\) of
\cref{sec:prop-syntax-semantics-04}, so the formula set
\(\set{\operatorname{form}_{\CPL}(C):C\in S}\) is not satisfiable, and
\cref{thm:prop-compactness} gives a finite subset of it that is not
satisfiable.  Selecting for each of its finitely many members one clause
of \(S\) interpreting it yields a finite \(S_0\subseteq S\) that is not
satisfiable.  \Cref{thm:prop-resolution-complete} applied to \(S_0\)
gives a resolution derivation of the empty clause from \(S_0\), and every
resolution derivation from \(S_0\) is one from \(S\).  Since the
soundness theorem above is stated without finiteness of the clause set,
the equivalence of \cref{thm:prop-resolution-complete} therefore holds
for every clause set, finite or not.
\end{remark}

Resolution is a refutation calculus: it derives the empty clause from an
unsatisfiable clause set.  Conversion to conjunctive normal form and
resolution must therefore be kept distinct.  The former supplies clauses; the
latter derives clauses.  The original resolution method is due to Robinson
\cite{Robinson1965}; normal-form transformations and their preservation
properties are treated systematically in \cite{BaazEglyLeitsch2001}.

\section{Common errors}
\label{sec:prop-metatheory-05}

\begin{warning}[Propositional ambiguity]
The following replacements are invalid unless separately justified:
\begin{enumerate}
  \item replacing \(\varphi\to\psi\) by \(\psi\to\varphi\);
  \item treating \(\varphi\lor\psi\) as exclusive;
  \item inferring either disjunct from a disjunction;
  \item inferring \(\varphi\) from \(\varphi\to\psi\) and \(\psi\);
  \item using a truth table for some variables while ignoring another variable
        occurring in the formula;
  \item citing soundness in the completeness direction; and
  \item calling a clause set satisfiable because resolution search has not yet
        found the empty clause.
\end{enumerate}
\end{warning}

\section{Exercises}
\label{sec:prop-metatheory-06}

\begin{exercise}
\label{ex:prop-metatheory-01}
Parse
\(\neg P\land Q\to R\lor S\) using the declared precedence and associativity
conventions.  Restore every parenthesis required by the grammar and list all
subformulas.
\end{exercise}

\begin{exercise}
\label{ex:prop-metatheory-02}
Prove by structural induction that, in the syntax tree of every formula, the
number of leaves---occurrences of propositional variables or \(\bottom\)---is
one greater than the number of binary-connective occurrences.
\end{exercise}

\begin{exercise}
\label{ex:prop-metatheory-03}
Construct a countervaluation to
\(P\to Q,Q\entails P\).  Explain why this refutes semantic consequence but
does not by itself inspect any derivation.
\end{exercise}

\begin{exercise}
\label{ex:prop-metatheory-04}
Translate ``The process terminates only if either the input is finite or the
search is bounded, but not both.''  Supply a key and state whether ``or'' is
inclusive before adding the final restriction.
\end{exercise}

\begin{exercise}
\label{ex:prop-metatheory-05}
Find conjunctive and disjunctive normal forms for
\((P\to Q)\land(Q\to R)\).  Verify both by a truth table.
\end{exercise}

\begin{exercise}
\label{ex:prop-metatheory-06}
Give a \(\ND_{\CPL}\)-derivation of
\((P\to R)\to((Q\to R)\to(P\lor Q\to R))\).
Mark every discharged assumption.
\end{exercise}

\begin{exercise}
\label{ex:prop-metatheory-07}
Derive \(\neg(P\land\neg P)\) without using \(\mathsf{RAA}\).  Then
identify the discharging rule used in place of \(\mathsf{RAA}\), and explain
why every rule in the derivation is available in both classical and
intuitionistic natural deduction.
\end{exercise}

\begin{exercise}
\label{ex:prop-metatheory-08}
Give a cut-free \(\LK_{\CPL}\)-derivation of
\(P\to Q, Q\to R\Rightarrow P\to R\).
\end{exercise}

\begin{exercise}
\label{ex:prop-metatheory-09}
Prove that logical equivalence is an equivalence relation and a congruence for
every primitive connective.
\end{exercise}

\begin{exercise}
\label{ex:prop-metatheory-10}
Use the maximal-consistent-set construction to find a countervaluation to a
chosen underivable formula.  Show explicitly how membership in the
set determines the valuation.
\end{exercise}

\begin{exercise}
\label{ex:prop-metatheory-11}
Convert
\[
 (P\lor Q)\land(\neg P\lor R)\land(\neg Q\lor R)\land\neg R
\]
to a clause set and give a resolution refutation.
\end{exercise}

\begin{exercise}
\label{ex:prop-metatheory-12}
Show that deleting a tautological clause preserves the set of satisfying
valuations.  Show by example that deleting a clause that some valuation
satisfies, but that is not tautological, can change the set of valuations
satisfying the clause set.
\end{exercise}

\part{Classical First-Order Logic}
\label{part:classical-first-order}

Classical first-order logic adds variables ranging over a nonempty domain,
symbols for functions and relations, and quantifiers.  That additional
expressive power makes the syntax of binding and substitution part of the
formal-logical development.  The definitions in this part follow the set-based semantics of
\cite{EbbinghausFlumThomas1994,Hodges1993}; equality is logical identity, not
an unspecified binary relation.

\chapter{First-Order Syntax}
\label{ch:fol-syntax-semantics}

Parallel textbook developments of the material in this part are
\cite{ChiswellHodges2007} and \cite{EbbinghausFlumThomas1994}.

\section{Signatures and terms}
\label{sec:fol-syntax-semantics-01}

Fix the countably infinite set
\(\Var\defeq\set{x_n:n\in\N}\) of object-language variables.  The letters
\(x,y,z\) denote members of \(\Var\), while \(r,t,u\) are metavariables
ranging over terms.  The letter \(s\) is reserved for variable assignments.

\begin{definition}[First-order signature]
A \emph{first-order signature} is a tuple
\[
 \Lang=\tuple{\mathcal F,\mathcal R,\arity_{\mathcal F},
                    \arity_{\mathcal R}},
\]
where \(\mathcal F\) and \(\mathcal R\) are sets of function and relation
symbols and
\[
 \arity_{\mathcal F}:\mathcal F\to\N,
 \qquad
 \arity_{\mathcal R}:\mathcal R\to\N
\]
assign each symbol its arity.  A function symbol of arity zero is a
\emph{constant symbol}.  The equality sign \(=\) is logical and is not a
member of \(\mathcal R\).  The sets \(\mathcal F\), \(\mathcal R\),
\(\Var\), and the set of reserved logical and punctuation tokens
\[
 \{=,\bottom,\neg,\land,\lor,\to,\forall,\exists,
   \texttt{(},\texttt{)},\texttt{,}\}
\]
are pairwise disjoint.  Consequently no signature symbol can also be a
variable, logical symbol, parenthesis, or comma.
\end{definition}

The cardinality of the signature is
\(\abs{\Lang}\defeq\abs{\mathcal F\cup\mathcal R}\).  Arity data add no
new symbols.

The signature is a syntactic object.  It assigns arities but assigns no
functions, relations, or elements.  A structure supplies those semantic
objects later.

\begin{definition}[Term]
For a signature \(\Lang\), the set \(\Term(\Lang)\) of \emph{terms} is the
least set satisfying:
\begin{enumerate}
  \item every \(x\in\Var\) belongs to \(\Term(\Lang)\); and
  \item if \(f\in\mathcal F\) has arity \(n\) and
        \(t_1,\ldots,t_n\in\Term(\Lang)\), then
        \(f(t_1,\ldots,t_n)\in\Term(\Lang)\).
\end{enumerate}
For \(n=0\), the expression is written \(f\), not \(f()\).  No other string
is a term.
\end{definition}

\begin{example}
Let \(\Lang_{\mathrm{ar}}\) have function symbols
\(\set{0,S,+,\cdot}\) of arities \(0,1,2,2\), respectively, and relation
symbol \(<\) of arity \(2\).  The strings
\[
 x_0,\qquad S(0),\qquad +(S(0),x_2),\qquad
 \cdot(+(x_0,x_1),S(x_2))
\]
are terms.  The string \(<(x_0,x_1)\) is not a term because \(<\) is a
relation symbol.  Conventional infix displays such as \(x_0+S(0)\) are
typographical abbreviations for the prefix grammar.
\end{example}

\section{Formulas and unique readability}
\label{sec:fol-syntax-semantics-02}

\begin{definition}[First-order formula]
For a signature \(\Lang\), the set \(\Fm(\Lang)\) of \emph{formulas} is the
least set satisfying:
\begin{enumerate}
  \item if \(r,t\in\Term(\Lang)\), then \(r=t\in\Fm(\Lang)\);
  \item if \(R\in\mathcal R\) has arity \(n\) and
        \(t_1,\ldots,t_n\in\Term(\Lang)\), then
        \(R(t_1,\ldots,t_n)\in\Fm(\Lang)\);
  \item \(\bottom\in\Fm(\Lang)\);
  \item if \(\varphi\in\Fm(\Lang)\), then
        \(\neg\varphi\in\Fm(\Lang)\);
  \item if \(\varphi,\psi\in\Fm(\Lang)\) and
        \(\circ\in\set{\land,\lor,\to}\), then
        \((\varphi\circ\psi)\in\Fm(\Lang)\); and
  \item if \(x\in\Var\) and \(\varphi\in\Fm(\Lang)\), then
        \(\forall x\,\varphi\) and \(\exists x\,\varphi\) belong to
        \(\Fm(\Lang)\).
\end{enumerate}
No other string is a formula.  A formula introduced by clause 1 or 2 is
\emph{atomic}.  When a relation symbol has arity zero, its atomic formula is
written \(R\), not \(R()\).
\end{definition}

The scope of a quantifier in the official grammar is the formula immediately
following its variable.  In abbreviated displays, this document gives a
quantifier tighter binding than every binary connective; parentheses extend
its scope.  Thus
\(\forall x\,P(x)\to Q(x)\) abbreviates
\((\forall x\,P(x))\to Q(x)\) under the declared connective precedence,
whereas \(\forall x\,(P(x)\to Q(x))\) has the implication inside the
quantifier's scope.

\begin{theorem}[Unique readability for first-order expressions]
\label{thm:fol-unique-readability}
Every \(\Lang\)-term has exactly one outermost formation clause, and its
immediate subterms are uniquely determined.  Every \(\Lang\)-formula has
exactly one outermost formation clause, and its principal symbol, bound
variable when present, and immediate subformulas are uniquely determined.
\end{theorem}
\begin{proof}
Function symbols and relation symbols are disjoint and carry fixed arities;
variables, logical symbols, punctuation, and signature symbols are pairwise
distinguishable tokens.  A variable term is one token.  A compound term begins
with its unique function symbol, and its fixed arity determines the number of
comma-separated argument positions.  Balanced-parenthesis scanning uniquely
separates those arguments.  Induction on term length therefore gives unique
term construction.

An atomic formula either has the uniquely located equality sign between two
terms or begins with a relation symbol whose arity uniquely determines its
term arguments.  The remaining formula constructors have distinguishable
initial forms: \(\bottom\), \(\neg\), an initial left parenthesis for a binary
connective, and \(\forall\) or \(\exists\) followed by one variable.  The
depth-zero scan used in \cref{thm:prop-unique-readability} uniquely separates a
binary formula.  Removing a unary logical symbol leaves its unique immediate
formula.  Induction on formula length completes the proof.
\end{proof}

\section{Free variables and sentences}
\label{sec:fol-syntax-semantics-03}

\begin{definition}[Variables occurring in a term]
The sans-serif symbol \(\Var\) without an argument denotes the fixed set
of all object-language variables; the roman operator
\(\operatorname{Var}\) with an argument denotes an occurrence set.  The
two roles are distinguished by the presence of the argument, and both are
registered here.  The set \(\operatorname{Var}(t)\) of variables occurring
in a term is defined by
\[
 \operatorname{Var}(x)=\set{x},
 \qquad
 \operatorname{Var}(f(t_1,\ldots,t_n))
 =\bigcup_{i=1}^{n}\operatorname{Var}(t_i).
\]
\end{definition}

\begin{definition}[Free variables]
The set \(\FV(\varphi)\) of \emph{free variables} of a formula is defined by
\[
\begin{aligned}
 \FV(r=t)&=\operatorname{Var}(r)\cup\operatorname{Var}(t),\\
 \FV(R(t_1,\ldots,t_n))&=\bigcup_{i=1}^{n}\operatorname{Var}(t_i),\\
 \FV(\bottom)&=\varnothing,\\
 \FV(\neg\varphi)&=\FV(\varphi),\\
 \FV(\varphi\circ\psi)&=\FV(\varphi)\cup\FV(\psi),\\
 \FV(\forall x\,\varphi)&=\FV(\varphi)\setminus\set{x},\\
 \FV(\exists x\,\varphi)&=\FV(\varphi)\setminus\set{x}.
\end{aligned}
\]
A formula with no free variable is a \emph{sentence}; the set of
\(\Lang\)-sentences is \(\Sent(\Lang)\).
\end{definition}

\begin{definition}[First-order formula complexity]
For \(\varphi\in\Fm(\Lang)\), its \emph{complexity}
\(c(\varphi)\in\N\) is defined by
\[
\begin{aligned}
 c(r=t)&=0,&c(R(t_1,\ldots,t_n))&=0,&c(\bottom)&=0,\\
 c(\neg\varphi)&=c(\varphi)+1,
 &c(\varphi\circ\psi)&=c(\varphi)+c(\psi)+1,\\
 c(Qx\,\varphi)&=c(\varphi)+1,
\end{aligned}
\]
where \(\circ\in\set{\land,\lor,\to}\) and
\(Q\in\set{\forall,\exists}\).  Every immediate subformula has smaller
complexity.  This is the measure used by every later induction or recursion
described as being on first-order formula complexity.
\end{definition}

For a set of formulas \(\Gamma\), extend the notation by
\[
  \FV(\Gamma)\defeq\bigcup_{\gamma\in\Gamma}\FV(\gamma).
\]
Thus \(x\notin\FV(\Gamma)\) means exactly that \(x\) is free in no member of
\(\Gamma\).

Let \(\operatorname{Var}(\varphi)\) denote the set of all variables having a
free or bound occurrence in \(\varphi\).  It is defined by the same recursive
clauses as \(\FV\), except that a quantifier clause adds its bound variable
instead of removing it.

An occurrence of \(x\) is \emph{bound} when it lies in the scope of a
quantifier \(\forall x\) or \(\exists x\) that governs that occurrence.  The
quantifier binds occurrences, not the variable symbol throughout the formula.
In
\[
  P(x)\land\exists x\,Q(x,y),
\]
the first occurrence of \(x\) is free, the second is bound, and the occurrence
of \(y\) is free.

\section{Capture-avoiding substitution}
\label{sec:fol-syntax-semantics-04}

\begin{definition}[Substitution in a term]
For terms \(r,t\) and a variable \(x\), the term \(r[t/x]\) is defined by
\[
 x[t/x]=t,\qquad y[t/x]=y\ (y\ne x),
\]
and
\[
 f(r_1,\ldots,r_n)[t/x]
 =f(r_1[t/x],\ldots,r_n[t/x]).
\]
Only occurrences of \(x\) are replaced.
\end{definition}

Fix the order \(x_0,x_1,\ldots\) on \(\Var\).  This order makes every choice
of a fresh variable below deterministic.

\begin{definition}[Capture-avoiding substitution in a formula]
For a formula \(\varphi\), term \(t\), and variable \(x\), the formula
\(\varphi[t/x]\) is defined by recursion on formula complexity.  In atomic
formulas, substitute in every term.  Substitution commutes with
\(\bottom,\neg,\land,\lor,\to\).  For
\(Q\in\set{\forall,\exists}\), define
\[
(Qy\,\psi)[t/x]=
\begin{cases}
 Qy\,\psi,&y=x,\\
 Qy\,(\psi[t/x]),&y\ne x\text{ and }y\notin\operatorname{Var}(t),\\
 Qz\,((\psi[z/y])[t/x]),&y\ne x\text{ and }y\in\operatorname{Var}(t),
\end{cases}
\]
where in the last clause \(z\) is the least variable not in
\[
 \operatorname{Var}(\psi)\cup\operatorname{Var}(t)\cup\set{x,y}.
\]
The recursion is on complexity, and it is well founded by a simultaneous
strong induction on complexity: the calls on \(\psi\) concern a formula of
lower complexity than \(Qy\,\psi\), and variable-for-variable
substitution preserves complexity---each clause preserves the outer
constructor---so in the last clause \(c(\psi[z/y])=c(\psi)<c(Qy\,\psi)\)
and the outer call is also at lower complexity.  Complexity preservation
for substitution by an arbitrary term is recorded as
\cref{lem:substitution-complexity} below.
\end{definition}

The last clause first renames the binder and then substitutes.  Its purpose is
not stylistic: it prevents a free variable of \(t\) from becoming bound.

\begin{lemma}[Substitution preserves formula complexity]
\label{lem:substitution-complexity}
For every formula \(\varphi\), term \(t\), and variable \(x\),
\[
  c(\varphi[t/x])=c(\varphi).
\]
\end{lemma}
\begin{proof}
Induct on \(c(\varphi)\).  Term substitution changes no logical constructor
in an atomic formula.  The connective clauses preserve the displayed outer
constructor and use the induction hypotheses.  At a quantifier, substitution
either stops, commutes with the quantifier, or renames its bound variable
before recurring on the matrix.  Bound-variable renaming changes no logical
constructor, so every quantifier case also preserves complexity.
\end{proof}

\begin{definition}[Alpha-equivalence]
The relation of \emph{alpha-equivalence}, written \(\equiv_\alpha\), is the
least equivalence relation on formulas that is a congruence for every formula
constructor and contains
\[
 Qx\,\varphi\equiv_\alpha Qy\,(\varphi[y/x])
\]
for \(Q\in\set{\forall,\exists}\) whenever
\(y\notin\operatorname{Var}(\varphi)\).  Thus a generating step renames one
bound variable and exactly the occurrences governed by its quantifier.
\end{definition}

Alpha-equivalence is syntactic and is not literal identity of strings.  It
will be shown below to preserve satisfaction.

\begin{example}[Avoiding capture]
Let \(\varphi=\forall y\,R(x,y)\) and substitute the term \(y\) for \(x\).
Direct textual replacement would produce \(\forall y\,R(y,y)\), incorrectly
binding the inserted occurrence.  If \(z\) is the least permitted fresh
variable, the defined result is
\[
  \varphi[y/x]=\forall z\,R(y,z).
\]
The free occurrence introduced by substitution remains free.
\end{example}

\begin{lemma}[Free variables under substitution]
\label{lem:substitution-fv}
For all terms \(r,t\) and every variable \(x\): if
\(x\notin\operatorname{Var}(r)\), then \(r[t/x]=r\); and
\[
 \operatorname{Var}(r[t/x])\subseteq
 (\operatorname{Var}(r)\setminus\set{x})\cup\operatorname{Var}(t).
\]
For every formula \(\varphi\), term \(t\), and variables \(x,y\):
\begin{enumerate}
  \item \(\FV(\varphi[t/x])\subseteq
        (\FV(\varphi)\setminus\set{x})\cup\operatorname{Var}(t)\); and
  \item if \(y\notin\operatorname{Var}(\varphi)\), then
        \(\FV(\varphi[y/x])=(\FV(\varphi)\setminus\set{x})\cup\set{y}\)
        when \(x\in\FV(\varphi)\), and
        \(\FV(\varphi[y/x])=\FV(\varphi)\) when
        \(x\notin\FV(\varphi)\).
\end{enumerate}
A term's variables are measured by \(\operatorname{Var}\), exactly as in
the \(\FV\) clauses for atomic formulas; no separate free-variable
operator on terms is defined in this document.
\textit{Status: complete proof.}
\end{lemma}
\begin{proof}
For the term claims, induct on \(r\).  A variable other than \(x\) is
unchanged, the variable \(x\) itself is replaced by \(t\), and the compound
clause applies the induction hypotheses argumentwise.

For the formula claims, proceed by strong induction on formula complexity
\(c\).  The atomic cases follow from the term claims and the \(\FV\)
clauses for atomic formulas; the \(\bottom\) case is immediate, and the
connective cases combine the induction hypotheses through the \(\FV\)
clauses.

For clause 2, the renaming clause of the substitution definition never
fires: every binder \(w\) of a subformula of \(\varphi\) lies in
\(\operatorname{Var}(\varphi)\), so \(w\ne y\), and the substitution is
textual replacement of the free occurrences of \(x\) by \(y\).  At a
quantifier \(Qw\,\psi\): if \(w=x\), substitution stops, \(x\) is not free,
and both sides equal \(\FV(Qw\,\psi)\); if \(w\ne x\), the substitution
commutes with the quantifier, and the induction hypothesis for \(\psi\)
passes through the clause
\(\FV(Qw\,(\psi[y/x]))=\FV(\psi[y/x])\setminus\set{w}\), using
\(w\ne y\) and the fact that \(x\in\FV(Qw\,\psi)\) iff
\(x\in\FV(\psi)\).

For clause 1 at a quantifier \(Qy\,\psi\): if \(y=x\), substitution stops
and \(x\notin\FV(Qy\,\psi)\); if \(y\ne x\) and
\(y\notin\operatorname{Var}(t)\), apply the induction hypothesis to
\(\psi\) and remove \(y\).  In the renaming case
\((Qy\,\psi)[t/x]=Qz\,((\psi[z/y])[t/x])\), apply the induction hypothesis
to \(\psi\) with \([z/y]\)---clause 2---and to \(\psi[z/y]\) with
\([t/x]\); the latter application is at the equal complexity
\(c(\psi[z/y])=c(\psi)<c(Qy\,\psi)\), licensed by
\cref{lem:substitution-complexity}.  Since
\(z\notin\operatorname{Var}(t)\cup\set{x}\), computing with the \(\FV\)
clauses gives
\[
 \FV\bigl(Qz\,((\psi[z/y])[t/x])\bigr)\subseteq
 (\FV(Qy\,\psi)\setminus\set{x})\cup\operatorname{Var}(t).
 \qedhere
\]
\end{proof}

\begin{lemma}
If \(x\notin\FV(\varphi)\), then
\(\varphi[t/x]\equiv_\alpha\varphi\).
\textit{Status: complete proof.}
\end{lemma}
\begin{proof}
Proceed by strong induction on formula complexity \(c\).  In the atomic
case, \(x\) occurs in no term of \(\varphi\), so every term is unchanged by
the term clause of \cref{lem:substitution-fv} and
\(\varphi[t/x]=\varphi\); the \(\bottom\) case is likewise immediate.  In
the connective cases, \(x\) is free in no immediate subformula; the
induction hypotheses and the congruence property of \(\equiv_\alpha\)
conclude.

At a quantifier \(Qy\,\psi\) with \(y=x\), substitution stops, and
reflexivity concludes.  If \(y\ne x\) and
\(y\notin\operatorname{Var}(t)\), then \(x\notin\FV(\psi)\); the induction
hypothesis gives \(\psi[t/x]\equiv_\alpha\psi\), and the quantifier
congruence concludes.  If \(y\ne x\) and \(y\in\operatorname{Var}(t)\),
the definition returns \(Qz\,((\psi[z/y])[t/x])\), where \(z\) is the
least variable outside
\(\operatorname{Var}(\psi)\cup\operatorname{Var}(t)\cup\set{x,y}\).  By
\cref{lem:substitution-fv}(2), \(x\notin\FV(\psi[z/y])\), because
\(x\notin\FV(\psi)\) and \(x\ne z\).  By
\cref{lem:substitution-complexity},
\(c(\psi[z/y])=c(\psi)<c(Qy\,\psi)\), so the strong induction hypothesis
applies to \(\psi[z/y]\) and gives
\((\psi[z/y])[t/x]\equiv_\alpha\psi[z/y]\).  The quantifier congruence
gives \(Qz\,((\psi[z/y])[t/x])\equiv_\alpha Qz\,(\psi[z/y])\); the
generating clause, licensed by \(z\notin\operatorname{Var}(\psi)\), gives
\(Qz\,(\psi[z/y])\equiv_\alpha Qy\,\psi\); transitivity concludes.
\end{proof}

\chapter{Structures, Satisfaction, and Translation}
\label{ch:fol-semantics}

\section{Structures and variable assignments}
\label{sec:fol-semantics-01}

\begin{definition}[First-order structure]
Let \(\Lang\) be a signature.  A \emph{\(\Lang\)-structure}
\(\Meta\) consists of:
\begin{enumerate}
  \item a nonempty set \(\abs{\Meta}\), called its domain;
  \item for each \(n\)-ary function symbol \(f\), a function
        \(f^{\Meta}:\abs{\Meta}^{n}\to\abs{\Meta}\); and
  \item for each \(n\)-ary relation symbol \(R\), a relation
        \(R^{\Meta}\subseteq\abs{\Meta}^{n}\).
\end{enumerate}
For a constant symbol \(c\), the zero-ary function is identified with its
value \(c^{\Meta}\in\abs{\Meta}\).  The equality sign is always interpreted
as actual identity on \(\abs{\Meta}\).
\end{definition}

\begin{definition}[Signature extension, expansion, and reduct]
Let \(\Lang=\tuple{\mathcal F,\mathcal R,\arity_{\mathcal F},
\arity_{\mathcal R}}\) and
\(\Lang'=\tuple{\mathcal F',\mathcal R',\arity_{\mathcal F'},
\arity_{\mathcal R'}}\).  The signature \(\Lang'\) \emph{extends}
\(\Lang\) iff \(\mathcal F\subseteq\mathcal F'\),
\(\mathcal R\subseteq\mathcal R'\), and both arity maps of \(\Lang'\)
restrict to the corresponding arity maps of \(\Lang\).

If \(\Meta\) is a \(\Lang\)-structure, a \(\Lang'\)-structure
\(\Meta'\) is an \emph{expansion} of \(\Meta\) iff it has the same domain
and gives every symbol of \(\Lang\) the same interpretation.  If
\(\Meta'\) is a \(\Lang'\)-structure, its \emph{reduct} to \(\Lang\),
written \(\Meta'\restr\Lang\), is the \(\Lang\)-structure with the same
domain and the interpretations of the symbols in \(\Lang\).  Expansion adds
interpretations and reduct forgets them; neither operation changes the domain,
so neither term means substructure.  These operations are required in the
Henkin, L\"owenheim--Skolem, and Skolemization arguments below.
\end{definition}

\begin{definition}[Variable assignment and update]
A \emph{variable assignment} in \(\Meta\) is a function
\(s:\Var\to\abs{\Meta}\).  For \(a\in\abs{\Meta}\), define the updated
assignment \(s[x\mapsto a]\) by
\[
 s[x\mapsto a](y)=
 \begin{cases}
   a,&y=x,\\
   s(y),&y\ne x.
 \end{cases}
\]
\end{definition}

\begin{definition}[Term value]
The \emph{value} of a term \(t\) in \(\Meta\) under \(s\), written
\(\sem{t}^{\Meta}_s\), is defined by
\[
 \sem{x}^{\Meta}_s=s(x),
 \qquad
 \sem{f(t_1,\ldots,t_n)}^{\Meta}_s
 =f^{\Meta}(\sem{t_1}^{\Meta}_s,\ldots,\sem{t_n}^{\Meta}_s).
\]
\end{definition}

\begin{lemma}[Term coincidence]
If assignments \(s\) and \(s'\) agree on \(\operatorname{Var}(t)\), then
\(\sem{t}^{\Meta}_s=\sem{t}^{\Meta}_{s'}\).
\end{lemma}
\begin{proof}
Induct on the construction of \(t\).  The variable case is the hypothesis.
In a function case, apply the induction hypothesis to each argument and then
apply the function \(f^{\Meta}\).
\end{proof}

\begin{lemma}[Term-substitution lemma]
\label{lem:term-substitution}
For all terms \(r,t\), assignments \(s\), and variables \(x\),
\[
 \sem{r[t/x]}^{\Meta}_s
 =\sem{r}^{\Meta}_{s[x\mapsto\sem{t}^{\Meta}_s]}.
\]
\end{lemma}
\begin{proof}
Induct on \(r\).  If \(r=x\), both sides equal
\(\sem{t}^{\Meta}_s\).  If \(r=y\ne x\), both sides equal \(s(y)\).  The
function case follows componentwise from the induction hypotheses and the
definition of term value.
\end{proof}

\section{Satisfaction}
\label{sec:fol-semantics-02}

\begin{definition}[Satisfaction]
For a structure \(\Meta\), assignment \(s\), and formula \(\varphi\), the
relation \(\Meta,s\entails\varphi\) is defined by structural recursion:
\[
\begin{aligned}
 \Meta,s\entails r=t
 &\quad\text{iff}\quad
 \sem{r}^{\Meta}_s=\sem{t}^{\Meta}_s,\\
 \Meta,s\entails R(t_1,\ldots,t_n)
 &\quad\text{iff}\quad
 (\sem{t_1}^{\Meta}_s,\ldots,\sem{t_n}^{\Meta}_s)\in R^{\Meta},\\
 \Meta,s\notentails\bottom,
 &\\
 \Meta,s\entails\neg\varphi
 &\quad\text{iff}\quad\Meta,s\notentails\varphi,\\
 \Meta,s\entails\varphi\land\psi
 &\quad\text{iff}\quad
 \Meta,s\entails\varphi\text{ and }\Meta,s\entails\psi,\\
 \Meta,s\entails\varphi\lor\psi
 &\quad\text{iff}\quad
 \Meta,s\entails\varphi\text{ or }\Meta,s\entails\psi,\\
 \Meta,s\entails\varphi\to\psi
 &\quad\text{iff}\quad
 \Meta,s\notentails\varphi\text{ or }\Meta,s\entails\psi,\\
 \Meta,s\entails\forall x\,\varphi
 &\quad\text{iff}\quad
 \Meta,s[x\mapsto a]\entails\varphi
 \text{ for every }a\in\abs{\Meta},\\
 \Meta,s\entails\exists x\,\varphi
 &\quad\text{iff}\quad
 \Meta,s[x\mapsto a]\entails\varphi
 \text{ for some }a\in\abs{\Meta}.
\end{aligned}
\]
\end{definition}

\begin{lemma}[Formula coincidence]
\label{lem:fol-coincidence}
If assignments \(s\) and \(s'\) agree on \(\FV(\varphi)\), then
\[
 \Meta,s\entails\varphi
 \quad\text{iff}\quad
 \Meta,s'\entails\varphi.
\]
\end{lemma}
\begin{proof}
Proceed by structural induction.  Atomic cases use term coincidence.
Connective cases follow from the induction hypotheses.  For
\(\forall x\,\psi\), the assignments
\(s[x\mapsto a]\) and \(s'[x\mapsto a]\) agree on \(\FV(\psi)\) for each
\(a\), because \(s\) and \(s'\) agree on
\(\FV(\psi)\setminus\set{x}\).  Apply the induction hypothesis for every
\(a\).  The existential case is identical with ``some'' in place of
``every.''
\end{proof}

\begin{theorem}[Substitution lemma]
\label{thm:fol-substitution}
For every formula \(\varphi\), term \(t\), variable \(x\), structure
\(\Meta\), and assignment \(s\),
\[
 \Meta,s\entails\varphi[t/x]
 \quad\text{iff}\quad
 \Meta,s[x\mapsto\sem{t}^{\Meta}_s]\entails\varphi.
\]
\end{theorem}
\begin{proof}
Proceed by induction on formula complexity.  Atomic cases follow from
\cref{lem:term-substitution}.  Connective cases follow immediately from the
induction hypotheses.

Let \(\varphi=Qy\,\psi\), where
\(Q\in\set{\forall,\exists}\).  If \(y=x\), substitution stops at the
quantifier.  Updating \(x\) before evaluating \(Qx\,\psi\) is immaterial,
because the quantified update of \(x\) overrides it.  If \(y\ne x\) and
\(y\notin\operatorname{Var}(t)\), then
\(\sem{t}^{\Meta}_{s[y\mapsto a]}=\sem{t}^{\Meta}_s\) for every \(a\).
Apply the induction hypothesis to \(\psi\) under each updated assignment and
then apply the truth clause for \(Q\).

If \(y\ne x\) and \(y\in\operatorname{Var}(t)\), the definition first replaces \(y\) by a
fresh \(z\).  The variable \(z\) occurs neither in \(t\) nor free in
\(\psi\).  The induction hypothesis applied to the renamed body---
legitimate because renaming preserves complexity, by
\cref{lem:substitution-complexity}---together with formula coincidence,
shows directly that
\(Qy\,\psi\) and \(Qz\,(\psi[z/y])\) have the same satisfaction condition.
The preceding quantifier argument then applies to the renamed body, whose
complexity equals \(c(\psi)\) by \cref{lem:substitution-complexity}.  Thus
the equivalence holds in all cases.
\end{proof}

\begin{corollary}[Alpha-invariance]
\label{cor:alpha-invariance}
If \(\varphi\equiv_\alpha\psi\), then
\(\Meta,s\entails\varphi\) iff \(\Meta,s\entails\psi\) for every structure
\(\Meta\) and assignment \(s\).
\end{corollary}
\begin{proof}
Define a relation \(R\) on formulas by \(\varphi R\psi\) iff, for every
structure \(\Meta\) and assignment \(s\),
\(\Meta,s\entails\varphi\) iff \(\Meta,s\entails\psi\).  The relation
\(R\) is reflexive, symmetric, and transitive.  It is a congruence for
\(\bottom,\neg,\land,\lor,\to\) by the corresponding satisfaction clauses.
It is also a congruence for each quantifier: if \(\varphi R\psi\), then the
equivalence holds under every assignment \(s[x\mapsto a]\), and the
universal or existential satisfaction clause transfers it to
\(Qx\,\varphi\) and \(Qx\,\psi\).

It remains to check a generating renaming.  Suppose
\(y\notin\operatorname{Var}(\varphi)\).  By
\cref{thm:fol-substitution}, satisfaction of \(\varphi[y/x]\) under
\(s[y\mapsto a]\) is equivalent to satisfaction of \(\varphi\) under
\(s[y\mapsto a][x\mapsto a]\).  Since \(y\notin\FV(\varphi)\), formula
coincidence identifies the latter satisfaction condition with that of
\(\varphi\) under \(s[x\mapsto a]\).  Applying the satisfaction clause for
\(Q\in\set{\forall,\exists}\) gives
\(Qx\,\varphi\;R\;Qy\,(\varphi[y/x])\).  Thus \(R\) is an equivalence
relation and a congruence containing every generating renaming.  By the
leastness clause in the definition of \(\equiv_\alpha\), alpha-equivalence is
contained in \(R\).
\end{proof}

\begin{corollary}[Assignment independence for sentences]
If \(\sigma\in\Sent(\Lang)\), then
\(\Meta,s\entails\sigma\) iff \(\Meta,s'\entails\sigma\) for all assignments
\(s,s'\) in \(\Meta\).
\end{corollary}
\begin{proof}
The assignments agree on \(\FV(\sigma)=\varnothing\); apply
\cref{lem:fol-coincidence}.
\end{proof}

\begin{definition}[Truth, validity, and semantic consequence]
For a sentence \(\sigma\), write \(\Meta\entails\sigma\) iff
\(\Meta,s\entails\sigma\) for one, equivalently every, assignment \(s\).
For \(\Gamma\subseteq\Fm(\Lang)\), write
\(\Gamma\entails\varphi\) iff, for every \(\Lang\)-structure \(\Meta\) and
assignment \(s\),
\[
 \bigl(\Meta,s\entails\gamma\text{ for every }\gamma\in\Gamma\bigr)
 \quad\Longrightarrow\quad
 \Meta,s\entails\varphi.
\]
The formula \(\varphi\) is \emph{valid} iff
\(\varnothing\entails\varphi\).
Two \(\Lang\)-formulas \(\varphi\) and \(\psi\) are \emph{logically
equivalent}, written \(\varphi\equiv\psi\), iff, for every
\(\Lang\)-structure \(\Meta\) and every assignment
\(s\),
\[
  \Meta,s\entails\varphi
  \quad\Longleftrightarrow\quad
  \Meta,s\entails\psi.
\]
This is logical equivalence over all \(\Lang\)-structures, not
alpha-equivalence and not literal identity of formula strings.
\end{definition}

\section{Translation with quantifiers}
\label{sec:fol-semantics-03}

Every first-order translation must state a signature and the intended range of
the variables.  Restricted quantification is an abbreviation:
\[
 \forall x\in A\,\varphi\quad\text{means}\quad
 \forall x\,(A(x)\to\varphi),
\]
and
\[
 \exists x\in A\,\varphi\quad\text{means}\quad
 \exists x\,(A(x)\land\varphi).
\]

\begin{example}[Quantifier order]
Let the variables range over people, with unary relation symbols
\(S\) and \(P\) for student and professor, and binary relation symbol \(A\)
for admires.  Then
\[
 \forall x\bigl(S(x)\to
   \exists y(P(y)\land A(x,y))\bigr)
\]
says that every student admires some professor.  The formula
\[
 \exists y\bigl(P(y)\land
   \forall x(S(x)\to A(x,y))\bigr)
\]
says that one professor is admired by every student.  The second formula
implies the first: its existentially quantified professor serves as a witness
for each student.  The first formula does not in general imply the second.
Quantifier order records dependence of a witness on earlier
universally quantified objects.
\end{example}

\begin{example}[Negated universal]
In the same signature, ``No professor admires every student'' is
\[
 \forall y\bigl(P(y)\to
   \neg\forall x(S(x)\to A(y,x))\bigr).
\]
By classical quantifier duality, this is equivalent to
\(\forall y(P(y)\to\exists x(S(x)\land\neg A(y,x)))\).  If there are no
students, the inner universal in the first formula is vacuously true and the
existential in the second formula is false; both formulas therefore require
that there be no professors.
\end{example}

\chapter{Theories and Axiomatization}
\label{ch:fol-theories}

\begin{definition}[Theory, model, satisfiability, and structure theory]
\label{def:fol-theory-satisfiability}
A \emph{theory} in \(\Lang\) is a set
\(T\subseteq\Sent(\Lang)\).  A \(\Lang\)-structure \(\Meta\) is a
\emph{model} of \(T\), written \(\Meta\entails T\), iff
\(\Meta\entails\sigma\) for every \(\sigma\in T\).  The class of models of
\(T\) is \(\Mod(T)\).

For any \(\Gamma\subseteq\Fm(\Lang)\), say that \(\Gamma\) is
\emph{satisfiable} iff there are a \(\Lang\)-structure \(\Meta\) and an
assignment \(s\) such that \(\Meta,s\entails\gamma\) for every
\(\gamma\in\Gamma\).  Because sentences are assignment-independent, a theory
is satisfiable iff it has a model.

For a \(\Lang\)-structure \(\Meta\), define its \emph{structure theory} by
\[
 \Th_{\Lang}(\Meta)
 \defeq\set{\sigma\in\Sent(\Lang):\Meta\entails\sigma}.
\]
When the signature is fixed by the surrounding statement, \(\Th(\Meta)\) is
the metalinguistic abbreviation for \(\Th_{\Lang}(\Meta)\).  A theory is an
arbitrary set of sentences, whereas a structure theory contains exactly all
sentences true in one specified structure.  Both notions are required below:
the first axiomatizes classes, and the second states elementary equivalence and
forms elementary diagrams.
\end{definition}

A theory is not required to be closed under semantic consequence.  Its
semantic closure is
\[
 \Cn(T)\defeq\set{\sigma\in\Sent(\Lang):T\entails\sigma}.
\]

\begin{definition}[Axiomatization]
A theory \(T\) \emph{axiomatizes} a class \(\mathcal K\) of
\(\Lang\)-structures iff \(\Mod(T)=\mathcal K\).  A class is
\emph{finitely axiomatizable} iff some finite theory axiomatizes it.
\end{definition}

\begin{example}[Groups]
Let \(\Lang_{\mathrm{grp}}\) contain a binary function symbol \(\cdot\), a
unary function symbol \(i\), and a constant symbol \(e\).  The following
finite theory axiomatizes groups:
\[
\begin{aligned}
&\forall x\forall y\forall z\,
  ((x\cdot y)\cdot z=x\cdot(y\cdot z)),\\
&\forall x\,(e\cdot x=x\land x\cdot e=x),\\
&\forall x\,(i(x)\cdot x=e\land x\cdot i(x)=e).
\end{aligned}
\]
The symbols are interpreted by the structure; the equations constrain those
interpretations.  Equality itself remains actual identity.
\end{example}

\begin{example}[Dense linear orders without endpoints]
In the signature containing one binary relation symbol \(<\), the following
sentences axiomatize dense linear orders without endpoints:
\[
\begin{aligned}
&\forall x\,\neg(x<x),\\
&\forall x\forall y\forall z\,
  ((x<y\land y<z)\to x<z),\\
&\forall x\forall y\,(x<y\lor x=y\lor y<x),\\
&\forall x\forall y\,(x<y\to\exists z(x<z\land z<y)),\\
&\forall x\,\exists y\exists z\,(y<x\land x<z).
\end{aligned}
\]
Each sentence contributes a distinct condition: irreflexivity, transitivity,
linearity, density, and absence of endpoints.
\end{example}

\begin{example}[Infinite structures]
For \(n\geq2\), let
\[
 \iota_n\defeq
 \exists x_0\cdots\exists x_{n-1}
 \bigwedge_{0\leq i<j<n}x_i\ne x_j,
\]
where the pairs \((i,j)\) are ordered lexicographically and the finite
conjunction uses the global right-association convention.  The notation
\(r\ne t\) abbreviates \(\neg(r=t)\).  A structure satisfies every
member of \(T_\infty=\set{\iota_n:n\geq2}\) iff its domain is
infinite.  This is an axiomatization by a theory, not by one sentence.
\end{example}

\begin{example}[Definitions inside a theory]
Suppose \(T\) is a theory in \(\Lang\), \(R\) is an \(n\)-ary relation
symbol not in \(\Lang\), and \(\varphi(x_1,\ldots,x_n)\) is an
\(\Lang\)-formula whose free variables are among the \(n\) distinct
displayed variables.  Add the \emph{defining axiom}
\[
 \delta_R\defeq
 \forall x_1\cdots\forall x_n
 \bigl(R(x_1,\ldots,x_n)\leftrightarrow
       \varphi(x_1,\ldots,x_n)\bigr).
\]
The passage from \(T\) to \(T\cup\set{\delta_R}\) is a
\emph{definitional extension}, not a new axiom about the old signature:
by \cref{thm:definitional-conservativity} below, every \(\Lang\)-model of
\(T\) has exactly one expansion satisfying \(\delta_R\), and the extended
theory has exactly the old \(\Lang\)-consequences.
\end{example}

\begin{theorem}[Conservativity of definitional extensions]
\label{thm:definitional-conservativity}
Let \(T\), \(R\), \(\varphi\), and \(\delta_R\) be as in the preceding
example, and let \(\Lang_R\) denote the signature extending \(\Lang\) by
the single relation symbol \(R\).
\begin{enumerate}
  \item Every \(\Lang\)-model \(\Meta\) of \(T\) has exactly one
        \(\Lang_R\)-expansion satisfying \(\delta_R\), and that expansion
        is a model of \(T\cup\set{\delta_R}\).
  \item The extension is \emph{conservative}: for every
        \(\sigma\in\Sent(\Lang)\),
        \(T\cup\set{\delta_R}\entails\sigma\) iff
        \(T\entails\sigma\).
\end{enumerate}
\textit{Status: complete proof.}
\end{theorem}
\begin{proof}
Both clauses use invariance under reduct
(\cref{lem:reduct-invariance}); that lemma is stated in
\cref{ch:fol-deduction}, and its proof uses only the satisfaction clauses,
not any material of this chapter.  Every member of
\(T\cup\set{\delta_R,\sigma}\) is a sentence, so satisfaction is
assignment-independent and it suffices to reason with models.

For clause 1, let \(\Meta\) be an \(\Lang\)-model of \(T\).  Since
\(\FV(\varphi)\subseteq\set{x_1,\ldots,x_n}\), formula coincidence
(\cref{lem:fol-coincidence}) shows that whether
\(\Meta,s[x_1\mapsto a_1]\cdots[x_n\mapsto a_n]\entails\varphi\) holds
depends only on the tuple \((a_1,\ldots,a_n)\), not on \(s\); define
\(R^*\subseteq\abs{\Meta}^n\) as the set of tuples for which it holds.
Now let \(\Meta'\) be any \(\Lang_R\)-expansion of \(\Meta\), and let
\(s\) be any assignment.  Unwinding the satisfaction clause for
\(\forall\) at each of the \(n\) distinct quantified variables and then
the clauses for \(\leftrightarrow\), \(\Meta'\) satisfies \(\delta_R\)
iff for every tuple \((a_1,\ldots,a_n)\in\abs{\Meta}^n\), writing
\(s'=s[x_1\mapsto a_1]\cdots[x_n\mapsto a_n]\),
\[
 \Meta',s'\entails R(x_1,\ldots,x_n)
 \quad\text{iff}\quad
 \Meta',s'\entails\varphi.
\]
The atomic clause evaluates the left side as
\((a_1,\ldots,a_n)\in R^{\Meta'}\), and because \(\varphi\) is an
\(\Lang\)-formula, invariance under reduct evaluates the right side as
\((a_1,\ldots,a_n)\in R^*\).  Hence \(\Meta'\) satisfies \(\delta_R\)
iff \(R^{\Meta'}=R^*\): the expansion interpreting \(R\) as \(R^*\) is
the unique one satisfying \(\delta_R\).  It also satisfies every member
of \(T\), because each is an \(\Lang\)-sentence true in its reduct
\(\Meta\).

For clause 2, first suppose \(T\entails\sigma\), and let \(\mathcal N\)
be any model of \(T\cup\set{\delta_R}\) in the signature \(\Lang_R\).
By invariance under reduct, \(\mathcal N\restr\Lang\) is an
\(\Lang\)-model of \(T\); hence it satisfies \(\sigma\), and invariance
under reduct returns \(\mathcal N\entails\sigma\).  Thus
\(T\cup\set{\delta_R}\entails\sigma\).  Conversely, suppose
\(T\cup\set{\delta_R}\entails\sigma\), and let \(\Meta\) be any
\(\Lang\)-model of \(T\).  Clause 1 expands \(\Meta\) to a model of
\(T\cup\set{\delta_R}\), which therefore satisfies \(\sigma\); since
\(\sigma\) is an \(\Lang\)-sentence, invariance under reduct gives
\(\Meta\entails\sigma\).  Thus \(T\entails\sigma\).
\end{proof}

Once soundness (\cref{thm:fol-soundness}) and completeness
(\cref{thm:fol-completeness}, under its Axiom-of-Choice hypothesis) are
available, clause 2 transfers to derivability: for every
\(\sigma\in\Sent(\Lang)\),
\(T\cup\set{\delta_R}\derives_{\ND_{\FOL}}\sigma\) iff
\(T\derives_{\ND_{\FOL}}\sigma\), since the left-to-right direction
composes soundness, clause 2, and completeness, and the converse holds
because a derivation from a set is a derivation from any superset.

\section{Axiomatization errors}
\label{sec:fol-theories-01}

\begin{warning}
When writing an axiomatization, check all of the following:
\begin{enumerate}
  \item every displayed symbol belongs to the declared signature;
  \item the arity of every symbol is respected;
  \item every axiom is a sentence;
  \item the nonempty-domain convention is accounted for;
  \item intended restrictions are written with implications for universal
        quantification and conjunctions for existential quantification;
  \item a schema is identified as a set of sentences, not printed as though it
        were one first-order sentence; and
  \item both directions of the claim \(\Mod(T)=\mathcal K\) are proved or
        explained.
\end{enumerate}
\end{warning}

\chapter{First-Order Deduction}
\label{ch:fol-deduction}

\section{Natural-deduction rules for quantifiers and equality}
\label{sec:fol-deduction-01}

The calculus \(\ND_{\FOL}\) contains every rule of \(\ND_{\CPL}\) and the
following rules.  Every displayed substitution is the capture-avoiding
substitution of \cref{ch:fol-syntax-semantics}.
\[
\begin{array}{c@{\qquad}c}
\Rule[\(\forall E\)]{\forall x\,\varphi}{\varphi[t/x]}
&
\Rule[\(\exists I\)]{\varphi[t/x]}{\exists x\,\varphi}
\end{array}
\]
\[
\Rule[\(\forall I\)]{\varphi}{\forall x\,\varphi}
\side{provided \(x\) is not free in any open assumption}.
\]
\[
\Rule[\(\exists E\)]{
 \exists x\,\varphi\qquad
 \begin{array}{c}[\varphi[y/x]]^i\\ \vdots\\ \psi\end{array}}
 {\psi}
\side{\(y\) is absent from \(\psi\), from every other open assumption,
and from \(\exists x\,\varphi\)}.
\]
A variable subject to a freshness side condition in a rule instance is the
\emph{eigenvariable} of that instance; in \(\exists E\) the eigenvariable
is \(y\), and in \(\forall I\) it is the generalized variable \(x\).  The
rule \(\exists E\) discharges
all assumptions marked \(i\).  Its freshness condition prevents the temporary
witness from appearing in the conclusion or surviving assumptions.

Because formulas are strings rather than alpha-equivalence classes, the
calculus also contains the explicit conversion rule
\[
  \Rule[\(\alpha\)]{\varphi}{\psi}
  \side{provided \(\varphi\equiv_\alpha\psi\)}.
\]
The rule's premise and conclusion are formulas, and its side condition is the
syntactic alpha-equivalence relation of
\cref{sec:fol-syntax-semantics-04}.  It changes only bound-variable names, so
it is neither identity elimination nor replacement by an arbitrary
logically equivalent formula.  The rule is required below when a fresh
eigenvariable of \(\forall I\) produces an alpha-equivalent, but not
literally identical, quantified formula.

Equality is governed by identity introduction and substitutivity:
\[
\begin{array}{c@{\qquad}c}
\Rule[\(=I\)]{}{t=t}
&
\Rule[\(=E\)]{r=t\qquad\varphi[r/x]}{\varphi[t/x]}.
\end{array}
\]
The formula \(\varphi\) may contain zero or more free occurrences of \(x\).
The rule uses actual identity in the semantics; it does not merely postulate
that an arbitrary relation is reflexive.

\begin{definition}[First-order derivability]
For \(\Gamma\subseteq\Fm(\Lang)\), write
\(\Gamma\derives_{\ND_{\FOL}}\varphi\) iff there is a
\(\ND_{\FOL}\)-derivation of \(\varphi\) whose open assumptions all belong to
\(\Gamma\).
\end{definition}

\begin{lemma}[Invariance under reduct]
\label{lem:reduct-invariance}
Let \(\Lang'\) extend \(\Lang\), let \(\Meta'\) be a
\(\Lang'\)-structure, and let \(s\) be a variable assignment in its domain.
For every \(\Lang\)-term \(t\) and \(\Lang\)-formula \(\varphi\),
\[
 \sem{t}^{\Meta'}_s=\sem{t}^{\Meta'\restr\Lang}_s,
 \qquad
 \Meta',s\entails\varphi
 \quad\Longleftrightarrow\quad
 \Meta'\restr\Lang,s\entails\varphi.
\]
\end{lemma}
\begin{proof}
Induct on terms for the first equality.  Every function symbol in the term has
the same interpretation in the expansion and reduct.  Then induct on formulas.
The atomic cases use the term result and equality of the old relation-symbol
interpretations.  The clauses for \(\bottom\) and the connectives are
immediate from the induction hypotheses.  In a quantifier case the two
structures have the same domain, so they range over the same elements; apply
the induction hypothesis under each updated assignment.
\end{proof}

\begin{example}[Universal reasoning]
The premises
\(\forall x(P(x)\to Q(x))\) and \(\forall xP(x)\) derive
\(\forall xQ(x)\):
\[
\begin{array}{rll}
1.&\forall x(P(x)\to Q(x))&\text{premise}\cr
2.&\forall xP(x)&\text{premise}\cr
3.&P(y)\to Q(y)&\forall E,1\cr
4.&P(y)&\forall E,2\cr
5.&Q(y)&\to E,3,4\cr
6.&\forall yQ(y)&\forall I,5\cr
7.&\forall xQ(x)&\alpha,6.
\end{array}
\]
The variable \(y\) is not free in either premise, so the side condition at
line 6 is met.  Line 7 applies the \(\alpha\) rule, because
\(\forall yQ(y)\) and \(\forall xQ(x)\) are alpha-equivalent but not
identical strings.
\end{example}

\begin{example}[Existential elimination]
From \(\exists x(P(x)\land Q(x))\) derive \(\exists xP(x)\):
\[
\begin{array}{rll}
1.&\exists x(P(x)\land Q(x))&\text{premise}\cr
2.&\quad[P(y)\land Q(y)]^1&\text{assumption, \(y\) fresh}\cr
3.&\quad P(y)&\land E_1,2\cr
4.&\quad\exists xP(x)&\exists I,3\cr
5.&\exists xP(x)&\exists E,1,2\text{--}4.
\end{array}
\]
The eigenvariable \(y\) does not occur in the conclusion.
\end{example}

\begin{example}[Symmetry of equality]
The rule \(=E\) derives \(t=r\) from \(r=t\).  Choose a variable \(x\)
with \(x\notin\operatorname{Var}(r)\); this is possible because \(\Var\)
is infinite and \(r\) is a finite string.  From \(=I\) obtain
\(r=r\).  In the formula \(\varphi\defeq x=r\), the instance
\(\varphi[r/x]\) is \(r=r\) and the instance \(\varphi[t/x]\) is
\(t=r\).  Apply \(=E\) to \(r=t\) and \(r=r\).
\end{example}

\section{First-order sequent rules}
\label{sec:fol-deduction-02}

The first-order sequent calculus \(\LK_{\FOL}\) extends
\(\LK_{\CPL}\).  A sequent is satisfied by \(\Meta,s\) under the same
left-conjunctive and right-disjunctive reading used for propositional
sequents.  Its additional rules are
\[
\begin{array}{c@{\qquad}c}
\Rule[\(\forall L\)]{\Gamma,\varphi[t/x]\Rightarrow\Delta}
 {\Gamma,\forall x\,\varphi\Rightarrow\Delta}
&
\Rule[\(\exists R\)]{\Gamma\Rightarrow\Delta,\varphi[t/x]}
 {\Gamma\Rightarrow\Delta,\exists x\,\varphi}
\\[2.8ex]
\Rule[\(\forall R\)]{\Gamma\Rightarrow\Delta,\varphi[y/x]}
 {\Gamma\Rightarrow\Delta,\forall x\,\varphi}
&
\Rule[\(\exists L\)]{\Gamma,\varphi[y/x]\Rightarrow\Delta}
 {\Gamma,\exists x\,\varphi\Rightarrow\Delta}.
\end{array}
\]
In \(\forall R\), the eigenvariable \(y\) occurs free in no formula of
\(\Gamma,\Delta,\forall x\,\varphi\).  In \(\exists L\), it occurs free in
no formula of \(\Gamma,\Delta,\exists x\,\varphi\).

Equality may be included by the initial sequent
\(\Rightarrow t=t\) and the replacement rule
\[
 \Rule[\(=L\)]{
   \Gamma,\varphi[r/x]\Rightarrow\Delta}
  {\Gamma,r=t,\varphi[t/x]\Rightarrow\Delta}.
\]
The symmetric replacement instance---from
\(\Gamma,\varphi[t/x]\Rightarrow\Delta\) infer
\(\Gamma,r=t,\varphi[r/x]\Rightarrow\Delta\)---is derivable.  First,
\(r=t\Rightarrow t=r\) is derivable: apply \(=L\) with the auxiliary
formula \(t=x\), for a variable \(x\) absent from \(r\) and \(t\), to
the identity sequent \(t=r\Rightarrow t=r\), obtaining
\(r=t,\,t=t\Rightarrow t=r\), and cut with the initial sequent
\(\Rightarrow t=t\).  Second, the displayed rule with the roles of the two
terms exchanged turns the premise into
\(\Gamma,t=r,\varphi[r/x]\Rightarrow\Delta\), and a cut with
\(r=t\Rightarrow t=r\) yields the instance.  A first-order sequent
derivation is a
finite tree generated by the propositional and displayed first-order rules.

\section{Soundness}
\label{sec:fol-deduction-03}

\begin{theorem}[Soundness of first-order natural deduction]
\label{thm:fol-soundness}
For every set \(\Gamma\cup\set{\varphi}\subseteq\Fm(\Lang)\),
\[
 \Gamma\derives_{\ND_{\FOL}}\varphi
 \quad\Longrightarrow\quad
 \Gamma\entails\varphi.
\]
\end{theorem}
\begin{proof}
Induct on the construction of the derivation, using the strengthened statement
that any structure and assignment satisfying all open assumptions satisfies
the conclusion.  The propositional cases are those of
\cref{thm:prop-nd-soundness}.

For \(\forall E\), if \(\Meta,s\entails\forall x\,\varphi\), then it
satisfies \(\varphi\) under
\(s[x\mapsto\sem{t}^{\Meta}_s]\).  The substitution lemma gives
\(\Meta,s\entails\varphi[t/x]\).  The \(\exists I\) case reverses this
argument: a satisfying value \(\sem{t}^{\Meta}_s\) supplies the existential
witness.

For \(\forall I\), let \(\Meta,s\) satisfy all open assumptions and choose an
arbitrary \(a\in\abs{\Meta}\).  Because \(x\) is free in no open assumption,
formula coincidence shows that \(s[x\mapsto a]\) also satisfies every open
assumption.  The induction hypothesis gives
\(\Meta,s[x\mapsto a]\entails\varphi\).  Since \(a\) was arbitrary,
\(\Meta,s\entails\forall x\,\varphi\).

For \(\exists E\), let \(\Meta,s\) satisfy the surviving open assumptions and
\(\exists x\,\varphi\).  Choose \(a\) such that
\(\Meta,s[x\mapsto a]\entails\varphi\), and put
\(s'=s[y\mapsto a]\).  Freshness of \(y\), formula coincidence, and the
substitution lemma give \(\Meta,s'\entails\varphi[y/x]\), while freshness also
ensures that \(s'\) satisfies every other open assumption.  The induction
hypothesis for the subsidiary derivation yields \(\Meta,s'\entails\psi\).
Because \(y\notin\FV(\psi)\), formula coincidence gives
\(\Meta,s\entails\psi\).

The \(\alpha\) rule is sound by alpha-invariance.  For \(=I\), actual
identity is reflexive.  For \(=E\), satisfaction of
\(r=t\) gives
\(\sem{r}^{\Meta}_s=\sem{t}^{\Meta}_s\).  The substitution lemma says that
the two instances of \(\varphi\) are evaluated under updates of \(s\) by
these equal elements, so satisfaction of one instance implies satisfaction of
the other.  This treats every rule of \(\ND_{\FOL}\) that is not a rule
of \(\ND_{\CPL}\), completing the induction.
\end{proof}

\begin{theorem}[Soundness of the first-order sequent calculus]
\label{thm:fol-lk-soundness}
Every sequent derivable in \(\LK_{\FOL}\) is satisfied by every
\(\Lang\)-structure and assignment.
\end{theorem}
\begin{proof}
Extend the induction in \cref{thm:prop-lk-soundness}: for each rule,
if every premise is satisfied by every structure and assignment, so is the
conclusion.  The rules \(\forall L\) and \(\exists R\) use the term value
and the substitution lemma.

For \(\forall R\), suppose toward a contradiction that some \(\Meta,s\)
does not satisfy the conclusion
\(\Gamma\Rightarrow\Delta,\forall x\,\varphi\) while the premise
\(\Gamma\Rightarrow\Delta,\varphi[y/x]\) is satisfied by every structure
and assignment.  Then \(\Meta,s\) satisfies every formula in
\(\operatorname{supp}(\Gamma)\), no formula in
\(\operatorname{supp}(\Delta)\), and not \(\forall x\,\varphi\).  The
universal satisfaction clause supplies an element \(a\in\abs{\Meta}\)
with \(\Meta,s[x\mapsto a]\notentails\varphi\).  Put
\(s'=s[y\mapsto a]\), where \(y\) is the eigenvariable.

First, \(\Meta,s'\notentails\varphi[y/x]\).  The substitution lemma
(\cref{thm:fol-substitution}) gives
\(\Meta,s'\entails\varphi[y/x]\) iff
\(\Meta,s'[x\mapsto\sem{y}^{\Meta}_{s'}]\entails\varphi\), and
\(\sem{y}^{\Meta}_{s'}=s'(y)=a\).  If \(y=x\), then
\(s'[x\mapsto a]=s[x\mapsto a]\) and the right side fails directly.  If
\(y\ne x\), then the eigenvariable condition---\(y\) occurs free in no
formula of \(\Gamma,\Delta,\forall x\,\varphi\)---gives
\(y\notin\FV(\varphi)\), so \(s'[x\mapsto a]\) and
\(s[x\mapsto a]\) agree on \(\FV(\varphi)\), and formula coincidence
(\cref{lem:fol-coincidence}) reduces the right side to
\(\Meta,s[x\mapsto a]\entails\varphi\), which fails.

Second, the same eigenvariable condition puts \(y\) outside the free
variables of every formula of \(\Gamma\) and \(\Delta\), so \(s\) and
\(s'\) agree on those free variables, and formula coincidence transfers
the pattern above to \(s'\): \(\Meta,s'\) satisfies every formula in
\(\operatorname{supp}(\Gamma)\) and no formula in
\(\operatorname{supp}(\Delta)\).  Together with
\(\Meta,s'\notentails\varphi[y/x]\), this means \(\Meta,s'\) does not
satisfy the premise, a contradiction.  The eigenvariable condition is what
licenses both uses of formula coincidence; without it, updating \(y\)
could change the truth values of the contexts or of \(\varphi\) itself.

The \(\exists L\) case is dual.  The equality initial sequent
\(\Rightarrow t=t\) is satisfied by every structure and assignment
because actual identity is reflexive, and \(=L\) preserves satisfaction by
every structure and assignment by the same equal-update argument used for
\(=E\) in \cref{thm:fol-soundness}.
\end{proof}

\section{Quantifier-rule errors}
\label{sec:fol-deduction-04}

\begin{warning}
None of the following inference steps is licensed by the rules of
\(\ND_{\FOL}\), and each associated semantic-consequence claim fails:
\begin{enumerate}
  \item from the open assumption \(P(x)\), infer \(\forall xP(x)\);
  \item from \(\exists xP(x)\), introduce a permanent constant \(c\) and
        conclude \(P(c)\);
  \item in \(\exists E\), allow the eigenvariable to occur in the conclusion;
  \item from \(\forall x\exists yR(x,y)\), infer
        \(\exists y\forall xR(x,y)\); and
  \item replace a free variable by a term using textual substitution that
        permits capture.
\end{enumerate}
The first and third violate explicit side conditions.  The second treats an
existential witness as a named element outside its subsidiary derivation.  The
fourth reverses witness dependence.  The fifth changes formula meaning.
\end{warning}

\section{Exercises}
\label{sec:fol-deduction-05}

\begin{exercise}
\label{ex:fol-deduction-01}
Compute \(\FV(\varphi)\) for
\[
 \varphi=P(x,z)\land\forall x\bigl(Q(x,y)\to\exists yR(x,y,z)\bigr).
\]
Identify the quantifier binding each bound occurrence.
\end{exercise}

\begin{exercise}
\label{ex:fol-deduction-02}
Compute capture-avoiding substitutions
\[
 (\exists yR(x,y))[y/x]
 \quad\text{and}\quad
 (\forall x(P(x)\to Q(y)))[f(x)/y].
\]
Record every bound-variable renaming.
\end{exercise}

\begin{exercise}
\label{ex:fol-deduction-03}
Prove term coincidence and formula coincidence simultaneously for a language
with one binary function and one binary relation, then explain why the same
proof covers arbitrary signatures.
\end{exercise}

\begin{exercise}
\label{ex:fol-deduction-04}
Give a \(\ND_{\FOL}\)-derivation of
\[
 \forall x(P(x)\to Q(x)),\ \exists xP(x)
 \derives_{\ND_{\FOL}}\exists xQ(x).
\]
Mark the eigenvariable and every discharged assumption.
\end{exercise}

\begin{exercise}
\label{ex:fol-deduction-05}
Derive transitivity of equality from \(=I\) and \(=E\).
\end{exercise}

\begin{exercise}
\label{ex:fol-deduction-06}
Translate each sentence using one declared signature: ``Every reviewer read
exactly one paper''; ``One paper was read by every reviewer''; ``No reviewer
read only papers that cited themselves.''  State the domain.
\end{exercise}

\begin{exercise}
\label{ex:fol-deduction-07}
Prove both directions of the claim that the displayed group theory
axiomatizes groups in the signature \(\Lang_{\mathrm{grp}}\).
\end{exercise}

\begin{exercise}
\label{ex:fol-deduction-08}
Write a theory whose models are simple undirected graphs.  Treat equality as
logical identity, and state whether loops are permitted.
\end{exercise}

\begin{exercise}
\label{ex:fol-deduction-09}
Give a cut-free \(\LK_{\FOL}\)-derivation of
\(\forall xP(x)\Rightarrow\exists xP(x)\).  Identify where the nonempty-domain
convention is used.
\end{exercise}

\chapter{Henkin Completeness and Compactness}
\label{ch:fol-metatheory}

The completeness theorem for classical first-order logic is G\"odel's
\cite{Godel1930}.  This chapter develops the construction introduced by
Henkin \cite{Henkin1949}, which extracts a term structure from a maximal
consistent, witness-saturated set.

\section{Syntactic preparation}
\label{sec:fol-metatheory-01}

\begin{definition}[First-order consistency]
A set \(\Gamma\subseteq\Fm(\Lang)\) is \emph{consistent} iff
\(\Gamma\notderives_{\ND_{\FOL}}\bottom\).  It is \emph{maximal consistent}
iff it is consistent and \(\Gamma\cup\set{\varphi}\) is inconsistent for
every \(\varphi\notin\Gamma\).
\end{definition}

The closure properties in \cref{lem:prop-mcs} also hold for a maximal
consistent \(\Delta\subseteq\Fm(\Lang)\), but their first-order proof does not
silently graft derivations across eigenvariable rules.  For deductive closure,
suppose \(\Delta\derives_{\ND_{\FOL}}\varphi\) and
\(\varphi\notin\Delta\).  Maximality supplies a derivation of \(\bottom\) from
\(\Delta\cup\set{\varphi}\).  Because \(\Delta\) is consistent, at least one
open occurrence of the added assumption remains in that derivation.  Apply
\(\neg I\) at its root, discharging every such occurrence, to obtain
\(\Delta\derives_{\ND_{\FOL}}\neg\varphi\).  Combining this derivation with
the given derivation of \(\varphi\) by \(\neg E\) derives \(\bottom\) from
\(\Delta\), contradicting consistency.  Clauses 2--5 of
\cref{lem:prop-mcs} then use this deductive-closure result and only the
displayed propositional rules, so their proofs transfer without an
eigenvariable-renaming step.

\begin{lemma}[Classical quantifier duality]
\label{lem:quantifier-duality}
The calculus \(\ND_{\FOL}\) derives
\[
 \neg\forall x\,\varphi\leftrightarrow\exists x\,\neg\varphi
 \quad\text{and}\quad
 \neg\exists x\,\varphi\leftrightarrow\forall x\,\neg\varphi.
\]
\end{lemma}
\begin{proof}
For \(\neg\exists x\,\varphi\to\forall x\neg\varphi\), take a variable
\(y\) absent from \(\varphi\) and free in no open assumption.  Assume \(\varphi[y/x]\); existential
introduction contradicts \(\neg\exists x\varphi\).  Negation introduction
gives \(\neg\varphi[y/x]\), and universal introduction gives
\(\forall y\,\neg\varphi[y/x]\).  Since \(y\) is fresh for \(\varphi\), the
\(\alpha\) rule gives \(\forall x\neg\varphi\).  The reverse
implication follows by existential elimination: a witness to \(\varphi\)
contradicts the corresponding instance of \(\forall x\neg\varphi\).

For \(\exists x\neg\varphi\to\neg\forall x\varphi\), eliminate the
existential and contradict the corresponding universal instance.  For the
remaining direction, assume \(\neg\forall x\varphi\) and, for reductio,
\(\neg\exists x\neg\varphi\).  The already proved implication gives
\(\forall x\neg\neg\varphi\).  For a variable \(y\) absent from
\(\varphi\) and free in no open assumption, double-negation
elimination gives \(\varphi[y/x]\); universal introduction gives
\(\forall y\,\varphi[y/x]\), and the \(\alpha\) rule gives
\(\forall x\varphi\), contradicting the first assumption.  Reductio yields
\(\exists x\neg\varphi\).  Every eigenvariable is chosen absent from the
surviving assumptions and conclusion.
\end{proof}

\begin{lemma}[Fresh-constant lemma]
\label{lem:fresh-constant}
Let \(c\) be a constant symbol occurring in neither \(\Gamma\) nor
\(\varphi\).  If
\(\Gamma\derives_{\ND_{\FOL}}\varphi[c/x]\), then
\(\Gamma\derives_{\ND_{\FOL}}\forall x\varphi\).
\end{lemma}
\begin{proof}
Fix a derivation witnessing the premise.  Because the derivation is finite,
only a finite set \(\Gamma_0\subseteq\Gamma\) occurs among its open
assumptions.  Choose a variable \(y\) absent from \(\Gamma_0\), from
\(\varphi\), and from the finite derivation.  Uniformly replace \(c\) by \(y\)
in that derivation.  Induction on the derivation shows that every rule instance
remains an instance: \(c\) has no binding behavior, it occurs in no open
assumption, and \(y\) occurred in no original rule instance, so every
eigenvariable condition is preserved.  Hence
\(\Gamma_0\derives_{\ND_{\FOL}}\varphi[y/x]\).  Since \(y\) is free in no
member of \(\Gamma_0\), \(\forall I\) gives
\(\Gamma_0\derives_{\ND_{\FOL}}\forall y\,\varphi[y/x]\).  Because
\(y\notin\operatorname{Var}(\varphi)\), this conclusion is alpha-equivalent
to \(\forall x\varphi\); the \(\alpha\) rule therefore gives the latter.  The
same finite derivation also witnesses
\(\Gamma\derives_{\ND_{\FOL}}\forall x\varphi\).
\end{proof}

\begin{definition}[Witness property]
A set \(\Delta\subseteq\Fm(\Lang)\) has the \emph{witness property} iff, for
every formula \(\exists x\varphi\in\Delta\), there is a constant symbol \(c\)
of \(\Lang\) such that \(\varphi[c/x]\in\Delta\).
\end{definition}

The witness property is necessary in the term-structure construction: membership
of an existential formula must supply a term naming a witness.

\begin{lemma}[Witness-extension lemma]
\label{lem:witness-extension}
Let \(\Sigma\) be consistent, let \(c\) occur in neither \(\Sigma\) nor
\(\exists x\varphi\), and put
\[
 H_{x,\varphi,c}\defeq
 \exists x\varphi\to\varphi[c/x].
\]
Then \(\Sigma\cup\set{H_{x,\varphi,c}}\) is consistent.
\end{lemma}
\begin{proof}
Suppose the extension is inconsistent.  Negation introduction gives
\(\Sigma\derives\neg H_{x,\varphi,c}\).  Classical propositional derivations
then give
\[
 \Sigma\derives\exists x\varphi
 \quad\text{and}\quad
 \Sigma\derives\neg\varphi[c/x],
\]
because \(\neg(A\to B)\derives A\land\neg B\).  The fresh-constant lemma
applied to the second derivation gives
\(\Sigma\derives\forall x\neg\varphi\).  Quantifier duality gives
\(\Sigma\derives\neg\exists x\varphi\), contradicting the first derivation.
Thus the extension is consistent.
\end{proof}

\begin{lemma}[Eliminating newly added constants]
\label{lem:new-constant-elimination}
Let \(\Lang'\) extend \(\Lang\) only by constant symbols, and let
\(\Gamma\cup\set{\varphi}\subseteq\Fm(\Lang)\).  If there is an
\(\ND_{\FOL}\)-derivation of \(\varphi\) from \(\Gamma\) in the expanded
signature \(\Lang'\), then there is such a derivation in \(\Lang\).
\end{lemma}
\begin{proof}
Fix an expanded-signature derivation.  It is finite, so only finitely many
constants from \(\Lang'\setminus\Lang\), say
\(d_1,\ldots,d_n\), occur in it.  Choose distinct variables
\(z_1,\ldots,z_n\) absent from every formula and rule annotation in the
derivation.  Simultaneously replace each \(d_i\) by \(z_i\) throughout the
derivation.

Induction on the derivation verifies that every inference remains an instance
of the same rule.  Replacement commutes with the propositional constructors
and with term formation.  Because the \(z_i\) occurred in no original binder,
replacement causes no capture, and for every quantifier instance it commutes
with capture-avoiding substitution.  The eigenvariables of
\(\forall I\) and of \(\exists E\) are original variables,
distinct from all \(z_i\); replacing constants therefore preserves their
absence conditions.  Replacement also preserves alpha-equivalence and every
identity-rule instance.  The end formula and open assumptions contain none of
the \(d_i\), so they are unchanged.  The transformed tree is consequently a
\(\Lang\)-derivation of \(\varphi\) from \(\Gamma\).
\end{proof}

\begin{lemma}[Transfinite Lindenbaum lemma]
\label{lem:transfinite-lindenbaum}
Let \(\Lang'\) be a signature, let \(\lambda\) be an ordinal, and let
\(\langle\varphi_\alpha:\alpha<\lambda\rangle\) be an enumeration of
\(\Fm(\Lang')\) in which every formula occurs at least once.  Then every
set \(\Sigma\subseteq\Fm(\Lang')\) that is consistent for
\(\ND_{\FOL}\) is contained in a maximal consistent
\(\Delta\subseteq\Fm(\Lang')\).  The construction is choice-free given the
enumeration; the Axiom of Choice enters only when such an enumeration must
first be produced.
\textit{Status: complete proof.}
\end{lemma}
\begin{proof}
By \cref{thm:transfinite-recursion} with domain \(\lambda+1\), define a
sequence \(\langle\Gamma_\alpha:\alpha\le\lambda\rangle\) by
\[
 \Gamma_0=\Sigma,\qquad
 \Gamma_{\alpha+1}=
 \begin{cases}
   \Gamma_\alpha\cup\set{\varphi_\alpha},&\text{if this set is
     consistent},\\
   \Gamma_\alpha,&\text{otherwise},
 \end{cases}
 \qquad
 \Gamma_\beta=\bigcup_{\alpha<\beta}\Gamma_\alpha
 \ \text{ for limit }\beta.
\]

\emph{Chain claim:} if \(\alpha\le\beta\le\lambda\), then
\(\Gamma_\alpha\subseteq\Gamma_\beta\).  Proceed by transfinite induction
on \(\beta\): the claim is trivial at \(\beta=\alpha\), both successor
clauses extend the preceding stage, and a limit stage is the union of all
earlier stages.

\emph{Consistency:} every \(\Gamma_\alpha\) is consistent, by transfinite
induction on \(\alpha\).  Stage zero is the hypothesis on \(\Sigma\), and
the successor clause preserves consistency by its case split.  At a limit
\(\beta\), suppose \(\Gamma_\beta\derives_{\ND_{\FOL}}\bottom\).  The
derivation is finite, so its open assumptions form a finite
\(F\subseteq\Gamma_\beta\).  Each \(\psi\in F\) lies in some
\(\Gamma_{\alpha_\psi}\) with \(\alpha_\psi<\beta\).  If
\(F\ne\varnothing\), put
\(\alpha^{*}=\max_{\psi\in F}\alpha_\psi<\beta\), which exists because a
finite nonempty set of ordinals has a maximum; if \(F=\varnothing\), put
\(\alpha^{*}=0<\beta\).  The chain claim locates
\(F\subseteq\Gamma_{\alpha^{*}}\), so \(\Gamma_{\alpha^{*}}\) is
inconsistent, contradicting the induction hypothesis at \(\alpha^{*}\).

\emph{Maximality:} put \(\Delta=\Gamma_\lambda\); it is consistent, and it
contains \(\Sigma\) by the chain claim.  Suppose \(\varphi\notin\Delta\)
and pick \(\alpha<\lambda\) with \(\varphi_\alpha=\varphi\).  Rejection at
stage \(\alpha+1\) means that \(\Gamma_\alpha\cup\set{\varphi}\) derives
\(\bottom\).  A derivation from a set is a derivation from any superset,
because derivability from a set requires only that every open assumption
belong to it; hence \(\Delta\cup\set{\varphi}\) derives \(\bottom\).
Therefore \(\Delta\) is maximal consistent.
\end{proof}

\begin{lemma}[Henkin--Lindenbaum lemma]
\label{lem:henkin-lindenbaum}
Let \(\Lang\) be a signature and assume the Axiom of Choice in the
ambient set theory.  Every consistent
\(\Sigma\subseteq\Fm(\Lang)\) is contained in a maximal consistent set
\(\Delta\) in an expanded signature \(\Lang^H\), and \(\Delta\) has the
witness property.
\textit{Status: complete proof.}
\end{lemma}
\begin{proof}
 Let \(\kappa=\max(\abs{\Lang},\aleph_0)\).  Add new constant symbols
 \(\langle d_\xi:\xi<\kappa\rangle\).  Regard \(\Sigma\) as a set of formulas
 in the expanded signature.
It remains consistent there: otherwise
\cref{lem:new-constant-elimination} would transform an expanded-signature
derivation of \(\bottom\) into an original-signature derivation from
\(\Sigma\).  By \cref{cor:syntax-counting}, the expanded
language has at most \(\kappa\) formulas.  The \(\kappa\) formulas
\(\exists x_0\,(x_0=d_\xi)\), for \(\xi<\kappa\), are pairwise distinct, so
it has exactly \(\kappa\) formulas and exactly \(\kappa\) existential formulas.
Since a cardinal
is an ordinal by definition, fix a bijective enumeration
\[
  \langle E_\alpha:\alpha<\kappa\rangle,
  \qquad
  E_\alpha=\exists x_{n_\alpha}\varphi_\alpha
  \text{ with }n_\alpha\in\N,
\]
 of its existential formulas.  By the transfinite-recursion theorem
\cref{thm:transfinite-recursion}, at stage
\(\alpha<\kappa\) add a witness implication
\(H_{x_{n_\alpha},\varphi_\alpha,c_\alpha}\), where \(c_\alpha\) is the
least \(d_\xi\), in the displayed index order, that is absent from the current
theory and from \(E_\alpha\).  This least-available convention makes the
recursive clause functional.  Before stage \(\alpha\), fewer than
\(\kappa\) finite witness formulas have been added.  By
\cref{thm:infinite-cardinal-finite-coding}, their union contains fewer than
\(\kappa\) constants; \(E_\alpha\) contains only finitely many constants.
Therefore one of the \(\kappa\) new constants remains available.  The
witness-extension lemma
preserves consistency at successor stages.  At a limit stage take the union;
any derivation of \(\bottom\) would use finitely many assumptions and would
already occur at an earlier stage.  Let \(\Sigma^H\) be the resulting
consistent set.

The set \(\Fm(\Lang^H)\) has cardinality exactly \(\kappa\), as
established above from \cref{cor:syntax-counting} and the
formulas \(\exists x_0\,(x_0=d_\xi)\).  By
\cref{cor:cardinal-assignment}, under the Axiom of Choice already assumed,
fix a bijective enumeration
\(\langle\varphi_\alpha:\alpha<\kappa\rangle\) of \(\Fm(\Lang^H)\).
Apply \cref{lem:transfinite-lindenbaum} to \(\Sigma^H\) under this
enumeration to obtain a maximal consistent
\(\Delta\supseteq\Sigma^H\supseteq\Sigma\).  Every witness implication remains
in \(\Delta\).  If \(\exists x\varphi\in\Delta\), implication elimination
and deductive closure put its designated instance \(\varphi[c/x]\) in
\(\Delta\).  Hence \(\Delta\) has the witness property.
\end{proof}

\section{The term structure}
\label{sec:fol-metatheory-02}

Fix a maximal consistent \(\Delta\subseteq\Fm(\Lang^H)\) with the witness
property.  Define a relation on \(\Term(\Lang^H)\) by
\[
 r\sim_\Delta t
 \quad\text{iff}\quad
 \Delta\derives_{\ND_{\FOL}}r=t.
\]

\begin{lemma}[Equality congruence]
\label{lem:equality-congruence}
The relation \(\sim_\Delta\) is an equivalence relation.  If
\(n\in\N\), \(r_i,t_i\in\Term(\Lang^H)\), and
\(r_i\sim_\Delta t_i\) for \(1\leq i\leq n\), then for every
\(n\)-ary \(f\in\mathcal F\),
\[
 f(r_1,\ldots,r_n)\sim_\Delta f(t_1,\ldots,t_n),
\]
and for every \(n\)-ary \(R\in\mathcal R\),
\[
 R(r_1,\ldots,r_n)\in\Delta
 \quad\text{iff}\quad
 R(t_1,\ldots,t_n)\in\Delta.
\]
\end{lemma}
\begin{proof}
Reflexivity is \(=I\).  Symmetry was derived in
\cref{ch:fol-deduction}; transitivity follows by applying \(=E\) to the
formula \(r=x\).  Repeated applications of \(=E\), one argument at a time,
give compatibility with every function and relation symbol.  Maximal
consistent sets are deductively closed, so each derived replacement belongs
to \(\Delta\).  Symmetry supplies the reverse relational implication.
\end{proof}

\begin{definition}[Term structure]
The \emph{term structure} \(\Meta_\Delta\) has domain
\[
 \abs{\Meta_\Delta}
 =\Term(\Lang^H)/{\sim_\Delta}.
\]
Write \([t]\) for the class of \(t\).  Interpret symbols by
\[
 f^{\Meta_\Delta}([t_1],\ldots,[t_n])
 =[f(t_1,\ldots,t_n)]
\]
and
\[
 R^{\Meta_\Delta}([t_1],\ldots,[t_n])
 \quad\text{iff}\quad
 R(t_1,\ldots,t_n)\in\Delta.
\]
Let \(s_0\) be the variable assignment \(s_0(x)=[x]\).
\end{definition}

\Cref{lem:equality-congruence} proves that these interpretations are
independent of the chosen representatives.  The domain is nonempty
because \(\Lang^H\) has constant symbols.

\begin{lemma}[Term evaluation in the term structure]
\label{lem:term-model-evaluation}
For every term \(t\),
\[
 \sem{t}^{\Meta_\Delta}_{s_0}=[t].
\]
\end{lemma}
\begin{proof}
Induct on \(t\).  For a variable this is the definition of \(s_0\); the
function case is the definition of \(f^{\Meta_\Delta}\) applied to the
induction hypotheses.
\end{proof}

\begin{lemma}[Quantifier membership]
\label{lem:quantifier-membership}
For every formula \(\varphi\) and variable \(x\):
\begin{enumerate}
  \item \(\exists x\varphi\in\Delta\) iff
        \(\varphi[t/x]\in\Delta\) for some term \(t\); and
  \item \(\forall x\varphi\in\Delta\) iff
        \(\varphi[t/x]\in\Delta\) for every term \(t\).
\end{enumerate}
\end{lemma}
\begin{proof}
If an existential formula belongs to \(\Delta\), the witness property supplies
a constant instance.  Conversely, \(\exists I\) derives the existential from
any term instance, and deductive closure places it in \(\Delta\).

Universal elimination proves the forward implication in 2.  Conversely,
suppose \(\forall x\varphi\notin\Delta\).  Maximal consistency gives
\(\neg\forall x\varphi\in\Delta\); quantifier duality and deductive closure
give \(\exists x\neg\varphi\in\Delta\).  By the witness property,
\(\neg\varphi[c/x]\in\Delta\) for some constant \(c\), so consistency gives
\(\varphi[c/x]\notin\Delta\).  Therefore not every term instance belongs to
\(\Delta\).
\end{proof}

\begin{lemma}[First-order truth lemma]
\label{lem:fol-truth}
For every \(\Lang^H\)-formula \(\varphi\),
\[
 \Meta_\Delta,s_0\entails\varphi
 \quad\text{iff}\quad
 \varphi\in\Delta.
\]
\end{lemma}
\begin{proof}
Proceed by strong induction on \(c(\varphi)\).  For equality, term evaluation gives
\[
 \Meta_\Delta,s_0\entails r=t
 \quad\text{iff}\quad [r]=[t]
 \quad\text{iff}\quad \Delta\derives r=t
 \quad\text{iff}\quad r=t\in\Delta,
\]
where the last equivalence is deductive closure and the definition of
\(\sim_\Delta\).  The relation case follows directly from the interpretation
of \(R^{\Meta_\Delta}\), and \(\bottom\notin\Delta\) by consistency.  The
Boolean cases use the propositional closure properties of maximal consistent
sets exactly as in \cref{lem:prop-truth}.

For the existential case, term-structure evaluation and the substitution
lemma (\cref{thm:fol-substitution}) give, for every term \(t\),
\[
 \Meta_\Delta,s_0[x\mapsto[t]]\entails\varphi
 \quad\text{iff}\quad
 \Meta_\Delta,s_0\entails\varphi[t/x].
\]
By \cref{lem:substitution-complexity}, every instance
\(\varphi[t/x]\) has complexity \(c(\varphi)\), which is smaller than the
complexity of \(\exists x\varphi\).  The strong induction hypothesis and
\cref{lem:quantifier-membership}(1) therefore show that
some class \([t]\) satisfies the matrix iff \(\exists x\varphi\in\Delta\).
The universal case uses ``every'' and
\cref{lem:quantifier-membership}(2).  Every element of the term structure has the
form \([t]\), so all semantic quantifier cases are covered.
\end{proof}

\section{Completeness}
\label{sec:fol-metatheory-03}

\begin{theorem}[Strong completeness of classical first-order logic]
\label{thm:fol-completeness}
Let \(\Lang\) be a signature and assume the Axiom of Choice in the
ambient set theory.  For every
\(\Gamma\cup\set{\varphi}\subseteq\Fm(\Lang)\),
\[
 \Gamma\entails\varphi
 \quad\Longrightarrow\quad
 \Gamma\derives_{\ND_{\FOL}}\varphi.
\]
\end{theorem}
\begin{proof}
We prove the contrapositive.  If \(\Gamma\notderives\varphi\), then
\(\Sigma=\Gamma\cup\set{\neg\varphi}\) is consistent.  For suppose a
derivation \(D\) of \(\bottom\) has open assumptions in \(\Sigma\).  If no
open occurrence of \(\neg\varphi\) appeared in \(D\), then \(D\) would
witness \(\Gamma\derives\bottom\), and \(\bottom E\) would give
\(\Gamma\derives\varphi\); otherwise mark all open occurrences of
\(\neg\varphi\) in \(D\) with a common fresh label and close \(D\) by
\(\mathsf{RAA}\) at the root, discharging every such occurrence, to
derive \(\varphi\) from \(\Gamma\).  Either way
\(\Gamma\derives\varphi\), contrary to supposition.  Apply the
Henkin--Lindenbaum lemma to obtain a maximal consistent witness set
\(\Delta\supseteq\Sigma\).  Its term structure and canonical assignment satisfy
every member of \(\Delta\) by the truth lemma.  In particular,
\(\Meta_\Delta,s_0\entails\Gamma\) and
\(\Meta_\Delta,s_0\entails\neg\varphi\).  Thus
the reduct \(\Meta_\Delta\restr\Lang\), with the same assignment \(s_0\),
satisfies \(\Gamma\) and falsifies \(\varphi\) by
\cref{lem:reduct-invariance}.  It is a \(\Lang\)-structure, so
\(\Gamma\notentails\varphi\) under the semantic consequence relation stated
in the theorem.
\end{proof}

Together with soundness, completeness gives
\[
 \Gamma\derives_{\ND_{\FOL}}\varphi
 \quad\text{iff}\quad
 \Gamma\entails\varphi.
\]
This equivalence concerns the specified calculus and the set-based semantics
with nonempty domains.  It does not identify syntactic derivability
with semantic truth as objects.

\begin{corollary}[Completeness for countable signatures, without choice]
\label{cor:fol-completeness-countable}
Let \(\Lang\) be a signature whose set of function and relation symbols
is countable in the sense of \cref{sec:set-prerequisites-05}, finite
signatures included.  Then, for every
\(\Gamma\cup\set{\varphi}\subseteq\Fm(\Lang)\),
\[
 \Gamma\entails\varphi
 \quad\Longrightarrow\quad
 \Gamma\derives_{\ND_{\FOL}}\varphi,
\]
and no use is made of the Axiom of Choice.  This is the
countable-language version of the completeness theorem promised in
\cref{sec:set-prerequisites-05}.
\textit{Status: complete proof.}
\end{corollary}
\begin{proof}
Run the proof of \cref{thm:fol-completeness} unchanged, replacing each
input that the Axiom of Choice supplied with an explicit one.  Choice
enters that proof only through \cref{lem:henkin-lindenbaum}, at exactly
two points: the cardinal bookkeeping at
\(\kappa=\max(\abs{\Lang},\aleph_0)\) through
\cref{cor:cardinal-assignment}, \cref{thm:infinite-cardinal-finite-coding},
and \cref{cor:syntax-counting}; and the bijective enumeration of
\(\Fm(\Lang^H)\) handed to \cref{lem:transfinite-lindenbaum}, whose own
construction is choice-free given the enumeration, as its statement
records.  The enumeration of the existential formulas
\(\langle E_\alpha\rangle\) is an instance of the same second point.

Both inputs are explicit here.  Fix an injection \(\iota\) of the symbol
set of \(\Lang\) into \(\N\); one exists by the definition of countable
set, and fixing a single witness of one existential statement is not an
application of the Axiom of Choice.  Take \(\kappa=\aleph_0\) and add the
witness constants \(\langle d_k:k\in\N\rangle\).  The alphabet of
\(\Lang^H\) is the union of three explicitly enumerated families: the
finitely many logical symbols and punctuation marks together with the
variables in the fixed order \(x_0,x_1,\ldots\) declared in
\cref{ch:fol-syntax-semantics}; the symbols of \(\Lang\) in the order
given by \(\iota\); and the constants \(d_k\) in index order.  Sending
the \(k\)th member of the three families to \(3k\), \(3k+1\), and
\(3k+2\) respectively injects the alphabet into \(\N\), and the string
coding of \cref{ch:strings}---iterating \(\operatorname{pair}_{\N}\) as
in the countability-of-finite-strings lemma---then assigns every formula
of \(\Lang^H\) a distinct natural number code.  Define \(\varphi_n\) to
be the formula whose code is the \(n\)th smallest member of the set of
formula codes; recursion on \(\N\) with the least-number principle
produces this bijective enumeration of \(\Fm(\Lang^H)\) without choice,
and its restriction to the existential formulas enumerates them likewise.
The transfinite recursions of \cref{lem:henkin-lindenbaum} and
\cref{lem:transfinite-lindenbaum} then run along \(\omega\) and
\(\omega+1\) with functional clauses.  At \(\kappa=\aleph_0\) the
cardinal-arithmetic appeals reduce to finite counting and to the displayed
coding: before any stage \(n\) of the witness recursion, only finitely
many witness implications have been added, so they contain only finitely
many of the \(d_k\) and a least available one exists; and the code
injection directly witnesses that \(\Fm(\Lang^H)\) is countably
infinite.

No other ingredient of the completeness proof invokes a choice principle:
\cref{lem:new-constant-elimination} and
\cref{lem:witness-extension} manipulate single finite derivations; the
term structure interprets symbols on equivalence classes with
representative independence proved by
\cref{lem:equality-congruence}; the canonical assignment is
\(s_0(x)=[x]\); and the truth lemma and
\cref{lem:reduct-invariance} are inductions all of whose clauses are
explicit.  The
contrapositive argument of \cref{thm:fol-completeness} therefore goes
through verbatim for countable \(\Lang\) with no appeal to the Axiom of
Choice.
\end{proof}

\begin{theorem}[Completeness of the first-order sequent calculus]
\label{thm:fol-lk-completeness}
Assume the Axiom of Choice in the ambient set theory.  For all finite
multisets \(\Gamma,\Delta\) of \(\Lang\)-formulas, the sequent
\(\Gamma\Rightarrow\Delta\) is derivable in \(\LK_{\FOL}\) iff for every
\(\Lang\)-structure \(\Meta\) and assignment \(s\): if
\(\Meta,s\entails\gamma\) for every
\(\gamma\in\operatorname{supp}(\Gamma)\), then
\(\Meta,s\entails\delta\) for some
\(\delta\in\operatorname{supp}(\Delta)\).
\textit{Status: proof sketch.}
\end{theorem}
\begin{proof}[Proof sketch]
The exact scope is the cut-admitting calculus of
\cref{ch:fol-deduction}.  Left to right is \cref{thm:fol-lk-soundness},
whose satisfaction condition on a sequent is exactly the one displayed
here; at first order the semantic data are a structure together with an
assignment.

For right to left, first record the translation subroutine.  Induction on
a \(\ND_{\FOL}\)-derivation translates the
propositional rules as in the propositional completeness proof, translates
\(\forall I\) and \(\exists E\) using the corresponding eigenvariable rules,
translates \(\forall E\), \(\exists I\), and equality rules directly, and
translates the \(\alpha\) rule by an induction on the generating
alpha-equivalence steps using the sequent quantifier rules.  Cut composes
translated elimination steps.  Bound-variable alpha-equivalence and the
separate eigenvariable freshness conditions are preserved by the relevant
subinductions.  The induction establishes that if
\(\Gamma_0\derives_{\ND_{\FOL}}\varphi\) for a finite set \(\Gamma_0\),
then the sequent \(\Gamma_0\Rightarrow\varphi\) is derivable in
\(\LK_{\FOL}\).  The individual translation trees and those subinductions
are omitted.

Now suppose the semantic condition holds, and list
\(\operatorname{supp}(\Delta)=\set{\delta_1,\ldots,\delta_m}\) with
\(m\geq0\).
\begin{enumerate}
  \item Let \(S\) be the set of members of
        \(\operatorname{supp}(\Gamma)\) together with
        \(\neg\delta_1,\ldots,\neg\delta_m\).  A structure and
        assignment satisfying all of \(S\) would satisfy every member
        of \(\operatorname{supp}(\Gamma)\), hence some \(\delta_j\)
        together with \(\neg\delta_j\).  So no structure and
        assignment satisfies \(S\), and \(S\entails\bottom\).
  \item \Cref{thm:fol-completeness} gives
        \(S\derives_{\ND_{\FOL}}\bottom\).  The open assumptions of
        that derivation form a finite subset of \(S\), say
        \(\gamma_1,\ldots,\gamma_k\in\operatorname{supp}(\Gamma)\) and
        \(\neg\delta_{j_1},\ldots,\neg\delta_{j_r}\).  The translation
        subroutine turns it into an \(\LK_{\FOL}\)-derivation of
        \(\gamma_1,\ldots,\gamma_k,
        \neg\delta_{j_1},\ldots,\neg\delta_{j_r}\Rightarrow\bottom\).
  \item Cut with the initial \(\bottom\)-sequent
        \(\bottom\Rightarrow\), the instance of
        \(\Gamma,\bottom\Rightarrow\Delta\) with both contexts empty,
        empties the succedent.
  \item For each \(j\), the sequent
        \(\Rightarrow\delta_j,\neg\delta_j\) is derivable: start from
        the identity sequent \(\delta_j\Rightarrow\delta_j\) and apply
        \(\neg R\).  Cutting it against the current sequent, with cut
        formula \(\neg\delta_j\),
        \[
         \Rule[Cut]{\Rightarrow\delta_j,\neg\delta_j\qquad
          \Gamma',\neg\delta_j\Rightarrow\Pi}
          {\Gamma'\Rightarrow\delta_j,\Pi},
        \]
        moves \(\delta_j\) into the succedent; after \(r\) such cuts
        the derived sequent is
        \(\gamma_1,\ldots,\gamma_k\Rightarrow
        \delta_{j_1},\ldots,\delta_{j_r}\).
  \item Weakening (\(W_L\), \(W_R\)) adds the unused members of
        \(\Gamma\) and \(\Delta\), and contraction (\(C_L\),
        \(C_R\)) merges duplicates, giving exactly the multiset
        sequent \(\Gamma\Rightarrow\Delta\).
\end{enumerate}
When \(m=0\), the hypothesis asserts that no structure and assignment
satisfies every member of \(\operatorname{supp}(\Gamma)\), step 4
performs no cuts, and the same route derives
\(\Gamma\Rightarrow\Delta\) with empty succedent; an empty \(\Gamma\) is
likewise covered.  The route uses cut, which is primitive in the
displayed calculus, and makes no cut-free claim; extending the
subroutine's translation induction to multi-succedent targets is
unnecessary, because the reduction reuses it verbatim.  The status is a
sketch because the individual translation trees and subinductions of the
subroutine remain omitted; the argument establishes completeness but not
cut elimination.
\end{proof}

\section{Compactness}
\label{sec:fol-metatheory-04}

\begin{theorem}[First-order compactness]
\label{thm:fol-compactness}
Assume the Axiom of Choice in the ambient set theory.  For every
\(T\subseteq\Sent(\Lang)\) and every \(\sigma\in\Sent(\Lang)\):
\begin{enumerate}
  \item \(T\) has a model iff every finite subset of \(T\) has a model;
        and
  \item if \(T\entails\sigma\), then some finite \(T_0\subseteq T\)
        satisfies \(T_0\entails\sigma\).
\end{enumerate}
Clause 1 is the model form of compactness and clause 2 is the consequence
form; they parallel the two clauses of \cref{thm:prop-compactness}.
\end{theorem}
\begin{proof}
For clause 1, left to right is immediate: a model of \(T\) is a model of
each of its finite subsets.  For right to left, suppose every finite
subset of \(T\) has a model.  If \(T\) were inconsistent, a finite
derivation of \(\bottom\) would use only a finite set
\(T_0\subseteq T\) of assumptions.  Soundness (\cref{thm:fol-soundness})
would make \(T_0\) unsatisfiable, contrary to hypothesis.  Thus \(T\) is
consistent.  The Henkin term-structure construction gives an
\(\Lang^H\)-structure satisfying \(T\).  Its \(\Lang\)-reduct is a model
of \(T\) by \cref{lem:reduct-invariance}.

For clause 2, suppose \(T\entails\sigma\).  Completeness
(\cref{thm:fol-completeness}) gives a finite derivation of \(\sigma\)
from \(T\).  Let \(T_0\) be the finite set of assumptions used; soundness
gives \(T_0\entails\sigma\).
\end{proof}

\begin{corollary}
Assume the Axiom of Choice in the ambient set theory.
The class of finite structures in the empty signature is not axiomatizable by
a first-order theory.
\end{corollary}
\begin{proof}
Suppose \(T\) had exactly the finite nonempty structures as models.  Add the
sentences \(\iota_n\) from \cref{ch:fol-theories}.  Every finite subset of
\(T\cup T_\infty\) has a sufficiently large finite model of \(T\).  By
compactness, the union has a model.  Its satisfaction of every \(\iota_n\)
makes it infinite, while its satisfaction of \(T\) makes it finite, a
contradiction.
\end{proof}

\chapter{Cardinality, Normal Forms, and Resolution}
\label{ch:fol-cardinality-normalforms}

\section{Elementary substructures}
\label{sec:fol-cardinality-normalforms-01}

\begin{definition}[Substructure]
Let \(\Meta\) and \(\mathcal N\) be \(\Lang\)-structures.  Then
\(\mathcal N\) is a \emph{substructure} of \(\Meta\), written
\(\mathcal N\subseteq\Meta\), iff:
\begin{enumerate}
  \item \(\abs{\mathcal N}\subseteq\abs{\Meta}\);
  \item every constant has the same value in both structures;
  \item \(\abs{\mathcal N}\) is closed under every function of \(\Meta\),
        and each function of \(\mathcal N\) is the corresponding restriction;
        and
  \item every relation of \(\mathcal N\) is the restriction of the
        corresponding relation of \(\Meta\) to \(\abs{\mathcal N}\).
\end{enumerate}
\end{definition}

\begin{definition}[Elementary substructure]
A substructure \(\mathcal N\subseteq\Meta\) is an \emph{elementary
substructure}, written \(\mathcal N\preccurlyeq\Meta\), iff for every formula
\(\varphi(x_1,\ldots,x_n)\) and every
\(a_1,\ldots,a_n\in\abs{\mathcal N}\),
\[
 \mathcal N\entails\varphi(a_1,\ldots,a_n)
 \quad\text{iff}\quad
 \Meta\entails\varphi(a_1,\ldots,a_n).
\]
The notation with parameters abbreviates satisfaction under assignments that
map the displayed variables to the displayed elements.
\end{definition}

\begin{theorem}[Tarski--Vaught test]
\label{thm:tarski-vaught-fol}
Let \(\mathcal N\subseteq\Meta\).  Then
\(\mathcal N\preccurlyeq\Meta\) iff, for every formula
\(\varphi(y,x_1,\ldots,x_n)\) and every tuple
\(\bar a\in\abs{\mathcal N}^{n}\),
\[
 \Meta\entails\exists y\,\varphi(y,\bar a)
 \quad\Longrightarrow\quad
 \Meta\entails\varphi(b,\bar a)
 \text{ for some }b\in\abs{\mathcal N}.
\]
\end{theorem}
\begin{proof}
If \(\mathcal N\preccurlyeq\Meta\) and the existential formula with
parameters is satisfied in \(\Meta\), it is satisfied in
\(\mathcal N\); its witness
belongs to \(\abs{\mathcal N}\).

Conversely, assume the witness condition and prove elementarity by structural
induction on formulas.  Term values of tuples from \(\mathcal N\) agree in
the two structures because \(\mathcal N\) is closed under the functions.
Atomic formulas therefore agree because relations are restricted.  The
Boolean cases follow from the induction hypotheses.  If an existential formula
with parameters from \(\mathcal N\) is true in \(\mathcal N\), its witness is
also in \(\Meta\), and the induction hypothesis transfers the matrix.  If it
is true in \(\Meta\), the assumed condition supplies a witness in
\(\mathcal N\), and the induction hypothesis transfers the matrix back.
For a universal formula \(\forall y\,\varphi\) with parameters from
\(\mathcal N\): if it is true in \(\Meta\), then every instance at an
element of \(\abs{\mathcal N}\) is true in \(\Meta\), and the induction
hypothesis for \(\varphi\) transfers each instance to \(\mathcal N\).  If
it is false in \(\Meta\), the satisfaction clauses give
\(\Meta\entails\exists y\,\neg\varphi\) at the same parameters; the
witness condition supplies a witness in \(\abs{\mathcal N}\), the
induction hypothesis for \(\varphi\) transfers the negated instance to
\(\mathcal N\), and \(\forall y\,\varphi\) is false in
\(\mathcal N\).
\end{proof}

\section{The L\"owenheim--Skolem theorems}
\label{sec:fol-cardinality-normalforms-02}

\begin{theorem}[Downward L\"owenheim--Skolem]
\label{thm:downward-ls}
Assume the Axiom of Choice in the ambient set theory.  Let \(\Meta\) be a
\(\Lang\)-structure and let
\(A\subseteq\abs{\Meta}\).  There is
\(\mathcal N\preccurlyeq\Meta\) such that \(A\subseteq\abs{\mathcal N}\) and
\[
 \abs{\abs{\mathcal N}}
 \leq\max(\abs{A},\abs{\Lang},\aleph_0).
\]
\end{theorem}
\begin{proof}
Let \(\mu=\max(\abs{A},\abs{\Lang},\aleph_0)\), and choose
\(m_0\in\abs{\Meta}\).  For each formula
\(\exists y\varphi(y,\bar x)\), where
\(\bar x=(x_1,\ldots,x_k)\) and \(k\in\N\), choose a function
\(F_\varphi:\abs{\Meta}^{k}\to\abs{\Meta}\) such that, whenever
\(\Meta\entails\exists y\varphi(y,\bar a)\),
\(F_\varphi(\bar a)\) is a witness; otherwise let its value be \(m_0\).

Define
\[
 A_0=A\cup\{m_0\}\cup
 \{c^{\Meta}:c\text{ is a constant symbol of }\Lang\}.
\]
Given \(A_n\), define \(A_{n+1}\) to be exactly
\[
\begin{aligned}
 A_n
 &\cup\{f^{\Meta}(\bar a): k\geq1, f\text{ is a }k\text{-ary function
       symbol},\ \bar a\in A_n^k\}\\
 &\cup\{F_\varphi(\bar a): k\in\N,
       \ \exists y\,\varphi(y,x_1,\ldots,x_k)\text{ is a formula},
       \ \bar a\in A_n^k\}.
\end{aligned}
\]
By
\cref{cor:syntax-counting}, there are at most \(\mu\) function
symbols, formulas, and finite tuples at each stage, so
\(\abs{A_n}\leq\mu\) by induction.  Put
\(N=\bigcup_{n\in\N}A_n\).  \Cref{thm:infinite-cardinal-finite-coding} gives
\(\abs{N}\leq\mu\), and \(N\) is nonempty,
contains \(A\), and is closed under every signature function.  Let
\(\mathcal N\) be the induced substructure on \(N\).

If \(\Meta\entails\exists y\varphi(y,\bar a)\) with \(\bar a\in N^k\), all
coordinates of \(\bar a\) lie together in some \(A_n\).  The chosen witness
\(F_\varphi(\bar a)\) lies in \(A_{n+1}\subseteq N\).  The Tarski--Vaught test
therefore gives \(\mathcal N\preccurlyeq\Meta\).
\end{proof}

\begin{theorem}[Upward L\"owenheim--Skolem]
\label{thm:upward-ls}
Assume the Axiom of Choice in the ambient set theory.  Let \(T\) be a
\(\Lang\)-theory with an infinite model.  For every cardinal
\[
 \kappa\geq\max(\abs{\Lang},\aleph_0),
\]
the theory \(T\) has a model of cardinality exactly \(\kappa\).  In
particular, if \(\Meta\) is infinite, then for each such \(\kappa\) there is
a structure \(\mathcal N\) of cardinality \(\kappa\) with
\(\mathcal N\equiv\Meta\), where elementary equivalence means satisfaction of
the same \(\Lang\)-sentences.
\end{theorem}
\begin{proof}
Add constants \(c_\alpha\) for \(\alpha<\kappa\) and add all sentences
\(c_\alpha\ne c_\beta\) for \(\alpha<\beta<\kappa\).  Every finite subset of
the expanded theory has a model: expand a fixed infinite model of \(T\) by
assigning distinct elements to the finitely many constants involved.
Compactness gives a model \(\Meta^*\) of the entire expanded theory, so the
set
\(A=\set{c_\alpha^{\Meta^*}:\alpha<\kappa}\) has cardinality \(\kappa\).

Apply downward L\"owenheim--Skolem in the expanded signature to obtain
\(\mathcal N^*\preccurlyeq\Meta^*\) containing \(A\) and of cardinality at
most \(\kappa\).  Containment of \(A\) gives cardinality at least \(\kappa\),
so \(\abs{\abs{\mathcal N^*}}=\kappa\).  Its reduct to \(\Lang\) is a model
 of \(T\).  For the final assertion take \(T=\Th_{\Lang}(\Meta)\), as
defined in \cref{def:fol-theory-satisfiability}.
\end{proof}

The downward theorem constructs a small elementary substructure of a given
structure.  The upward theorem constructs a new model of the same theory.  It
does not assert that the larger model contains the original structure as an
elementary substructure.

\section{Prenex form and Skolemization}
\label{sec:fol-cardinality-normalforms-03}

\begin{definition}[First-order literal and negation normal form]
A \emph{first-order literal} is an atomic formula or the negation of an
atomic formula.  A first-order formula is in \emph{negation normal form} iff
it is generated from first-order literals, \(\bottom\), and
\(\neg\bottom\) by \(\land,\lor,\forall,\exists\).  Thus every negation in
such a formula occurs only in a first-order literal or immediately before
\(\bottom\).  This extends the earlier propositional definition by adding
the two quantifier constructors and supplies the induction class used in the
prenex theorem.
\end{definition}

\begin{definition}[Prenex formula]
A formula is in \emph{prenex form} iff it has the shape
\[
 Q_1x_1\cdots Q_nx_n\,\theta,
\]
where each \(Q_i\) is \(\forall\) or \(\exists\), and \(\theta\) contains no
quantifier.  The formula is \emph{rectified} iff no variable has both a free
and a bound occurrence and distinct quantifiers bind distinct variables.
\end{definition}

\begin{lemma}[First-order negation normal form]
\label{lem:fol-nnf}
There exist functions
\(\operatorname{nnf},\operatorname{nnf}':\Fm(\Lang)\to\Fm(\Lang)\) such
that, for every formula \(\varphi\):
\begin{enumerate}
  \item \(\operatorname{nnf}(\varphi)\) and \(\operatorname{nnf}'(\varphi)\)
        are in negation normal form;
  \item \(\operatorname{nnf}(\varphi)\equiv\varphi\) and
        \(\operatorname{nnf}'(\varphi)\equiv\neg\varphi\); and
  \item \(\FV\) and \(\operatorname{Var}\) of
        \(\operatorname{nnf}(\varphi)\) and of
        \(\operatorname{nnf}'(\varphi)\) equal \(\FV(\varphi)\) and
        \(\operatorname{Var}(\varphi)\).
\end{enumerate}
\textit{Status: complete proof.}
\end{lemma}
\begin{proof}
Define the two functions by simultaneous recursion.  For \(\alpha\)
atomic and \(Q\in\set{\forall,\exists}\),
\[
\begin{aligned}
 \operatorname{nnf}(\alpha)&=\alpha,
 &\operatorname{nnf}'(\alpha)&=\neg\alpha,\\
 \operatorname{nnf}(\bottom)&=\bottom,
 &\operatorname{nnf}'(\bottom)&=\neg\bottom,\\
 \operatorname{nnf}(\neg\varphi)&=\operatorname{nnf}'(\varphi),
 &\operatorname{nnf}'(\neg\varphi)&=\operatorname{nnf}(\varphi),\\
 \operatorname{nnf}(\varphi\land\psi)&=
   \operatorname{nnf}(\varphi)\land\operatorname{nnf}(\psi),
 &\operatorname{nnf}'(\varphi\land\psi)&=
   \operatorname{nnf}'(\varphi)\lor\operatorname{nnf}'(\psi),\\
 \operatorname{nnf}(\varphi\lor\psi)&=
   \operatorname{nnf}(\varphi)\lor\operatorname{nnf}(\psi),
 &\operatorname{nnf}'(\varphi\lor\psi)&=
   \operatorname{nnf}'(\varphi)\land\operatorname{nnf}'(\psi),\\
 \operatorname{nnf}(\varphi\to\psi)&=
   \operatorname{nnf}'(\varphi)\lor\operatorname{nnf}(\psi),
 &\operatorname{nnf}'(\varphi\to\psi)&=
   \operatorname{nnf}(\varphi)\land\operatorname{nnf}'(\psi),\\
 \operatorname{nnf}(Qx\,\varphi)&=Qx\,\operatorname{nnf}(\varphi),
 &\operatorname{nnf}'(\forall x\,\varphi)&=
   \exists x\,\operatorname{nnf}'(\varphi),\\
 &&\operatorname{nnf}'(\exists x\,\varphi)&=
   \forall x\,\operatorname{nnf}'(\varphi).
\end{aligned}
\]
Every recursive call is on an immediate subformula, so the complexity
\(c\) strictly decreases and the recursion is legitimate.

The three claims follow by simultaneous strong induction on
\(c(\varphi)\).  For claim 1, every clause of \(\operatorname{nnf}\) and
\(\operatorname{nnf}'\) produces a first-order literal, \(\bottom\), or
\(\neg\bottom\)---each admitted by the definition of negation normal
form---or combines values of the two functions by
\(\land,\lor,\forall,\exists\); the induction hypotheses close each case.

For claim 2, check each case against the satisfaction clauses.  The
atomic, \(\bottom\), and negation cases are immediate; the binary cases
for \(\operatorname{nnf}'\) are the recursive truth clauses for
\(\neg(\varphi\land\psi)\), \(\neg(\varphi\lor\psi)\), and
\(\neg(\varphi\to\psi)\), and the \(\to\) case of \(\operatorname{nnf}\)
is the truth clause for \(\to\).  The quantifier cases of
\(\operatorname{nnf}'\) are the pointwise computations
\(\Meta,s\entails\neg\forall x\,\varphi\) iff some \(a\in\abs{\Meta}\) has
\(\Meta,s[x\mapsto a]\entails\neg\varphi\) iff
\(\Meta,s\entails\exists x\,\neg\varphi\), and dually for
\(\exists\)---computed directly from the satisfaction clauses, with the
induction hypothesis applied under each updated assignment.

For claim 3, every clause preserves the multiset of atomic subformulas
and the quantifier occurrences with their variables, so the recursive
\(\FV\) and \(\operatorname{Var}\) clauses compute the same sets on both
sides; the induction hypotheses close each case.
\end{proof}

\begin{lemma}[Rectified variant]
\label{lem:rectified-variant}
For every formula \(\varphi\) and every finite set
\(V\supseteq\operatorname{Var}(\varphi)\) of variables there is a formula
\(r(\varphi,V)\) such that:
\begin{enumerate}
  \item \(r(\varphi,V)\equiv_\alpha\varphi\), hence
        \(r(\varphi,V)\equiv\varphi\);
  \item \(r(\varphi,V)\) is rectified, its binders are pairwise distinct
        and lie outside \(V\), and
        \(\FV(r(\varphi,V))=\FV(\varphi)\);
  \item alpha-equivalent formulas have equal free-variable sets; and
  \item \(r\) preserves the constructor tree; in particular, if
        \(\varphi\) is in negation normal form, then so is
        \(r(\varphi,V)\).
\end{enumerate}
\textit{Status: complete proof.}
\end{lemma}
\begin{proof}
Define \(r\) by recursion: atomic formulas and \(\bottom\) are unchanged;
\(r(\neg\varphi,V)=\neg r(\varphi,V)\);
\(r(\varphi\circ\psi,V)=r(\varphi,V)\circ
r(\psi,\,V\cup\operatorname{Var}(r(\varphi,V)))\) for
\(\circ\in\set{\land,\lor,\to}\); and
\(r(Qx\,\varphi,V)=Qz\,r(\varphi[z/x],\,V\cup\set{z})\), where \(z\) is
the least variable outside \(V\).  The recursion is well founded by the
pattern of the substitution definition: since
\(\operatorname{Var}(\varphi)\subseteq V\) and \(z\notin V\), we have
\(z\notin\operatorname{Var}(\varphi)\cup\set{x}\), so no binder of
\(\varphi\) equals \(z\), the substitution \(\varphi[z/x]\) runs on the
first two quantifier clauses only and is textual replacement---whence
\(\operatorname{Var}(\varphi[z/x])\subseteq
\operatorname{Var}(\varphi)\cup\set{z}\), so the invariant
\(\operatorname{Var}\subseteq V\) persists into every recursive
call---and \(c(\varphi[z/x])=c(\varphi)<c(Qx\,\varphi)\) by
\cref{lem:substitution-complexity}.

Prove claim 3 first.  The relation ``has the same free-variable set''
is an equivalence relation and a congruence for every formula
constructor, by the recursive \(\FV\) clauses.  It contains each
generating pair \(Qx\,\varphi\equiv_\alpha Qy\,(\varphi[y/x])\) with
\(y\notin\operatorname{Var}(\varphi)\): both cases of
\cref{lem:substitution-fv}(2) give
\(\FV(Qy\,(\varphi[y/x]))=\FV(Qx\,\varphi)\), because when
\(x\in\FV(\varphi)\) the sets
\((\FV(\varphi)\setminus\set{x})\cup\set{y}\) minus \(\set{y}\) and
\(\FV(\varphi)\setminus\set{x}\) coincide (\(y\notin\FV(\varphi)\)), and
when \(x\notin\FV(\varphi)\) both sides equal \(\FV(\varphi)\).  By
leastness of \(\equiv_\alpha\), alpha-equivalence is contained in the
relation.

Claims 1, 2, and 4 follow by strong induction on \(c(\varphi)\) under the
invariant \(\operatorname{Var}(\varphi)\subseteq V\).  For claim 1, the
atomic and \(\bottom\) cases are reflexivity, and the connective cases
are the congruence property of \(\equiv_\alpha\).  In the quantifier
case, the induction hypothesis applies to \(\varphi[z/x]\) at complexity
\(c(\varphi[z/x])=c(\varphi)<c(Qx\,\varphi)\)---licensed by
\cref{lem:substitution-complexity}---and gives
\(r(\varphi[z/x],V\cup\set{z})\equiv_\alpha\varphi[z/x]\); the quantifier
congruence gives
\(Qz\,r(\varphi[z/x],V\cup\set{z})\equiv_\alpha Qz\,(\varphi[z/x])\); the
generating clause \(Qx\,\varphi\equiv_\alpha Qz\,(\varphi[z/x])\) is
licensed by \(z\notin\operatorname{Var}(\varphi)\); transitivity and
symmetry conclude.  Then \(r(\varphi,V)\equiv\varphi\) by
\cref{cor:alpha-invariance}.

For claim 2, the binders of \(r(\varphi,V)\) lie outside \(V\) and are
pairwise distinct: the quantifier clause chooses \(z\notin V\) and
recurses with \(V\cup\set{z}\), so all inner binders avoid \(z\) and
\(V\), and the binary clause forbids every variable of the left output to
the right recursion, so the right binders avoid the left binders.  By
claims 1 and 3, \(\FV(r(\varphi,V))=\FV(\varphi)\subseteq
\operatorname{Var}(\varphi)\subseteq V\), so no binder is also free and
\(r(\varphi,V)\) is rectified.

For claim 4, every clause of \(r\) preserves the outer constructor, and
in the quantifier case the textual replacement \(\varphi[z/x]\) changes
only variables inside atomic formulas; induction on the clauses shows
that the constructor tree, and hence membership in the
negation-normal-form class, is preserved.
\end{proof}

\begin{lemma}[Prenex extraction]
\label{lem:prenex-extraction}
For every rectified formula \(\varphi\) in negation normal form there is
a formula \(p(\varphi)\) such that:
\begin{enumerate}
  \item \(p(\varphi)\equiv\varphi\);
  \item \(p(\varphi)\) is rectified and in prenex form, and its matrix is
        quantifier-free and in negation normal form; and
  \item \emph{(invariant)} \(\FV\) and \(\operatorname{Var}\) are
        unchanged, and the prefix of \(p(\varphi)\) consists of exactly
        the quantifier occurrences of \(\varphi\), with the same
        quantifier and the same variable; no new free or bound variables
        occur.
\end{enumerate}
\textit{Status: complete proof.}
\end{lemma}
\begin{proof}
Record two facts.  (C1) If \(y\notin\FV(\chi)\), then
\((Qy\,\pi)\circ\chi\equiv Qy\,(\pi\circ\chi)\) and
\(\chi\circ(Qy\,\pi)\equiv Qy\,(\chi\circ\pi)\) for
\(Q\in\set{\forall,\exists}\) and \(\circ\in\set{\land,\lor}\).  For the
\(\forall\)/\(\lor\) case: if
\(\Meta,s\entails(\forall y\,\pi)\lor\chi\), then either
\(\Meta,s[y\mapsto a]\entails\pi\) for every \(a\), or
\(\Meta,s\entails\chi\) and, by \cref{lem:fol-coincidence},
\(\Meta,s[y\mapsto a]\entails\chi\) for every \(a\); in both cases
\(\Meta,s\entails\forall y\,(\pi\lor\chi)\).  Conversely, if
\(\Meta,s\entails\forall y\,(\pi\lor\chi)\) and
\(\Meta,s\notentails\chi\), then \cref{lem:fol-coincidence} gives
\(\Meta,s[y\mapsto a]\notentails\chi\) for every \(a\), so every
\(s[y\mapsto a]\) satisfies \(\pi\), and
\(\Meta,s\entails\forall y\,\pi\).  The remaining seven cases follow by
the same computation from the satisfaction clauses and
\cref{lem:fol-coincidence}.  (C2) \(\equiv\) is a congruence for
\(\land,\lor\) and for each quantifier, immediately from the
satisfaction clauses.

Define \(p\) by structural induction on the negation-normal generation
of \(\varphi\); every subformula of a rectified formula is rectified, so
the induction hypotheses apply.  On first-order literals, \(\bottom\),
and \(\neg\bottom\), let \(p\) be the identity.  For a quantifier case
\(Qx\,\alpha\), put \(p(Qx\,\alpha)=Qx\,p(\alpha)\): claim 1 is (C2) with
the induction hypothesis, and claims 2 and 3 hold because the prefix of
\(p(Qx\,\alpha)\) consists of \(Qx\) followed by the prefix of
\(p(\alpha)\), whose entries are exactly the quantifier occurrences of
\(\alpha\), pairwise distinct from \(x\) by rectifiedness.

For a binary case \(\alpha\circ\beta\) with
\(\circ\in\set{\land,\lor}\), let
\(p(\alpha)=Q_1y_1\cdots Q_ky_k\,\alpha_0\) and
\(p(\beta)=R_1z_1\cdots R_mz_m\,\beta_0\) with quantifier-free
\(\alpha_0,\beta_0\), and put
\[
 p(\alpha\circ\beta)=
 Q_1y_1\cdots Q_ky_k\,R_1z_1\cdots R_mz_m\,(\alpha_0\circ\beta_0).
\]
The required side conditions hold by rectifiedness of
\(\alpha\circ\beta\) and the invariant: each \(y_i\) is bound in
\(\alpha\), so \(y_i\notin\FV(\beta)=\FV(p(\beta))\); and each
\(z_j\) is bound in \(\beta\), so
\(z_j\notin\FV(\alpha_0)\subseteq\FV(\alpha)\cup\set{y_1,\ldots,y_k}\),
using binder distinctness for \(z_j\ne y_i\).  An inner induction on
\(k+m\) then proves
\(p(\alpha\circ\beta)\equiv p(\alpha)\circ p(\beta)\): while a
quantifier heads either side of \(\circ\), (C1) moves it outward across
\(\circ\)---its variable is not free in the other side by the side
conditions just checked---and (C2) transports the equivalence under the
already-extracted prefix.  With the induction hypotheses
\(p(\alpha)\equiv\alpha\) and \(p(\beta)\equiv\beta\) and (C2), claim 1
follows.  The output is prenex with quantifier-free negation-normal
matrix \(\alpha_0\circ\beta_0\); it is rectified and satisfies the
invariant because nothing is renamed or introduced: its prefix lists
exactly the quantifier occurrences of \(\alpha\) followed by those of
\(\beta\), which are pairwise distinct and disjoint from the free
variables by rectifiedness of \(\alpha\circ\beta\).  The base,
quantifier, and binary cases exhaust the constructors of a first-order
formula in negation normal form.
\end{proof}

\begin{theorem}[Prenex normal form]
\label{thm:prenex-normal-form}
Every first-order formula is logically equivalent to a rectified formula in
prenex form in the same signature.
\textit{Status: complete proof.}
\end{theorem}
\begin{proof}
Let \(\varphi\in\Fm(\Lang)\).  Put
\(\psi=r\bigl(\operatorname{nnf}(\varphi),
\operatorname{Var}(\operatorname{nnf}(\varphi))\bigr)\) and
\(\chi=p(\psi)\), with the three operations supplied by
\cref{lem:fol-nnf,lem:rectified-variant,lem:prenex-extraction}.  By
\cref{lem:fol-nnf}, \(\operatorname{nnf}(\varphi)\equiv\varphi\) and
\(\operatorname{nnf}(\varphi)\) is in negation normal form.  By
\cref{lem:rectified-variant}, \(\psi\) is rectified, remains in negation
normal form, and satisfies \(\psi\equiv\operatorname{nnf}(\varphi)\).  By
\cref{lem:prenex-extraction}, \(\chi\) is rectified and in prenex form
and satisfies \(\chi\equiv\psi\).  Hence
\(\chi\equiv\varphi\), and all three operations stay inside
\(\Fm(\Lang)\), so \(\chi\) is in the same signature.
\end{proof}

\begin{definition}[Skolemization]
Let
\[
 \sigma=Q_1x_1\cdots Q_nx_n\,\theta
\]
be a rectified prenex sentence.  Traverse its quantifier prefix from left to
right.  At an existential quantifier \(\exists x_i\), let
\(u_1,\ldots,u_m\) be the variables governed by preceding universal
quantifiers, introduce a new \(m\)-ary function symbol \(f_i\), replace the
occurrences bound by \(\exists x_i\) with
\(f_i(u_1,\ldots,u_m)\), and delete that existential quantifier.  Retain every
universal quantifier.  The resulting universal sentence in the expanded
signature is the \emph{Skolemization} \(\operatorname{Sk}(\sigma)\).
\end{definition}

An existential quantifier with no preceding universal quantifier produces a
new constant symbol.  A different new function symbol is used for each
existential-quantifier occurrence.

\begin{example}
Skolemization transforms
\[
 \forall x\exists y\forall z\exists w\,R(x,y,z,w)
\]
into
\[
 \forall x\forall z\,R(x,f(x),z,g(x,z)).
\]
The second witness may depend on \(x\) and \(z\); its indirect dependence on
the first witness is already determined by \(f(x)\).
\end{example}

\begin{theorem}[Equisatisfiability of Skolemization]
\label{thm:skolem-equisatisfiable}
Let \(\sigma\) be a rectified prenex \(\Lang\)-sentence.
\begin{enumerate}
  \item \emph{(Axiom of Choice.)}  Every \(\Lang\)-structure satisfying
        \(\sigma\) has an expansion to the Skolem signature satisfying
        \(\operatorname{Sk}(\sigma)\).
  \item \emph{(No choice.)}  The \(\Lang\)-reduct of every model of
        \(\operatorname{Sk}(\sigma)\) satisfies \(\sigma\).
\end{enumerate}
Hence, assuming the Axiom of Choice in the ambient set theory, \(\sigma\)
has a \(\Lang\)-model iff \(\operatorname{Sk}(\sigma)\) has a model in its
expanded signature; the right-to-left implication holds without choice.
\textit{Status: complete proof.}
\end{theorem}
\begin{proof}
Write \(\sigma=Q_1x_1\cdots Q_nx_n\,\theta\) with \(\theta\)
quantifier-free.  For \(0\le i\le n\), let
\(\sigma_i=Q_{i+1}x_{i+1}\cdots Q_nx_n\,\theta\), so that
\(\sigma_0=\sigma\), \(\sigma_n=\theta\), and
\(\FV(\sigma_i)\subseteq\set{x_1,\ldots,x_i}\).  Let \(q(i)\) be the
number of universal positions among the first \(i\), let \(\bar u(j)\)
list the variables of the universal positions before position \(j\), and,
for each existential position \(j\), let \(t_j=f_j(\bar u(j))\) be its
Skolem term.  Then
\(\operatorname{Sk}(\sigma)=\forall u_1\cdots\forall u_p\,\theta^{*}\),
where \(u_1,\ldots,u_p\) are the universal variables in prefix order,
\(p=q(n)\), and \(\theta^{*}=\theta[t_j/x_j]\) simultaneously for all
existential positions \(j\).  Because \(\theta\) is quantifier-free, each
substitution is plain replacement, and because no \(t_j\) contains an
existential variable, the simultaneous substitution equals the
composition of the single substitutions in any order.

Given a structure and a tuple \(\bar a\) of values for the universal
variables among \(x_1,\ldots,x_i\), the \emph{Skolem-term assignment}
\(s_{\bar a}\) assigns to a universal \(x_j\) its entry of \(\bar a\) and
to an existential \(x_j\) the value \(v_j\) of \(t_j\) computed from the
entries for \(\bar u(j)\); the values of all other variables are
irrelevant by \cref{lem:fol-coincidence}.

\emph{Clause 1.}  Let \(\Meta\entails\sigma\).  By finite recursion on
\(i\), choose functions
\(F_j:\abs{\Meta}^{\,q(j-1)}\to\abs{\Meta}\) for the existential
positions \(j\le i\) while proving the invariant
\[
 \mathrm{Inv}(i):\quad
 \Meta,s_{\bar a}\entails\sigma_i
 \quad\text{for every }\bar a\in\abs{\Meta}^{\,q(i)},
\]
where \(s_{\bar a}\) computes each \(v_j\) as
\(F_j\) applied to the entries of \(\bar a\) for \(\bar u(j)\).
\(\mathrm{Inv}(0)\) is the hypothesis \(\Meta\entails\sigma\).

Universal step (\(Q_{i+1}=\forall\)): from
\(\Meta,s_{\bar a}\entails\forall x_{i+1}\,\sigma_{i+1}\), the universal
satisfaction clause gives
\(\Meta,s_{\bar a}[x_{i+1}\mapsto b]\entails\sigma_{i+1}\) for every
\(b\in\abs{\Meta}\), which is \(\mathrm{Inv}(i+1)\) for the tuple
\((\bar a,b)\).

Existential step (\(Q_{i+1}=\exists\)): by \(\mathrm{Inv}(i)\) and the
existential satisfaction clause, each witness set
\[
 W(\bar a)=\set{b\in\abs{\Meta}:
 \Meta,s_{\bar a}[x_{i+1}\mapsto b]\entails\sigma_{i+1}}
\]
is nonempty.  Apply the Axiom of Choice to the set
\(\set{\set{\bar a}\times W(\bar a):\bar a\in\abs{\Meta}^{\,q(i)}}\) of
nonempty sets to obtain a function \(F_{i+1}\) with
\(F_{i+1}(\bar a)\in W(\bar a)\) for every \(\bar a\).  This is one
application of choice per existential quantifier, and choice is used
nowhere else.  With \(v_{i+1}=F_{i+1}(\bar a)\), the Skolem-term
assignment satisfies \(\mathrm{Inv}(i+1)\).

Define the expansion \(\Meta^{*}\) of \(\Meta\) to the Skolem signature
by \(f_j^{\Meta^{*}}=F_j\) for every existential position \(j\).  Fix
\(\bar a\in\abs{\Meta}^{p}\).  By \(\mathrm{Inv}(n)\),
\(\Meta,s_{\bar a}\entails\theta\), hence
\(\Meta^{*},s_{\bar a}\entails\theta\) by \cref{lem:reduct-invariance}.
The term-value clauses give each \(t_j\) the value \(v_j\) under
\([\bar u\mapsto\bar a]\) in \(\Meta^{*}\), so iterating
\cref{thm:fol-substitution} once per existential position turns
\(\Meta^{*},s_{\bar a}\entails\theta\) into
\(\Meta^{*},[\bar u\mapsto\bar a]\entails\theta^{*}\).  Since \(\bar a\)
was arbitrary, the \(p\)-fold universal satisfaction clause yields
\(\Meta^{*}\entails\operatorname{Sk}(\sigma)\).

\emph{Clause 2.}  Let \(\mathcal N\entails\operatorname{Sk}(\sigma)\) in
the expanded signature; now \(v_j\) is the value of the term \(t_j\) in
\(\mathcal N\) at the entries for \(\bar u(j)\), and no choice is needed
to produce it.  By downward induction from \(i=n\) to \(i=0\), prove
\[
 \mathrm{Inv}'(i):\quad
 \mathcal N,s_{\bar a}\entails\sigma_i
 \quad\text{for every }\bar a\in\abs{\mathcal N}^{\,q(i)}.
\]
Base \(i=n\): instantiating the universal prefix of
\(\operatorname{Sk}(\sigma)\) at \(\bar a\) gives
\(\mathcal N,[\bar u\mapsto\bar a]\entails\theta^{*}\), and reversing the
iterated substitution-lemma computation of clause 1, again by
\cref{thm:fol-substitution}, gives
\(\mathcal N,s_{\bar a}\entails\theta\).  Universal step
(\(Q_{i+1}=\forall\)): for every \(b\), \(\mathrm{Inv}'(i+1)\) at
\((\bar a,b)\) gives
\(\mathcal N,s_{\bar a}[x_{i+1}\mapsto b]\entails\sigma_{i+1}\)---the
values \(v_j\) for the existential \(j\le i\) depend only on the entries
of \(\bar a\)---and the universal satisfaction clause gives
\(\mathrm{Inv}'(i)\).  Existential step (\(Q_{i+1}=\exists\)):
\(q(i+1)=q(i)\), and \(\mathrm{Inv}'(i+1)\) at \(\bar a\) exhibits the
witness \(v_{i+1}\), so the existential satisfaction clause gives
\(\mathrm{Inv}'(i)\).  Then \(\mathrm{Inv}'(0)\) says
\(\mathcal N\entails\sigma\) in the expanded signature, and
\cref{lem:reduct-invariance} yields
\(\mathcal N\restr\Lang\entails\sigma\).  No choice principle is used.
\end{proof}

Skolemization generally does not preserve logical equivalence: its output is
in an expanded signature and fixes witness functions that the original
sentence did not name.  The theorem states equisatisfiability and the precise
expansion--reduct relation, not equivalence of formulas.

\section{First-order resolution}
\label{sec:fol-cardinality-normalforms-04}

\begin{definition}[First-order clause and formula interpretation]
Using the first-order literal defined before the prenex theorem, a
\emph{clause} is a finite set of first-order literals, with
set equality as its equality relation.  Thus literal order and multiplicity
are not clause data.

For an exact map back to formulas, take as presentation data a well-order
\(\prec_{\Lang}\) of the symbol tokens of \(\Lang\).  Supplying this order is
an explicit well-ordering assumption; for an arbitrary signature,
its existence follows from the well-ordering theorem and hence from the Axiom
of Choice.  The semantic results below do not depend on which supplied
well-order is used.  Combine \(\prec_{\Lang}\) with the order
\(x_0,x_1,\ldots\) and a fixed order of
the finitely many logical and punctuation tokens, and order formula strings
first by length and then lexicographically.  For a clause \(C\), let
\(k=\abs{C}\); when \(k>0\), list its members as
\(L_1\prec\cdots\prec L_k\) in the induced order and define, by downward
recursion,
\[
 D_k\defeq L_k,
 \qquad
 D_i\defeq(L_i\lor D_{i+1})\quad(1\leq i<k).
\]
Then define
\[
 \operatorname{form}_{\Lang}(C)\defeq
 \begin{cases}
  \bottom,&C=\varnothing,\\
  D_1,&C\ne\varnothing.
 \end{cases}
\]
Let \(x_{i_1},\ldots,x_{i_m}\), with
\(i_1<\cdots<i_m\), be exactly the variables occurring in \(C\).  Its
formula interpretation is
\[
 \operatorname{ucl}_{\Lang}(C)\defeq
 \forall x_{i_1}\cdots\forall x_{i_m}\,
 \operatorname{form}_{\Lang}(C),
\]
with the quantifier prefix empty when \(m=0\).  The functions
\(\operatorname{form}_{\Lang}\) and \(\operatorname{ucl}_{\Lang}\) are
metalinguistic, not new object-language constructors.  They are required
because a clause is a set rather than a formula;
\(\operatorname{ucl}_{\Lang}\) differs from \(\operatorname{form}_{\Lang}\)
by closing every clause variable.  A structure satisfies \(C\) iff it
satisfies \(\operatorname{ucl}_{\Lang}(C)\), and it satisfies a clause set
iff it satisfies every member.

A clause is \emph{equality-free} iff no literal in it contains \(=\); that is,
the atom underlying each of its literals has the form
\(R(t_1,\ldots,t_n)\), not \(r=t\).  A clause set is equality-free iff each
of its clauses is equality-free.  This condition concerns the literals in
the clauses, not the signature, because logical equality is not a signature
member.
\end{definition}

\begin{definition}[Finite-support simultaneous term substitution]
A \emph{term substitution} is a total function
\(\theta:\Var\to\Term(\Lang)\) whose support
\[
 \operatorname{supp}(\theta)
 \defeq\set{x\in\Var:\theta(x)\ne x}
\]
is finite.  A finite display such as
\(\set{x_1\mapsto t_1,\ldots,x_n\mapsto t_n}\) denotes the substitution
shown on those variables and the identity on every other variable.

Its simultaneous homomorphic action is defined by
\[
 x\theta=\theta(x),
 \qquad
 f(t_1,\ldots,t_n)\theta
 =f(t_1\theta,\ldots,t_n\theta),
\]
\[
 \begin{aligned}
 R(t_1,\ldots,t_n)\theta
   &\defeq R(t_1\theta,\ldots,t_n\theta),\\
 (r=t)\theta&\defeq(r\theta=t\theta),\\
 (\neg A)\theta&\defeq\neg(A\theta),\\
 C\theta&\defeq\set{L\theta:L\in C}.
 \end{aligned}
\]
Here \(A\) is atomic, \(L\) is a literal, and \(C\) is a clause.  The
application is simultaneous: when \(\theta(x)\) is inserted for \(x\), the
same application does not apply \(\theta\) again inside that inserted term.
The clause action may identify two literals, in which case set equality
removes the duplicate.  This quantifier-free action is the one needed by
resolution; it neither replaces nor restricts the total capture-avoiding
single-variable formula substitution \(\varphi[t/x]\) defined earlier.
\end{definition}

\begin{definition}[Composition and instance preorder]
The identity substitution \(\mathrm{id}\) maps every variable to itself.  For
substitutions \(\theta\) and \(\eta\), their composition \(\theta\eta\) is
oriented by
\[
 (\theta\eta)(x)\defeq(\theta(x))\eta.
\]
Its support is contained in
\(\operatorname{supp}(\theta)\cup\operatorname{supp}(\eta)\), so it is
again a finite-support term substitution.
Thus \(t(\theta\eta)=(t\theta)\eta\), and the same equation holds for
atoms, literals, and clauses.  Define the \emph{instance preorder} on term
substitutions by
\[
 \theta\preceq_{\mathrm{inst}}\rho
 \quad\Longleftrightarrow\quad
 \text{there is a term substitution \(\eta\) such that }
 \rho=\theta\eta.
\]
This is a relation among substitutions, not equality: it records that
\(\rho\) results from further simultaneous instantiation of \(\theta\).
It is reflexive by \(\mathrm{id}\) and transitive by associativity of the
displayed composition, so it is a preorder.  Its use below is precisely to
state what ``most general'' means.
\end{definition}

\begin{lemma}[Composition of simultaneous substitutions]
For every term \(t\), atom \(A\), literal \(L\), and clause \(C\),
\[
 \begin{gathered}
 t(\theta\eta)=(t\theta)\eta,
 \qquad A(\theta\eta)=(A\theta)\eta,\\
 L(\theta\eta)=(L\theta)\eta,
 \qquad C(\theta\eta)=(C\theta)\eta.
 \end{gathered}
\]
\end{lemma}
\begin{proof}
Induct on terms for the first equation.  The atomic and literal equations then
follow from their displayed homomorphic clauses.  Taking the pointwise image
of a finite set proves the clause equation.  The term equation also proves
associativity of substitution composition by extensionality.
\end{proof}

\begin{lemma}[Satisfaction under simultaneous term substitution]
\label{lem:clause-substitution}
If a \(\Lang\)-structure \(\Meta\) satisfies a clause \(C\), then it
satisfies \(C\theta\) for every term substitution \(\theta\).
\end{lemma}
\begin{proof}
For an assignment \(s\), define another assignment by
\[
 s^{\theta}(x)\defeq\sem{\theta(x)}^{\Meta}_{s}.
\]
Induction on terms gives
\(\sem{t\theta}^{\Meta}_{s}=\sem{t}^{\Meta}_{s^{\theta}}\), and the atomic
truth clauses extend this equality to literals.  Because \(\Meta\) satisfies
the universal closure of \(C\), the assignment \(s^{\theta}\) satisfies some
literal \(L\in C\).  Hence \(s\) satisfies \(L\theta\in C\theta\).  Since
\(s\) was arbitrary, \(\Meta\) satisfies the universal closure of
\(C\theta\).
\end{proof}

\begin{definition}[Renaming apart]
A \emph{variable renaming} is a finite-support term substitution induced by a
bijection \(\pi:\Var\to\Var\); hence every \(\pi(x)\) is a variable.  For a
clause \(C\), put
\[
 \operatorname{Var}(C)\defeq
 \bigcup_{L\in C}\operatorname{Var}(L).
\]
Clauses \(C\pi\) and \(D\rho\) are \emph{renamed apart} iff
\(\operatorname{Var}(C\pi)\cap\operatorname{Var}(D\rho)=\varnothing\).
Before a resolution inference, each parent is replaced by such a renamed
copy.  The required renamings exist because both variable sets are finite and
\(\Var\) is infinite.  Universal closure makes a renamed copy equivalent to
the original clause; the operation changes bound variable names in its
formula interpretation, not its satisfaction conditions.
\end{definition}

\begin{definition}[Unifier and most general unifier]
A term substitution \(\theta\) is a \emph{unifier} of atomic formulas
\(A\) and \(B\) iff \(A\theta\) and \(B\theta\) are literally identical.
A unifier \(\theta\) is a \emph{most general unifier} iff
\(\theta\preceq_{\mathrm{inst}}\rho\) for every unifier \(\rho\); equivalently, every
unifier \(\rho\) has the form \(\theta\eta\) for some term substitution
\(\eta\).
\end{definition}

\begin{definition}[First-order unification procedure]
Below, \emph{atom} abbreviates \emph{atomic formula}.  To unify atoms
\(A\) and \(B\), first fail if their outer atomic symbols or
arities differ.  Otherwise form a finite set \(E\) of metalinguistic term
equations \(r\doteq t\) by pairing corresponding arguments (the two ordered
arguments when the common atomic symbol is \(=\)), and initialize
\(\theta\) to \(\mathrm{id}\).  Repeatedly apply the following transformations
to a selected equation:
\begin{enumerate}
  \item delete \(t\doteq t\);
  \item replace
        \(f(r_1,\ldots,r_n)\doteq f(t_1,\ldots,t_n)\) by the equations
        \(r_i\doteq t_i\) for \(1\leq i\leq n\);
  \item fail on
        \(f(r_1,\ldots,r_m)\doteq g(t_1,\ldots,t_n)\) when
        \(f\ne g\) or \(m\ne n\);
  \item orient \(t\doteq x\) as \(x\doteq t\) when \(x\) is a variable
        and \(t\) is not;
  \item for \(x\doteq t\) with \(x\ne t\), perform the \emph{occurs
        check}: if \(x\in\operatorname{Var}(t)\), fail; otherwise put
        \(\delta=\set{x\mapsto t}\), apply \(\delta\) simultaneously to
        every remaining equation, and replace the accumulator by
        \(\theta\delta\).
\end{enumerate}
If the equation set becomes empty, return \(\theta\).  The occurs check is
part of this procedure rather than part of the definition of a term
substitution: it rejects the impossible demand that a finite term equal a
proper term containing itself.
\end{definition}

\begin{theorem}[First-order unification]
The unification procedure terminates.  It returns failure iff its
input atoms are not unifiable; otherwise it returns a most general unifier.
\end{theorem}
\begin{proof}[Proof sketch]
Maintain the invariant that the unifiers of the original problem are exactly
the substitutions \(\theta\eta\) for substitutions \(\eta\) solving every
equation still in \(E\).  Deletion, decomposition, and orientation preserve
this invariant directly.  Elimination preserves it by the composition lemma.
The clash rule has no solution.  If \(x\) occurs properly in \(t\), no finite
term can satisfy \(x\eta=t\eta\), which justifies occurs-check failure.

An elimination removes its variable from the remaining problem; a
decomposition removes matching outer function symbols; deletion removes an
equation, and orientation is immediately followed by one of these steps.
The corresponding lexicographic measure therefore decreases, proving
termination.  At success the remaining equation set is empty, so the
invariant says both that \(\theta\) unifies the original atoms and that every
other unifier is an instance \(\theta\eta\).  At either failure rule the
invariant says that no unifier exists.  The full lexicographic-measure check
and the transformation-by-transformation invariant calculations are omitted.
\end{proof}

After the two parent clauses \(E\) and \(F\) have been renamed apart, choose
\(A\in E\) and \(\neg B\in F\), and let \(\theta\) be a most general
unifier of the atoms \(A\) and \(B\).  The first-order resolution rule is
\[
 \Rule[Resolution]{E\qquad F}
 {\bigl((E\setminus\set{A})\cup(F\setminus\set{\neg B})\bigr)\theta}.
\]
Factoring applies to two distinct literals of the same polarity in one
clause \(E\).  If their underlying atoms are \(A\) and \(B\), and
\(\theta\) is a most general unifier of those atoms, its rule is
\[
 \Rule[Factoring]{E}{E\theta}.
\]
Because clauses are sets, the two selected literals become one member of
\(E\theta\).  A \emph{first-order resolution derivation} from a clause set
\(S\) is a finite sequence in which each clause belongs to \(S\) or follows
by resolution or factoring from earlier clauses, using renamed-apart copies
of two parents before every resolution step.

\begin{theorem}[Soundness of first-order resolution and factoring]
Every structure satisfying the parent clauses of a resolution or factoring
inference satisfies its conclusion.  Consequently, if a first-order
resolution derivation from \(S\) contains the empty clause, then \(S\) is not
satisfiable.
\end{theorem}
\begin{proof}
Renaming apart preserves the universally closed formula interpretations: one
direction follows from \cref{lem:clause-substitution}, and the other follows
by applying the inverse variable renaming.  By the same lemma, a structure
satisfying \(E\) and \(F\) satisfies \(E\theta\) and \(F\theta\).  If an
assignment falsified every literal in the displayed resolvent, satisfaction
of the substituted first parent would force \(A\theta\), while satisfaction
of the substituted second parent would force \(\neg(B\theta)\).  But
\(A\theta=B\theta\) literally, a contradiction.  Factoring is sound directly
by \cref{lem:clause-substitution}.  Induction on a derivation now preserves
satisfaction, whereas the formula interpretation of the empty clause is
\(\bottom\).
\end{proof}

\begin{definition}[Ground term and ground instance]
A term \(t\) is \emph{ground} iff \(\operatorname{Var}(t)=\varnothing\).  Let
\(C\) be a first-order clause.  A \emph{ground instance} of \(C\) is a clause
\(C\theta\) where \(\theta\) maps every variable occurring in \(C\) to a
ground term.  These terms classify terms and substitution instances,
respectively; neither means a closed formula.  They are required in the two
cited results and the resolution-completeness proof below.
\end{definition}

\begin{theorem}[Herbrand theorem for clauses]
\label{thm:herbrand-clauses}
Let \(S\) be a finite equality-free set of first-order clauses in a signature
having at least one ground term.  Then \(S\) is not satisfiable iff some finite
set of ground instances of members of \(S\), regarded as propositional
clauses whose distinct ground atoms are propositional variables, is not
satisfiable.
\textit{Status: cited theorem.}
\end{theorem}

\begin{theorem}[Lifting lemma]
\label{thm:resolution-lifting}
Let \(C,D\) be first-order clauses with disjoint variable sets, and let \(G\)
be the ground resolvent produced by one ground resolution step from ground
instances of \(C\) and \(D\).  Factor each parent literal family that its
ground substitution identifies, and let \(R\) be the first-order resolvent of
the resulting parents obtained with a most general unifier.  Then there is a
substitution \(\eta\) that maps every variable occurring in \(R\) to a ground
term and satisfies \(R\eta=G\).  The analogous equality holds for a ground
factoring step and its lifted first-order factoring step.  Consequently, for
every clause set \(S\), every
 finite ground resolution refutation of ground instances of \(S\) lifts, step by
step, to a first-order resolution derivation with factoring from \(S\).
\textit{Status: cited theorem.}
\end{theorem}

See \cite{Robinson1965,BachmairGanzinger2001} for complete proofs of
\cref{thm:herbrand-clauses,thm:resolution-lifting}.

\begin{example}[Unification and resolution]
Resolve
\[
 \set{P(f(x)),Q(x)}
 \quad\text{and}\quad
 \set{\neg P(y),R(y)}.
\]
The substitution \(\theta=\set{y\mapsto f(x)}\) is a most general unifier of
\(P(f(x))\) and \(P(y)\).  The resolvent is
\[
 \set{Q(x),R(f(x))}.
\]
\end{example}

\begin{theorem}[Refutational completeness of first-order resolution]
\label{thm:fol-resolution}
Let \(S\) be a finite equality-free set of first-order clauses.  After adding
a constant if the signature has no ground term, \(S\) is not satisfiable iff
resolution with factoring derives the empty clause.
\end{theorem}
\begin{proof}[Proof sketch]
The preceding soundness theorem proves the reverse implication.  For the
forward implication, \cref{thm:herbrand-clauses} yields a finite propositionally
unsatisfiable set of ground instances of clauses in \(S\).  Propositional
refutational completeness gives a ground resolution refutation.
\Cref{thm:resolution-lifting} replaces each ground resolution step by a first-order resolution or
factoring step whose instance is that ground step.  Induction through the
finite ground refutation yields a first-order derivation of the empty clause.
This proof depends on the two cited theorems, the existence of most general
unifiers for unifiable first-order terms, and propositional refutational
completeness; it is the
precise equality-free scope of the classical result in
\cite{Robinson1965,BachmairGanzinger2001}.  The cited complete proofs and the
step-by-step lifted derivation are not reproduced here.
\end{proof}

Bare resolution is not complete for logical equality.  One must add
an equality-complete treatment, such as equality axioms together with
resolution or a sound and complete paramodulation or superposition calculus.
The theorem above deliberately excludes every clause having a literal that
contains logical equality.

\section{Common errors}
\label{sec:fol-cardinality-normalforms-05}

\begin{longtable}{@{}p{0.25\textwidth}p{0.32\textwidth}p{0.35\textwidth}@{}}
\toprule
Error & Invalid move & Required correction \\ \midrule
\endhead
Capture & Replace \(x\) by \(y\) in \(\forall y\,R(x,y)\) textually. &
Rename the binder, then perform capture-avoiding substitution. \\
\addlinespace
Structure/model confusion & Call every \(\Lang\)-structure a model without
naming a theory. & Use ``structure'' until a set of sentences has been stated
and satisfied. \\
\addlinespace
Formula/sentence confusion & Write \(\Meta\entails\varphi(x)\) without an
assignment or parameter convention. & Supply an assignment, or close the
formula. \\
\addlinespace
Witness dependence & Reverse \(\forall x\exists y\) to
\(\exists y\forall x\). & Preserve quantifier order or prove a uniform witness
exists. \\
\addlinespace
Skolem equivalence & Claim that a sentence and its Skolemization are
equivalent in one signature. & State equisatisfiability and the
expansion--reduct directions. \\
\addlinespace
Cardinality omission & State a L\"owenheim--Skolem theorem without a language
size or infinite-model hypothesis. & State the exact cardinal bound and the
direction of the theorem. \\
\addlinespace
Resolution with equality & Apply bare resolution and claim completeness for
logical equality. & Add an equality-complete calculus or restrict the theorem.
\\
\bottomrule
\end{longtable}

\section{Exercises}
\label{sec:fol-cardinality-normalforms-06}

\begin{exercise}
\label{ex:fol-cardinality-normalforms-01}
Prove the Tarski--Vaught test in full detail, handling both quantifiers
directly in the structural induction and justifying every appeal to the
induction hypothesis.
\end{exercise}

\begin{exercise}
\label{ex:fol-cardinality-normalforms-02}
Let \(\Meta\) be a structure in a countable signature and let
\(A\subseteq\abs{\Meta}\) be countable.  Write out the first three stages
\(A_0,A_1,A_2\) of the downward L\"owenheim--Skolem construction and prove
that the union \(\bigcup_{n\in\N}A_n\) of all stages is closed under every
signature function.
\end{exercise}

\begin{exercise}
\label{ex:fol-cardinality-normalforms-03}
Use compactness to prove that if a theory has arbitrarily large finite models,
then it has an infinite model.  State why this conclusion alone does not say
that any finite model embeds into the infinite model.
\end{exercise}

\begin{exercise}
\label{ex:fol-cardinality-normalforms-04}
Prove that no first-order sentence in the empty signature is true exactly in
the finite nonempty structures.
\end{exercise}

\begin{exercise}
\label{ex:fol-cardinality-normalforms-05}
Convert
\[
 \neg\forall x\bigl(P(x)\to\exists y(Q(y)\land R(x,y))\bigr)
\]
to rectified prenex normal form.  Justify each equivalence and side condition.
\end{exercise}

\begin{exercise}
\label{ex:fol-cardinality-normalforms-06}
Skolemize
\[
 \exists u\forall x\exists y\forall z\exists w\,
 R(u,x,y,z,w).
\]
Give the arity and argument list of every new symbol.
\end{exercise}

\begin{exercise}
\label{ex:fol-cardinality-normalforms-07}
Give two structures showing why a sentence and its Skolemization cannot be
compared for logical equivalence before an expanded signature and a reduct map
have been specified.
\end{exercise}

\begin{exercise}
\label{ex:fol-cardinality-normalforms-08}
Find a most general unifier, or prove none exists, for each pair:
\[
 P(f(x),x),P(y,g(y));
 \qquad
 P(f(x),x),P(f(y),y);
 \qquad
 P(x,x),P(y,f(y)).
\]
The first and final pairs require the occurs check.
\end{exercise}

\begin{exercise}
\label{ex:fol-cardinality-normalforms-09}
Clausify and refute by first-order resolution:
\[
 \forall x(P(x)\to Q(x)),
 \qquad
 \exists xP(x),
 \qquad
 \forall x\neg Q(x).
\]
Record the Skolem constant, every substitution, and the empty clause.
\end{exercise}

\begin{exercise}
\label{ex:fol-cardinality-normalforms-10}
Explain exactly why failure to find a resolution refutation within
a fixed time bound does not establish satisfiability.
\end{exercise}
\part{Model Theory}
\label{part:model-theory}

\chapter{Maps, Substructures, and Elementarity}
\label{ch:maps-elementarity}

This part assumes the Axiom of Choice in the ambient set theory.  A signature means a
single-sorted first-order signature with equality, and all structures have a
nonempty domain.  The definitions extend to many-sorted signatures by requiring
each map to preserve sorts.  The main references are
\cite{ChangKeisler1990,Hodges1993,Marker2002}.

\begin{convention}[Status of results]
Each major result in this part carries one of three status statements.
``Complete proof'' means that all substantive steps are supplied here.
``Proof sketch'' identifies the construction and the missing verification.
``Cited theorem'' records a theorem with exact hypotheses but refers to the
given source for its proof.  These descriptions concern this text only; they do
not rank the formal-logical importance of the results.
\end{convention}

\section{Homomorphisms, embeddings, and isomorphisms}
\label{sec:maps-elementarity-01}

Let \(\Lang\) be a signature, and let \(\mathcal M\) and \(\mathcal N\) be
\(\Lang\)-structures with domains \(M\) and \(N\).

\begin{definition}[Homomorphism]
A \emph{homomorphism} from \(\mathcal M\) to \(\mathcal N\) is a function
\(h:M\to N\) such that:
\begin{enumerate}
  \item \(h(c^{\mathcal M})=c^{\mathcal N}\) for every constant symbol \(c\);
  \item
  \[
    h\bigl(f^{\mathcal M}(a_1,\ldots,a_n)\bigr)
    =f^{\mathcal N}\bigl(h(a_1),\ldots,h(a_n)\bigr)
  \]
  for every \(n\)-ary function symbol \(f\) and all
  \(a_1,\ldots,a_n\in M\); and
  \item if
  \(\mathcal M\entails R(a_1,\ldots,a_n)\), then
  \(\mathcal N\entails R(h(a_1),\ldots,h(a_n))\), for every relation symbol
  \(R\).
\end{enumerate}
The third clause preserves relations only in the forward direction.
\end{definition}

\begin{definition}[Embedding]
An \emph{embedding} \(h:\mathcal M\to\mathcal N\) is an injective
homomorphism that also reflects every relation symbol:
\[
  \mathcal M\entails R(\bar a)
  \quad\Longleftrightarrow\quad
  \mathcal N\entails R(h\bar a).
\]
Here \(h\bar a\) denotes the coordinatewise image of the finite tuple
\(\bar a\).  Because equality is logical identity, injectivity is exactly the
reflection condition for equality.
\end{definition}

\begin{definition}[Isomorphism and automorphism]
An \emph{isomorphism} is a surjective embedding.  An \emph{automorphism} of
\(\mathcal M\) is an isomorphism from \(\mathcal M\) to itself.
\end{definition}

\begin{lemma}[Preservation of terms]
\label{lem:term-preservation}
If \(h:\mathcal M\to\mathcal N\) is a homomorphism, \(t(\bar x)\) is a
\(\Lang\)-term, and \(\bar a\in M\), then
\[
  h\bigl(t^{\mathcal M}(\bar a)\bigr)
  =t^{\mathcal N}(h\bar a).
\]
\textit{Status: complete proof.}
\end{lemma}

\begin{proof}
Induct on the construction of \(t\).  A variable is immediate, and a constant
uses the constant clause in the definition of homomorphism.  If
\(t=f(t_1,\ldots,t_n)\), apply the induction hypothesis to each \(t_i\), then
apply preservation of \(f\).
\end{proof}

\begin{warning}
A homomorphism need not preserve negated atomic formulas, and an embedding need
not preserve formulas with quantifiers.  Elementarity is the additional
condition that preserves all first-order formulas.
\end{warning}

\section{Substructures}
\label{sec:maps-elementarity-02}

Substructure, elementary substructure, and the Tarski--Vaught test were
introduced for the first-order metatheory in
\cref{sec:fol-cardinality-normalforms-01}.  This chapter restates them once
in its map-based framework; the definitions agree, and the restatement is
recorded here rather than at each item.

\begin{definition}[Substructure]
The structure \(\mathcal M\) is a \emph{substructure} of \(\mathcal N\),
written \(\mathcal M\subseteq\mathcal N\), iff \(M\subseteq N\), the
inclusion map is an embedding, and therefore:
\begin{enumerate}
  \item every constant of \(\mathcal N\) lies in \(M\);
  \item \(M\) is closed under every function of \(\mathcal N\); and
  \item every relation of \(\mathcal M\) is the restriction to \(M\) of the
        corresponding relation of \(\mathcal N\).
\end{enumerate}
\end{definition}

Suppose that \(A\subseteq N\) is nonempty, or that \(\Lang\) has a constant
symbol.  The \emph{substructure generated by \(A\)} is the intersection of all
substructures of \(\mathcal N\) whose domains contain \(A\).  Its domain
consists exactly of values \(t^{\mathcal N}(\bar a)\), where \(t\) is a term
and \(\bar a\) is a finite tuple from \(A\), including the values of constant
terms.  If structures are required to be nonempty, \(A=\varnothing\), and the
signature has no constant, a least substructure need not exist; the phrase
``generated by \(\varnothing\)'' then requires a separate convention.

\begin{example}[A substructure that is not elementary]
In the relational signature \(\Lang=\{<\}\), the standard order
\((\N,<)\) is a substructure of \((\Z,<)\).  It is not elementary.  With the
parameter \(0\in\N\),
\[
  (\Z,<)\entails\exists y\,(y<0),
  \qquad
  (\N,<)\notentails\exists y\,(y<0).
\]
This example also shows why agreement on atomic formulas does not imply
agreement on quantified formulas.
\end{example}

\begin{lemma}[Isomorphism invariance of satisfaction]
\label{lem:iso-invariance}
Let \(h:\mathcal M\to\mathcal N\) be an isomorphism.  Then, for every
formula \(\varphi(x_1,\ldots,x_k)\) and all \(a_1,\ldots,a_k\in M\),
\[
  \mathcal M\entails\varphi(\bar a)
  \quad\Longleftrightarrow\quad
  \mathcal N\entails\varphi(h\bar a).
\]
In particular, isomorphic structures satisfy the same sentences.
\textit{Status: complete proof.}
\end{lemma}
\begin{proof}
Induct on \(\varphi\).  For an equality atom \(t_1=t_2\),
\cref{lem:term-preservation} gives
\(h(t_i^{\mathcal M}(\bar a))=t_i^{\mathcal N}(h\bar a)\) for both terms,
and injectivity of \(h\) yields
\(t_1^{\mathcal M}(\bar a)=t_2^{\mathcal M}(\bar a)\) iff
\(t_1^{\mathcal N}(h\bar a)=t_2^{\mathcal N}(h\bar a)\).  For a relation
atom \(R(t_1,\ldots,t_l)\), \cref{lem:term-preservation} reduces the claim
to the preservation and reflection of \(R\) required of an embedding.
The Boolean cases follow from the induction hypotheses and the
satisfaction clauses.  For \(\exists y\,\psi\): a witness \(a'\in M\) on
the left maps to the witness \(h(a')\) on the right by the induction
hypothesis; conversely, a witness \(b\in N\) equals \(h(a')\) for some
\(a'\in M\) by surjectivity, and the induction hypothesis returns the
witness \(a'\).  The universal case is dual, using the same surjectivity.
\end{proof}

\begin{lemma}[Image of an embedding]
\label{lem:embedding-image}
Let \(h:\mathcal M\to\mathcal N\) be an embedding.  Then \(h(M)\)
contains \(c^{\mathcal N}\) for every constant symbol \(c\) and is closed
under \(f^{\mathcal N}\) for every function symbol \(f\).  The structure
\(h[\mathcal M]\) with domain \(h(M)\), with the constants
\(c^{\mathcal N}\), with the restrictions of the functions
\(f^{\mathcal N}\), and with the relations \(R^{\mathcal N}\) restricted
to \(h(M)\) is a substructure of \(\mathcal N\), and \(h\), taken with
codomain \(h(M)\), is an isomorphism from \(\mathcal M\) to
\(h[\mathcal M]\).
\textit{Status: complete proof.}
\end{lemma}
\begin{proof}
By the homomorphism clauses,
\(c^{\mathcal N}=h(c^{\mathcal M})\in h(M)\), and
\(f^{\mathcal N}(h\bar a)=h(f^{\mathcal M}(\bar a))\in h(M)\) for every
tuple \(h\bar a\) from \(h(M)\); so \(h(M)\) is closed as claimed and the
definition of \(h[\mathcal M]\) is legitimate.  Its interpretations are by
construction the restrictions of those of \(\mathcal N\), so the
inclusion map preserves and reflects atomic formulas, and
\(h[\mathcal M]\subseteq\mathcal N\) in the sense of the substructure
definition above.  The map \(h\) with codomain \(h(M)\) is injective,
because \(h\) is, and surjective by construction.  It is a homomorphism
into \(h[\mathcal M]\), since the interpretations there are the
restrictions of those of \(\mathcal N\) to a set containing every
required value, and it reflects every relation symbol \(R\) because the
relation of \(h[\mathcal M]\) interpreting \(R\) is the restriction of
\(R^{\mathcal N}\) to tuples from \(h(M)\) while \(h\) reflects
\(R^{\mathcal N}\).  Hence the corestricted map is a surjective
embedding, which is an isomorphism by definition.
\end{proof}

\section{Elementary maps and elementary substructures}
\label{sec:maps-elementarity-03}

\begin{definition}[Elementary map]
Let \(A\subseteq M\).  A map \(h:A\to N\) is \emph{elementary} iff, for every
\(\Lang\)-formula \(\varphi(x_1,\ldots,x_n)\) and every
\(a_1,\ldots,a_n\in A\),
\[
  \mathcal M\entails\varphi(a_1,\ldots,a_n)
  \quad\Longleftrightarrow\quad
  \mathcal N\entails\varphi(h(a_1),\ldots,h(a_n)).
\]
When \(A=M\), the map is an \emph{elementary embedding}.
\end{definition}

\begin{definition}[Elementary substructure]
If \(\mathcal M\subseteq\mathcal N\) and the inclusion map is elementary,
then \(\mathcal M\) is an \emph{elementary substructure} of \(\mathcal N\),
written \(\mathcal M\preccurlyeq\mathcal N\).  The structure
\(\mathcal N\) is then an \emph{elementary extension} of \(\mathcal M\).
\end{definition}

\begin{remark}
The symbols \(\subseteq\) and \(\preccurlyeq\) express different claims.
The first concerns the interpretations of the signature symbols; the second
concerns all formulas with parameters from the smaller domain.
\end{remark}

\begin{theorem}[Tarski--Vaught test]
\label{thm:tarski-vaught}
Let \(\mathcal M\subseteq\mathcal N\).  Then
\(\mathcal M\preccurlyeq\mathcal N\) iff, for every formula
\(\varphi(y,\bar x)\) and every finite tuple \(\bar a\) from \(M\),
\[
  \mathcal N\entails\exists y\,\varphi(y,\bar a)
  \quad\Longrightarrow\quad
  \mathcal N\entails\varphi(b,\bar a)
  \text{ for some }b\in M.
\]
\textit{Status: complete proof.}
\end{theorem}

This restates \cref{thm:tarski-vaught-fol} in the map-based framework of
this chapter; the proof is repeated so that the present part remains
self-contained.

\begin{proof}
Suppose first that \(\mathcal M\preccurlyeq\mathcal N\) and that the
displayed existential formula is satisfied in \(\mathcal N\).  By
elementarity it is satisfied in
\(\mathcal M\), so a witness \(b\in M\) exists.  Elementarity again gives
\(\mathcal N\entails\varphi(b,\bar a)\).

Conversely, assume the witness condition.  We prove by induction on formulas
\(\psi(\bar x)\) that, for every \(\bar a\) from \(M\),
\[
  \mathcal M\entails\psi(\bar a)
  \quad\Longleftrightarrow\quad
  \mathcal N\entails\psi(\bar a).
\]
Atomic formulas agree because \(\mathcal M\subseteq\mathcal N\), using
\cref{lem:term-preservation}.  Boolean connectives follow from the induction
hypothesis.  For the existential step, if
\(\mathcal M\entails\exists y\,\varphi(y,\bar a)\), a witness in \(M\) and the
induction hypothesis give the same existential statement in \(\mathcal N\).
In the reverse direction, the assumed witness condition supplies \(b\in M\)
with \(\mathcal N\entails\varphi(b,\bar a)\); the induction hypothesis yields
\(\mathcal M\entails\varphi(b,\bar a)\).  For a universal formula
\(\forall y\,\varphi(y,\bar x)\): if
\(\mathcal N\entails\forall y\,\varphi(y,\bar a)\), then for each
\(b\in M\) the instance \(\varphi(b,\bar a)\) is satisfied in
\(\mathcal N\), and the induction hypothesis transfers it to
\(\mathcal M\).  Conversely, if
\(\mathcal N\notentails\forall y\,\varphi(y,\bar a)\), the satisfaction
clauses give \(\mathcal N\entails\exists y\,\neg\varphi(y,\bar a)\); the
witness condition supplies \(b\in M\) with
\(\mathcal N\entails\neg\varphi(b,\bar a)\), the induction hypothesis for
\(\varphi\) gives \(\mathcal M\entails\neg\varphi(b,\bar a)\), and hence
\(\mathcal M\notentails\forall y\,\varphi(y,\bar a)\).
\end{proof}

\begin{example}[Applying the test]
Let \(K\subseteq L\) be algebraically closed fields in the language of rings.
The later quantifier-elimination theorem for algebraically closed fields
implies \(K\preccurlyeq L\).  Without that theorem, closure under the ring
operations establishes only \(K\subseteq L\); algebraic closure alone does
not yield the Tarski--Vaught witness condition without further argument.
\end{example}

\section{Diagrams and their embedding tests}
\label{sec:maps-elementarity-04}
\label{sec:diagrams}

\begin{definition}[Constants for parameters]
\label{def:parameter-expansion}
If \(A\subseteq M\), let
\[
  \Lang(A)=\Lang\cup\{c_a:a\in A\},
\]
where the \(c_a\) are pairwise distinct new constant symbols, and let
\(\mathcal M_A\) be the \(\Lang(A)\)-expansion of \(\mathcal M\) in which
\(c_a^{\mathcal M_A}=a\).  The notation is needed to distinguish formulas
that may name members of \(A\) from formulas in the original signature.  It
names a signature expansion and an interpreted structure, not a theory, and it
will be used both for diagrams and for types over parameters.
\end{definition}

\begin{definition}[Atomic and elementary diagrams]
\label{def:diagrams}
For a \(\Lang\)-structure \(\mathcal M\), its \emph{atomic diagram}
\(\operatorname{Diag}(\mathcal M)\) is the set of all atomic and negated
atomic \(\Lang(M)\)-sentences true in \(\mathcal M_M\).  Its
\emph{elementary diagram} is
\[
  \operatorname{ElDiag}(\mathcal M)
  =\Th_{\Lang(M)}(\mathcal M_M),
\]
the set of all \(\Lang(M)\)-sentences true in \(\mathcal M_M\).  Both
objects are sets of sentences.  The atomic diagram records exactly the
quantifier-free atomic information needed to test embeddings, whereas the
elementary diagram records every formula needed to test elementary
embeddings; the two names are therefore not interchangeable, and the two
tests below use the distinction.
\end{definition}

\begin{theorem}[Diagram tests for embeddings]
\label{thm:diagram-tests}
Let \(\mathcal M\) and \(\mathcal N\) be \(\Lang\)-structures.
\begin{enumerate}
  \item There is an embedding \(h:\mathcal M\to\mathcal N\) iff
        \(\mathcal N\) has a \(\Lang(M)\)-expansion \(\mathcal N'\) such
        that \(\mathcal N'\entails\operatorname{Diag}(\mathcal M)\).
  \item There is an elementary embedding
        \(h:\mathcal M\to\mathcal N\) iff
        \(\mathcal N\) has a \(\Lang(M)\)-expansion \(\mathcal N'\) such
        that \(\mathcal N'\entails\operatorname{ElDiag}(\mathcal M)\).
\end{enumerate}
\textit{Status: complete proof.}
\end{theorem}

\begin{proof}
Given an embedding \(h\), expand \(\mathcal N\) by interpreting
\(c_a^{\mathcal N'}=h(a)\).  Preservation and reflection of atomic formulas
then make every sentence in \(\operatorname{Diag}(\mathcal M)\) true in
\(\mathcal N'\).  Conversely, suppose such an expansion is given and define
\(h(a)=c_a^{\mathcal N'}\).  For \(a\ne b\), the sentence
\(c_a\ne c_b\) belongs to the atomic diagram, so \(h\) is injective.  The
atomic sentences specifying constants and function values show that \(h\)
preserves those symbols, and the positive and negated relation instances show
that it preserves and reflects relations.  Hence \(h\) is an embedding.

The same expansion constructed from an elementary embedding satisfies every
sentence in the elementary diagram.  Conversely, if
\(\mathcal N'\entails\operatorname{ElDiag}(\mathcal M)\), the map just
defined is an embedding by the first part.  For every formula
\(\varphi(\bar x)\) and tuple \(\bar a\) from \(M\), exactly one of
\(\varphi(\bar c_{\bar a})\) and
\(\neg\varphi(\bar c_{\bar a})\) belongs to the elementary diagram.
Satisfaction of that sentence by \(\mathcal N'\) gives
\[
  \mathcal M\entails\varphi(\bar a)
  \quad\Longleftrightarrow\quad
  \mathcal N\entails\varphi(h\bar a),
\]
so \(h\) is elementary.
\end{proof}

For diagrams and the two embedding tests, see
\cite{Hodges1993,Marker2002}.

\begin{example}[The two diagrams make different tests]
For the inclusion \((\N,<)\subseteq(\Z,<)\), interpret each new constant
\(c_n\) by the integer \(n\).  The resulting expansion of \((\Z,<)\)
satisfies the atomic diagram of \((\N,<)\), in agreement with the fact that
the inclusion is an embedding.  It does not satisfy the elementary diagram:
\(\neg\exists y\,(y<c_0)\) is true in \((\N,<)_{\N}\) and false in the
expanded integer order.
\end{example}

\begin{exercise}
\label{ex:maps-elementarity-01}
\label{ex:diagram-test}
Let \(\mathcal M\) be a substructure of \(\mathcal N\).  Prove directly that
the canonical \(\Lang(M)\)-expansion of \(\mathcal N\) satisfies
\(\operatorname{Diag}(\mathcal M)\).  Show that it satisfies
\(\operatorname{ElDiag}(\mathcal M)\) iff
\(\mathcal M\preccurlyeq\mathcal N\).
\end{exercise}

\section{Elementary chains}
\label{sec:maps-elementarity-05}

\begin{definition}[Elementary chain]
Let \(\delta\) be a limit ordinal.  A sequence
\(\langle\mathcal M_i:i<\delta\rangle\) is an \emph{elementary chain} iff
\(\mathcal M_i\preccurlyeq\mathcal M_j\) whenever \(i<j<\delta\).
Its union \(\bigcup_{i<\delta}\mathcal M_i\) has domain
\(\bigcup_{i<\delta}M_i\), with each symbol interpreted by the compatible
union of its interpretations in the chain.
\end{definition}

\begin{theorem}[Elementary chain theorem]
\label{thm:elementary-chain}
If \(\langle\mathcal M_i:i<\delta\rangle\) is an elementary chain and
\(\mathcal M=\bigcup_{i<\delta}\mathcal M_i\), then
\(\mathcal M_i\preccurlyeq\mathcal M\) for every \(i<\delta\).
\textit{Status: complete proof.}
\end{theorem}

\begin{proof}
We prove simultaneously, by induction on formulas \(\psi(\bar x)\), that for
every \(i<\delta\) and every tuple \(\bar a\) from \(M_i\),
\[
  \mathcal M_i\entails\psi(\bar a)
  \quad\Longleftrightarrow\quad
  \mathcal M\entails\psi(\bar a).
\]
Atomic formulas agree because all interpretations are compatible; Boolean
cases follow from the induction hypothesis.  For the existential case, a
witness in \(M_i\) remains a witness in the union by the induction hypothesis.
Conversely, suppose
\(\mathcal M\entails\exists y\,\varphi(y,\bar a)\), and choose a witness
\(b\in M\).  Some \(j\geq i\) contains \(b\).  Applying the induction
hypothesis to the strictly simpler formula \(\varphi\), now at stage \(j\),
gives \(\mathcal M_j\entails\varphi(b,\bar a)\).  Therefore
\(\mathcal M_j\entails\exists y\,\varphi(y,\bar a)\), and
\(\mathcal M_i\preccurlyeq\mathcal M_j\) yields the existential statement in
\(\mathcal M_i\).

For a universal formula \(\forall y\,\varphi(y,\bar x)\): if
\(\mathcal M\entails\forall y\,\varphi(y,\bar a)\), then for every
\(b\in M_i\) the induction hypothesis at stage \(i\) transfers the
instance \(\varphi(b,\bar a)\) from \(\mathcal M\) to
\(\mathcal M_i\), so \(\mathcal M_i\entails\forall y\,\varphi(y,\bar a)\).
If \(\mathcal M\notentails\forall y\,\varphi(y,\bar a)\), choose
\(b\in M\) with \(\mathcal M\notentails\varphi(b,\bar a)\) and
\(j\geq i\) with \(b\in M_j\); the induction hypothesis at stage \(j\)
gives \(\mathcal M_j\notentails\varphi(b,\bar a)\), so
\(\mathcal M_j\notentails\forall y\,\varphi(y,\bar a)\), and
\(\mathcal M_i\preccurlyeq\mathcal M_j\) transfers this failure to
\(\mathcal M_i\).  This completes the induction and proves elementarity of
every stage in the union.
\end{proof}

\begin{warning}[A frequent chain error]
The union of an increasing chain of substructures is a substructure, but it
need not be an elementary extension of every stage.  Elementarity of the
transition maps is the hypothesis that supports
\cref{thm:elementary-chain}.
\end{warning}

\begin{exercise}
\label{ex:maps-elementarity-02}
Let \(\mathcal M\preccurlyeq\mathcal N\) and let \(A\subseteq M\).  Prove that
the expansion \((\mathcal M,a)_{a\in A}\) is an elementary substructure of
\((\mathcal N,a)_{a\in A}\) in the language with a constant symbol for every
member of \(A\).
\end{exercise}

\begin{exercise}
\label{ex:maps-elementarity-03}
Give an increasing chain of substructures
\(\mathcal M_0\subseteq\mathcal M_1\subseteq\cdots\) such that
\(\mathcal M_0\not\preccurlyeq\bigcup_n\mathcal M_n\).  Identify a formula
and parameters witnessing failure of elementarity.
\end{exercise}

\chapter{Types, Saturation, and Ultraproducts}
\label{ch:types-ultraproducts}

\section{Types over parameter sets}
\label{sec:types-ultraproducts-01}

Let \(\mathcal M\) be a \(\Lang\)-structure and \(A\subseteq M\).  Use the
parameter expansion \(\Lang(A)\) and \(\mathcal M_A\) fixed in
\cref{def:parameter-expansion}.

\begin{definition}[Type]
Let \(\bar x=(x_1,\ldots,x_n)\) be a finite tuple of distinct variables.  An
\emph{\(n\)-type over \(A\) consistent with \(\mathcal M\)} is a set
\(p(\bar x)\) of \(\Lang(A)\)-formulas whose free variables are among
\(\bar x\) such that
\[
  \Th_{\Lang(A)}(\mathcal M_A)\cup
  \set{\varphi(\bar c):\varphi(\bar x)\in p}
\]
is satisfiable, where \(\bar c\) is a new tuple of constants.  The type is
\emph{complete} iff, for every \(\Lang(A)\)-formula \(\varphi(\bar x)\),
exactly one of \(\varphi\) and \(\neg\varphi\) belongs to \(p\).  The
phrase \emph{consistent with} is the standard model-theoretic term and is
defined here semantically, by satisfiability; by soundness and completeness
of the fixed first-order calculus, this sense agrees with consistency in
the registered derivability sense.  That departure is declared once here
and governs types over parameter sets and over theories.
\end{definition}

\begin{definition}[Types over a theory and type spaces]
\label{def:type-space}
Let \(T\) be a satisfiable \(\Lang\)-theory.  An \emph{\(n\)-type
consistent with \(T\)} is a set \(p(\bar x)\) of \(\Lang\)-formulas whose
free variables are among the distinct variables in \(\bar x\), such that
\[
  T\cup\{\varphi(\bar c):\varphi(\bar x)\in p\}
\]
is satisfiable for a new \(n\)-tuple of constants \(\bar c\).  It is complete
when it contains exactly one of \(\varphi\) and \(\neg\varphi\) for every
\(\Lang\)-formula \(\varphi(\bar x)\).  Let \(S_n(T)\) be the set of all
complete \(n\)-types consistent with \(T\).  For each formula
\(\varphi(\bar x)\), put
\[
  [\varphi]=\{p\in S_n(T):\varphi\in p\}.
\]
The sets \([\varphi]\) form a basis of clopen sets because
\([\varphi]\cap[\psi]=[\varphi\land\psi]\) and
\(S_n(T)\setminus[\varphi]=[\neg\varphi]\).  The set \(S_n(T)\) with this
topology is the \emph{\(n\)-type space of \(T\)}.  The topology is required
to express isolation and compactness of complete descriptions; it is not the
set of tuples realized in one chosen model.  This distinction is used below
in the omitting-types theorem and the Ryll--Nardzewski theorem.
\end{definition}

\begin{theorem}[Compactness of type spaces]
\label{thm:type-space-compact}
Assume the Axiom of Choice in the ambient set theory.  If \(T\) is satisfiable, then
\(S_n(T)\) is compact, Hausdorff, and has a
basis of clopen sets.
\textit{Status: complete proof.}
\end{theorem}

\begin{proof}
The displayed basis in \cref{def:type-space} consists of clopen sets.  If
\(p\ne q\), choose \(\varphi\in p\setminus q\); completeness gives
\(\neg\varphi\in q\), so the disjoint neighborhoods \([\varphi]\) and
\([\neg\varphi]\) separate them.

For compactness it suffices to consider a cover by basic open sets: an
arbitrary open cover refines to one by basic sets, because every point of an
open set lies in a basic set contained in that open set, and a finite
subfamily of the basic refinement yields a finite subcover of the original
cover upon replacing each basic set by a member of the original cover
containing it.

So let \(\set{[\varphi_i]:i\in I}\) cover \(S_n(T)\) and suppose, toward a
contradiction, that no finite subfamily covers \(S_n(T)\).  Let \(\bar c\)
be a new \(n\)-tuple of constants and consider the set
\[
  T\cup\set{\neg\varphi_i(\bar c):i\in I}.
\]
Every finite subset of the displayed set is contained in
\(T\cup\set{\neg\varphi_j(\bar c):j\in J}\) for some finite \(J\subseteq I\).
Since \(\set{[\varphi_j]:j\in J}\) is a finite subfamily, it does not cover
\(S_n(T)\), so there is a complete type \(p\in S_n(T)\) with
\(\varphi_j\notin p\) for every \(j\in J\); completeness of \(p\) gives
\(\neg\varphi_j\in p\).  By \cref{def:type-space},
\(T\cup\set{\psi(\bar c):\psi\in p}\) is satisfiable, and any model of it is
a model of \(T\cup\set{\neg\varphi_j(\bar c):j\in J}\).

By first-order compactness (\cref{thm:fol-compactness}, clause 1, applied in
the signature expanded by the constants \(\bar c\)), the displayed set has a
model \(\mathcal N\); let \(\bar b\) be the interpretation of \(\bar c\).
The complete-type step, which here plays the role of the Lindenbaum lemma,
extracts a point of \(S_n(T)\) from \(\mathcal N\): let \(q\) be the set of
all \(\Lang\)-formulas \(\psi(\bar x)\) with free variables among \(\bar x\)
such that \(\mathcal N\entails\psi(\bar b)\).  For every \(\Lang\)-formula
\(\psi(\bar x)\), exactly one of \(\psi\) and \(\neg\psi\) is satisfied by
\(\bar b\) in \(\mathcal N\), so \(q\) contains exactly one of them; and
\(\mathcal N\) itself is a model of \(T\cup\set{\psi(\bar c):\psi\in q}\),
because it satisfies \(T\) and interprets \(\bar c\) as \(\bar b\).  Hence
\(q\) is a complete \(n\)-type consistent with \(T\), that is,
\(q\in S_n(T)\).  For every \(i\in I\), the model \(\mathcal N\) satisfies
\(\neg\varphi_i(\bar b)\), so \(\varphi_i\notin q\) and
\(q\notin[\varphi_i]\).  This contradicts the assumption that
\(\set{[\varphi_i]:i\in I}\) covers \(S_n(T)\).
\end{proof}

\begin{example}[A finite type space]
Let \(T=\Th(\mathcal M)\), where \(\mathcal M\) is a two-element structure in
the equality-only signature.  Every model of \(T\) has exactly two
elements and is isomorphic to \(\mathcal M\), so every complete type
consistent with \(T\) is realized in \(\mathcal M\).  Its one-type space
has one point, since the two
elements are exchanged by an automorphism.  Its two-type space has exactly two
points: the complete type containing \(x_1=x_2\) and the complete type
containing \(x_1\ne x_2\).  Both points are clopen, so this type space is
discrete.
\end{example}

\begin{exercise}
\label{ex:types-ultraproducts-01}
\label{ex:type-space-basis}
For formulas \(\varphi(\bar x)\) and \(\psi(\bar x)\), verify directly that
\([\varphi]\cap[\psi]=[\varphi\land\psi]\) and that
\(S_n(T)\setminus[\varphi]=[\neg\varphi]\).  Compute the one- and two-type
spaces in the preceding example and exhibit every singleton as a basic open
set.
\end{exercise}

\begin{definition}[Realization]
A tuple \(\bar b\in M^n\) \emph{realizes} \(p(\bar x)\) in
\(\mathcal M\) iff \(\mathcal M\entails\varphi(\bar b)\) for every
\(\varphi(\bar x)\in p\).  The structure \(\mathcal M\) \emph{omits}
\(p\) iff no tuple in \(M^n\) realizes \(p\) in \(\mathcal M\).
\end{definition}

\begin{example}[A finitely satisfiable type not realized]
Work in the parameter expansion of \((\N,<)\) having the previously defined
constant \(c_n\) for each \(n\in\N\), and let
\[
  p(x)=\set{c_n<x:n\in\N}.
\]
Thus every displayed member is a formula of the fixed parameter language.
Every finite subset is realized in \(\N\), but
the whole type is not.  Assuming the Axiom of Choice in the ambient set theory,
compactness shows that \(p\) is realized in an elementary
extension of \((\N,<)\).  Such a realization is greater than every standard
natural number.
\end{example}

\begin{definition}[Saturation]
Let \(\kappa\) be an infinite cardinal.  The structure \(\mathcal M\) is
\emph{\(\kappa\)-saturated} iff, whenever \(A\subseteq M\) has
\(\abs{A}<\kappa\), every complete type in finitely many variables over
\(A\) consistent with \(\mathcal M\) is realized in \(\mathcal M\).
An infinite structure is \emph{saturated} iff it is
\(\abs{M}\)-saturated.  This reference does not apply the unqualified term
``saturated'' to a finite structure, because \(\kappa\)-saturation has just
been defined only for infinite \(\kappa\).
\end{definition}

\begin{remark}
Some sources allow tuples of length less than \(\kappa\) in the definition of
\(\kappa\)-saturation.  This reference uses finite tuples.  The convention must
be restated before applying a theorem quoted from a source using the stronger
version.
\end{remark}

\begin{theorem}[Existence of saturated elementary extensions]
\label{thm:saturated-extension}
Assume the Axiom of Choice in the ambient set theory.  Let \(\mathcal M\) be a
structure and let \(\kappa\) be an infinite cardinal.
There is a \(\kappa\)-saturated elementary extension
\(\mathcal N\succcurlyeq\mathcal M\).
\textit{Status: cited theorem.}
\end{theorem}

See \cite{Marker2002} for a complete proof.
The construction iterates elementary extensions realizing all types over
small parameter sets, and takes unions at limit stages using
\cref{thm:elementary-chain}.  Producing a specified cardinal bound requires
additional cardinal-arithmetic bookkeeping and is not asserted here.

\section{Homogeneity}
\label{sec:types-ultraproducts-02}
\label{sec:homogeneity}

\begin{definition}[Homogeneity]
\label{def:homogeneity}
Let \(\kappa\) be an infinite cardinal.  A structure \(\mathcal M\) is
\emph{\(\kappa\)-homogeneous} iff, whenever \(A\subseteq M\) has
\(\abs{A}<\kappa\), \(h:A\to M\) is an elementary map, and \(b\in M\),
there is an elementary map \(h':A\cup\{b\}\to M\) extending \(h\).  This is
a property of a structure, not a new kind of map.  It records extension of
partial elementary self-maps, whereas saturation records realization of
types; the notions are therefore not definitionally redundant.  The theorem
below relates them, and the later back-and-forth method uses the extension
property in both directions.
\end{definition}

\begin{theorem}[Saturation supplies homogeneity]
\label{thm:saturation-homogeneity}
Assume the Axiom of Choice in the ambient set theory.
Every \(\kappa\)-saturated structure is \(\kappa\)-homogeneous.
\textit{Status: complete proof.}
\end{theorem}

\begin{proof}
Let \(h:A\to M\) be elementary with \(\abs{A}<\kappa\) and let \(b\in M\).
If \(b\in A\), take \(h'=h\).  So suppose \(b\notin A\), and let \(p(x)\)
be the complete type of \(b\) over \(A\): the set of
\(\Lang(A)\)-formulas \(\varphi(x)\) with
\(\mathcal M_A\entails\varphi(b)\).  For an \(\Lang(A)\)-formula
\(\varphi\), let \(h(\varphi)\) be the \(\Lang(h[A])\)-formula obtained by
replacing each parameter constant \(c_a\) by \(c_{h(a)}\), and put
\(q(x)=\set{h(\varphi):\varphi\in p}\).  Every \(\Lang(h[A])\)-formula
\(\psi(x)\) has the form \(h(\varphi)\), and
\(h(\neg\varphi)=\neg h(\varphi)\), so \(q\) contains \(\psi\) or
\(\neg\psi\) for every such \(\psi\).

The set \(q\) is a type over \(h[A]\) consistent with \(\mathcal M\); this
is the compactness appeal.  Fix a new constant \(c\).  By the type
definition, the claim is that
\[
  \Th_{\Lang(h[A])}(\mathcal M_{h[A]})\cup\set{h(\varphi)(c):\varphi\in p}
\]
is satisfiable.  Every finite subset of the displayed set is contained in
\(\Th_{\Lang(h[A])}(\mathcal M_{h[A]})\cup
\set{h(\varphi_1)(c),\ldots,h(\varphi_k)(c)}\) for some
\(\varphi_1,\ldots,\varphi_k\in p\).  Since \(b\) satisfies each
\(\varphi_i\), the structure \(\mathcal M_A\) satisfies
\(\exists x(\varphi_1\land\cdots\land\varphi_k)\).  Write this sentence as
\(\exists x\,\chi(x,c_{a_1},\ldots,c_{a_m})\) with \(\chi\) an
\(\Lang\)-formula and \(a_1,\ldots,a_m\in A\); then
\(\mathcal M\entails\exists x\,\chi(x,a_1,\ldots,a_m)\), and elementarity
of \(h\) gives
\(\mathcal M\entails\exists x\,\chi(x,h(a_1),\ldots,h(a_m))\), that is,
\(\mathcal M_{h[A]}\entails
\exists x(h(\varphi_1)\land\cdots\land h(\varphi_k))\).  Interpreting
\(c\) as such a witness expands \(\mathcal M_{h[A]}\) to a model of the
finite subset.  By first-order compactness (\cref{thm:fol-compactness},
clause 1, applied in the signature \(\Lang(h[A])\) expanded by \(c\)), the
displayed set is satisfiable; this appeal is the step that uses the Axiom
of Choice assumed in the statement.  In particular \(q\) never contains
both \(\psi\) and \(\neg\psi\), so with the previous paragraph \(q\) is a
complete type over \(h[A]\) consistent with \(\mathcal M\).

Since \(\abs{h[A]}\leq\abs{A}<\kappa\), saturation supplies a realization
\(d\in M\) of \(q\).  Define \(h'(b)=d\) and let \(h'\) agree with \(h\) on
\(A\).  For every \(\Lang\)-formula \(\chi(x,\bar y)\) and tuple \(\bar a\)
from \(A\), exactly one of \(\chi(x,c_{a_1},\ldots,c_{a_m})\) and its
negation belongs to \(p\), and \(d\) satisfies its image under \(h\);
hence \(\mathcal M\entails\chi(b,\bar a)\) iff
\(\mathcal M\entails\chi(d,h(a_1),\ldots,h(a_m))\).  A tuple from
\(A\cup\{b\}\) arises from a tuple from \(A\) by inserting \(b\) into some
coordinates, and replacing those coordinates by one fresh variable reduces
the elementarity condition for \(h'\) to the displayed equivalence.  Thus
\(h'\) is elementary.
\end{proof}

See \cite{Marker2002} for the general saturation--homogeneity development.

\begin{example}[What homogeneity adds]
The cited saturated-extension theorem produces, for any structure
\(\mathcal M\) and infinite \(\kappa\), an elementary extension
\(\mathcal N\) in which every type over fewer than \(\kappa\) parameters is
realized.  By \cref{thm:saturation-homogeneity}, every elementary map from
fewer than \(\kappa\) elements of \(N\) can then be extended to include any
one further domain element.  This extension claim is the homogeneity content;
it is not part of the definition of saturation.
\end{example}

\begin{exercise}
\label{ex:types-ultraproducts-02}
\label{ex:homogeneity-back-step}
Let \(\mathcal M\) be \(\kappa\)-homogeneous, let \(h:A\to M\) be elementary
with \(\abs{A}<\kappa\), and let \(d\in M\).  Apply homogeneity to the inverse
map on \(h[A]\) to prove that \(h\) has an elementary extension whose range
contains \(d\).
\end{exercise}

\section{Omitting types}
\label{sec:types-ultraproducts-03}
\label{sec:omitting-types}

\begin{definition}[Principal and nonprincipal types]
\label{def:principal-type}
Let \(p(\bar x)\) be a type consistent with a theory \(T\).  The type is
\emph{principal over \(T\)} iff some formula \(\theta(\bar x)\) satisfies
\begin{enumerate}
  \item \(T\cup\{\exists\bar x\,\theta(\bar x)\}\) is satisfiable; and
  \item for every \(\varphi(\bar x)\in p\),
        \(T\entails\forall\bar x(\theta\to\varphi)\).
\end{enumerate}
If no such formula exists, \(p\) is \emph{nonprincipal over \(T\)}.  This is
a property of a type relative to a theory, not a claim that a selected model
realizes or omits it.  It distinguishes a type forced by one consistent
formula from a merely finitely satisfiable collection, and that distinction is
the hypothesis of the theorem below.  When \(p\) is complete, it is principal
iff its point is isolated in \(S_n(T)\): an isolating formula
\(\theta\) gives \(\{p\}=[\theta]\), and a basic open singleton gives such a
formula.
\end{definition}

\begin{theorem}[Omitting types theorem]
\label{thm:omitting-types}
Let \(T\) be a satisfiable theory in a countable signature, and let
\(p_0,p_1,\ldots\) be a countable family of nonprincipal types consistent
with \(T\), each in a finite tuple of variables.  Then \(T\) has a countable
model that omits every \(p_n\).
\textit{Status: cited theorem.}
\end{theorem}

See \cite{Hodges1993,Marker2002} for complete proofs.  The Henkin construction
enumerates both witness requirements and pairs consisting of a term tuple and
a type \(p_n\).  At the latter requirement, nonprincipality permits adding the
negation of some formula in \(p_n\) at that tuple without destroying finite
satisfiability.  The term model therefore witnesses the theory while no tuple
realizes every formula of any listed type.

\begin{example}[A type that a countable model omits]
Take a signature with constants \(c_0,c_1,\ldots\), and let
\[
  T=\{c_m\ne c_n:m\ne n\},
  \qquad
  p(x)=\{x\ne c_n:n\in\N\}.
\]
The type is consistent and nonprincipal over \(T\).  The model whose domain
consists exactly of the distinct interpretations of the constants satisfies
\(T\) and omits \(p\), while a model with one additional unnamed element
realizes \(p\).  Thus omission is a property of a model, whereas
nonprincipality is the syntactic-semantic condition on the type used by
\cref{thm:omitting-types}.
\end{example}

\begin{exercise}
\label{ex:types-ultraproducts-03}
\label{ex:omit-constant-type}
Prove the two assertions in the example.  For nonprincipality, use the fact
that a single formula contains only finitely many of the constant symbols; for
omission, show that every element of the displayed countable model is the
interpretation of some \(c_n\).
\end{exercise}

\section{Filters and ultrafilters}
\label{sec:types-ultraproducts-04}

\begin{definition}[Filter]
Let \(I\) be a nonempty set.  A \emph{filter} on \(I\) is a set
\(\mathcal F\subseteq\Pow(I)\) such that:
\begin{enumerate}
  \item \(I\in\mathcal F\) and \(\varnothing\notin\mathcal F\);
  \item if \(X,Y\in\mathcal F\), then \(X\cap Y\in\mathcal F\); and
  \item if \(X\in\mathcal F\) and \(X\subseteq Y\subseteq I\), then
        \(Y\in\mathcal F\).
\end{enumerate}
An \emph{ultrafilter} is a filter \(\mathcal U\) such that, for every
\(X\subseteq I\), exactly one of \(X\) and \(I\setminus X\) belongs to
\(\mathcal U\).  The requirement that exactly one of \(X\) and
\(I\setminus X\) belong to \(\mathcal U\) is cited below as the
\emph{ultrafilter alternative}.
\end{definition}

An ultrafilter is \emph{principal} iff it has the form
\(\set{X\subseteq I:i_0\in X}\) for some \(i_0\in I\).  A subset
\(X\subseteq I\) is \emph{cofinite} iff \(I\setminus X\) is finite.

\begin{lemma}[Existence of nonprincipal ultrafilters]
\label{lem:nonprincipal-ultrafilter}
Assume the Ultrafilter Lemma in the ambient set theory.  Let \(I\) be an
infinite set.
\begin{enumerate}
  \item The cofinite subsets of \(I\) form a filter on \(I\), the
        \emph{Fr\'echet filter} on \(I\).
  \item Some ultrafilter on \(I\) contains the Fr\'echet filter, and every
        ultrafilter containing it is nonprincipal.
  \item An ultrafilter on \(I\) is nonprincipal iff it contains every
        cofinite subset of \(I\).
\end{enumerate}
In particular, there is a nonprincipal ultrafilter on \(\N\).
\textit{Status: complete proof.}
\end{lemma}

\begin{proof}
For clause 1, check the three filter conditions.  \(I\setminus
I=\varnothing\) is finite, so \(I\) is cofinite, while
\(I\setminus\varnothing=I\) is infinite, so \(\varnothing\) is not.  If
\(X\) and \(Y\) are cofinite, then
\(I\setminus(X\cap Y)=(I\setminus X)\cup(I\setminus Y)\) is a union of two
finite sets, hence finite, so \(X\cap Y\) is cofinite.  If \(X\) is
cofinite and \(X\subseteq Y\subseteq I\), then
\(I\setminus Y\subseteq I\setminus X\) is finite, so \(Y\) is cofinite.

For clause 2, the Fr\'echet filter satisfies the two displayed conditions
of the Ultrafilter Lemma (\cref{def:ultrafilter-lemma}); these are exactly
the filter conditions just verified.  The assumed principle therefore
supplies an ultrafilter containing the Fr\'echet filter.  Now let
\(\mathcal U\) be any ultrafilter containing the Fr\'echet filter and let
\(i_0\in I\).  The set \(I\setminus\set{i_0}\) is cofinite, so it belongs
to \(\mathcal U\); it omits \(i_0\), so
\(\mathcal U\ne\set{X\subseteq I:i_0\in X}\).  Since \(i_0\) was
arbitrary, \(\mathcal U\) is nonprincipal.

For clause 3, right to left, an ultrafilter containing every cofinite
subset contains the Fr\'echet filter and is nonprincipal by clause 2.
Left to right, argue contrapositively in three steps that an ultrafilter
\(\mathcal U\) containing some finite set is principal.  First, if
\(X\cup Y\in\mathcal U\), then \(X\in\mathcal U\) or \(Y\in\mathcal U\):
otherwise the ultrafilter alternative would put \(I\setminus X\) and
\(I\setminus Y\) in \(\mathcal U\), hence also
\((I\setminus X)\cap(I\setminus Y)=I\setminus(X\cup Y)\), and the
alternative would then exclude \(X\cup Y\), contrary to assumption.
Second, by induction on \(m\geq 1\): if \(F\in\mathcal U\) is finite with
\(\abs{F}=m\), then \(\set{i}\in\mathcal U\) for some \(i\in F\).  For
\(m=1\) the set \(F\) is itself a singleton.  For the step from \(m\) to
\(m+1\), choose \(i\in F\) and write
\(F=(F\setminus\set{i})\cup\set{i}\); the first step gives
\(\set{i}\in\mathcal U\), in which case the claim holds, or
\(F\setminus\set{i}\in\mathcal U\), a finite member of size \(m\) to which
the induction hypothesis applies.  Third, if \(\set{i_0}\in\mathcal U\),
then \(\mathcal U=\set{X\subseteq I:i_0\in X}\): every \(X\) with
\(i_0\in X\) includes \(\set{i_0}\) and belongs to \(\mathcal U\) by
upward closure, and no \(X\) with \(i_0\notin X\) belongs to
\(\mathcal U\), since otherwise \(X\cap\set{i_0}=\varnothing\) would lie
in \(\mathcal U\).  Contrapositively, a nonprincipal ultrafilter contains
no finite set, and the ultrafilter alternative then puts the complement of
every finite set, that is every cofinite set, into it.

The final claim is clause 2 applied to the infinite set \(\N\).
\end{proof}

\section{Ultraproducts}
\label{sec:types-ultraproducts-05}

Throughout this section, assume the Axiom of Choice in the ambient
metatheory.  It supplies an element of every Cartesian product of the
nonempty domains below and supplies the coordinatewise witness functions used
in the quantifier cases of \L o\'s's theorem.  This is a metatheoretic
assumption, not a first-order axiom of any factor structure.

Let \(\langle\mathcal M_i:i\in I\rangle\) be \(\Lang\)-structures and let
\(\mathcal U\) be an ultrafilter on \(I\).  On the Cartesian product
\(\prod_{i\in I}M_i\), define
\[
  a\sim_{\mathcal U}b
  \quad\Longleftrightarrow\quad
  \set{i\in I:a(i)=b(i)}\in\mathcal U.
\]

\begin{definition}[Ultraproduct]
The \emph{ultraproduct}
\[
  \prod_{i\in I}\mathcal M_i/\mathcal U
\]
has as its domain the set of \(\sim_{\mathcal U}\)-equivalence classes
\([a]_{\mathcal U}\).  Constants and functions are interpreted
coordinatewise.  For an \(n\)-ary relation symbol \(R\), define
\[
  R([a_1]_{\mathcal U},\ldots,[a_n]_{\mathcal U})
  \quad\Longleftrightarrow\quad
  \set{i\in I:\mathcal M_i\entails R(a_1(i),\ldots,a_n(i))}
  \in\mathcal U.
\]
When all \(\mathcal M_i\) are the same structure \(\mathcal M\), the result is
the \emph{ultrapower} \(\mathcal M^I/\mathcal U\).
\end{definition}

\begin{lemma}
The relation \(\sim_{\mathcal U}\) is an equivalence relation, and the
coordinatewise interpretations in the ultraproduct do not depend on the
chosen representatives.
\textit{Status: complete proof.}
\end{lemma}

\begin{proof}
Reflexivity and symmetry are immediate.  For transitivity, intersect the two
\(\mathcal U\)-large equality sets and use upward closure.  For
well-definedness, if representatives agree on \(\mathcal U\)-large sets, then
their common intersection is \(\mathcal U\)-large.  On that intersection the
corresponding function values agree, and the truth values of each relation
instance agree.  Upward closure completes both verifications.
\end{proof}

\begin{lemma}[Term evaluation in an ultraproduct]
\label{lem:los-terms}
Let \(t(x_1,\ldots,x_n)\) be a \(\Lang\)-term and let
\(a_1,\ldots,a_n\in\prod_{i\in I}M_i\).  Define
\(b_t\in\prod_{i\in I}M_i\) by
\[
  b_t(i)=t^{\mathcal M_i}(a_1(i),\ldots,a_n(i)).
\]
Then
\[
  t^{\prod_i\mathcal M_i/\mathcal U}
  ([a_1]_{\mathcal U},\ldots,[a_n]_{\mathcal U})
  =[b_t]_{\mathcal U}.
\]
\textit{Status: complete proof.}
\end{lemma}

\begin{proof}
Induct on \(t\).  For a variable \(x_j\), both sides are
\([a_j]_{\mathcal U}\).  For a constant, the claim is the coordinatewise
definition of its ultraproduct interpretation.  If
\(t=f(t_1,\ldots,t_m)\), apply the induction hypothesis to every \(t_j\) and
then use the coordinatewise definition of \(f\) in the ultraproduct.
\end{proof}

\begin{theorem}[\L o\'s's theorem]
\label{thm:los}
Assume the Axiom of Choice in the ambient set theory, as fixed for this
section.  Let \(\varphi(x_1,\ldots,x_n)\) be a first-order
\(\Lang\)-formula and let
\(a_1,\ldots,a_n\in\prod_{i\in I}M_i\).  Then
\[
  \prod_{i\in I}\mathcal M_i/\mathcal U
  \entails\varphi([a_1]_{\mathcal U},\ldots,[a_n]_{\mathcal U})
\]
iff
\[
  \set{i\in I:\mathcal M_i\entails
  \varphi(a_1(i),\ldots,a_n(i))}\in\mathcal U.
\]
\textit{Status: complete proof.}
\end{theorem}

The original source is \cite{Los1955}.

\begin{proof}
We prove the displayed equivalence by structural induction for every formula
and every tuple of representatives for its free variables.  Brackets on a
tuple are coordinatewise, so
\([\bar a]_{\mathcal U}=([a_1]_{\mathcal U},\ldots,[a_n]_{\mathcal U})\).
For such a formula \(\theta\), write
\[
  X_\theta(\bar a)=
  \set{i\in I:\mathcal M_i\entails\theta(\bar a(i))}.
\]

If \(\theta\) is \(t=s\), \cref{lem:los-terms} and the definition of the
quotient give
\[
  \prod_i\mathcal M_i/\mathcal U\entails
  t([\bar a]_{\mathcal U})=s([\bar a]_{\mathcal U})
  \quad\Longleftrightarrow\quad
  \set{i:t^{\mathcal M_i}(\bar a(i))
               =s^{\mathcal M_i}(\bar a(i))}\in\mathcal U.
\]
For \(R(t_1,\ldots,t_m)\), the term lemma followed by the definition of the
ultraproduct relation gives the required equivalence.  For \(\bottom\), both
sides are false because \(X_{\bottom}=\varnothing\notin\mathcal U\).

For negation,
\(X_{\neg\alpha}=I\setminus X_\alpha\), and the ultrafilter alternative says
that this complement is large iff \(X_\alpha\) is not large.  For
conjunction,
\(X_{\alpha\land\beta}=X_\alpha\cap X_\beta\); a filter contains this
intersection iff it contains both factors.  For disjunction,
\(X_{\alpha\lor\beta}=X_\alpha\cup X_\beta\), and an ultrafilter contains a
finite union iff it contains at least one member.  Indeed, if it
contained neither, it would contain both complements and hence the complement
of the union.  Finally,
\[
  X_{\alpha\to\beta}=(I\setminus X_\alpha)\cup X_\beta.
\]
The union property and the ultrafilter alternative make this set large iff
\(X_\alpha\) is not large or \(X_\beta\) is large.  In each of these four
cases, the induction hypotheses and the corresponding truth clause in the
ultraproduct give the desired equivalence.

Suppose \(\theta(\bar x)=\exists y\,\psi(y,\bar x)\).  If the ultraproduct
has a witness \([b]_{\mathcal U}\), the induction hypothesis makes
\(X_\psi(b,\bar a)\) large.  This set is contained in
\(X_\theta(\bar a)\), so the latter is large.  Conversely, if
\(X_\theta(\bar a)\in\mathcal U\), choose for each index in that set a
witness \(b(i)\in M_i\), and choose any \(b(i)\in M_i\) at every remaining
index.  Then \(X_\psi(b,\bar a)\) contains \(X_\theta(\bar a)\) and is large;
the induction hypothesis makes \([b]_{\mathcal U}\) a witness in the
ultraproduct.

Suppose \(\theta(\bar x)=\forall y\,\psi(y,\bar x)\).  If
\(X_\theta(\bar a)\) is large, then for every \(b\in\prod_iM_i\) it is
contained in \(X_\psi(b,\bar a)\).  Upward closure and the induction
hypothesis show that every \([b]_{\mathcal U}\) satisfies \(\psi\), so the
ultraproduct satisfies \(\theta\).  Conversely, suppose the ultraproduct
satisfies \(\theta\) but \(X_\theta(\bar a)\notin\mathcal U\).  Its complement
is large.  At each index in that complement choose \(b(i)\) for which
\(\psi(b(i),\bar a(i))\) is false, and choose any value elsewhere.  If
\(X_\psi(b,\bar a)\) were large, its intersection with the large complement
of \(X_\theta(\bar a)\) would be empty and large, a contradiction.  Thus the
induction hypothesis says that \(\psi([b]_{\mathcal U},[\bar a]_{\mathcal U})\)
is false in the ultraproduct, contradicting its universal statement.  The two
quantifier cases use the stated Axiom of Choice and the stipulated
nonemptiness of every \(M_i\).
\end{proof}

\begin{corollary}[Diagonal elementary embedding]
\label{cor:diagonal-ultrapower}
The map \(d:M\to M^I/\mathcal U\) sending \(a\) to the class of the constant
function with value \(a\) is an elementary embedding.
\textit{Status: complete proof.}
\end{corollary}

\begin{proof}
For a formula with parameters from \(M\), its truth set under constant
representatives is either all of \(I\) or empty.  Apply
\cref{thm:los}.
\end{proof}

\begin{example}[An infinite ultraproduct of finite structures]
Let \(\mathcal M_n\) be the finite linear order
\(\{0,\ldots,n-1\}\), and let \(\mathcal U\) be a nonprincipal ultrafilter on
\(\N\setminus\{0\}\).  The ultraproduct is infinite: for each \(k\), the
sentence asserting that at least \(k\) distinct elements exist holds in all
\(\mathcal M_n\) with \(n\geq k\), a cofinite set of indices.  By
\cref{thm:los}, it holds in the ultraproduct.
\end{example}

\begin{warning}[First-order scope]
\L o\'s's theorem applies to first-order formulas.  The property ``the domain
is finite'' is not first-order definable over all structures, and the example
above must not be read as a contradiction.
\end{warning}

\begin{exercise}
\label{ex:types-ultraproducts-04}
Let \(\mathcal U\) be principal at \(i_0\).  Construct an explicit isomorphism
\[
  \prod_{i\in I}\mathcal M_i/\mathcal U\cong\mathcal M_{i_0}.
\]
\end{exercise}

\begin{exercise}
\label{ex:types-ultraproducts-05}
Use an ultrapower of \((\N,0,S,+,\times,<)\) by a nonprincipal ultrafilter to
produce an element greater than every diagonal image \(d(n)\).  Verify the
claim with \cref{thm:los}.
\end{exercise}

\chapter{Quantifier Elimination and Categoricity}
\label{ch:qe-categoricity}

\section{Quantifier elimination}
\label{sec:qe-categoricity-01}

\begin{definition}[Quantifier elimination]
A theory \(T\) in a signature \(\Lang\) \emph{eliminates quantifiers} iff,
for every \(\Lang\)-formula \(\varphi(\bar x)\), there is a quantifier-free
\(\Lang\)-formula \(\psi(\bar x)\) such that
\[
  T\entails\forall\bar x\,
  \bigl(\varphi(\bar x)\leftrightarrow\psi(\bar x)\bigr).
\]
The equivalence is relative to the models of \(T\); it is not asserted in all
\(\Lang\)-structures.
\end{definition}

\begin{method}[One-quantifier reduction]
To prove quantifier elimination, it suffices to give an explicit reduction of
every formula
\(\exists y\,\theta(y,\bar x)\), with \(\theta\) quantifier-free, to a
quantifier-free formula modulo \(T\).  Boolean induction, replacement of
\(\forall y\,\theta\) by \(\neg\exists y\,\neg\theta\), and renaming of
bound variables then remove all quantifiers from arbitrary formulas.
\end{method}

\begin{theorem}[Dense linear orders]
\label{thm:dlo-qe}
Let \(\mathsf{DLO}\) be the first-order theory of nonempty dense linear orders
without endpoints in the signature \(\{<\}\).  Then \(\mathsf{DLO}\)
eliminates quantifiers.
\textit{Status: cited theorem.}
\end{theorem}

\begin{remark}[Construction underlying the cited theorem]
Put a quantifier-free matrix in disjunctive normal form and distribute the
existential quantifier over the finite disjunction.  In one conjunction,
separate literals involving \(y\) into lower bounds \(t_i<y\), upper bounds
\(y<r_j\), equalities \(y=u_k\), and disequalities.  An equality substitutes
its term for \(y\).  Without an equality, density and the absence of endpoints
show that a witness exists iff every required lower bound is below
every required upper bound, subject to the finitely many equalities and
disequalities among parameters.  Writing those finite compatibility conditions
gives a quantifier-free formula.  A complete proof verifies all inconsistent
literal combinations and then applies the one-quantifier reduction method; see
\cite{Marker2002}.
\end{remark}

\begin{example}[Eliminating one quantifier in \(\mathsf{DLO}\)]
Modulo \(\mathsf{DLO}\),
\[
  \exists y\,(x<y\land y<z)
  \quad\longleftrightarrow\quad x<z.
\]
The forward implication uses transitivity.  The reverse implication is exactly
density.
\end{example}

\begin{definition}[Deductive completeness, semantic completeness, and model completeness]
\label{def:theory-completeness}
A theory \(T\) is \emph{deductively complete} iff, for every sentence
\(\sigma\) in its signature, either \(T\derives\sigma\) or
\(T\derives\neg\sigma\).  It is \emph{semantically complete} iff, for every
such \(\sigma\), either \(T\entails\sigma\) or
\(T\entails\neg\sigma\).  Because derivability in this reference is taken
in the
first-order calculus fixed by
\cref{thm:fol-soundness,thm:fol-completeness},
\[
  T\derives\sigma\quad\Longleftrightarrow\quad T\entails\sigma,
\]
so these two properties are equivalent here.  The fixed term for the
derivability property is ``deductively complete,'' and the fixed term for the
semantic property is ``semantically complete.''

A theory is \emph{model complete} iff every embedding between its models is
elementary.  This is not either decision property above: it concerns formulas
with parameters along embeddings, and later results require that distinction.
\end{definition}

\begin{lemma}[Reduction to substructure inclusions]
\label{lem:model-completeness-reduction}
Let \(T\) be a theory such that
\(\mathcal M\preccurlyeq\mathcal N\) whenever
\(\mathcal M\subseteq\mathcal N\) are both models of \(T\).  Then \(T\)
is model complete.
\textit{Status: complete proof.}
\end{lemma}
\begin{proof}
Let \(h:\mathcal M\to\mathcal N\) be an embedding between models of
\(T\).  By \cref{lem:embedding-image}, \(h\) factors as the isomorphism
\(h_0:\mathcal M\to h[\mathcal M]\) obtained by corestriction, followed
by the inclusion \(h[\mathcal M]\subseteq\mathcal N\).  By
\cref{lem:iso-invariance}, \(h[\mathcal M]\) satisfies the same sentences
as \(\mathcal M\), so \(h[\mathcal M]\) is a model of \(T\).  The
hypothesis applied to \(h[\mathcal M]\subseteq\mathcal N\) gives
\(h[\mathcal M]\preccurlyeq\mathcal N\).  Then, for every formula
\(\varphi(x_1,\ldots,x_k)\) and every tuple \(\bar a\) from \(M\),
\[
  \mathcal M\entails\varphi(\bar a)
  \quad\Longleftrightarrow\quad
  h[\mathcal M]\entails\varphi(h\bar a)
  \quad\Longleftrightarrow\quad
  \mathcal N\entails\varphi(h\bar a),
\]
the first equivalence by \cref{lem:iso-invariance} applied to \(h_0\) and
the second because the inclusion is elementary.  Hence \(h\) is
elementary, and \(T\) is model complete.
\end{proof}

\begin{corollary}
The theory \(\mathsf{DLO}\) is deductively complete (equivalently,
semantically complete) and model complete.
\textit{Status: proof sketch.}
\end{corollary}

\begin{proofplan}
A sentence is equivalent modulo \(\mathsf{DLO}\) to a quantifier-free sentence.
The signature has no constant symbols, so there are no closed atomic
formulas; every quantifier-free sentence is a truth-functional combination
of \(\bottom\), and its truth value is fixed by the semantics of the
connectives alone.  Hence any two models of
\(\mathsf{DLO}\) satisfy the same sentences.  Thus the theory is semantically
complete and therefore deductively complete by
\cref{def:theory-completeness}.
If \(\mathcal M\subseteq\mathcal N\) are models of \(\mathsf{DLO}\), they agree
on quantifier-free formulas with parameters from \(M\).  Quantifier elimination
therefore gives agreement on all formulas, so
\(\mathcal M\preccurlyeq\mathcal N\).  By
\cref{lem:model-completeness-reduction}, every embedding between models
of \(\mathsf{DLO}\) is elementary; this is model completeness.
The omitted literal-case verification is the one described in the remark
following \cref{thm:dlo-qe}.
\end{proofplan}

\begin{theorem}[Algebraically closed fields]
For each characteristic \(p\), where \(p=0\) or \(p\) is prime, the theory
\(\mathsf{ACF}_p\) of algebraically closed fields of characteristic \(p\)
eliminates quantifiers in the language of rings.
\textit{Status: cited theorem.}
\end{theorem}

See \cite{Marker2002} for a complete proof.

\section{Categoricity}
\label{sec:qe-categoricity-02}

\begin{definition}[Categoricity]
Let \(T\) be a theory and let \(\kappa\) be an infinite cardinal.  The theory
\(T\) is \emph{\(\kappa\)-categorical} iff \(T\) has a model of cardinality
\(\kappa\) and any two models of \(T\) of cardinality \(\kappa\) are
isomorphic.
\end{definition}

\begin{warning}
Categoricity at one cardinal is not categoricity at every cardinal.  The
L\"owenheim--Skolem theorems imply that a first-order theory with an infinite
model has models in multiple infinite cardinalities, subject to the language
bound.  They do not preclude uniqueness up to isomorphism separately within
each cardinality.  For example, the equality-only theory asserting that there
are at least \(n\) elements for every \(n\geq1\) is
\(\kappa\)-categorical for every infinite cardinal \(\kappa\): any two pure
equality structures of size \(\kappa\) are isomorphic.  A transfer theorem
such as Morley's theorem requires additional hypotheses before categoricity at
one cardinal implies categoricity at others.
\end{warning}

\begin{definition}[Finite partial isomorphism]
A \emph{finite partial isomorphism} from \(\mathcal M\) to \(\mathcal N\) is a
finite injective map \(f:A\to N\), with \(A\subseteq M\), that preserves and
reflects all atomic formulas on tuples from \(A\).
\end{definition}

\begin{method}[Back-and-forth]
\label{method:back-and-forth}
For countable structures, build an isomorphism as the union of an increasing
sequence of finite partial isomorphisms.  At each ``forth'' stage extend the
domain to contain the next enumerated element of the first structure; at each
``back'' stage extend the range to contain the next enumerated element of the
second.  The name denotes this alternating proof method, not a new map or
relation.  It differs from the one-sided extension property in
\cref{def:homogeneity} because it enforces totality and surjectivity together.
The next theorem gives the exact hypotheses, and the dense-order proof then
uses them.
\end{method}

\begin{theorem}[Back-and-forth test]
\label{thm:back-and-forth-test}
Let \(\mathcal M\) and \(\mathcal N\) be countable structures.  Suppose
\(\mathcal F\) is a nonempty set of finite partial isomorphisms from
\(\mathcal M\) to \(\mathcal N\) such that, for every \(f\in\mathcal F\):
\begin{enumerate}
  \item for every \(a\in M\), some \(g\in\mathcal F\) extends \(f\) and has
        \(a\) in its domain; and
  \item for every \(b\in N\), some \(g\in\mathcal F\) extends \(f\) and has
        \(b\) in its range.
\end{enumerate}
Then, for every \(f_0\in\mathcal F\), there is an isomorphism
\(h:\mathcal M\to\mathcal N\) with \(f_0\subseteq h\), extension taken as
sets of ordered pairs; in particular \(\mathcal M\cong\mathcal N\).
\textit{Status: complete proof.}
\end{theorem}

\begin{proof}
Enumerate \(M=\{a_0,a_1,\ldots\}\) and
\(N=\{b_0,b_1,\ldots\}\), and let \(f_0\) be the given member of
\(\mathcal F\).
Recursively use clause 1 at stage \(2n\) to put \(a_n\) into the domain and
clause 2 at stage \(2n+1\) to put \(b_n\) into the range, obtaining an
increasing sequence \(f_0\subseteq f_1\subseteq f_2\subseteq\cdots\) in
\(\mathcal F\).  The union \(h=\bigcup_{k\in\N}f_k\) of the
resulting increasing sequence is a function because the maps extend one
another, and it is injective because every stage is.  The alternating
requirements make it total and surjective.  Every atomic instance involves
only finitely many arguments and is therefore already preserved and reflected
at some stage.  The union is an isomorphism.  Since
\(f_0\subseteq f_1\subseteq\cdots\) and \(h=\bigcup_{k\in\N}f_k\), the
isomorphism extends \(f_0\).
\end{proof}

See \cite{Hodges1993,Marker2002} for further back-and-forth applications.

\begin{corollary}[Countable saturated structures]
Assume the Axiom of Choice in the ambient set theory.
If \(\mathcal M\) is countable and \(\aleph_0\)-saturated, every finite
elementary map from \(\mathcal M\) to itself extends to an automorphism of
\(\mathcal M\).
\textit{Status: complete proof.}
\end{corollary}

\begin{proof}
Let \(f\) be the given finite elementary map, and let \(\mathcal F\) be
the set of finite elementary maps from \(\mathcal M\) to itself that
extend \(f\).  The set is nonempty, because \(f\in\mathcal F\).  Every
member is a finite partial isomorphism from \(\mathcal M\) to
\(\mathcal M\): elementary maps preserve and reflect all atomic formulas,
and they reflect equality, hence are injective.  Clause 1 of
\cref{thm:back-and-forth-test} holds within \(\mathcal F\): given
\(g\in\mathcal F\) and \(a\in M\), \cref{thm:saturation-homogeneity}
extends \(g\) to an elementary map whose domain also contains \(a\), and
this extension still extends \(f\).  Clause 2 holds by applying the same
theorem to \(g^{-1}\), which is elementary because the inverse of a
finite elementary self-map is elementary: extend \(g^{-1}\) to an
elementary map whose domain contains the prescribed element \(b\) and
invert the result, obtaining a member of \(\mathcal F\) that extends
\(g\supseteq f\) and has \(b\) in its range.
\Cref{thm:back-and-forth-test}, applied with \(f_0\defeq f\), yields an
isomorphism \(h:\mathcal M\to\mathcal M\) with \(f\subseteq h\): an
automorphism extending \(f\).
\end{proof}

\begin{theorem}[Countable categoricity of dense linear order]
\label{thm:dlo-countable-categorical}
Any two countable dense linear orders without endpoints are isomorphic.
Consequently, \(\mathsf{DLO}\) is \(\aleph_0\)-categorical.
\textit{Status: complete proof.}
\end{theorem}

\begin{proof}
Let \(\mathcal M\) and \(\mathcal N\) be countable models, with enumerations
\(M=\{m_0,m_1,\ldots\}\) and \(N=\{n_0,n_1,\ldots\}\).  Construct a chain of
finite order-preserving bijections
\[
  f_0\subseteq f_1\subseteq f_2\subseteq\cdots.
\]
This is the back-and-forth method.  At an even stage, take the first \(m_i\)
not in the current domain.  Its order
relations to the finitely many domain points determine a cut in the finite
range.  Density and the absence of endpoints supply an unused point of \(N\)
in the corresponding interval; extend the map by that pair.  At an odd stage,
perform the symmetric construction for the first \(n_i\) not in the range.
The union \(f=\bigcup_{n\in\N} f_n\) is total and surjective by the alternating
construction, and it preserves and reflects order.  Hence it is an
isomorphism.
\end{proof}

\begin{example}[Uncountable categoricity of algebraically closed fields]
Fix \(p=0\) or a prime.  For every uncountable cardinal \(\kappa\), any
algebraically closed field of characteristic \(p\) and cardinality \(\kappa\)
has transcendence degree \(\kappa\) over its prime field.  Transcendence bases
of the same cardinality yield isomorphic algebraic closures.  Thus
\(\mathsf{ACF}_p\) is \(\kappa\)-categorical for every uncountable
\(\kappa\).  It is not \(\aleph_0\)-categorical: countable models can have
different finite or countable transcendence degree.
\end{example}

\begin{theorem}[Ryll--Nardzewski]
Let \(T\) be a deductively complete theory in a countable signature with an infinite
model.  The following are equivalent:
\begin{enumerate}
  \item \(T\) is \(\aleph_0\)-categorical;
  \item for every \(n\geq1\), \(T\) has only finitely many complete
        \(n\)-types over \(\varnothing\); and
  \item for every \(n\geq1\), there are, up to equivalence modulo
        \(T\), only finitely many formulas whose free variables are among
        \(x_1,\ldots,x_n\).
\end{enumerate}
\textit{Status: cited theorem.}
\end{theorem}

See \cite{Hodges1993} for a complete proof.

\begin{theorem}[Morley's categoricity theorem]
Let \(T\) be a deductively complete theory in a countable signature.  If \(T\) is
categorical in one uncountable cardinal, then \(T\) is categorical in every
uncountable cardinal.
\textit{Status: cited theorem.}
\end{theorem}

See \cite{Marker2002} for a complete proof.

\begin{warning}[Completeness versus categoricity]
A deductively complete theory may have many nonisomorphic models of the same
cardinality.  By \cref{def:theory-completeness}, its equivalent semantic
condition says that its models agree on sentences; categoricity says that
models of a specified size are isomorphic.
\end{warning}

\begin{exercise}
\label{ex:qe-categoricity-01}
Prove directly that every countable model of \(\mathsf{DLO}\) is isomorphic to
\((\Q,<)\).
\end{exercise}

\begin{exercise}
\label{ex:qe-categoricity-02}
Let \(T\) eliminate quantifiers, and let \(\mathcal M\subseteq\mathcal N\) be
models of \(T\).  Prove that \(\mathcal M\preccurlyeq\mathcal N\).  Explain
why the assumption that both structures model \(T\) cannot be omitted.
\end{exercise}

\begin{exercise}
\label{ex:qe-categoricity-03}
\label{ex:back-and-forth-equivalence}
Show that the theory of an equivalence relation with infinitely many
classes, all infinite---axiomatized by the equivalence axioms together with,
for every \(n\), the sentences asserting that there are at least \(n\)
pairwise inequivalent elements and that every element is equivalent to at
least \(n\) distinct elements---is
\(\aleph_0\)-categorical.  Give an explicit back-and-forth argument
rather than invoking Ryll--Nardzewski.
\end{exercise}

\part{Computability}
\label{part:computability}

\chapter{Machines and Effective Procedures}
\label{ch:turing-machines}

This part uses deterministic single-tape Turing machines as the primary
formalism for computation.  Other standard formalisms are mentioned only
after the primary definitions have been fixed.  Historical sources are Church and Turing
\cite{Church1936,Turing1937}; systematic references include
\cite{Shoenfield1967,BoolosBurgessJeffrey2007}.

\section{Deterministic single-tape Turing machines}
\label{sec:turing-machines-01}

\begin{definition}[Turing machine]
A \emph{deterministic single-tape Turing machine} is a tuple
\[
  \mathcal T=(Q,\mathcal A_{\mathrm{in}},\mathcal A_{\mathrm{tape}},
  \delta,q_0,q_{\mathrm{acc}},q_{\mathrm{rej}})
\]
such that:
\begin{enumerate}
  \item \(Q\) is a finite set of states;
  \item \(\mathcal A_{\mathrm{in}}\) is a finite input alphabet not containing
        the blank symbol \(\sqcup\);
  \item \(\mathcal A_{\mathrm{tape}}\) is a finite tape alphabet with
        \(\mathcal A_{\mathrm{in}}\cup\{\sqcup\}
        \subseteq\mathcal A_{\mathrm{tape}}\);
  \item \(q_0,q_{\mathrm{acc}},q_{\mathrm{rej}}\in Q\), with
        \(q_{\mathrm{acc}}\neq q_{\mathrm{rej}}\); and
  \item the transition function is
  \[
    \delta:(Q\setminus\{q_{\mathrm{acc}},q_{\mathrm{rej}}\})
    \times\mathcal A_{\mathrm{tape}}
    \longrightarrow Q\times\mathcal A_{\mathrm{tape}}\times\{L,R\}.
  \]
\end{enumerate}
The two \(\mathcal A\)-symbols name components of this machine tuple only;
they are not formula contexts or arithmetical-hierarchy class names.
The tape is infinite in both directions, has only finitely many nonblank cells
at every finite time, and the head scans one cell.
\end{definition}

\begin{definition}[Configuration and computation]
A \emph{configuration} of \(\mathcal T\) is a tuple
\[
  C=(q,T,h),
\]
where \(q\in Q\), \(h\in\Z\), and
\(T:\Z\to\mathcal A_{\mathrm{tape}}\) has finite nonblank support:
\[
  \set{i\in\Z:T(i)\ne\sqcup}
  \quad\text{is finite}.
\]
Suppose \(q\notin\set{q_{\mathrm{acc}},q_{\mathrm{rej}}}\) and
\[
  \delta(q,T(h))=(q',a,d).
\]
Let \(T':\Z\to\mathcal A_{\mathrm{tape}}\) be given by
\[
  T'(i)=
  \begin{cases}
    a,&i=h,\\
    T(i),&i\ne h,
  \end{cases}
\]
and let \(h'=h-1\) when \(d=L\), while \(h'=h+1\) when \(d=R\).  Then, and
only then,
\[
  (q,T,h)\longrightarrow_{\mathcal T}(q',T',h').
\]
A configuration in either halting state has no successor.  Write
\(\longrightarrow_{\mathcal T}^{*}\) for the reflexive transitive closure of
this one-step relation.

For an input word \(w\in\mathcal A_{\mathrm{in}}^*\) of length \(n\), define
\[
  T_w(i)=
  \begin{cases}
    w(i),&0\leq i<n,\\
    \sqcup,&\text{otherwise}.
  \end{cases}
\]
The initial configuration is \((q_0,T_w,0)\).  Thus on empty input the head is
at cell \(0\), which contains a blank.  A computation on \(w\) is the unique
maximal finite or infinite sequence beginning with this initial configuration
and following \(\longrightarrow_{\mathcal T}\) at each successor step.  It
\emph{halts} iff the sequence reaches either halting state, and it
\emph{accepts} or \emph{rejects} according to which halting state it reaches.
\end{definition}

Fix the canonical binary word \(\operatorname{bin}(n)\) for each
\(n\in\N\): \(\operatorname{bin}(0)\) is the one-symbol word \(0\), and for
\(n>0\) it is the ordinary base-two expansion without leading zeros, with the
most significant bit first.  Use the numerical pairing map fixed in the
finite-string chapter:
\[
  \operatorname{pair}_{\N}(a,b)
  =\frac{(a+b)(a+b+1)}2+b.
\]
Its arguments and value are natural numbers; it is distinct in notation and
type from the set-theoretic ordered pair \(\tuple{a,b}\).  For each fixed
arity \(k\), define an injective tuple code by
\[
 \operatorname{code}_0()\defeq0,
 \qquad
 \operatorname{code}_1(n)\defeq n,
\]
and, for \(k\geq2\),
\[
 \operatorname{code}_k(n_0,\ldots,n_{k-1})
 \defeq
 \operatorname{pair}_{\N}\bigl(n_0,
   \operatorname{pair}_{\N}(n_1,\ldots,
     \operatorname{pair}_{\N}(n_{k-2},n_{k-1})\ldots)\bigr).
\]
The arity is fixed by the type \(\N^k\), so it need not be stored inside the
code.  A numerical input for a \(k\)-ary function is the word
\(\operatorname{bin}(\operatorname{code}_k(\bar n))\), placed with its most
significant bit in cell \(0\).

For machines used to compute numerical functions, the input alphabet contains
the symbols \(0,1\).  If a halting configuration has tape function \(T\), let
\[
  m\defeq\min\set{i\in\N:T(i)=\sqcup}.
\]
Finite nonblank support makes \(m\) defined.  The tape prefix of length \(m\)
is a \emph{valid numerical output} iff it is exactly
\(\operatorname{bin}(y)\) for a unique \(y\in\N\).  Thus the empty word, a
word containing a symbol other than \(0,1\), and a nonzero binary word with a
leading zero are invalid.  Cells outside this prefix and the final head
position do not affect the output.  A machine used to compute a numerical
function is required to halt in \(q_{\mathrm{acc}}\) with a valid numerical
output.

\begin{definition}[Partial computable function]
\label{def:partial-computable}
A partial function \(f:\N^k\rightharpoonup\N\) is \emph{partial computable}
iff some Turing machine satisfies both of the following conditions for every
\(\bar n\in\N^k\): when \(f(\bar n)\) is defined, its computation on
\(\operatorname{bin}(\operatorname{code}_k(\bar n))\) reaches
\(q_{\mathrm{acc}}\) with output word \(\operatorname{bin}(f(\bar n))\); and
when \(f(\bar n)\) is undefined, that computation is infinite.  A total
partial computable function is called \emph{computable}.
\end{definition}

The arrow \(\rightharpoonup\) permits undefined values.  Divergence is not a
third numerical output.

The fixed binary and tuple encodings have computable encoding and decoding
maps.  For pairing, given \(z\), boundedly search for the greatest \(s\leq z\)
with \(s(s+1)/2\leq z\), put
\(b=z-s(s+1)/2\), and put \(a=s-b\).  When the arity \(k\) is supplied,
iterating these projections decodes the \(k\) tuple entries; the code alone
does not encode \(k\).  Repeated division by two produces the binary
word, and the left-to-right recurrence \(r\mapsto2r+b\) decodes it.  Each
description is a terminating Turing-machine procedure, so the uses of
``code'' in the definition above do not add an unexamined computation model.

\begin{example}[Unary successor]
A machine computing \(S(n)=n+1\) reads the arity-one input word
\(\operatorname{bin}(n)\), since \(\operatorname{code}_1(n)=n\).  On that block it
scans right to the end, propagates a carry left through consecutive \(1\)s,
 and writes \(1\) at the first \(0\) encountered.  If the carry passes the most
 significant bit of an \(m\)-bit input, all old bits are now \(0\).  The
 machine shifts that \(m\)-cell zero block one cell to the right, writes \(1\)
 in cell \(0\), and leaves a blank in cell \(m+1\).  Thus the output prefix in
 cells \(0,\ldots,m\) changes \(111\) to \(1000\).  It then treats that prefix as the output
word, not as another tuple code, because the output convention requires the
binary word for the numerical value.  The loop invariant is that every
processed suffix is the correct low-order suffix of \(n+1\), while the state
records whether a carry remains.  Decoding and block relocation are justified
by the computable coding operations established above.
\end{example}

\begin{warning}[Machine formalism versus complexity measure]
Changing from one tape to several tapes does not change which partial functions
are computable, but it can change running time by a polynomial factor.  A
computability claim and a complexity claim therefore require different
invariance arguments; see \cite{BoolosBurgessJeffrey2007} for the
multi-tape invariance argument.
\end{warning}

\section{Primitive recursion}
\label{sec:turing-machines-02}

\begin{definition}[Primitive recursive functions]
The class of \emph{primitive recursive functions} is the least class of total
functions on finite powers of \(\N\) containing:
\begin{enumerate}
  \item for every \(k\in\N\), the zero function
        \(Z_k:\N^k\to\N\) given by \(Z_k(\bar n)=0\), including
        \(Z_0()=0\);
  \item the successor function \(S(n)=n+1\); and
  \item every projection \(P_i^k(n_1,\ldots,n_k)=n_i\), for
        \(k\geq1\) and \(1\leq i\leq k\),
\end{enumerate}
and closed under the following two operations.  Composition forms
\(f:\N^k\to\N\) from \(g:\N^m\to\N\) and
\(h_1,\ldots,h_m:\N^k\to\N\) by
\[
  f(\bar x)=g(h_1(\bar x),\ldots,h_m(\bar x)).
\]
For every \(k\in\N\), primitive recursion forms
\(f:\N^{k+1}\to\N\) from \(g:\N^k\to\N\) and
\(h:\N^{k+2}\to\N\) by
\[
  f(\bar x,0)=g(\bar x),
  \qquad
  f(\bar x,n+1)=h(\bar x,n,f(\bar x,n)).
\]
\end{definition}

\begin{definition}[Primitive recursive relation]
A relation \(R\subseteq\N^k\) is \emph{primitive recursive} iff its
characteristic function \(\chi_R:\N^k\to\{0,1\}\), with value \(1\) exactly
on tuples in \(R\), is primitive recursive.  This term classifies a relation,
not a representing arithmetic formula.  It is required below because the
coding predicates for formulas and derivations are relations, whereas the
preceding definition applies only to functions.
\end{definition}

\begin{theorem}
Every primitive recursive function is computable and total.
\textit{Status: proof sketch.}
\end{theorem}

\begin{proofplan}
The initial functions have direct machines.  Turing machines are closed under
composition by running machines sequentially while retaining the required
inputs.  For primitive recursion, store \(\bar x\), a counter \(i\), and the
current value \(f(\bar x,i)\).  Initialize with the machine for \(g\), and run
the machine for \(h\) exactly \(n\) times while incrementing \(i\).  The loop
has a prescribed finite number of iterations, so it halts.  Induction on the
generation of the primitive recursive class proves the claim.
The omitted verification gives concrete transition tables and invariants for
the closure constructions.
\end{proofplan}

\begin{example}[Addition and multiplication]
Addition and multiplication are primitive recursive because
\[
\begin{aligned}
  x+0&=x, & x+(y+1)&=S(x+y),\\
  x\times0&=0, & x\times(y+1)&=(x\times y)+x.
\end{aligned}
\]
The second definition legitimately uses addition because addition was
generated at an earlier stage.
\end{example}

\begin{remark}
The Ackermann function is computable and total but not primitive recursive.
Thus primitive recursion is a proper subclass of total computability.  The
proof requires a domination argument and is omitted here; see
\cite{BoolosBurgessJeffrey2007}.
\end{remark}

\section{Decidable and computably enumerable sets}
\label{sec:turing-machines-03}

\begin{definition}[Decidable set]
A set \(A\subseteq\N\) is \emph{decidable} iff its characteristic function
\(\chi_A:\N\to\{0,1\}\) is computable.
\end{definition}

\begin{definition}[Computably enumerable set]
A set \(A\subseteq\N\) is \emph{computably enumerable} iff it is the domain
of a partial computable function.
\end{definition}

\begin{definition}[Decidable and computably enumerable relations]
For \(R\subseteq\N^k\), the relation \(R\) is \emph{decidable} iff its
characteristic function is computable.  It is \emph{computably enumerable}
iff
\[
  \set{\operatorname{code}_k(\bar n):\bar n\in R}
\]
is a computably enumerable subset of \(\N\).  These terms classify relations;
for \(k=1\) they reproduce the preceding definitions for sets.  The extension
to relations is needed by the computation-history and proof predicates below.
\end{definition}

\begin{definition}[Machine decision and recognition]
Let \(A\subseteq\N\).  A Turing machine \emph{decides} \(A\) iff it halts on
every unary numerical input \(\operatorname{bin}(n)\), accepts exactly those inputs
whose decoded number \(n\) belongs to \(A\), and rejects exactly those whose
decoded number lies outside \(A\).  A Turing machine \emph{recognizes} \(A\)
iff it accepts exactly the unary numerical inputs whose decoded numbers belong
to \(A\), with no halting requirement on inputs outside \(A\).  A machine with
the first property is a \emph{decider} for \(A\); a machine with the second is
a \emph{recognizer} for \(A\).

These terms classify machine behavior, not sets or numerical functions.  They
are required to state the machine forms of decidability and computable
enumerability without ambiguity.  Every decider is a recognizer, but a
recognizer may diverge outside \(A\), so the two terms are not redundant.  The
next lemma and theorem are their uses.
\end{definition}

\begin{lemma}[Machine form of decidability]
A set \(A\subseteq\N\) is decidable iff some Turing machine decides \(A\).
\textit{Status: complete proof.}
\end{lemma}

\begin{proof}
Given a machine computing \(\chi_A\), simulate it and accept when its output is
\(1\), while rejecting when its output is \(0\).  Totality makes this machine
halt on every input.  Conversely, simulate a decider and write \(1\) after an
accepting halt and \(0\) after a rejecting halt.  The resulting machine
computes \(\chi_A\).
\end{proof}

\begin{theorem}[Equivalent forms of computable enumerability]
For \(A\subseteq\N\), the following are equivalent:
\begin{enumerate}
  \item \(A\) is computably enumerable;
  \item some Turing machine recognizes \(A\);
  \item either \(A=\varnothing\), or \(A\) is the range of a computable
        function \(e:\N\to\N\); and
  \item there is a decidable relation \(R(x,k)\) such that
        \(x\in A\) iff \(\exists k\,R(x,k)\).
\end{enumerate}
\textit{Status: proof sketch.}
\end{theorem}

\begin{proofplan}
Convert a domain machine into a recognizer.  Dovetail all computations to
enumerate the accepted inputs.  Record a finite accepting computation as the
witness \(k\) for a decidable computation-history relation.  Conversely,
search successive \(k\) for a witness to \(R(x,k)\).  The empty-set exception
in the range formulation is necessary because a total function on \(\N\) has
nonempty range.
\end{proofplan}

\begin{theorem}
A set \(A\subseteq\N\) is decidable iff both \(A\) and
\(\N\setminus A\) are computably enumerable.
\textit{Status: complete proof.}
\end{theorem}

\begin{proof}
If \(A\) is decidable, a decider recognizes \(A\), and reversing its
two halting answers gives a recognizer for the complement.  Conversely, run a
recognizer for \(A\) and a recognizer for its complement in dovetailing fashion.
Exactly one eventually accepts, and its identity determines membership in
\(A\).
\end{proof}

\begin{exercise}
\label{ex:turing-machines-01}
Prove that the decidable sets are closed under complement, finite union, and
finite intersection.  Prove that computably enumerable sets are closed under
finite union and finite intersection, but do not use closure under complement.
\end{exercise}

\chapter{Universal Computation, Reductions, and Fixed Points}
\label{ch:universal-reductions}

\section{An effective numbering of programs}
\label{sec:universal-reductions-01}

Normalize every machine description as follows.  A description is
\emph{normalized} iff its tape symbols are the numbers
\(0,\ldots,\abs{\mathcal A_{\mathrm{tape}}}-1\) with the designated
assignments \(\sqcup\mapsto0\), digit \(0\mapsto1\), digit \(1\mapsto2\),
and its states are the numbers \(0,\ldots,\abs{Q}-1\) with the designated
assignments \(q_0\mapsto0\), \(q_{\mathrm{acc}}\mapsto1\),
\(q_{\mathrm{rej}}\mapsto2\); the remaining numbers are assigned
arbitrarily, and the row list is the lexicographic listing, by numbered key
\((q,a)\), of the full transition table over its non-halting domain.  (A
clause fixing an ``order of first appearance'' instead would be circular:
the lexicographic row order presupposes the numbering of states and
symbols.  Without the designated assignments, ``valid code'' below would
not be well defined.)  Encode a table tuple of length \(k\) as
\(\operatorname{pair}_{\N}(k,\operatorname{code}_k(\bar n))\).  No
canonical machine-to-code map is fixed; none is needed, since indices are
not required to be unique.  A natural number that is not the code of a
valid normalized numerical-machine description denotes a machine that
diverges on every input.  Such a description must in particular satisfy
\(\{0,1\}\subseteq\mathcal A_{\mathrm{in}}\), read at the designated digit
numbers \(1,2\) together with the input and output conventions fixed
above; a normalized description whose input alphabet omits either binary
symbol is treated as invalid by this numbering.  In a valid code the three
designated states are pairwise distinct; a machine whose start state is
halting, which the machine definition permits by requiring only
\(q_{\mathrm{acc}}\neq q_{\mathrm{rej}}\), is captured up to behavior by
the surrogate-start step in the proof of \cref{thm:enumeration} below.
For a valid code \(e\), define
\(\varphi_e(x)\) to be the numerical output when that machine accepts with a
valid numerical output, and leave \(\varphi_e(x)\) undefined when it rejects,
diverges, or halts without a valid numerical output.  A fixed wrapper simulates the decoded
machine, copies a canonical numerical output after a valid accepting halt, and
enters an infinite loop after either rejection or an accepting halt without a
valid numerical output.

The number \(e\) is an \emph{index} for \(\varphi_e\).  Different indices may
compute the same partial function.  For a finite tuple \(\bar x\), the notation
\(\varphi_e(\bar x)\), with \(\bar x\in\N^k\), denotes the behavior of indexed
program \(e\) on the input word
\(\operatorname{bin}(\operatorname{code}_k(\bar x))\).  This is
arity-dependent application notation for the one fixed program numbering, not
a second numbering and not a second application of the unary tuple code.

For expressions denoting possibly undefined values, \(a\simeq b\) means that
both sides are undefined or both are defined and equal.  For partial
functions \(f,g:\N\rightharpoonup\N\), the notation \(f\simeq g\)
abbreviates: \(f(x)\simeq g(x)\) for every \(x\in\N\).

\begin{lemma}[Renaming invariance]
\label{lem:renaming-invariance}
Let \(\mathcal T\) be a Turing machine used to compute a numerical function
and let \(\iota=(\iota_Q,\iota_A)\) bijectively rename its states and tape
symbols, preserving the roles of
\(q_0,q_{\mathrm{acc}},q_{\mathrm{rej}}\) and of \(\sqcup\) and the digits
\(0,1\).  Write \(\mathcal T^{\iota}\) for the machine obtained by applying
\(\iota\) to every component of the machine tuple, and \(C^{\iota}\) for the
corresponding renaming of a configuration \(C\).  Then
\(C\longrightarrow_{\mathcal T}C'\) iff
\(C^{\iota}\longrightarrow_{\mathcal T^{\iota}}(C')^{\iota}\).  Hence, on
every pair of input words corresponding letter by letter under \(\iota_A\),
the machines \(\mathcal T\) and \(\mathcal T^{\iota}\) accept, reject, or
diverge together, with corresponding tapes at every step.
\textit{Status: complete proof.}
\end{lemma}

\begin{proof}
For the one-step claim, read the configuration definition on both sides:
\(\delta(q,T(h))=(q',a,d)\) holds in \(\mathcal T\) iff
\(\delta^{\iota}(\iota_Q(q),\iota_A(T(h)))
=(\iota_Q(q'),\iota_A(a),d)\) holds in \(\mathcal T^{\iota}\), since
\(\delta^{\iota}\) is by construction the \(\iota\)-conjugate of \(\delta\);
the rewritten cell and the moved head then correspond position by position.
The claim for whole computations follows by induction on the number of
steps.  Initial configurations correspond because \(\iota_A\) preserves
\(\sqcup\), and halting, acceptance, and rejection correspond because
\(\iota_Q\) preserves the designated states.  Valid numerical outputs
correspond because \(\iota_A\) preserves the digits.
\end{proof}

\begin{theorem}[Adequacy of the numbering: enumeration theorem]
\label{thm:enumeration}
\begin{enumerate}[label=(\alph*)]
\item Every \(\varphi_e\) is partial computable.
\item For every \(k\geq1\) and every partial computable
      \(f:\N^k\rightharpoonup\N\) there is an index \(e\) with
      \(\varphi_e(\bar x)\simeq f(\bar x)\) for all \(\bar x\in\N^k\),
      under the arity-\(k\) application convention fixed above.
\end{enumerate}
\textit{Status: complete proof.}
\end{theorem}

\begin{proof}
(a) For a valid code \(e\), the fixed wrapper described in the numbering
paragraph simulates the decoded machine, accepts with the canonical output
word exactly when that machine accepts with a valid numerical output, and
otherwise runs forever; for an invalid code, \(\varphi_e\) is the
everywhere-undefined function, which a machine with an immediate infinite
loop computes.  In both cases \(\varphi_e\) satisfies
\cref{def:partial-computable}.

(b) Let \(\mathcal T\) witness that \(f\) is partial computable in the sense
of \cref{def:partial-computable}.  First,
\(\{0,1\}\subseteq\mathcal A_{\mathrm{in}}\) holds automatically:
\(\operatorname{code}_1\) is the identity and, for \(k\geq2\),
\(\operatorname{code}_k\) is built from the pairing bijection, so
\(\operatorname{code}_k\) is surjective for every \(k\geq1\) and the input
words \(\operatorname{bin}(\operatorname{code}_k(\bar x))\) exhaust all
binary numerals.  An initial configuration exists only for a word over
\(\mathcal A_{\mathrm{in}}\), and \cref{def:partial-computable} demands an
accepting or infinite computation of \(\mathcal T\) on each of these words;
since some required numeral contains both digits, both digits lie in
\(\mathcal A_{\mathrm{in}}\), and no alphabet extension is needed.

Second, since the machine definition requires only
\(q_{\mathrm{acc}}\neq q_{\mathrm{rej}}\), the start state of
\(\mathcal T\) may coincide with a halting state.  In that case replace it
by a fresh start state \(q_0'\) with \(\delta(q_0',a)=(q_0,a,R)\) for every
tape symbol \(a\): the one extra step leaves the tape unchanged, and the
output convention ignores the final head position and the cells outside the
output prefix, so the accepted outputs and the divergence behavior are
unchanged.  After this step the three designated states are pairwise
distinct, as are the three designated symbols \(\sqcup,0,1\).

Now choose any bijective renaming \(\iota\) of states and tape symbols onto
the designated-slot numbering that preserves the six designated roles, and
let \(\mathcal T^*\) be the resulting normalized machine;
\cref{lem:renaming-invariance} makes \(\mathcal T^*\) behave exactly as
\(\mathcal T\) on corresponding inputs, and the binary numerals correspond
to themselves read at the designated digit numbers \(1,2\).  List the rows
of \(\mathcal T^*\) lexicographically by numbered key \((q,a)\); the row
list is the full graph of the transition function, which is total on
non-halting pairs by the machine definition.  Encode the resulting
description tuple of length \(k'\) as
\(\operatorname{pair}_{\N}(k',\operatorname{code}_{k'}(\bar n))\), obtaining
a number \(e\).  This \(e\) is a valid code by construction.  On each input
word \(\operatorname{bin}(\operatorname{code}_k(\bar x))\), the decoded
machine \(\mathcal T^*\) accepts with output word
\(\operatorname{bin}(f(\bar x))\) when \(f(\bar x)\) is defined and runs
forever when \(f(\bar x)\) is undefined, so
\(\varphi_e(\bar x)\simeq f(\bar x)\) for all \(\bar x\in\N^k\).
\end{proof}

Kleene's treatment of enumeration and normal form over standard numberings
\cite{Kleene1952} is parallel reading; the in-document proof is primary
because this numbering is bespoke: its normalization, its validity
constraint \(\{0,1\}\subseteq\mathcal A_{\mathrm{in}}\), and its wrapper are
fixed here rather than taken from the cited source.

\begin{theorem}[Universal computation]
\label{thm:universal-computation}
There is a partial computable function \(U(e,x)\) such that
\[
  U(e,x)\simeq\varphi_e(x)
\]
for all \(e,x\in\N\).
\textit{Status: proof sketch.}
\end{theorem}

\begin{proofplan}
Decode \(e\) as a machine description.  If decoding fails, diverge.  Otherwise
store an encoded configuration and repeatedly apply the machine's finite
 transition table.  Return an output only when the simulated machine reaches
its accepting state with a valid numerical output; copy that output into
canonical form.  If it rejects or accepts without a valid numerical output,
enter an
 infinite loop.  Each simulation step is effective, and every finite computation is
eventually reproduced.  Turing's original universal-machine construction is
the historical source \cite{Turing1937}.
\end{proofplan}

\begin{theorem}[Parameter theorem]
\label{thm:parameter}
For every \(m,n\geq1\), there is a total computable function
\(\mathsf{par}_n^m:\N^{m+1}\to\N\) such that, for every index \(e\), every \(\bar a\in\N^m\), and every
\(\bar x\in\N^n\),
\[
  \varphi_{\mathsf{par}_n^m(e,\bar a)}(\bar x)
  \simeq \varphi_e(\bar a,\bar x).
\]
\textit{Status: cited theorem.}
\end{theorem}

From \(e\) and the constants \(\bar a\), construct a normalized transition
table for a machine that decodes its runtime input as the \(n\)-tuple
\(\bar x\), computes
\(\operatorname{code}_{m+n}(\bar a,\bar x)\), writes the binary word for that code,
and then simulates program \(e\).  The established tuple encoding and decoding
maps, state renumbering, and finite table concatenation are computable, so the
map from \((e,\bar a)\) to the new index is total and computable.  The omitted
verification writes the compiler's transition table and proves its simulation
invariant.  See \cite{Kleene1952} for a complete proof of the parameter
theorem for effective program numberings.

\section{Many-one reductions}
\label{sec:universal-reductions-02}

\begin{definition}[Many-one reduction]
For \(A,B\subseteq\N\), a \emph{many-one reduction} from \(A\) to \(B\) is a
total computable function \(f:\N\to\N\) satisfying
\[
  x\in A\quad\Longleftrightarrow\quad f(x)\in B
  \qquad\text{for every }x\in\N.
\]
Write \(A\leq_m B\) when such an \(f\) exists.
\end{definition}

\begin{lemma}
If \(A\leq_m B\) and \(B\) is decidable, then \(A\) is decidable.  If
\(A\leq_m B\) and \(B\) is computably enumerable, then \(A\) is computably
enumerable.
\textit{Status: complete proof.}
\end{lemma}

\begin{proof}
Compute \(f(x)\).  In the first case run the decider for \(B\); in the second
case run the recognizer for \(B\).  The defining biconditional transfers the
answer or acceptance condition to \(A\).
\end{proof}

\begin{warning}[Direction of a reduction]
To prove that \(B\) is at least as hard as \(A\), construct
\(A\leq_m B\).  Reversing the reduction does not establish the desired lower
bound.
\end{warning}

\section{The self-application domain problem}
\label{sec:universal-reductions-03}

Define the \emph{self-application domain set}
\[
  K=\set{e\in\N:\varphi_e(e)\text{ is defined}}.
\]

\begin{theorem}[Undecidability of the self-application domain set]
\label{thm:halting-undecidable}
The set \(K\) is computably enumerable but not decidable.
\textit{Status: complete proof.}
\end{theorem}

\begin{proof}
The universal machine recognizes \(K\): simulate \(U(e,e)\) and accept if it
returns a value.  Suppose, for contradiction, that \(K\) is decidable,
and let \(h\) be its total computable characteristic function.  Define a
partial computable function
\[
  d(e)\simeq
  \begin{cases}
    0, & h(e)=0,\\
    \text{undefined}, & h(e)=1.
  \end{cases}
\]
By \cref{thm:enumeration}, fix an index \(k\) with \(\varphi_k\simeq d\).
If \(k\in K\), then \(h(k)=1\), so
\(\varphi_k(k)\) is undefined and \(k\notin K\).  If \(k\notin K\), then
\(h(k)=0\), so \(\varphi_k(k)=0\) and \(k\in K\).  Both alternatives are
contradictions.
\end{proof}

\begin{corollary}
The complement \(\N\setminus K\) is not computably enumerable.
\end{corollary}

\begin{proof}
If both \(K\) and its complement were computably enumerable, then \(K\)
would be decidable.
\end{proof}

\begin{example}[Reduction to the binary application domain set]
Let
\[
  H=\set{\operatorname{pair}_{\N}(e,x):e,x\in\N
       \text{ and }\varphi_e(x)\text{ is defined}}.
\]
The map \(e\mapsto\operatorname{pair}_{\N}(e,e)\) reduces \(K\) to \(H\), so
\(H\) is undecidable.  Conversely, decode an input to \(H\) as
\(\operatorname{pair}_{\N}(e,x)\).  Fix an index \(b\) satisfying
\[
 \varphi_b(e,x,y)\simeq\varphi_e(x)
 \qquad\text{for all }e,x,y\in\N;
\]
such an index exists because omission of the third input followed by
universal simulation is partial computable.  Put
\(q=\mathsf{par}_1^2(b,e,x)\).  The parameter theorem makes
\((e,x)\mapsto q\) total computable, and
\(\varphi_q(q)\simeq\varphi_e(x)\).  Thus \(\varphi_q(q)\) is defined iff
\(\varphi_e(x)\) is defined, so this map reduces \(H\) to \(K\).
\end{example}

\section{The recursion theorem}
\label{sec:universal-reductions-04}

\begin{theorem}[Kleene recursion theorem]
\label{thm:recursion-theorem}
If \(f:\N\to\N\) is total computable, then there is an index \(e\) such that
\[
  \varphi_e\simeq\varphi_{f(e)}.
\]
\textit{Status: proof sketch.}
\end{theorem}

\begin{proofplan}
By universal computation and the parameter theorem, there is a total
computable function \(p\) such that
\[
  \varphi_{p(x)}(y)\simeq\varphi_{\varphi_x(x)}(y),
\]
with the right side undefined when \(\varphi_x(x)\) is undefined.  By
\cref{thm:enumeration}, let \(a\)
be an index of the total computable function \(x\mapsto f(p(x))\).  Then
\(\varphi_a(a)=f(p(a))\).  Set \(e=p(a)\).  We obtain
\[
  \varphi_e
  =\varphi_{p(a)}
  \simeq\varphi_{\varphi_a(a)}
  =\varphi_{f(p(a))}
  =\varphi_{f(e)}.
\]
This completes the fixed-point calculation using the cited universal and
parameter theorems.  The direct construction of the program transformation
\(p\) from the fixed encodings is omitted; this omission is why the recursion
theorem retains proof-sketch status.
\end{proofplan}

\begin{example}[Self-reference without a self-inspection primitive]
Suppose a total computable transformation \(f\) adds a fixed output header to
the behavior of a program whose index it receives.  The recursion theorem
supplies a program \(e\) whose behavior is exactly that of the transformed
program \(f(e)\).  The theorem does not assert \(e=f(e)\); it asserts equality
of the computed partial functions.
\end{example}

\section{Rice's theorem}
\label{sec:universal-reductions-05}

\begin{definition}[Extensional index property]
A set \(A\subseteq\N\) is \emph{extensional} iff
\[
  \varphi_e\simeq\varphi_d
  \quad\Longrightarrow\quad
  (e\in A\Longleftrightarrow d\in A)
  \qquad\text{for all }e,d\in\N.
\]
It is \emph{nontrivial} iff it is neither empty nor all of \(\N\).
\end{definition}

\begin{theorem}[Rice]
\label{thm:rice}
Every nontrivial extensional set of indices of unary partial computable
functions is undecidable.
\textit{Status: complete proof.}
\end{theorem}

\begin{proof}
Let \(A\) be nontrivial and extensional.  First suppose an index of the
everywhere-undefined function does not belong to \(A\).  Choose \(a\in A\).
Fix an index \(b\) for the binary partial computable function
\[
 F(e,x)\simeq
 \begin{cases}
   \varphi_a(x),&\text{if }\varphi_e(e)\text{ returns a value},\\
   \text{undefined},&\text{otherwise}.
 \end{cases}
\]
It is partial computable by first simulating \(U(e,e)\) and, only after that
simulation returns, simulating \(U(a,x)\); \cref{thm:enumeration} then
supplies the index \(b\).  Define the total computable map
\(g(e)=\mathsf{par}_1^1(b,e)\).  The parameter theorem gives
\(\varphi_{g(e)}(x)\simeq F(e,x)\) for every \(x\).  If \(e\notin K\), the
new program is everywhere undefined, so \(g(e)\notin A\).  If \(e\in K\), the
new program computes \(\varphi_a\), so \(g(e)\in A\).  Hence
\(K\leq_m A\), and \(A\) is undecidable.

If the everywhere-undefined function has an index in \(A\), apply the first
case to the nontrivial extensional complement \(\N\setminus A\).  Decidability
of \(A\) would imply decidability of its complement, which is impossible.
\end{proof}

\begin{warning}
Rice's theorem concerns semantic properties of the partial function computed
by a program.  Syntactic properties such as ``the source contains ten states''
can be decidable and are not extensional.
\end{warning}

\begin{exercise}
\label{ex:universal-reductions-01}
Use Rice's theorem to prove undecidable each of the following index sets:
programs computing a total function; programs whose range contains \(0\); and
programs computing a constant partial function, where this last phrase means a
partial function with nonempty domain that takes one fixed value on every
member of its domain.  Verify nontriviality in each case.
\end{exercise}

\chapter{Logical Decision Problems and the Arithmetical Hierarchy}
\label{ch:decision-hierarchy}

\section{The classical decision problem}
\label{sec:decision-hierarchy-01}

Fix an effective coding of first-order formulas and finite derivations.

\begin{theorem}[Negative solution to the classical Entscheidungsproblem]
\label{thm:fol-validity-undecidable}
Under the fixed effective coding, no Turing machine decides the set of
codes of logically valid first-order sentences.
\textit{Status: cited theorem.}
\end{theorem}

Church and Turing supply the historical negative solution
\cite{Church1936,Turing1937}.  The theorem records only that historically
sourced general decision problem.

\begin{corollary}[Enumerability without decidability]
Under the fixed effective coding, the set of logically valid first-order
sentences is computably enumerable but not decidable.
\textit{Status: complete proof.}
\end{corollary}

\begin{proof}
Enumerate all finite derivations in the fixed effective first-order
calculus that have no open assumptions and whose conclusion is a sentence,
and output their conclusions.  Soundness for the empty premise set ensures
that every output is valid, and completeness ensures that every valid
sentence eventually appears as such a conclusion, so the
set is computably enumerable.  A decider for this set would solve the general
decision problem forbidden by \cref{thm:fol-validity-undecidable}.
\end{proof}

\begin{warning}[Restricted signatures]
The theorem concerns the general decision problem.  The monadic fragment
without function symbols is decidable, so undecidability may not be transferred
to an individual fixed signature or fragment without a theorem controlling
that restriction.
\end{warning}

\begin{theorem}[Trakhtenbrot]
For every finite relational signature containing a binary relation symbol, the
set of sentences having a finite model is computably enumerable but
undecidable, and the set of sentences true in every finite structure is not
computably enumerable.
\textit{Status: cited theorem.}
\end{theorem}

See \cite{Trakhtenbrot1950} for a complete proof.

The contrast with unrestricted validity is important: unrestricted
validities are computably enumerable by completeness of the fixed effective
first-order calculus, while finite validities admit no effective calculus
both sound and complete for truth in all finite
structures.

\section{Quantifier complexity over arithmetic}
\label{sec:decision-hierarchy-02}

\begin{definition}[Language and standard structure of arithmetic]
\label{def:language-arithmetic}
The \emph{language of arithmetic} is
\[
  \Lang_A=\{0,S,+,\times\},
\]
where \(0\) is a constant, \(S\) is unary, and \(+\) and \(\times\) are
binary function symbols.  The \emph{standard structure of arithmetic} is
\[
  \mathcal N=(\N,0,S,+,\times)
\]
with the usual operations.  These objects supply the syntax and semantics for
every hierarchy classification below.
\end{definition}

For \(n\in\N\), the \emph{numeral} for \(n\) is the closed term
\(\bar n=S^n0\).  For \(\bar a=(a_1,\ldots,a_k)\), write
\(\overline{\bar a}=(\bar a_1,\ldots,\bar a_k)\).  The bar is
metalinguistic: it maps a natural number or finite tuple of natural numbers to
an arithmetic term or tuple of arithmetic terms.  When a bar sits over a
single letter carrying a numeric subscript, as in \(\bar a_1\), it denotes
the numeral of that tuple entry; a bar is never applied to a subscripted
tuple name.

For arithmetic terms \(r,t\), write
\(r<t\) as the metalinguistic abbreviation
\(\exists z\,(r+Sz=t)\), where \(z\) is the least variable absent from
\(r,t\).  If \(x\notin\operatorname{Var}(t)\), the marked bounded quantifiers
abbreviate
\[
\begin{aligned}
  \exists x<t\,\varphi &\defeq \exists x\,(x<t\land\varphi),\\
  \forall x<t\,\varphi &\defeq \forall x\,(x<t\to\varphi).
\end{aligned}
\]
Their satisfaction clauses therefore quantify over exactly the natural
numbers smaller than the value of \(t\).  The hierarchy counts each marked
bounded quantifier as bounded before this eliminable expansion.

Alongside \(<\), write \(r\leq t\) as the metalinguistic abbreviation
\(\exists z\,(z+r=t)\), where \(z\) is the least variable absent from
\(r,t\).  If \(x\notin\operatorname{Var}(t)\), the marked bounded
quantifiers
\[
\begin{aligned}
  \exists x\leq t\,\varphi &\defeq \exists x\,(x\leq t\land\varphi),\\
  \forall x\leq t\,\varphi &\defeq \forall x\,(x\leq t\to\varphi)
\end{aligned}
\]
are counted as bounded exactly as their \(<\)-marked counterparts.  In the
standard structure, \(\mathcal N\entails x\leq t\leftrightarrow x<St\);
provability of this equivalence in \(\mathsf Q\) is not claimed and is not
used below.  The witness \(z\) sits to the \emph{left} of \(+\) in
\(z+r=t\); this choice is what lets the comparison lemmas proved for
\(\leq\) in the incompleteness part rest only on the right-recursion
axioms \(x+0=x\) and \(x+Sy=S(x+y)\) for addition.

\begin{definition}[Syntactic arithmetical hierarchy]
A \emph{bounded arithmetic formula} is generated from atomic
\(\Lang_A\)-formulas by \(\bottom,\neg,\land,\lor,\to\) and marked bounded
quantifiers only.

To classify an arbitrary arithmetic formula, first eliminate implications and
push negations to atomic formulas by the De Morgan and quantifier-duality
rewrites, treating a marked bounded quantifier by the corresponding bounded
duality.  Rename bound variables by the fixed least-fresh-variable order when
needed to avoid capture.  This fixed conversion reaches negation normal form;
it does not move an unbounded quantifier across a bounded quantifier or a
connective.

At level zero, the \(\Sigma^0_0\) formulas and the \(\Pi^0_0\) formulas are
exactly the bounded arithmetic formulas in negation normal form.  Having
defined both classes at level \(n\), define the classes at level \(n+1\)
simultaneously as follows:
\begin{itemize}
  \item the \(\Sigma^0_{n+1}\) formulas in negation normal form form the least
        class containing every \(\Sigma^0_n\) and \(\Pi^0_n\) formula and
        closed under \(\land,\lor\), both marked bounded quantifiers, and
        unbounded existential quantification;
  \item the \(\Pi^0_{n+1}\) formulas in negation normal form form the least
        class containing every \(\Sigma^0_n\) and \(\Pi^0_n\) formula and
        closed under \(\land,\lor\), both marked bounded quantifiers, and
        unbounded universal quantification.
\end{itemize}
An arbitrary formula is a \emph{\(\Sigma^0_n\) formula} or a
\emph{\(\Pi^0_n\) formula} iff its fixed negation-normal-form conversion lies
in the corresponding recursively generated class.  The definition therefore
handles a bounded quantifier outside an unbounded quantifier without an
invalid quantifier exchange.  It depends on syntax and the fixed conversion,
not on whether the sentence happens to be true in \(\mathcal N\).
\end{definition}

\begin{definition}[Definability classes of relations]
Let \(k\in\N\).  A relation \(A\subseteq\N^k\) is
\emph{\(\Sigma^0_n\)-definable} iff there is a \(\Sigma^0_n\) formula
\(\varphi(\bar x)\) such that, for every \(\bar a\in\N^k\),
\[
  \bar a\in A
  \quad\Longleftrightarrow\quad
  \mathcal N\entails\varphi(\overline{\bar a}).
\]
Define \(\Pi^0_n\)-definable analogously, and call \(A\)
\emph{\(\Delta^0_n\)-definable} iff it is both
\(\Sigma^0_n\)-definable and \(\Pi^0_n\)-definable.  Write
\[
\begin{aligned}
 \boldsymbol{\Sigma}^0_n(k)&=\{A\subseteq\N^k:A\text{ is }
       \Sigma^0_n\text{-definable}\},\\
 \boldsymbol{\Pi}^0_n(k)&=\{A\subseteq\N^k:A\text{ is }
       \Pi^0_n\text{-definable}\},\\
 \boldsymbol{\Delta}^0_n(k)&=
       \boldsymbol{\Sigma}^0_n(k)\cap\boldsymbol{\Pi}^0_n(k).
\end{aligned}
\]
For unary relations, the argument \((1)\) may be omitted.  The unbold symbols
classify formulas syntactically; the bold symbols classify relations by
definability in \(\mathcal N\).  This typography records different standing
and prevents equivalence in the standard structure \(\mathcal N\) from
reclassifying every true or
false sentence as bounded.  Both standings are used in the results below.
\end{definition}

The semantic fact that a relation is decidable does not by itself provide an
object-language formula for that relation.  The lemmas and coding
definition below supply the particular formulas required in the hierarchy
arguments.  The coding of computation histories is built in-document from
the numerical pairing map of \cref{ch:strings} together with the G\"odel
beta function; every projection used is defined by a bounded formula,
proved as such here.

\begin{lemma}[Pairing graph]
\label{lem:pairing-graph}
Let
\[
  \mathsf P(x,y,z)\defeq(x+y)\times S(x+y)+(y+y)=z+z.
\]
Then \(\mathsf P\) is a quantifier-free \(\Lang_A\)-formula; for all
\(a,b,c\in\N\), \(\mathcal N\entails\mathsf P(\bar a,\bar b,\bar c)\) iff
\(c=\operatorname{pair}_{\N}(a,b)\); and both
\(a\leq\operatorname{pair}_{\N}(a,b)\) and
\(b\leq\operatorname{pair}_{\N}(a,b)\) hold for all \(a,b\in\N\).
\textit{Status: complete proof.}
\end{lemma}

\begin{proof}
The display is a single atomic equation, so \(\mathsf P\) is
quantifier-free.  In \(\mathcal N\) it evaluates to
\((a+b)(a+b+1)+2b=2c\), which holds iff
\(c=(a+b)(a+b+1)/2+b=\operatorname{pair}_{\N}(a,b)\); the doubled form
avoids division, which \(\Lang_A\) does not provide.  For the monotonicity
bounds: \(b\leq(a+b)(a+b+1)/2+b\) is immediate, and
\(a\leq a+b\leq(a+b)(a+b+1)/2\leq\operatorname{pair}_{\N}(a,b)\), where the
middle inequality holds because \(a+b=0\) forces \(a=0\), while
\(a+b\geq1\) makes the factor \((a+b+1)/2\) at least \(1\).
\end{proof}

In the next display, \(d\) abbreviates the term \(S((Si)\times m)\), and
\(u\) is a fresh variable with \(u\notin\operatorname{Var}(d)\).  The final
conjunct \(\exists u<d\,(v+Su=d)\) is the \emph{marked-bounded} form of the
comparison \(v<d\): it is \(\mathcal N\)-equivalent to the plain
abbreviation of \(v<d\), because any witness of \(v+Su=d\) satisfies
\(Su\leq d\), hence \(u<d\).  The marked form is required because the
hierarchy definition above counts only marked bounded quantifiers as
bounded, while the bare comparison \(v<d\) expands to an unbounded
existential.

\begin{lemma}[G\"odel beta function]
\label{lem:goedel-beta}
Define
\[
  \mathsf B(w,i,v)\defeq\exists t<Sw\,\exists m<Sw\,\exists q<St\,
  \bigl(\mathsf P(t,m,w)\land t=d\times q+v\land
  \exists u<d\,(v+Su=d)\bigr).
\]
Then \(\mathsf B\) is a bounded arithmetic formula, and:
\begin{enumerate}[label=(\alph*)]
\item for all \(w,i\in\N\) there is exactly one \(v\in\N\) with
      \(\mathcal N\entails\mathsf B(\bar w,\bar i,\bar v)\), and this
      \(v\) satisfies \(v\leq w\); write \(\beta(w,i)\) for it; and
\item for every finite sequence \(a_0,\ldots,a_{k-1}\) of natural numbers
      there is \(w\in\N\) with \(\beta(w,i)=a_i\) for every \(i<k\).
\end{enumerate}
\textit{Status: complete proof.}
\end{lemma}

\begin{proof}
Boundedness: every quantifier in the display is marked bounded, the final
conjunct is the marked-bounded form of \(v<d\) as recorded above, and the
remaining conjuncts are atomic.

(a) Since \(\operatorname{pair}_{\N}\) is a bijection, exactly one pair
\((t,m)\) satisfies \(\operatorname{pair}_{\N}(t,m)=w\), and
\cref{lem:pairing-graph} gives \(t\leq w\) and \(m\leq w\), so the
displayed bounds \(t<Sw\) and \(m<Sw\) do not exclude this pair, while
\(\mathsf P(t,m,w)\) excludes every other.  The term \(d\) evaluates to
\(1+(i+1)m\geq1\), so Euclidean division of \(t\) by that value determines
\(q\) and \(v\) uniquely via \(t=(1+(i+1)m)\,q+v\) with \(v<1+(i+1)m\);
moreover \(q\leq t<St\) and \(v\leq t\leq w\), so the bound on \(q\) is
respected and \(v\leq w\).

(b) Given \(a_0,\ldots,a_{k-1}\), put
\(N\defeq\max(k,a_0,\ldots,a_{k-1})+1\) and \(m\defeq N!\).  The moduli
\(1+(i+1)m\) for \(i<k\) are pairwise coprime: a prime \(p\) dividing both
\(1+(i+1)m\) and \(1+(j+1)m\) with \(i<j<k\) divides their difference
\((j-i)m\) but not \(m\), hence divides \(j-i\); yet \(0<j-i<k\leq N\)
gives \(p\leq N\), so \(p\) divides \(N!=m\), a contradiction.  Each
\(a_i<N\leq m<1+(i+1)m\).  By the Chinese remainder theorem for pairwise
coprime moduli there is \(t\) with \(t\equiv a_i\pmod{1+(i+1)m}\) for
every \(i<k\), and since \(a_i<1+(i+1)m\), the remainder of \(t\) on
division by \(1+(i+1)m\) is exactly \(a_i\).  Put
\(w\defeq\operatorname{pair}_{\N}(t,m)\); then \(\beta(w,i)=a_i\) for
every \(i<k\).
\end{proof}

G\"odel's original beta-function argument \cite{Godel1931} and the
treatment in \cite[ch.~16]{BoolosBurgessJeffrey2007} are parallel reading
for \cref{lem:goedel-beta}.

\begin{lemma}[Pairing contraction]
\label{lem:pair-contraction}
Let \(\theta\) be an \(\Lang_A\)-formula and let \(z\) be a variable not
free in \(\theta\) and distinct from \(x,y\).  Then \(\mathcal N\)
satisfies, under every variable assignment,
\[
  \exists x\,\exists y\,\theta
  \;\leftrightarrow\;
  \exists z\,\exists x<Sz\,\exists y<Sz\,(\mathsf P(x,y,z)\land\theta),
\]
and dually
\[
  \forall x\,\forall y\,\theta
  \;\leftrightarrow\;
  \forall z\,\forall x<Sz\,\forall y<Sz\,(\mathsf P(x,y,z)\to\theta).
\]
\textit{Status: complete proof.}
\end{lemma}

\begin{proof}
Fix an assignment.  If \(a,b\in\N\) satisfy \(\theta\) at \(x,y\), put
\(c\defeq\operatorname{pair}_{\N}(a,b)\); then \(\mathsf P(a,b,c)\) holds
in \(\mathcal N\) by \cref{lem:pairing-graph}, and its monotonicity bounds
give \(a,b\leq c<c+1\), so \(a\) and \(b\) witness the two marked bounded
quantifiers under \(z\mapsto c\).  Conversely, witnesses of the right side
satisfy \(\theta\) in particular.  Since \(z\) is not free in \(\theta\)
and differs from \(x,y\), no capture occurs and \(\theta\) keeps its
value.  The dual equivalence is the first applied to \(\neg\theta\),
negating both sides and moving negations inward.
\end{proof}

\begin{lemma}[Bounded-quantifier exchange]
\label{lem:bounded-collection}
Let \(\theta\) be an \(\Lang_A\)-formula, let \(t\) be an
\(\Lang_A\)-term with \(y,w\notin\operatorname{Var}(t)\), and let \(w\) be
fresh.  Then \(\mathcal N\) satisfies, under every variable assignment,
\[
  \forall y<t\,\exists x\,\theta
  \;\leftrightarrow\;
  \exists w\,\forall y<t\,\exists x<w\,\theta,
  \qquad
  \exists y<t\,\forall x\,\theta
  \;\leftrightarrow\;
  \forall w\,\exists y<t\,\forall x<w\,\theta.
\]
Both are equivalences of satisfaction in the standard structure, used
below for the boldface definability classes only; they are not
derivability claims in any formal theory of arithmetic.
\textit{Status: complete proof.}
\end{lemma}

\begin{proof}
Fix an assignment and consider the first equivalence.  Right to left is
immediate: a witness below \(w\) is in particular a witness.  Left to
right: the assignment gives \(t\) a value \(n\), and for each of the
finitely many \(b<n\) there is by hypothesis a witness \(x_b\in\N\) for
\(\theta\) at \(y\mapsto b\); let \(c\) be \(0\) when \(n=0\) and
\(\max_{b<n}x_b+1\) otherwise.  Then \(w\mapsto c\) makes
\(\forall y<t\,\exists x<w\,\theta\) true.  Freshness of \(w\) and
\(y,w\notin\operatorname{Var}(t)\) prevent capture and keep the value of
\(t\) fixed.  The second equivalence is the first applied to
\(\neg\theta\), negating both sides and moving negations inward.
\end{proof}

Call two \(\Lang_A\)-formulas \emph{\(\mathcal N\)-equivalent} iff exactly
the same variable assignments satisfy them in \(\mathcal N\).  Call a
formula \emph{prefix-form} \(\Sigma^0_n\) iff it consists of \(n\)
alternating blocks of unbounded quantifiers, the first block existential,
applied to a bounded arithmetic formula in negation normal form, and
\emph{prefix-form} \(\Pi^0_n\) dually, with the first block universal.
Every prefix-form formula lies in the corresponding syntactic class, since
the generating clauses allow each unbounded quantifier to be prepended at
the appropriate level.  The cited hierarchy sources define their classes
by prefix forms; the next theorem identifies the two presentations.

\begin{theorem}[Prefix normal form]
\label{thm:prefix-normal-form}
For every \(n\geq0\):
\begin{enumerate}[label=(\alph*)]
\item the relations defined in \(\mathcal N\) by prefix-form
      \(\Sigma^0_n\) formulas are closed, up to \(\mathcal N\)-equivalence
      of the defining formulas, under \(\land\), \(\lor\), and both marked
      bounded quantifiers, and for \(n\geq1\) also under unbounded
      existential quantification; dually for \(\Pi^0_n\) with unbounded
      universal quantification;
\item every \(\Sigma^0_n\) formula is \(\mathcal N\)-equivalent to a
      prefix-form \(\Sigma^0_n\) formula, and every \(\Pi^0_n\) formula to
      a prefix-form \(\Pi^0_n\) formula; and
\item the boldface classes \(\boldsymbol{\Sigma}^0_n(k)\),
      \(\boldsymbol{\Pi}^0_n(k)\), and \(\boldsymbol{\Delta}^0_n(k)\)
      coincide with the classes of relations defined in \(\mathcal N\) by
      prefix-form formulas of the corresponding kinds.
\end{enumerate}
\textit{Status: complete proof.}
\end{theorem}

\begin{proof}
Two preliminary observations.  \emph{Padding}: for \(m<n\), every
prefix-form formula of level \(m\) is \(\mathcal N\)-equivalent to a
prefix-form \(\Sigma^0_n\) formula and to a prefix-form \(\Pi^0_n\)
formula, obtained by prepending and interleaving quantifiers on fresh
variables; each added quantifier is vacuous because the domain \(\N\) is
nonempty.  \emph{Normal-form bookkeeping}: an implication
\(\mathsf P\to\theta\) introduced by the dual contraction of
\cref{lem:pair-contraction} is written \(\neg\mathsf P\lor\theta\), which
is in negation normal form because \(\mathsf P\) is atomic.

(a) Induction on \(n\), both classes at once; the \(\Pi^0_n\) case is dual
throughout, exchanging \(\exists\) with \(\forall\), \(\land\) with
\(\lor\) where the argument negates, and
\cref{lem:pair-contraction,lem:bounded-collection} with their dual halves.

Level \(n=0\): a prefix-form \(\Sigma^0_0\) or \(\Pi^0_0\) formula is a
bounded formula in negation normal form, and \(\land\), \(\lor\), and
marked bounded quantification produce bounded formulas in negation normal
form syntactically; no unbounded closure is claimed at this level.

Level \(n\geq1\): a prefix-form \(\Sigma^0_n\) formula is
\(\mathcal N\)-equivalent, by iterating the contraction of
\cref{lem:pair-contraction} inside its first block, to one of the form
\(\exists x\,B\) with \(B\) prefix-form \(\Pi^0_{n-1}\); the residual
bounded quantifiers and the conjuncts \(\mathsf P\) introduced by each
contraction are absorbed into \(B\) by the induction hypothesis for level
\(n-1\).  It therefore suffices to treat formulas \(\exists x\,B\) and
\(\exists x'\,B'\) with \(B,B'\) prefix-form \(\Pi^0_{n-1}\).
\begin{enumerate}[label=(\roman*)]
\item \(\land\), \(\lor\): rename the bound variables apart, so that
      \(x\notin\FV(B')\) and \(x'\notin\FV(B)\), and pull the existentials
      out: \(\exists x\,B\land\exists x'\,B'\) is
      \(\mathcal N\)-equivalent to \(\exists x\,\exists x'(B\land B')\),
      and \(\exists x\,B\lor\exists x'\,B'\) to
      \(\exists x\,\exists x'(B\lor B')\), the latter because the
      quantifier on the disjunct that did not supply the witness is
      vacuous over the nonempty domain.  Contract the two existentials by
      \cref{lem:pair-contraction}; the resulting matrix applies two marked
      bounded quantifiers and a conjunction with the quantifier-free
      \(\mathsf P\) to \(B\land B'\), respectively \(B\lor B'\), each
      \(\mathcal N\)-equivalent to a prefix-form \(\Pi^0_{n-1}\) formula
      by the induction hypothesis; prepending \(\exists z\) gives
      prefix-form \(\Sigma^0_n\).
\item Marked bounded \(\exists\): with \(x\) renamed so that \(x\neq y\)
      and \(x\notin\operatorname{Var}(t)\),
      \(\exists y<t\,\exists x\,B\) is \(\mathcal N\)-equivalent to
      \(\exists x\,\exists y<t\,B\) by pure logic; the induction
      hypothesis absorbs \(\exists y<t\) into \(B\), and prepending
      \(\exists x\) gives prefix-form \(\Sigma^0_n\).
\item Marked bounded \(\forall\): by \cref{lem:bounded-collection},
      \(\forall y<t\,\exists x\,B\) is \(\mathcal N\)-equivalent to
      \(\exists w\,\forall y<t\,\exists x<w\,B\); the induction hypothesis
      absorbs the two marked bounded quantifiers into \(B\), and
      prepending \(\exists w\) gives prefix-form \(\Sigma^0_n\).
\item Unbounded \(\exists\): contract
      \(\exists x'\,\exists x\,B\) by \cref{lem:pair-contraction} and
      absorb as in (i).
\end{enumerate}

(b) By the hierarchy definition, an arbitrary formula is \(\Sigma^0_n\)
or \(\Pi^0_n\) exactly when its fixed negation-normal-form conversion
lies in the corresponding generated class, and each conversion step ---
implication elimination, the De Morgan and quantifier-duality rewrites,
the bounded dualities, and capture-avoiding renaming --- preserves
satisfaction under every assignment, so every formula is
\(\mathcal N\)-equivalent to its conversion.  It therefore suffices to
prove the claim for the members of the generated classes, by induction on
\(n\) and, within each level, by induction on the generating clauses of
the syntactic hierarchy definition.  At level \(0\) the generated classes
consist of the bounded formulas in negation normal form, which are
exactly the prefix-form formulas of level \(0\).  At level \(n+1\), the
base formulas are the
\(\Sigma^0_n\) and \(\Pi^0_n\) formulas, which are
\(\mathcal N\)-equivalent to prefix-form formulas of level \(n\) by the
induction hypothesis and hence, by padding, to prefix-form
\(\Sigma^0_{n+1}\) formulas; each closure clause --- \(\land\), \(\lor\),
either marked bounded quantifier, unbounded \(\exists\) --- preserves
\(\mathcal N\)-equivalence to prefix form by (a) at level \(n+1\).  The
\(\Pi^0_{n+1}\) half is dual.

(c) By (b), a relation definable by a \(\Sigma^0_n\) formula is definable
by a prefix-form \(\Sigma^0_n\) formula, and conversely every prefix-form
\(\Sigma^0_n\) formula is a \(\Sigma^0_n\) formula; likewise for
\(\Pi^0_n\), hence for the intersections \(\boldsymbol{\Delta}^0_n(k)\).
\end{proof}

\begin{definition}[History coding]
\label{def:history-coding}
Fix an index scheme storing finitely many \emph{entry families} in a
single number \(s\): the member \(j\) of the family with tag \(t\) is the
entry \(\beta(s,\operatorname{pair}_{\N}(t,j))\), with finitely many
family tags fixed once for the clauses below.  A number \(s\) is a
\emph{history code} for the pair \((e,x)\) iff its entries comprise, under
this scheme:
\begin{enumerate}[label=(h\arabic*)]
\item the binary digits \(d_0,\ldots,d_{L-1}\) of \(x\), most significant
      digit first, with certified expansion values
      \(v_0,\ldots,v_{L-1}\) satisfying \(v_0=d_0\),
      \(v_{j+1}=2v_j+d_{j+1}\) for \(j+1<L\), and \(v_{L-1}=x\), together
      with the certificates \(v_j\leq x\) and \(d_j\leq1\) for all
      \(j<L\);
\item the transition table of \(e\): a decoding chain consisting of
      \(\mathsf P(k,c_0,e)\), \(\mathsf P(r_i,c_{i+1},c_i)\) for
      \(i<k-1\), and \(r_{k-1}=c_{k-1}\), unpacking the table tuple of
      length \(k\) into its entries \(r_0,\ldots,r_{k-1}\), each entry
      unpacked further by fixed-depth \(\mathsf P\)-steps, together with
      the validity checks of the fixed numbering under its
      designated-slot normalization;
\item configurations \(C_0,\ldots,C_T\) of the decoded machine, each tape
      coded in the tape-alphabet base with certified expansions as in
      (h1), the configuration \(C_0\) initial for the input word certified
      by (h1), and each successor configuration matching its predecessor
      and the decoded table by bounded digit comparisons; and
\item a final configuration whose state is the accepting state and whose
      tape prefix before the first blank is a valid binary numeral in the
      sense of the fixed output convention.
\end{enumerate}
Every auxiliary value occurring in these clauses is an entry of \(s\),
unique by \cref{lem:goedel-beta}(a), or a component of a pairing code,
bounded by the monotonicity clause of \cref{lem:pairing-graph}.
\end{definition}

\begin{theorem}[Bounded verification of computations]
\label{thm:bounded-verification}
There is a bounded \(\Lang_A\)-formula \(\mathsf{Comp}(e,x,s)\) such that:
\begin{enumerate}[label=(\alph*)]
\item for all \(e,x,s\in\N\),
      \(\mathcal N\entails\mathsf{Comp}(\bar e,\bar x,\bar s)\) iff
      \(s\) is a history code for \((e,x)\), so that \(s\) codes a valid
      accepting computation of program \(e\) on input \(x\) in the fixed
      machine numbering whose final tape has a valid numerical output;
      and
\item \(\varphi_e(x)\) is defined iff
      \(\mathcal N\entails\exists s\,\mathsf{Comp}(\bar e,\bar x,s)\).
\end{enumerate}
The symbol \(\mathsf{Comp}\) stands
for this arithmetic formula, not for the metalinguistic relation it defines;
the distinction is used in the theorem and exercises immediately below.
\textit{Status: proof sketch.}
\end{theorem}

\begin{proofplan}
Render each clause of \cref{def:history-coding} as a conjunction of
arithmetic conditions on the entries of \(s\).  Every auxiliary value is
an entry of \(s\) or a component of a pairing code, so every quantifier
can be bounded by a term in \(e,x,s\): entry extraction is the bounded
formula \(\mathsf B\) of \cref{lem:goedel-beta}, whose own witnesses are
bounded by \(s\), and component extraction is \(\mathsf P\) of
\cref{lem:pairing-graph} with its monotonicity bounds.  Every comparison
conjunct \(r<t\) or \(r\leq t\) is written in its marked-bounded form
\(\exists u<t\,(r+Su=t)\), respectively \(\exists u\leq t\,(u+r=t)\), with
\(u\) fresh, each \(\mathcal N\)-equivalent to the plain abbreviation; the
digit recurrence of (h1) is rendered with a numeral constant as
\(v_{j+1}=\bar2\times v_j+d_{j+1}\).  Under the counting rule for marked
bounded quantifiers, the resulting matrix is a bounded formula, which
gives boundedness.  Clause (a) then restates the history-code conditions.
For clause (b): if \(\varphi_e(x)\) is defined, the finite accepting
computation of the decoded machine, its table decoding, and the certified
expansions form one finite sequence of numbers, and
\cref{lem:goedel-beta}(b) supplies an \(s\) whose entries realize it;
if \(e\) is not a valid code, no \(s\) passes the validity checks of
(h2), matching the convention that invalid codes diverge everywhere.  The
omitted substantive step is the transcription of the full clause list
(h1)--(h4) into one displayed formula and the verification of the
step-matching digit comparisons of (h3); \cref{lem:pairing-graph} and
\cref{lem:goedel-beta} carry complete proofs.
\end{proofplan}

See \cite{Shoenfield1967,BoolosBurgessJeffrey2007} for complete coding proofs
using bounded checks of coded configurations and finite computation
histories; those sources state their hierarchy results for prefix-form
classes, and \cref{thm:prefix-normal-form} licenses transferring them to
the closure classes fixed here.

\begin{theorem}
A set is computably enumerable iff it belongs to
\(\boldsymbol{\Sigma}^0_1\).  A set is decidable iff it belongs to
\(\boldsymbol{\Delta}^0_1\).
\textit{Status: proof sketch.}
\end{theorem}

\begin{proofplan}
If \(A\) is computably enumerable, membership is equivalent to the existence
of a finite accepting computation.  The formula \(\mathsf{Comp}\) gives its
bounded arithmetic matrix, yielding a \(\Sigma^0_1\) definition.  Conversely,
search successive witnesses for the bounded matrix of a \(\Sigma^0_1\)
definition.
For the second claim, \(A\in\boldsymbol{\Delta}^0_1\) means that both \(A\) and its
complement are computably enumerable; apply the earlier characterization of
decidability.
The omitted substantive step constructs and verifies the numerical coding of
finite computation histories.
\end{proofplan}

\begin{definition}[Many-one hardness and completeness for a class]
Let \(\mathcal C\) be a class of subsets of \(\N\).  A set \(B\) is
\emph{\(\mathcal C\)-hard under many-one reductions} iff
\(A\leq_m B\) for every \(A\in\mathcal C\).  It is
\emph{\(\mathcal C\)-complete under many-one reductions} iff, in addition,
\(B\in\mathcal C\).  This term is required to distinguish an upper-bound
classification from the uniform lower bound used in the next theorem.  It is
a property of a set relative to \(\mathcal C\) and \(\leq_m\), and is not
semantic completeness, deductive completeness, or the status of a proof.
The theorem and the following exercise use exactly this distinction.
\end{definition}

\begin{theorem}[Hierarchy classifications of program domains and totality]
The set \(K\) is \(\boldsymbol{\Sigma}^0_1\)-complete under many-one
reductions, and its complement is \(\boldsymbol{\Pi}^0_1\)-complete under
many-one reductions.  The totality set
\[
  \mathsf{TOT}=\set{e:\forall x\,\exists k\,
  \text{program }e\text{ accepts on }x\text{ with a valid numerical output}
  \text{ within }k\text{ steps}}
\]
is \(\boldsymbol{\Pi}^0_2\)-complete under many-one reductions.
\textit{Status: cited theorem.}
\end{theorem}

See \cite{Shoenfield1967} for complete proofs.  The displayed prefix gives
the upper bound for \(\mathsf{TOT}\), read into the closure classes by
\cref{thm:prefix-normal-form}; the completeness conclusion additionally
requires a uniform many-one reduction.  The cited completeness proofs
classify relations by prefix-form formulas, and
\cref{thm:prefix-normal-form} licenses their transfer to the classes
fixed here.

\begin{theorem}[Strictness of the arithmetical hierarchy]
For every \(n\geq1\),
\[
  \boldsymbol{\Delta}^0_n\subsetneq\boldsymbol{\Sigma}^0_n,
  \qquad
  \boldsymbol{\Delta}^0_n\subsetneq\boldsymbol{\Pi}^0_n,
  \qquad
  \boldsymbol{\Sigma}^0_n\neq\boldsymbol{\Pi}^0_n,
\]
and
\(\boldsymbol{\Sigma}^0_n\cup\boldsymbol{\Pi}^0_n
\subsetneq\boldsymbol{\Delta}^0_{n+1}\), for unary relations.
\textit{Status: cited theorem.}
\end{theorem}

See \cite{Shoenfield1967} for a complete proof.  The cited proof works
with prefix-form classes; \cref{thm:prefix-normal-form} transfers the
statement to the closure classes fixed here.

\begin{warning}[Syntactic prefix versus semantic class]
The syntactic classification of a displayed formula supplies an upper bound
on the definability class of the relation it defines.  A completeness
reduction or a nondefinability argument is needed to show that the relation
does not lie in a lower boldface class.
\end{warning}

\begin{exercise}
\label{ex:decision-hierarchy-01}
Give a \(\Sigma^0_1\) formula defining the set of indices whose domain is
nonempty.  Prove the set is \(\boldsymbol{\Sigma}^0_1\)-complete under
many-one reductions by reducing \(K\) to it.
\end{exercise}

\begin{exercise}
\label{ex:decision-hierarchy-02}
Classify the following index sets by giving the best immediate upper bound in
the arithmetical hierarchy: finite domain, infinite domain, and cofinite domain.
Do not claim completeness without constructing reductions.
\end{exercise}

\begin{exercise}
\label{ex:decision-hierarchy-03}
Explain why first-order validity is computably enumerable even though validity
is defined by a universal quantification over all structures.  Identify
exactly where the completeness theorem is used.
\end{exercise}

\section{Relative computation and Post's theorem}
\label{sec:decision-hierarchy-03}

Many-one reducibility compares sets by a total computable translation of
questions.  Relative computation compares them by permitting one
computation to consult a fixed set during its run.

\begin{definition}[Oracle Turing machine and Turing reducibility]
An \emph{oracle Turing machine} with oracle \(B\subseteq\N\) is a
deterministic single-tape Turing machine augmented with a distinguished
query operation: at a query, the machine computes a number \(n\) and
control passes to one of two designated states according as \(n\in B\) or
\(n\notin B\).  The run is otherwise governed by the fixed machine
definition.  An oracle Turing machine \emph{computes} a partial numerical
function by the input and output discipline fixed for ordinary machines,
read verbatim with the oracle present: on the input word
\(\operatorname{bin}(x)\), its value is the numerical output when the run
halts accepting with a valid numerical output, and its value is undefined
when the run rejects, diverges, or accepts without a valid numerical
output.  A set \(A\subseteq\N\) is \emph{Turing reducible} to
\(B\), written \(A\leq_T B\), iff some oracle Turing machine with oracle
\(B\) decides \(A\).  A set is \emph{computably enumerable relative to}
\(B\) iff it is the domain of a partial function computed by some oracle
Turing machine with oracle \(B\).
\end{definition}

Every many-one reduction yields a Turing reduction: translate the input and
query the oracle once.  The converse fails: \(\N\setminus K\leq_T K\) by
a single query with the answer reversed, while
\(\N\setminus K\not\leq_m K\), because computable enumerability is
preserved downward by many-one reductions, \(K\) is computably enumerable,
and \(\N\setminus K\) is not.

\begin{definition}[Turing jump]
\label{def:turing-jump}
Fix an effective numbering of oracle Turing machines analogous to the
fixed machine numbering: the same designated-slot normalization, with the
query rows included in the row list, and with every number that is not a
valid code denoting an oracle machine that diverges on every input.  For
a valid code \(e\) and an oracle \(B\subseteq\N\), let \(\varphi^B_e(x)\)
be the numerical output when the decoded machine with oracle \(B\), on
the input word \(\operatorname{bin}(x)\), halts accepting with a valid
numerical output, and leave \(\varphi^B_e(x)\) undefined when that run
rejects, diverges, or accepts without a valid numerical output.  The
\emph{Turing jump} of \(B\subseteq\N\) is
\[
  B'\defeq\set{e\in\N:\varphi^B_e(e)\text{ is defined}}.
\]
Define \(B^{(0)}\defeq B\) and
\(B^{(n+1)}\defeq(B^{(n)})'\); in particular \(\varnothing^{(n)}\) is the
\(n\)-th jump of the empty set.
\end{definition}

\begin{lemma}[Relativized adequacy]
\label{lem:relativized-adequacy}
For every oracle \(B\subseteq\N\):
\begin{enumerate}[label=(\alph*)]
\item every partial function computed by an oracle Turing machine with
      oracle \(B\) occurs as \(\varphi^B_e\) for some valid code \(e\);
\item there is a fixed oracle machine computing a function
      \(U^B(e,x)\simeq\varphi^B_e(x)\) for all \(e,x\in\N\); and
\item for every \(m,n\geq1\) there is a total computable function
      \(\mathsf{par}_n^m:\N^{m+1}\to\N\), not depending on \(B\), with
      \(\varphi^B_{\mathsf{par}_n^m(e,\bar a)}(\bar x)
      \simeq\varphi^B_e(\bar a,\bar x)\) for every index \(e\), all
      argument tuples, and every oracle \(B\) simultaneously.
\end{enumerate}
\textit{Status: proof sketch.}
\end{lemma}

\begin{proofplan}
Clause (a) relativizes the proof of \cref{thm:enumeration}: renaming
commutes with the step relation with the query rows included, the
surrogate-start and automatic-alphabet arguments are unchanged, and the
designated-slot normalization applies to the extended row list.  Clause
(b) relativizes the universal simulation behind
\cref{thm:universal-computation}, passing each query of the simulated
machine through to the oracle.  Clause (c) relativizes the compiler
construction given for \cref{thm:parameter}, copying the query rows
unchanged; the compiled table transformation does not consult the oracle,
which gives the simultaneous claim.  The omitted verifications are
exactly those of the results being relativized.
\end{proofplan}

By \cref{lem:relativized-adequacy}(a), a set is computably enumerable
relative to \(B\) iff it is the domain of some \(\varphi^B_e\).

\begin{theorem}[Jump theorem]
\label{thm:jump}
For every \(B\subseteq\N\):
\begin{enumerate}[label=(\alph*)]
\item \(B'\) is computably enumerable relative to \(B\); and
\item \(B'\not\leq_T B\).
\end{enumerate}
\textit{Status: proof sketch.}
\end{theorem}

\begin{proofplan}
(a) \(B'\) is the domain of the partial function
\(e\mapsto U^B(e,e)\), which an oracle machine with oracle \(B\) computes
by \cref{lem:relativized-adequacy}(b).

(b) Relativized diagonal.  Suppose some oracle machine with oracle \(B\)
decides \(B'\), and let \(h^B\) be its total characteristic function,
with \(h^B(e)=1\) iff \(e\in B'\).  Define
\[
  d^B(e)\simeq
  \begin{cases}
    0, & h^B(e)=0,\\
    \text{undefined}, & h^B(e)=1.
  \end{cases}
\]
Then \(d^B\) is computed by an oracle machine with oracle \(B\): run the
decider, then write \(0\) and accept, or enter an infinite loop.  By
\cref{lem:relativized-adequacy}(a), fix \(b\) with
\(\varphi^B_b\simeq d^B\).  Then \(b\in B'\) iff \(\varphi^B_b(b)\) is
defined iff \(h^B(b)=0\) iff \(b\notin B'\), a contradiction.

The displayed diagonal calculation is complete, but clause (a) rests on
the relativized universal machine and clause (b) on the occurrence clause
of \cref{lem:relativized-adequacy}, both proof sketches; the status
policy of the front matter forbids complete-proof status over a sketch
dependency, which is why this theorem is a proof sketch.
\end{proofplan}

\begin{theorem}[Post]
For every \(n\in\N\) and every \(A\subseteq\N\):
\begin{enumerate}
  \item \(A\in\boldsymbol{\Sigma}^0_{n+1}\) iff \(A\) is computably
        enumerable relative to \(\varnothing^{(n)}\); and
  \item \(A\in\boldsymbol{\Delta}^0_{n+1}\) iff
        \(A\leq_T\varnothing^{(n)}\).
\end{enumerate}
\textit{Status: cited theorem.}
\end{theorem}

See \cite{Shoenfield1967} for a complete proof; the handbook treatment in
\cite{Barwise1977} develops the relativized machinery in full.  The cited
proofs classify relations through prefix-form formulas;
\cref{thm:prefix-normal-form} transfers those classifications to the
classes fixed here.  Post's
theorem ties the syntactic hierarchy to iterated jumps, and the first
three clauses of the strictness theorem above follow from it together with
the jump theorem, \cref{thm:jump}:
\(\varnothing^{(n+1)}\) is computably enumerable relative to
\(\varnothing^{(n)}\) but not Turing reducible to it, so
\(\boldsymbol{\Sigma}^0_{n+1}\ne\boldsymbol{\Delta}^0_{n+1}\).  The
final inclusion clause additionally uses the join of
\(\varnothing^{(n)}\) with its complement and the completeness of
\(\varnothing^{(n)}\) among sets computably enumerable relative to
\(\varnothing^{(n-1)}\); those steps are part of the cited proof.

\begin{remark}[Convention bridge for the jump]
\label{rem:jump-convention}
Write \(A\equiv_m B\), \emph{many-one equivalence}, iff
\(A\leq_m B\) and \(B\leq_m A\); the standing is a metalinguistic relation
between subsets of \(\N\), and it is used only in this remark and the
results that cite it.  Four computable table transformations align this
jump with the conventions of the cited sources.  (i)
\(K\equiv_m\varnothing'\): an
ordinary transition table embeds as a query-free oracle table, and an
oracle machine with oracle \(\varnothing\) is simulated by answering
every query negatively.  (ii) The halting-based jump of the cited sources
is \(\equiv_m\) the present \(B'\): the map \(e\mapsto e^*\), where
\(e^*\) simulates \(e\) and, on any halt of the simulation, writes \(0\)
and accepts, reduces halting to definedness; the map
\(e\mapsto e^{**}\), where \(e^{**}\) diverges unless the simulated halt
is accepting with a valid numerical output, reduces definedness to
halting.  All four maps are computable table transformations, made
uniform in their inputs by \cref{thm:parameter}.  Jump monotonicity ---
\(A\leq_T B\) implies \(A'\leq_m B'\) --- is used explicitly in this
alignment: replace each \(A\)-query of an oracle machine with oracle
\(A\) by a run of a fixed oracle-\(B\) decider for \(A\), then align the
self-application indices by the relativized parameter theorem,
\cref{lem:relativized-adequacy}(c), with padding.  Hence Post's theorem
is insensitive to the choice between the two conventions, and the
citations stand.
\textit{Status: proof sketch --- the four table transformations are
elementary, but the monotonicity step runs through
\cref{lem:relativized-adequacy}(c), itself a proof sketch.}
\end{remark}

\part{Formal Arithmetic and Incompleteness}
\label{part:incompleteness}

\chapter{Arithmetic, Representation, and Coding}
\label{ch:arithmetic-coding}

This part separates three levels throughout: arithmetic in the standard
structure \(\N\), formal derivability in a theory \(T\), and metalogical
reasoning about \(T\).  G\"odel's original paper is \cite{Godel1931}; modern
treatments include \cite{Shoenfield1967,BoolosBurgessJeffrey2007}.

\section{The language and standard structure of arithmetic}
\label{sec:arithmetic-coding-01}

The language \(\Lang_A\), the standard structure \(\mathcal N\), and the
numeral notation were fixed
in \cref{def:language-arithmetic}.  This part retains exactly that syntax and
semantics.

\begin{definition}[Robinson arithmetic]
Robinson arithmetic, written \(\mathsf Q\), is the finitely axiomatized
\(\Lang_A\)-theory consisting of the universal closures of:
\[
\begin{array}{ll}
  Sx\neq0, & Sx=Sy\to x=y,\\
  x\neq0\to\exists y\,x=Sy, & x+0=x,\\
  x+Sy=S(x+y), & x\times0=0,\\
  x\times Sy=(x\times y)+x. &
\end{array}
\]
It has no induction scheme.
\end{definition}

The incompleteness argument of Rosser evaluates bounded numerical
comparisons inside \(\mathsf Q\).  The two lemmas below supply exactly
those evaluations for the object-language abbreviation \(\leq\) registered
in the hierarchy chapter, whose expansion \(r\leq t\defeq\exists
z\,(z+r=t)\) places the witness on the left of \(+\).  Both proofs
therefore use, besides the successor axioms, only the right-recursion
axioms \(x+0=x\) and \(x+Sy=S(x+y)\).

\begin{lemma}[Bounded-case enumeration in \(\mathsf Q\)]
\label{lem:q-bounded-cases}
For every \(n\in\N\),
\[
  \mathsf Q\derives\forall y\,\bigl(y\leq\bar n\to
  (y=\bar0\lor y=\bar1\lor\dots\lor y=\bar n)\bigr).
\]
\textit{Status: complete proof.}
\end{lemma}

\begin{proof}
Metainduction on \(n\); each instance is a \(\mathsf Q\)-derivation using
only the axioms \(Sx\neq0\), \(Sx=Sy\to x=y\),
\(x\neq0\to\exists y\,x=Sy\), \(x+0=x\), and \(x+Sy=S(x+y)\).

Base \(n=0\).  Reason in \(\mathsf Q\) from the expansion of
\(y\leq\bar0\): fix \(z\) with \(z+y=\bar0\).  By the case axiom,
\(y=0\) or \(y=Sw\) for some \(w\).  In the second case
\(S(z+w)=z+Sw=\bar0\) contradicts \(Sx\neq0\).  Hence \(y=\bar0\).

Step from \(n\) to \(n+1\).  Reason in \(\mathsf Q\) from \(z+y=S\bar n\),
noting that \(\overline{n+1}\) is \(S\bar n\).  If \(y=0\), the disjunct
\(y=\bar0\) holds.  If \(y=Sw\), then \(S(z+w)=z+Sw=S\bar n\), so
\(S\)-injectivity gives \(z+w=\bar n\), hence \(w\leq\bar n\); the
already-derived instance for \(n\) gives
\(w=\bar0\lor\dots\lor w=\bar n\), and applying \(S\) to each disjunct
gives \(y=\bar1\lor\dots\lor y=\overline{n+1}\).  In both cases
\(y=\bar0\lor\dots\lor y=\overline{n+1}\).
\end{proof}

\begin{lemma}[Dichotomy in \(\mathsf Q\)]
\label{lem:q-dichotomy}
For every \(n\in\N\),
\[
  \mathsf Q\derives\forall y\,
  \bigl(y\leq\bar n\lor\bar n\leq y\bigr).
\]
\textit{Status: complete proof.}
\end{lemma}

\begin{proof}
Metainduction on \(n\), from the same axiom list.

Base \(n=0\).  For every \(y\), the witness \(z\defeq y\) gives
\(y+\bar0=y\) by the axiom \(x+0=x\), so \(\bar0\leq y\).

Step from \(n\) to \(n+1\).  Fix \(y\); by the case axiom, \(y=0\) or
\(y=Sw\) for some \(w\).  If \(y=0\), the witness \(z\defeq S\bar n\)
gives \(S\bar n+0=S\bar n\), so \(y\leq S\bar n\).  If \(y=Sw\), the
already-derived instance for \(n\), applied to \(w\), gives either some
\(z\) with \(z+w=\bar n\), whence \(z+Sw=S(z+w)=S\bar n\) and
\(y\leq S\bar n\), or some \(u\) with \(u+\bar n=w\), whence
\(u+S\bar n=S(u+\bar n)=Sw=y\) and \(S\bar n\leq y\).
\end{proof}

\begin{definition}[Peano arithmetic]
Peano arithmetic, written \(\PA\), is \(\mathsf Q\) together with every
instance of the induction scheme
\[
  \bigl(\varphi(0,\bar z)\land
  \forall x(\varphi(x,\bar z)\to\varphi(Sx,\bar z))\bigr)
  \to\forall x\,\varphi(x,\bar z),
\]
where \(\varphi(x,\bar z)\) is any \(\Lang_A\)-formula and the parameters
\(\bar z\) remain universally quantified in the closure.
\end{definition}

\section{Standard and nonstandard models of arithmetic}
\label{sec:arithmetic-coding-02}
\label{sec:nonstandard-arithmetic}

\begin{definition}[Standard and nonstandard models]
\label{def:nonstandard-model}
Let \(T\) be an \(\Lang_A\)-theory satisfied by \(\mathcal N\).  A model
\(\mathcal M\) of \(T\) is \emph{standard} iff
\(\mathcal M\cong\mathcal N\) as an \(\Lang_A\)-structure, and it is
\emph{nonstandard} otherwise.  These terms classify models of arithmetic by
isomorphism type; they do not classify sentences or proof predicates.  The
distinction is stronger than agreement on all sentences, as the theorem and
exercise below show, and it is needed to state what compactness produces.
\end{definition}

\begin{theorem}[Compactness produces nonstandard models]
\label{thm:nonstandard-model}
Assume the Axiom of Choice in the ambient set theory.  Every \(\Lang_A\)-theory
\(T\) satisfied by \(\mathcal N\) has a countable
nonstandard model.
\textit{Status: complete proof.}
\end{theorem}

\begin{proof}
Add one new constant \(c\) and consider
\[
  T\cup\{c\ne\bar n:n\in\N\}.
\]
Every finite subset has an expansion of \(\mathcal N\) as a model: interpret
\(c\) by a natural number different from the finitely many numerals mentioned.
By first-order compactness, the whole set has a model \(\mathcal M'\).  In its
\(\Lang_A\)-reduct \(\mathcal M\), the element \(c^{\mathcal M'}\) is unequal
to \(\bar n^{\mathcal M}\) for every standard \(n\).  If
\(\mathcal M\cong\mathcal N\), every element of \(M\) would be the value of
some numeral, a contradiction.  Thus the reduct is a nonstandard model of
\(T\).  Applying \cref{thm:downward-ls} in the expanded countable signature,
with \(c^{\mathcal M'}\) among the designated elements, makes the model
countable while preserving all the displayed sentences.
\end{proof}

See \cite{Hodges1993,Marker2002} for further compactness constructions of
nonstandard models.

\begin{example}[Nonstandard Peano arithmetic]
Since \(\mathcal N\entails\PA\), \cref{thm:nonstandard-model} gives a
countable model of \(\PA\) containing an element unequal to every numeral.
The theorem is an existence result: it does not identify that element with
the value of any numeral, and it does not make the resulting model
isomorphic to
\(\mathcal N\).
\end{example}

\begin{exercise}
\label{ex:arithmetic-coding-01}
\label{ex:nonstandard-true-arithmetic}
Apply \cref{thm:nonstandard-model} to \(T=\Th(\mathcal N)\).  Prove that the
resulting countable model agrees with \(\mathcal N\) on every
\(\Lang_A\)-sentence but is not isomorphic to \(\mathcal N\).  Explain why
elementary equivalence therefore does not imply isomorphism.
\end{exercise}

\begin{warning}[Standard truth versus formal proof]
The notation \(\mathcal N\entails\sigma\) says that \(\sigma\) is true in the
standard structure.  The notation \(T\derives\sigma\) says that a finite
derivation exists from the axioms of \(T\) in a fixed effective calculus.
Passage from \(T\derives\sigma\) to \(\mathcal N\entails\sigma\)
requires soundness of the fixed calculus together with the hypothesis
\(\mathcal N\entails T\).  Passage from \(\mathcal N\entails\sigma\) to
\(T\derives\sigma\) fails in general even though the fixed calculus is
complete---this failure is the content of the first incompleteness theorem
below---and holds only for restricted sentence classes, such as
\(\Sigma^0_1\) sentences over theories extending \(\mathsf Q\).
\end{warning}

\section{Representability}
\label{sec:arithmetic-coding-03}

\begin{definition}[Numeralwise representation of a relation]
Let \(T\supseteq\mathsf Q\), and let \(R\subseteq\N^k\).  A formula
\(\rho_R(x_1,\ldots,x_k)\) \emph{numeralwise represents} \(R\) in \(T\) if,
for every \(\bar n\in\N^k\),
\[
\begin{aligned}
  R(\bar n)&\quad\Longrightarrow\quad
  T\derives\rho_R(\overline{\bar n}),\\
  \neg R(\bar n)&\quad\Longrightarrow\quad
  T\derives\neg\rho_R(\overline{\bar n}).
\end{aligned}
\]
The two clauses concern each standard tuple of numerals; they do not assert
that \(T\) proves a universally quantified decision statement.
\end{definition}

\begin{definition}[Representation of a function]
Let \(T\supseteq\mathsf Q\).  A formula \(\gamma_f(\bar x,y)\)
\emph{represents} a total function
\(f:\N^k\to\N\) in \(T\) if, for every \(\bar n\in\N^k\),
\[
  T\derives\forall y\,
  \bigl(\gamma_f(\overline{\bar n},y)
  \leftrightarrow y=\overline{f(\bar n)}\bigr).
\]
\end{definition}

\begin{theorem}[Representability theorem]
\label{thm:representability}
Every primitive recursive relation is numeralwise representable in
\(\mathsf Q\) by a \(\Sigma^0_1\) formula, and every primitive recursive
function is representable in \(\mathsf Q\).  Every computably enumerable
relation is definable in \(\mathcal N\) by a \(\Sigma^0_1\) formula.
\textit{Status: cited theorem.}
\end{theorem}

See \cite{Shoenfield1967} for a complete proof.

The construction represents the initial functions directly and preserves
representability through composition and primitive recursion by quantifying
over coded finite computation sequences.  The sequence-coding operations and
the check that a code satisfies the recursion equations are primitive
recursive.  For a computably enumerable relation, a finite accepting
computation supplies the existential witness.  This paragraph records the
construction used below; the cited source supplies the complete coding proof.

\begin{example}[Representation of addition]
The formula
\[
  \mathsf{Add}(x,z,y)\defeq x+z=y
\]
represents the total function \((m,n)\mapsto m+n\) in \(\mathsf Q\).  For each
standard \(m,n\), finitely many uses of the two addition axioms give
\(\mathsf Q\derives\bar m+\bar n=\overline{m+n}\).  Equality reasoning then
gives
\[
  \mathsf Q\derives\forall y\,
  \bigl(\mathsf{Add}(\bar m,\bar n,y)
  \leftrightarrow y=\overline{m+n}\bigr).
\]
This numeral-by-numeral verification uses no induction scheme.
\end{example}

\section{G\"odel coding}
\label{sec:arithmetic-coding-04}

\begin{definition}[G\"odel numbering]
A \emph{G\"odel numbering} for the chosen formal language and calculus is an
injective effective map from symbols, finite strings, formulas, and finite
derivations to natural numbers such that the basic syntactic operations and
predicates required below are primitive recursive on codes.
\end{definition}

The exact numerical values are conventional; the effectiveness properties are
substantive.  Fix one such numbering.  Write \(\ulcorner E\urcorner\) for the
code of the expression or derivation \(E\).  In an arithmetic formula,
\(\overline{\ulcorner E\urcorner}\) is the numeral denoting that code.

\begin{definition}[Computably enumerable theory]
Relative to the fixed G\"odel numbering, a theory
\(T\subseteq\Sent(\Lang_A)\) is \emph{computably enumerable} iff
\[
  \set{\ulcorner\sigma\urcorner:\sigma\in T}
\]
is a computably enumerable set of natural numbers.  This term classifies the
given set of axioms \(T\); it does not mean that \(T\) is deductively closed,
and it does not define \(T\) as the set of its derivable consequences.  The
distinction is required because the proof predicate below verifies uses of
members of \(T\) by finite enumeration-stage certificates.
\end{definition}

\begin{example}[A finite-sequence code]
One illustrative component code assigns to
\((a_0,\ldots,a_{k-1})\) the natural number
\[
  \operatorname{pair}_{\N}\left(k,\prod_{i<k}p_i^{a_i+1}\right),
\]
where \(p_i\) is the \(i\)-th prime.  For \((2,0,1)\), the product is
\[
  2^3\cdot3^1\cdot5^2=600,
\]
so the complete code is \(\operatorname{pair}_{\N}(3,600)\).  The
exponent \(a_i+1\) rather than \(a_i\) is what distinguishes trailing
zeros; pairing with \(k\) additionally makes the length available by a
projection rather than by factoring.  A full arithmetization usually uses a coding with especially
simple primitive recursive projection and concatenation operations.
\end{example}

\begin{theorem}[Arithmetization of syntax]
For the fixed numbering, each of the following is primitive recursive: being a
symbol of a specified kind, being a term or formula, substitution of a numeral
for a free variable, being an axiom instance of a fixed finitely axiomatized
theory, and being a correct finite derivation from such axioms.
For a computably enumerable theory, derivations may carry finite certificates
identifying the stage at which each nonlogical axiom is enumerated, again
giving a primitive recursive proof predicate.
\textit{Status: proof sketch.}
\end{theorem}

\begin{proofplan}
Proceed by bounded search through the finite coded expression or derivation.
Formation trees, substitutions, and rule checks are finite data.  Replace an
unbounded appeal to computably enumerable axiom membership by a bounded
verification of the enumeration-stage certificate stored in the derivation
code.
\end{proofplan}

\begin{definition}[Proof and provability predicates]
\label{def:fixed-proof-predicate}
Let \(T\) be computably enumerable and contain \(\mathsf Q\).  Fix a primitive recursive
relation \(\mathsf{Proof}_T(p,x)\) meaning that \(p\) codes a correct
\(T\)-proof whose concluding sentence has code \(x\).  By the
\(\Sigma^0_1\) clause of \cref{thm:representability}, fix the same symbol
for a
specific \(\Sigma^0_1\) formula that numeralwise represents this relation in
\(\mathsf Q\).  Define
\[
  \mathsf{Prov}_T(x)\defeq\exists p\,\mathsf{Proof}_T(p,x).
\]
The additional existential quantifier merges with the initial existential
block of the chosen representation, so \(\mathsf{Prov}_T\) is a
\(\Sigma^0_1\) formula by the syntactic hierarchy definition.  The
coded-computation-sequence apparatus behind the \(\Sigma^0_1\) clause of
\cref{thm:representability} --- the construction paragraph following it
--- is the \(\beta\) coding of \cref{lem:goedel-beta}, so a single
in-document coding serves machine tables, computation histories, and
proof sequences.
\end{definition}

\begin{warning}[Mentioning a sentence inside arithmetic]
The expression \(\mathsf{Prov}_T(\varphi)\) is ill-typed if \(\varphi\) is a
formula rather than a number term.  The correct expression is
\(\mathsf{Prov}_T(\overline{\ulcorner\varphi\urcorner})\), often abbreviated
only after the coding convention has been declared.
\end{warning}

\begin{exercise}
\label{ex:arithmetic-coding-02}
Define a primitive recursive relation saying that a finite sequence is a
correct addition computation
\(a,a+1,\ldots,a+b\).  Explain how the relation supports representability of
addition.
\end{exercise}

\begin{exercise}
\label{ex:arithmetic-coding-03}
Distinguish the following three assertions by giving their types:
\(T\derives\sigma\),
\(\mathcal N\entails\mathsf{Prov}_T(\overline{\ulcorner\sigma\urcorner})\),
and
\(T\derives\mathsf{Prov}_T(\overline{\ulcorner\sigma\urcorner})\).
State one hypothesis that connects each adjacent pair.
\end{exercise}

\chapter{Diagonalization and the First Incompleteness Theorem}
\label{ch:first-incompleteness}

\section{The diagonal lemma}
\label{sec:first-incompleteness-01}

Let \(\mathsf{diag}(n)\) be the code of the result of substituting the numeral
\(\bar n\) for the distinguished free variable in the one-variable formula
with code \(n\).  On nonformula codes, define an arbitrary default value.  The
function \(\mathsf{diag}\) is primitive recursive.

\begin{theorem}[Diagonal lemma]
\label{thm:diagonal-lemma}
Let \(T\) be any theory extending \(\mathsf Q\).  For every
\(\Lang_A\)-formula \(\theta(x)\) with at most the displayed free variable,
there is a sentence \(G\) such that
\[
  T\derives G\leftrightarrow
  \theta(\overline{\ulcorner G\urcorner}).
\]
\textit{Status: cited theorem.}
\end{theorem}

See \cite{BoolosBurgessJeffrey2007} for a complete proof.  The following
construction explains the fixed point used by subsequent complete proofs.
Let \(D(x,y)\) represent the function \(\mathsf{diag}\) in \(\mathsf Q\), so
for each standard \(n\),
\[
  \mathsf Q\derives\forall y\,
  \bigl(D(\bar n,y)\leftrightarrow
  y=\overline{\mathsf{diag}(n)}\bigr).
\]
Define
\[
  \beta(x)\defeq\exists y\,(D(x,y)\land\theta(y)).
\]
Let \(b=\ulcorner\beta\urcorner\), and let
\(G\) be the sentence \(\beta(\bar b)\).  By construction,
\(\mathsf{diag}(b)=\ulcorner G\urcorner\).  Representability of
\(\mathsf{diag}\) gives, already in \(\mathsf Q\),
\[
\begin{aligned}
  G
  &\leftrightarrow \exists y\,(D(\bar b,y)\land\theta(y))\\
  &\leftrightarrow
  \theta(\overline{\mathsf{diag}(b)})\\
  &\leftrightarrow\theta(\overline{\ulcorner G\urcorner}).
\end{aligned}
\]
Since \(T\supseteq\mathsf Q\), \(T\) proves the same biconditional.
The cited proof supplies the primitive-recursive verification and numeralwise
representability of \(\mathsf{diag}\) on which this explanation depends.

\begin{remark}
The lemma does not create a sentence that literally contains its own
inscription.  It creates a sentence provably equivalent to the result of
applying \(\theta\) to the numeral for that sentence's code.
\end{remark}

\section{Consistency and soundness conditions}
\label{sec:first-incompleteness-02}

\begin{definition}[Consistency]
A theory \(T\) is \emph{consistent} iff there is no sentence \(\sigma\) such
that both \(T\derives\sigma\) and \(T\derives\neg\sigma\).  In the fixed
classical explosive calculus, this is equivalent to
\(T\notderives\bottom\).  If \(T\derives 0\ne S0\), then it is also
equivalent to \(T\notderives 0=S0\); in particular this specialization applies
when \(T\supseteq\mathsf Q\).  Without the displayed derivability hypothesis,
nonderivability of \(0=S0\) is not a general test for consistency.
\end{definition}

Recall from \cref{def:theory-completeness} that \(T\) is deductively complete
when, for every sentence \(\sigma\) in its language, either
\(T\derives\sigma\) or \(T\derives\neg\sigma\).  This property of a theory is
distinct from semantic completeness of a calculus, while soundness and
completeness of the fixed calculus identify it with the quantified semantic
condition stated there.

\begin{definition}[\(\Sigma^0_1\)-soundness]
A theory \(T\) is \emph{\(\Sigma^0_1\)-sound} iff, whenever \(\sigma\) is a
\(\Sigma^0_1\) sentence and \(T\derives\sigma\), then
\(\mathcal N\entails\sigma\).
\end{definition}

\begin{definition}[Omega-consistency]
A theory \(T\) is \emph{omega-consistent} iff there is no formula
\(\varphi\) with \(\FV(\varphi)\subseteq\{x\}\) such that
\[
  T\derives\exists x\,\varphi(x)
  \quad\text{and}\quad
  T\derives\neg\varphi(\bar n)
  \text{ for every }n\in\N.
\]
\end{definition}

Both omega-consistency and \(\Sigma^0_1\)-soundness imply consistency, but
consistency alone is weaker.  Rosser's theorem will remove the stronger
hypothesis from the incompleteness conclusion.

\section{The G\"odel sentence}
\label{sec:first-incompleteness-03}

Let \(T\) be a computably enumerable theory extending \(\mathsf Q\).  Apply
\cref{thm:diagonal-lemma}, in the complexity-preserving form discussed
below display (G), to
\(\theta(x)=\neg\mathsf{Prov}_T(x)\).  Obtain a \(\Pi^0_1\) sentence \(G_T\)
satisfying
\[
  T\derives G_T\leftrightarrow
  \neg\mathsf{Prov}_T(\overline{\ulcorner G_T\urcorner}).
  \tag{G}
\]
The fixed-point biconditional is in fact derivable already in \(\mathsf Q\).
The complexity-preservation clause is a standard refinement of the coding
construction: choose a suitable representation of substitution on codes and
put the resulting fixed point back into prenex form.  The fixed point of a
\(\Pi^0_1\) formula may then be chosen \(\Pi^0_1\).  This clause is cited
rather than proved here; see
\cite{Shoenfield1967}.

\begin{theorem}[G\"odel's first incompleteness theorem, original form]
\label{thm:godel-first}
Let \(T\) be a computably enumerable theory extending \(\mathsf Q\).
\begin{enumerate}
  \item If \(T\) is consistent, then \(T\notderives G_T\).
  \item If \(T\) is \(\Sigma^0_1\)-sound, then
        \(T\notderives\neg G_T\).
  \item Omega-consistency also suffices for
        \(T\notderives\neg G_T\).
\end{enumerate}
Consequently, if \(T\) is omega-consistent or \(\Sigma^0_1\)-sound, it is not
deductively complete.
\textit{Status: proof sketch.}
\end{theorem}

\begin{proofplan}
Suppose \(T\derives G_T\).  Because a concrete \(T\)-proof exists and the
proof predicate represents concrete proofs,
\(T\derives\mathsf{Prov}_T(\overline{\ulcorner G_T\urcorner})\).  But (G)
also gives
\(T\derives\neg\mathsf{Prov}_T(\overline{\ulcorner G_T\urcorner})\).
Thus \(T\) is inconsistent.  This proves (1).

If \(T\) is consistent, (1) says that no standard number codes a
\(T\)-proof of \(G_T\).  Hence
\[
  \mathcal N\entails
  \neg\mathsf{Prov}_T(\overline{\ulcorner G_T\urcorner}).
\]
The biconditional (G) is already a theorem of \(\mathsf Q\), not merely an
arbitrary theorem of \(T\); since \(\mathcal N\entails\mathsf Q\), it follows
that \(\mathcal N\entails G_T\).
Therefore \(\neg G_T\) is a false \(\Sigma^0_1\) sentence.  A
\(\Sigma^0_1\)-sound \(T\) cannot prove it, establishing (2).

For (3), suppose \(T\derives\neg G_T\).  By (G),
\[
  T\derives\exists p\,
  \mathsf{Proof}_T(p,\overline{\ulcorner G_T\urcorner}).
\]
For each standard \(n\), consistency and (1) imply that \(n\) is not the code
of a \(T\)-proof of \(G_T\).  Numeralwise representation of the primitive
recursive proof relation gives
\[
  T\derives\neg
  \mathsf{Proof}_T(\bar n,\overline{\ulcorner G_T\urcorner})
\]
for every \(n\).  These statements violate omega-consistency.
The argument is complete once the cited complexity-preserving fixed-point
construction and the cited representability theorem are supplied in full;
those dependencies are why the present status is a proof sketch.
\end{proofplan}

\begin{warning}[What the theorem does not say]
The theorem does not say that every formal theory is deductively incomplete,
that \(G_T\) is unprovable in every stronger theory, or that no deductively
incomplete theory can prove any consistency statement.  The hypotheses and
the subscript \(T\) are essential.
\end{warning}

\section{Rosser's strengthening}
\label{sec:first-incompleteness-04}

\begin{theorem}[Rosser]
\label{thm:rosser}
Let \(T\) be a consistent computably enumerable theory extending \(\mathsf Q\).
Then there is
a sentence \(R_T\) such that
\[
  T\notderives R_T
  \qquad\text{and}\qquad
  T\notderives\neg R_T.
\]
Thus consistency alone suffices for deductive incompleteness.
\textit{Status: proof sketch.}
\end{theorem}

The original source is \cite{Rosser1936}.

\begin{proofplan}
By \cref{thm:diagonal-lemma}, fix a sentence \(R_T\) with
\[
  \mathsf Q\derives R_T\leftrightarrow\forall y\,\bigl(
  \mathsf{Proof}_T(y,\overline{\ulcorner R_T\urcorner})\to
  \exists z\leq y\,
  \mathsf{Proof}_T(z,\overline{\ulcorner\neg R_T\urcorner})\bigr),
\]
so that \(R_T\) says: every proof of \(R_T\) has a proof of \(\neg R_T\)
with no larger code.  The bounded comparison uses the registered
object-language abbreviation \(\leq\).

First branch: suppose \(T\derives R_T\), and let \(p\) be the least code
of such a proof.  Numeralwise representation of the proof relation gives
\(T\derives\mathsf{Proof}_T(\bar p,\overline{\ulcorner R_T\urcorner})\),
so the fixed point yields
\(T\derives\exists z\leq\bar p\,
\mathsf{Proof}_T(z,\overline{\ulcorner\neg R_T\urcorner})\).
Consistency gives \(T\notderives\neg R_T\), so numeralwise representation
gives
\(T\derives\neg\mathsf{Proof}_T(\bar q,
\overline{\ulcorner\neg R_T\urcorner})\) for each \(q\leq p\).
\Cref{lem:q-bounded-cases} turns the bounded existential into the finite
disjunction over \(z=\bar0,\ldots,\bar p\), each disjunct refuted by the
preceding line, so \(T\) is inconsistent, contradicting consistency.  No
metatheoretic proof of \(\neg R_T\) is produced at any point.

Second branch: suppose \(T\derives\neg R_T\), and let \(q\) be the least
code of such a proof; numeralwise representation gives
\(T\derives\mathsf{Proof}_T(\bar q,
\overline{\ulcorner\neg R_T\urcorner})\).  Reason inside \(T\) from the
assumption \(\mathsf{Proof}_T(y,\overline{\ulcorner R_T\urcorner})\).
\Cref{lem:q-dichotomy} gives \(y\leq\bar q\lor\bar q\leq y\).  In the
first case, \cref{lem:q-bounded-cases} gives
\(y=\bar0\lor\dots\lor y=\bar q\); consistency gives
\(T\notderives R_T\), so numeralwise representation gives
\(T\derives\neg\mathsf{Proof}_T(\bar k,
\overline{\ulcorner R_T\urcorner})\) for each \(k\leq q\); every
disjunct is contradicted, and the bounded existential follows from the
contradiction.  In the second case, \(\bar q\leq y\) makes
\(z\defeq\bar q\) a witness for
\(\exists z\leq y\,
\mathsf{Proof}_T(z,\overline{\ulcorner\neg R_T\urcorner})\).
Discharging the assumption and applying the fixed point yields
\(T\derives R_T\), so \(T\) is inconsistent, again contradicting
consistency.

A complete proof supplies the cited representability theorem and diagonal
lemma in full; those dependencies are why the status remains a proof
sketch.  The bounded-comparison steps themselves are carried by
\cref{lem:q-bounded-cases,lem:q-dichotomy}, which have complete proofs.
\end{proofplan}

\begin{corollary}[Undecidability of theoremhood]
If \(T\) is a consistent computably enumerable extension of \(\mathsf Q\),
then the set of codes of theorems of \(T\) is computably enumerable but
undecidable.
\textit{Status: proof sketch.}
\end{corollary}

\begin{proofplan}
The theorem set is computably enumerable by enumeration of finite \(T\)-proofs.  If it were
decidable, then for each sentence \(\sigma\) one could decide whether
\(T\derives\sigma\).  Run this decider on both \(R_T\) and \(\neg R_T\); the
two negative answers are compatible with decidability, so this alone is not a
contradiction.  Instead enumerate the sentences as
\(\sigma_0,\sigma_1,\ldots\).  Starting with the empty finite set \(D_0\), let
\(D_{n+1}=D_n\cup\{\sigma_n\}\) if \(T\cup D_n\cup\{\sigma_n\}\) is
consistent, and otherwise let
\(D_{n+1}=D_n\cup\{\neg\sigma_n\}\).  The consistency test is decidable:
by the deduction theorem, inconsistency of the finite extension is equivalent
to \(T\) proving the negation of the conjunction of its finitely many added
sentences.  At least one choice preserves consistency, since otherwise
\(T\cup D_n\) would prove both \(\sigma_n\) and \(\neg\sigma_n\).

The theory \(T^*=T\cup\bigcup_nD_n\) is computably enumerable, consistent, extends
\(\mathsf Q\), and is deductively complete because a side of every sentence
is among its axioms.  Rosser's theorem applied to \(T^*\) gives an independent
sentence, a contradiction.
This derivation uses the preceding proof sketch of Rosser's theorem; a complete
proof supplies its bounded proof-comparison verification.
\end{proofplan}

\begin{remark}
An alternative proof reduces the self-application domain set to theoremhood by uniform
representability.  The proof above makes explicit why independence of one
fixed sentence does not by itself immediately contradict decidability of a
set.
\end{remark}

\begin{exercise}
\label{ex:first-incompleteness-01}
In the proof sketch of \cref{thm:godel-first}, identify the precise uses of
effectiveness, representability, consistency, and \(\Sigma^0_1\)-soundness.
\end{exercise}

\begin{exercise}
\label{ex:first-incompleteness-02}
Show that if \(T\) is consistent, then the G\"odel sentence \(G_T\) constructed
above is true in \(\mathcal N\).  Explain why this does not imply that
\(T\cup\{G_T\}\) proves every arithmetical truth.
\end{exercise}

\begin{exercise}
\label{ex:first-incompleteness-03}
Why can the proof sketch of Rosser's theorem compare proof codes with a
bounded
quantifier without assuming omega-consistency?  State the role of the least
proof code.
\end{exercise}

\chapter{Provability, the Second Incompleteness Theorem, and Truth}
\label{ch:second-incompleteness}

\section{Derivability conditions}
\label{sec:second-incompleteness-01}

Fix a computably enumerable theory \(T\supseteq\mathsf Q\) and use the proof predicate and
associated formula \(\mathsf{Prov}_T(x)\) previously fixed in
\cref{def:fixed-proof-predicate}.  To reduce clutter, write
\(\mathsf{Prov}_T(\ulcorner\varphi\urcorner)\) for
\(\mathsf{Prov}_T(\overline{\ulcorner\varphi\urcorner})\).

\begin{definition}[Hilbert--Bernays--L\"ob derivability conditions]
The provability predicate satisfies the \emph{derivability conditions} when,
for all sentences \(\varphi\) and \(\psi\):
\begin{enumerate}
  \item if \(T\derives\varphi\), then
        \(T\derives\mathsf{Prov}_T(\ulcorner\varphi\urcorner)\);
  \item
  \[
    T\derives
    \mathsf{Prov}_T(\ulcorner\varphi\to\psi\urcorner)
    \to
    \bigl(\mathsf{Prov}_T(\ulcorner\varphi\urcorner)
    \to\mathsf{Prov}_T(\ulcorner\psi\urcorner)\bigr);
  \]
  \item
  \[
    T\derives
    \mathsf{Prov}_T(\ulcorner\varphi\urcorner)
    \to
    \mathsf{Prov}_T
    (\ulcorner\mathsf{Prov}_T(\ulcorner\varphi\urcorner)\urcorner).
  \]
\end{enumerate}
The first clause is a metatheoretic implication; the latter two are theorem
schemes internal to \(T\).  The three clauses are cited below as D1, D2,
and D3.
\end{definition}

For \(T=\PA\) and its fixed proof predicate, all three derivability conditions
hold: clause~1 is the displayed metatheoretic implication, while \(\PA\)
derives the schemes in clauses~2 and~3 by arithmetizing proof concatenation and
proof verification \cite{HilbertBernays1939,Lob1955}.  For an arbitrary
computably enumerable \(T\supseteq\mathsf Q\), the results below state the
three conditions as hypotheses rather than inferring them from a vague
strength requirement.

\begin{definition}[Formal consistency sentence]
\label{def:formal-consistency-sentence}
Fix \(\bottom_A\) as the false arithmetic sentence \(0=S0\).  Define
\[
  \mathsf{Con}(T)\defeq
  \neg\mathsf{Prov}_T(\overline{\ulcorner\bottom_A\urcorner}).
\]
This formula expresses consistency relative to the fixed numbering and proof
predicate.
\end{definition}

\begin{warning}[External and internal consistency]
The metalinguistic claim that \(T\) is consistent is not literally the formula
\(\mathsf{Con}(T)\).  Correctness of the chosen proof predicate connects them
in the standard structure \(\mathcal N\); the second incompleteness
theorem concerns whether
\(T\) proves this particular internal sentence.
\end{warning}

\section{L\"ob's theorem}
\label{sec:second-incompleteness-02}

\begin{theorem}[L\"ob]
\label{thm:lob}
Let \(T\supseteq\mathsf Q\) be computably enumerable, and assume its previously
fixed provability predicate satisfies the three derivability conditions.  For
every sentence
\(\varphi\), if
\[
  T\derives
  \mathsf{Prov}_T(\ulcorner\varphi\urcorner)\to\varphi,
\]
then \(T\derives\varphi\).
\textit{Status: complete proof.}
\end{theorem}

\begin{proof}
By the diagonal lemma, choose a sentence \(\psi\) such that
\[
  T\derives
  \psi\leftrightarrow
  \bigl(\mathsf{Prov}_T(\ulcorner\psi\urcorner)\to\varphi\bigr).
  \tag{L1}
\]
From the left-to-right direction of (L1), derivability conditions D1 and D2
give
\[
  T\derives
  \mathsf{Prov}_T(\ulcorner\psi\urcorner)
  \to
  \mathsf{Prov}_T
  (\ulcorner\mathsf{Prov}_T(\ulcorner\psi\urcorner)\to\varphi\urcorner).
  \tag{L2}
\]
Condition D3 gives
\[
  T\derives
  \mathsf{Prov}_T(\ulcorner\psi\urcorner)
  \to
  \mathsf{Prov}_T
  (\ulcorner\mathsf{Prov}_T(\ulcorner\psi\urcorner)\urcorner).
  \tag{L3}
\]
Apply condition D2, internalized modus ponens, to (L2) and (L3).  Then
\[
  T\derives
  \mathsf{Prov}_T(\ulcorner\psi\urcorner)
  \to\mathsf{Prov}_T(\ulcorner\varphi\urcorner).
\]
By the theorem's hypothesis,
\[
  T\derives
  \mathsf{Prov}_T(\ulcorner\psi\urcorner)\to\varphi.
\]
The right-to-left direction of (L1) yields \(T\derives\psi\).  Condition D1
now gives \(T\derives\mathsf{Prov}_T(\ulcorner\psi\urcorner)\), and the last
display yields \(T\derives\varphi\).
\end{proof}

\section{The second incompleteness theorem}
\label{sec:second-incompleteness-03}

\begin{theorem}[G\"odel's second incompleteness theorem]
\label{thm:godel-second}
Let \(T\) be a consistent computably enumerable extension of \(\mathsf Q\) whose provability
predicate, previously fixed in \cref{def:fixed-proof-predicate}, satisfies the
derivability conditions.  Then
\[
  T\notderives\mathsf{Con}(T).
\]
\textit{Status: complete proof.}
\end{theorem}

\begin{proof}
Suppose \(T\derives\mathsf{Con}(T)\).  From
\(\neg\mathsf{Prov}_T(\overline{\ulcorner\bottom_A\urcorner})\), assume
\(\mathsf{Prov}_T(\overline{\ulcorner\bottom_A\urcorner})\); negation
elimination yields \(\bottom\), and bottom elimination yields \(\bottom_A\).
Implication introduction therefore gives
\[
  T\derives
  \mathsf{Prov}_T(\overline{\ulcorner\bottom_A\urcorner})\to\bottom_A.
\]
The hypothesis of L\"ob's theorem now holds with \(\varphi=\bottom_A\).
Therefore \(T\derives\bottom_A\), contradicting the assumed consistency of
\(T\).
\end{proof}

\begin{warning}[Scope of the second theorem]
The conclusion is not ``\(T\) cannot prove any sentence expressing any notion
of consistency.''  It concerns the arithmetized proof predicate fixed in
\cref{def:fixed-proof-predicate} and the corresponding sentence
\(\mathsf{Con}(T)\).  Other deliberately weak predicates can fail the
derivability conditions.
\end{warning}

\begin{example}[Relative consistency is not excluded]
A stronger theory \(T'\) may prove \(\mathsf{Con}(T)\); the second theorem
places no prohibition on such a relative result.  For example, sufficiently
strong standard set theories verify the existence of the standard natural
number structure and its satisfaction of the axioms of \(\PA\), from which
they formalize \(\mathsf{Con}(\PA)\).  This is a claim about what the stronger
theory proves, not a claim that its consistency alone implies
\(\mathsf{Con}(T)\).  The theorem says that a consistent \(T\) cannot prove its
own fixed consistency sentence under the stated hypotheses.
\end{example}

\section{Tarski's undefinability theorem}
\label{sec:second-incompleteness-04}

\begin{definition}[Arithmetical truth definition]
An \(\Lang_A\)-formula \(\mathsf{Tr}(x)\) is an \emph{arithmetical truth
definition} if, for every \(\Lang_A\)-sentence \(\sigma\),
\[
  \mathcal N\entails
  \mathsf{Tr}(\overline{\ulcorner\sigma\urcorner})
  \leftrightarrow\sigma.
\]
The biconditional is evaluated in the standard structure and ranges
metalinguistically over all arithmetic sentences.
\end{definition}

\begin{theorem}[Tarski undefinability]
\label{thm:tarski-undefinability}
There is no arithmetical truth definition for the full set of true
\(\Lang_A\)-sentences.
\textit{Status: complete proof.}
\end{theorem}

\begin{proof}
Suppose \(\mathsf{Tr}(x)\) were such a formula.  Apply the diagonal lemma to
\(\neg\mathsf{Tr}(x)\).  There is a sentence \(L\) such that
\[
  \mathsf Q\derives
  L\leftrightarrow
  \neg\mathsf{Tr}(\overline{\ulcorner L\urcorner}).
\]
Because \(\mathcal N\entails\mathsf Q\), the biconditional is true in
\(\mathcal N\).  But the assumed truth definition also gives
\[
  \mathcal N\entails
  \mathsf{Tr}(\overline{\ulcorner L\urcorner})\leftrightarrow L.
\]
Combining the two yields
\(\mathcal N\entails L\leftrightarrow\neg L\), impossible in the classical
standard structure.
\end{proof}

\begin{corollary}
The set of G\"odel numbers
\[
  \set{\ulcorner\sigma\urcorner:
  \sigma\text{ is an arithmetic sentence and }
  \mathcal N\entails\sigma}
\]
is not definable in \(\mathcal N\) by any \(\Lang_A\)-formula.  In
particular, it is not computably
enumerable.
\end{corollary}

\begin{proof}
A formula defining the displayed set in \(\mathcal N\) would be an
arithmetical truth
definition.  Every computably enumerable set is \(\Sigma^0_1\)-definable, so
nondefinability implies that the set is not computably enumerable.
\end{proof}

\begin{remark}[Partial truth definitions]
For each fixed level of the arithmetical hierarchy, arithmetic can define a
satisfaction predicate for formulas of that bounded complexity.  Tarski's
theorem excludes one arithmetic formula correct for all arithmetic sentences
at once.  This distinction is essential in reflection principles and proof
theory \cite{Tarski1983}.
\end{remark}

\section{A final audit table}
\label{sec:second-incompleteness-05}

\begin{longtable}{@{}p{0.24\textwidth}p{0.32\textwidth}p{0.34\textwidth}@{}}
\toprule
Claim & Required hypotheses & Common invalid strengthening \\
\midrule
\endhead
First incompleteness, G\"odel sentence & \(T\) computably enumerable,
\(T\supseteq\mathsf Q\); consistency for unprovability; omega-consistency or
\(\Sigma^0_1\)-soundness for independence of \(G_T\), that is, for
underivability of both \(G_T\) and \(\neg G_T\) & Every consistent theory
has the displayed G\"odel sentence independent under no effectiveness or
arithmetic-strength assumption. \\
\addlinespace
First incompleteness, Rosser & \(T\) consistent, computably enumerable, and
\(T\supseteq\mathsf Q\) & Every consistent theory whatever, with no
computable-enumerability or \(\mathsf Q\)-extension hypothesis, has an
independent sentence. \\
\addlinespace
Second incompleteness & \(T\supseteq\mathsf Q\) consistent and computably
enumerable; the previously fixed provability predicate satisfies all three
derivability conditions &
No theory can prove a consistency statement about itself or any weaker theory.
\\
\addlinespace
L\"ob's theorem & Diagonal lemma and all three derivability conditions & Any
informal predicate called ``provable'' satisfies L\"ob's theorem. \\
\addlinespace
Tarski undefinability & One arithmetic formula purporting to define truth for
all arithmetic sentences in \(\mathcal N\) & No restricted satisfaction or
partial truth predicate can be defined. \\
\bottomrule
\end{longtable}

\begin{exercise}
\label{ex:second-incompleteness-01}
Reprove \cref{thm:godel-second} directly from a G\"odel sentence by showing in
\(T\) that \(\mathsf{Con}(T)\to G_T\).  Identify the derivability conditions
used in the internal argument.
\end{exercise}

\begin{exercise}
\label{ex:second-incompleteness-02}
Let \(T\) satisfy the hypotheses of L\"ob's theorem.  Prove that if
\(T\derives\mathsf{Prov}_T(\ulcorner\varphi\urcorner)
\leftrightarrow\varphi\), then \(T\derives\varphi\).
\end{exercise}

\begin{exercise}
\label{ex:second-incompleteness-03}
Formulate a syntactic version of Tarski's theorem: show that no consistent
theory extending \(\mathsf Q\) can prove every biconditional
\[
  \mathsf{Tr}(\overline{\ulcorner\sigma\urcorner})
  \leftrightarrow\sigma
\]
for all arithmetic sentences \(\sigma\).  Use the diagonal lemma and state
where consistency is required.
\end{exercise}

\begin{exercise}
\label{ex:second-incompleteness-04}
Audit the sentence ``G\"odel proved that true statements cannot be proved.''
Replace it by three distinct correct statements concerning \(G_T\),
\(\mathsf{Con}(T)\), and \(\Th(\mathcal N)\), each with exact hypotheses.
\end{exercise}
\part{Axiomatic Set Theory}
\label{part:set-theory}

\chapter[The Axioms of ZF and ZFC]{The Axioms of \(\ZF\) and \(\ZFC\)}
\label{ch:zf-zfc}

\begin{convention}[Scoped semantic notation]
Throughout these advanced chapters, \(\mathcal M\) and \(\mathcal N\) are the
normative metavariables for semantic structures.  The letter \(\mathcal H\) is
reserved for the explicitly defined Henkin structure.  The word \emph{model}
is reserved for a structure together with an assertion that it satisfies a
specified sentence or set of sentences.  Letters assigned a different local
role, such as forcing names or type metavariables, are declared at the start of
that scope.
\end{convention}

Set theory supplies a common metatheory for syntax, structures, proofs, and
cardinality.  It can also be studied as an object theory.  These two roles must
remain separate.  In this part, \(\ZF\) and \(\ZFC\) are first-order theories in
the one-sorted language \(\Lang_{\in}=\set{\in}\), where equality is logical
identity and \(\in\) is a binary relation symbol.  Unless a displayed formula
is explicitly quoted, it is written in the metalanguage as a representative of
an object-language axiom or axiom instance.

The treatment follows the standard organization of \cite{Jech2003}.  A
relative-consistency assertion below is a metalogical implication between
arithmetized consistency statements.  It is not an assertion that the stronger
theory is absolutely consistent.

\section{The language and the axiom schemas}
\label{sec:zf-zfc-01}

\begin{convention}[Classes]
A formula \(\varphi(x,\bar p)\) may be used to abbreviate the collection
\[
  \set{x:\varphi(x,\bar p)}.
\]
This class notation is metalinguistic.  It does not assert that the collection
is a set, and it does not add class variables to \(\Lang_{\in}\).  A class is a
set iff an instance of an axiom or a theorem supplies a set having the
specified members.
\end{convention}

\begin{definition}[Extensionality]
The axiom of extensionality is
\[
  \forall a\,\forall b\,
  \bigl(\forall x\,(x\in a\leftrightarrow x\in b)\to a=b\bigr).
\]
It states that membership determines set identity.
\end{definition}

\begin{definition}[Empty set]
The empty-set axiom is
\[
  \exists a\,\forall x\,(x\notin a).
\]
Extensionality makes the witness unique; it is denoted by \(\varnothing\).
\end{definition}

\begin{definition}[Pairing, union, and power set]
The pairing axiom is
\[
  \forall a\,\forall b\,\exists c\,\forall x\,
  \bigl(x\in c\leftrightarrow(x=a\lor x=b)\bigr).
\]
The union axiom is
\[
  \forall a\,\exists u\,\forall x\,
  \bigl(x\in u\leftrightarrow\exists y\,(y\in a\land x\in y)\bigr).
\]
The power-set axiom is
\[
  \forall a\,\exists p\,\forall x\,
  \bigl(x\in p\leftrightarrow x\subseteq a\bigr),
  \qquad
  x\subseteq a\;:\Longleftrightarrow\;
  \forall z\,(z\in x\to z\in a).
\]
The unique witnesses are written \(\set{a,b}\), \(\bigcup a\), and \(\Pow(a)\),
respectively.  The singleton \(\set{a}\) abbreviates \(\set{a,a}\).
\end{definition}

\begin{definition}[Infinity]
Let \(S(x)\defeq x\cup\set{x}\).  The axiom of infinity is
\[
  \exists I\,
  \bigl(\varnothing\in I\land
        \forall x\,(x\in I\to S(x)\in I)\bigr).
\]
It asserts the existence of an inductive set.  The set \(\omega\) is
defined below as the least inductive set; that leastness is a theorem using
separation, not part of the axiom of infinity.
\end{definition}

\begin{definition}[Separation schema]
For every formula \(\varphi(x,\bar p)\) of \(\Lang_{\in}\) in which \(b\) is
not free, the corresponding instance of separation is
\[
  \forall a\,\exists b\,\forall x\,
  \bigl(x\in b\leftrightarrow(x\in a\land\varphi(x,\bar p))\bigr).
\]
The parameters \(\bar p\) are arbitrary sets.  Separation forms a subset of a
previously given set; it does not assert an unrestricted comprehension
principle.
\end{definition}

\begin{definition}[Replacement schema]
For every formula \(\varphi(x,y,\bar p)\) of \(\Lang_{\in}\) in which \(b\) is
not free, the corresponding instance of replacement is
\[
\begin{aligned}
  \forall a\,\bigl[&
    \forall x\in a\,\exists!y\,\varphi(x,y,\bar p)\to{}\\
    &\exists b\,\forall y\,
      \bigl(y\in b\leftrightarrow
            \exists x\in a\,\varphi(x,y,\bar p)\bigr)
  \bigr].
\end{aligned}
\]
The antecedent is a functionality requirement on the set \(a\); the
quantifier abbreviation \(\exists!y\,\psi(y)\) stands for
\(\exists y\,(\psi(y)\land\forall z\,(\psi(z)\to z=y))\) with \(z\)
fresh.  Replacement
then asserts that the image is a set.  The schema is not one first-order
sentence: each eligible formula supplies one axiom instance.
\end{definition}

\begin{definition}[The set \(\omega\)]
By infinity, fix an inductive set \(I_0\).  By separation, form
\[
  \omega\defeq\set{x\in I_0:
  x\text{ belongs to every inductive set}},
\]
where the membership condition abbreviates the \(\Lang_{\in}\)-formula
obtained from the displayed inductive condition of the axiom of infinity.
Every member of \(\omega\) belongs to every inductive set, so \(\omega\)
is included in every inductive set.  The set \(\omega\) is itself
inductive, because \(\varnothing\) belongs to every inductive set and
inductive sets are closed under \(S\).  Extensionality makes \(\omega\)
unique.  This discharges the promise recorded with the axiom of infinity.
\end{definition}

\begin{definition}[Foundation]
The axiom of foundation is
\[
  \forall a\,
  \bigl(a\ne\varnothing\to
        \exists x\in a\,\forall y\in x\,(y\notin a)\bigr).
\]
Thus every nonempty set has a member disjoint from it.
\end{definition}

\begin{definition}[Axiom of Choice]
This definition introduces the object-language axiom of \(\ZFC\); it is
distinct in standing from the metatheoretic Axiom of Choice of
\cref{sec:set-prerequisites-06}, which governs the ambient set theory.
Which of the two a later argument uses is determined by whether the
assertion is a formula of \(\Lang_{\in}\) or a metatheoretic principle.
The \emph{Axiom of Choice} states that for every set \(A\) of nonempty sets there is a
function \(f\) with domain \(A\) such that \(f(x)\in x\) for every \(x\in A\):
here the Kuratowski pair notation is the object-language abbreviation
\(\tuple{u,v}\defeq\set{\set u,\set{u,v}}\).  In the following display, the
other two abbreviations are object-language abbreviations for the indicated
membership formulas,
\[
\begin{aligned}
  \operatorname{Fun}(f)
  \quad:\Longleftrightarrow\quad
  &\bigl[\forall z\in f\,\exists u\,\exists v\,
          z=\tuple{u,v}\bigr]\ \land{}\\
  &\forall u\,\forall v\,\forall w\,
    \bigl((\tuple{u,v}\in f\land\tuple{u,w}\in f)\to v=w\bigr),\\
  \dom(f)=A
  \quad:\Longleftrightarrow\quad
  &\forall u\,
    \bigl(u\in A\leftrightarrow
          \exists v\,\tuple{u,v}\in f\bigr).
\end{aligned}
\]
Thus an exact object-language instance is
\[
  \forall A\,
  \left[
    \forall x\in A\,(x\ne\varnothing)
    \to
    \exists f\,
    \left(
      \operatorname{Fun}(f)\land\dom(f)=A\land
      \forall x\in A\,\exists y\,
        (\tuple{x,y}\in f\land y\in x)
    \right)
  \right].
\]
The first conjunct says that every member of \(f\) is an ordered pair and that
two pairs in \(f\) with the same first coordinate have the same second
coordinate.  The domain clause and the final conjunct then give exactly one
chosen member of each member of \(A\).  Pairwise disjointness is not required.
\end{definition}

\begin{definition}[The theories \(\ZF\) and \(\ZFC\)]
The theory \(\ZF\) consists of extensionality, empty set, pairing, union, power
set, infinity, foundation, every separation instance, and every replacement
instance.  The theory \(\ZFC\) is \(\ZF\) together with the Axiom of Choice.
\end{definition}

\begin{remark}
Some axiom lists derive the empty-set axiom from another axiom and separation.
That change is an alternative axiomatization, not a different set-theoretic
claim.  The active axiom list above is fixed for this document.
\end{remark}

\section{First constructions and exact uses of the axioms}
\label{sec:zf-zfc-02}

Recall that the Kuratowski ordered pair is
\[
  \tuple{a,b}\defeq\set{\set{a},\set{a,b}}.
\]
A relation is a set of ordered pairs.  A function is a relation having exactly
one output for each member of its domain.  These definitions are eliminable in
the language of membership.

\begin{lemma}[Ordered-pair identity]
For all sets \(a,b,c,d\),
\[
  \tuple{a,b}=\tuple{c,d}
  \quad\Longleftrightarrow\quad
  (a=c\land b=d).
\]
\end{lemma}

\begin{proof}
The right-to-left direction follows by substitution into the definition.  For
the converse, compare the members of the two Kuratowski pairs by
extensionality.  If \(a=b\), the left pair is \(\set{\set a}\), and equality
forces \(c=d=a\).  If \(a\ne b\), then \(\set a\) and \(\set{a,b}\) are
distinct; equality identifies the singleton member and then the two-element
member, yielding \(a=c\) and \(b=d\).
\end{proof}

\begin{example}[Foundation excludes a membership loop]
Suppose \(a\in a\).  Pairing supplies \(A=\set a\).  Foundation applied to
\(A\) requires an \(x\in A\) such that no member of \(x\) belongs to \(A\).
The only possible \(x\) is \(a\), but \(a\in a\) and \(a\in A\), a
contradiction.  Hence \(a\notin a\) for every set \(a\).

This argument proves one consequence of foundation.  It does not prove that
the nonexistence of membership loops by itself is equivalent to foundation.
\end{example}

\begin{theorem}[Set recursion]
Let \(G\) be a class operation specified by a formula such that, for every set
\(x\) and every function \(g\) with domain \(x\), there is exactly one set
\(G(x,g)\).  Then there is a unique class function \(F\) satisfying
\[
  F(x)=G(x,F\restr x)
\]
for every set \(x\), where the recursion is taken over the well-founded
membership relation.  Every restriction of \(F\) to a set is a set function.
\end{theorem}

\paragraph{Proof sketch.}
Foundation supplies well-founded induction.  The \emph{transitive closure}
of a set \(x\) is the least transitive set including \(x\), obtained as
\(\bigcup_{n\in\omega}u_n\) with \(u_0=x\) and
\(u_{n+1}=\bigcup u_n\); replacement and union make it a set, and it is
distinct from the reflexive transitive closure of a binary relation defined
in \cref{ch:strings}.  Define approximations on the
transitive closure of a given set and use replacement to collect the values.
Extensionality proves uniqueness by choosing a membership-minimal input at
which two proposed solutions differ.  A complete proof requires the previously
developed coding of functions and is therefore cited to \cite{Jech2003}; this
paragraph records the proof architecture, not a line-by-line proof.

\begin{warning}[Unrestricted comprehension]
The notation \(\set{x:\varphi(x)}\) does not guarantee a set.  In particular,
the class \(\set{x:x\notin x}\) cannot be a set.  Separation permits
\(\set{x\in a:x\notin x}\) only after a bounding set \(a\) has been supplied.
\end{warning}

\begin{exercise}
\label{ex:zf-zfc-01}
Identify the axioms used to construct \(a\times b\) as a set.  Your proof must
give a bounding set for the ordered pairs before applying separation.
\end{exercise}

\begin{exercise}
\label{ex:zf-zfc-02}
Show that foundation excludes a two-cycle \(a\in b\in a\).  Do not appeal to a
diagram; apply the axiom to a specifically constructed nonempty set.
\end{exercise}

\chapter[Hierarchy, Ordinals, and Cardinals]{The Cumulative Hierarchy, Ordinals, and Cardinals}
\label{ch:cumulative-hierarchy}

\section{Transitive sets and ordinals}
\label{sec:cumulative-hierarchy-01}

Transitive sets and ordinals were defined in
\cref{def:transitive-set,def:ordinal-cardinal}.  This chapter uses those exact
definitions inside set theory rather than introducing new terms for them.

For ordinals \(\alpha\) and \(\beta\), write \(\alpha<\beta\) iff
\(\alpha\in\beta\).  Define
\[
  \alpha+1\defeq\alpha\cup\set{\alpha}.
\]
An ordinal is a \emph{successor ordinal} iff it has this form; a nonzero ordinal
that is not a successor ordinal is a \emph{limit ordinal}.  These terms name
ordinal forms, not arithmetic operations on arbitrary sets.

\begin{theorem}[Trichotomy and transitivity of ordinals]
For ordinals \(\alpha\) and \(\beta\), exactly one of
\[
  \alpha\in\beta,\qquad \alpha=\beta,\qquad \beta\in\alpha
\]
holds.  Every member of an ordinal is an ordinal.
\end{theorem}

\paragraph{Proof sketch.}
The route runs through \(\alpha\cap\beta\).  First, \(\alpha\cap\beta\) is an
ordinal: it is transitive as an intersection of transitive sets, and \(\in\)
strictly well-orders it as a subset of \(\alpha\).  Second, an ordinal that
is a proper subset of an ordinal is a member of it: if
\(\gamma\subsetneq\delta\), the well-order of \(\delta\) supplies the
\(\in\)-least element \(\xi\) of \(\delta\setminus\gamma\); minimality of
\(\xi\) together with transitivity of \(\delta\) puts every member of
\(\xi\) in \(\gamma\), and linearity of \(\in\) on \(\delta\) together with
transitivity of \(\gamma\) puts every member of \(\gamma\) in \(\xi\), so
\(\gamma=\xi\in\delta\).  Now if \(\alpha\cap\beta\) were a proper subset of
both \(\alpha\) and \(\beta\), the second step would give
\(\alpha\cap\beta\in\alpha\cap\beta\), contradicting the irreflexivity of
the strict well-order on the ordinal \(\alpha\cap\beta\); hence
\(\alpha\cap\beta=\alpha\) or \(\alpha\cap\beta=\beta\), and the second step
converts the surviving inclusion, when proper, into \(\alpha\in\beta\) or
\(\beta\in\alpha\).  At most one alternative holds: any two of them combine,
using transitivity of the sets involved, into an ordinal that is a member of
itself.  For the final claim, a member \(x\) of an ordinal \(\alpha\) is a
subset of \(\alpha\) by transitivity of \(\alpha\), so \(\in\) strictly
well-orders \(x\); and \(x\) is transitive because transitivity of
\(\alpha\) puts the members of members of \(x\) inside \(\alpha\), where the
transitivity of the strict order \(\in\) on \(\alpha\) yields the required
memberships.  The omitted obligations are the two elementwise membership
calculations in the second step; they are supplied by \cite{Jech2003}.

\begin{theorem}[Transfinite induction]
Let \(\varphi(\alpha,\bar p)\) be a formula.  If
\[
  \forall\alpha\,
  \left[
    \bigl(\forall\beta<\alpha\,\varphi(\beta,\bar p)\bigr)
    \to\varphi(\alpha,\bar p)
  \right],
\]
then \(\varphi(\alpha,\bar p)\) holds for every ordinal \(\alpha\).
\end{theorem}

\begin{proof}
If the conclusion fails at some ordinal \(\delta\), separation applied to
\(\delta+1\) produces the nonempty set
\(\set{\beta\in\delta+1:\neg\varphi(\beta,\bar p)}\).  The ordinal
well-order supplies its least member \(\alpha\).  Every \(\beta<\alpha\) satisfies \(\varphi\), so
the displayed hypothesis gives \(\varphi(\alpha,\bar p)\), a contradiction.
\end{proof}

\begin{theorem}[Transfinite recursion]
Let \(G\) be a definable class operation assigning one set \(G(f)\) to every
function \(f\) whose domain is an ordinal.  There is a unique class function
\(F\) on the ordinals such that
\[
  F(\alpha)=G(F\restr\alpha)
\]
for every ordinal \(\alpha\).  For each ordinal \(\gamma\), the restriction
\(F\restr\gamma\) is a set.
\end{theorem}

\paragraph{Proof sketch.}
For each ordinal \(\gamma\), use transfinite induction to construct the unique
function \(f_\gamma\) with domain \(\gamma\) satisfying
\(f_\gamma(\alpha)=G(f_\gamma\restr\alpha)\).  At a successor stage append
the uniquely specified value; at a limit stage take the union of the coherent
earlier functions.  Replacement collects the earlier values and separation
forms the required graphs.  Transfinite induction proves
\(f_\beta=f_\gamma\restr\beta\) whenever \(\beta<\gamma\), so their union is
the asserted class function; the same induction proves uniqueness.  The
omitted Separation and Replacement instances are supplied in
\cite[Chapter 2]{Jech2003}.

Replacement is essential in the general theorem: it collects the values
already constructed at earlier stages.

\section{The cumulative hierarchy and rank}
\label{sec:cumulative-hierarchy-02}

\begin{definition}[Cumulative hierarchy]
By transfinite recursion define
\[
  V_0=\varnothing,\qquad
  V_{\alpha+1}=\Pow(V_\alpha),\qquad
  V_\lambda=\bigcup_{\beta<\lambda}V_\beta
\]
when \(\lambda\) is a limit ordinal.  The symbol
\[
  V=\bigcup_{\alpha\in\operatorname{Ord}}V_\alpha
\]
is class notation, not a set-forming union; \(\operatorname{Ord}\)
denotes the class of all ordinals, which the Burali-Forti argument shows to
be a proper class.
\end{definition}

\begin{theorem}[Hierarchy theorem]
Assume \(\ZF\).  Then:
\begin{enumerate}
  \item every \(V_\alpha\) is transitive;
  \item if \(\alpha\leq\beta\), then \(V_\alpha\subseteq V_\beta\);
  \item every set belongs to some \(V_\alpha\); and
  \item for every set \(x\), there is a least ordinal \(\alpha\) such that
        \(x\subseteq V_\alpha\).
\end{enumerate}
\end{theorem}

\paragraph{Proof sketch.}
The first two clauses follow by transfinite induction.  For the third,
suppose some set \(x\) belongs to no stage.  Foundation, applied to the
set of members of the transitive closure of \(\set{x}\) that belong to no
stage, supplies a membership-minimal counterexample.
All of its members occur at stages; replacement bounds their stage indices by
a set of ordinals.  Their union gives one stage containing every member, so the
counterexample appears at the next stage.  The fourth clause uses ordinal
well-ordering.  The dependency structure is recorded; routine coding details
are omitted.

\begin{definition}[Rank]
The rank of a set \(x\) is
\[
  \rank(x)\defeq
  \min\set{\alpha:x\subseteq V_\alpha}.
\]
Equivalently,
\[
  \rank(x)=\sup\set{\rank(y)+1:y\in x},
\]
where, for a set \(S\) of ordinals, \(\sup S\defeq\bigcup S\) is the
least ordinal that is at least every member of \(S\).
The equivalence is a theorem obtained from the hierarchy theorem and
transfinite recursion.
\end{definition}

\begin{example}[The first stages]
\[
  V_0=\varnothing,\qquad
  V_1=\set{\varnothing},\qquad
  V_2=\set{\varnothing,\set{\varnothing}},
\]
and \(V_3=\Pow(V_2)\) has four members.  Consequently
\(\rank(\varnothing)=0\), \(\rank(\set{\varnothing})=1\), and
\(\rank(V_2)=2\).  The index of the first stage containing \(x\) as a member
is \(\rank(x)+1\); it is not \(\rank(x)\).
\end{example}

\section{Well-founded relations and the Mostowski collapse}
\label{sec:cumulative-hierarchy-03}

\begin{definition}[Well-founded relation]
Let \(R\subseteq A\times A\).  The relation \(R\) is well-founded on \(A\)
iff every nonempty subset \(B\subseteq A\) has an \(R\)-minimal
member:
\[
  \forall B\subseteq A\,
  \bigl(B\ne\varnothing\to
        \exists b\in B\,\forall c\in B\,\neg(cRb)\bigr).
\]
It is extensional iff distinct elements have distinct sets of
\(R\)-predecessors.
\end{definition}

\begin{theorem}[Mostowski collapse]
Let \(A\) be a set and let \(R\subseteq A\times A\) be well-founded and
extensional.  There are a unique transitive set \(M\) and a unique isomorphism
\(\pi:(A,R)\to(M,\in)\).  They satisfy
\[
  \pi(a)=\set{\pi(b):bRa}
\]
for every \(a\in A\).
\end{theorem}

\paragraph{Proof sketch.}
Well-founded recursion---recursion along the well-founded set relation
\(R\), obtained from the set-recursion theorem via the rank function of
\(R\)---defines \(\pi\) by the displayed equation.  Replacement
makes each displayed collection a set.  An induction on \(R\) proves that
extensionality makes \(\pi\) injective.  Its range is transitive by the defining
equation, and the same equation proves preservation and reflection of the
relations.  A second well-founded induction proves uniqueness.  The
underlying recursion principle for arbitrary well-founded set relations and
the separation and replacement instances are supplied in
\cite[Chapter~6]{Jech2003}.

\section{Cardinals and the Axiom of Choice}
\label{sec:cumulative-hierarchy-04}

\begin{definition}[Equipotence]
Sets \(A\) and \(B\) are \emph{equipotent}, written \(A\approx B\), iff
a bijection \(A\to B\) exists.  Thus a cardinal, in the sense recalled
from \cref{def:ordinal-cardinal}, is an ordinal not equipotent
with any smaller ordinal.
\end{definition}

In \(\ZF\), a set \(A\) has an ordinal cardinal \(\abs A\) iff \(A\)
is well-orderable.  Do not write \(\abs A\) for an arbitrary set before
well-orderability has been established.  Under the Axiom of Choice, every set is
well-orderable and therefore has a unique cardinal.

\begin{theorem}[Hartogs]
For every set \(A\), there is a least ordinal \(h(A)\) for which no injection
\(h(A)\to A\) exists.  This theorem is provable in \(\ZF\), without the Axiom
of Choice.
\end{theorem}

\paragraph{Proof sketch.}
The well-orders carried by subsets of \(A\) form a set.  Collapse each such
well-order to its unique ordinal and use replacement to collect the resulting
ordinals.  Their supremum plus one cannot inject into \(A\).  Leastness
follows from ordinal well-ordering.  The set-coding of well-orders by
subsets of \(A\) and the replacement instance are supplied in
\cite{Jech2003}.

\begin{definition}[Order vocabulary]
A \emph{partial order} on a set \(A\) is a binary relation on \(A\) that
is reflexive, transitive, and \emph{antisymmetric}: whenever it relates
\(x\) to \(y\) and \(y\) to \(x\), then \(x=y\).  A \emph{partially
ordered set} is a set with a fixed partial order.  A \emph{chain} is a
subset on which the order is total.  An \emph{upper bound} of a subset is
an element that is greater than or equal to every member of that subset.  A
\emph{maximal element} is an element strictly below no element; a
\emph{greatest member} is an element greater than or equal to every
element.  One well-order \emph{end-extends} another
iff it extends it and every added element is greater than every original
element.
\end{definition}

\begin{theorem}[Equivalent forms of the Axiom of Choice]
Over \(\ZF\), each of the following implies the others:
\begin{enumerate}
  \item the Axiom of Choice;
  \item every set can be well-ordered;
  \item every partially ordered set in which every chain has an upper bound
        has a maximal element; and
  \item every product of a set-indexed family of nonempty sets is nonempty.
\end{enumerate}
\end{theorem}

\paragraph{Proof sketch.}
The implication from the Axiom of Choice to the maximal-element principle is proved by
choosing, whenever possible, a strict extension and applying transfinite
recursion; Hartogs' theorem prevents the construction from continuing through
all ordinals.  The maximal-element principle applied to partial well-orders
ordered by end-extension yields a well-order of an arbitrary set.  A well-order
of the union of a set-indexed family of nonempty sets selects the least member
of each factor, giving a member of the product.  Finally, a member of every
such product is a choice function.  This records the implication cycle and its
use of replacement; the set-coding details are imported from
\cite[Chapter 5]{Jech2003}.

The four clauses of the preceding equivalence theorem are equivalent over
the stated base theory; they are not definitions of one another.  Any use
of one clause in a proof must either cite that theorem or assume the clause
directly.

\begin{definition}[Aleph hierarchy, cardinal exponentiation, CH, and GCH]
By transfinite recursion, \(\aleph_0\defeq\omega\);
\(\aleph_{\alpha+1}\) is the least cardinal greater than
\(\aleph_\alpha\), which exists by Hartogs' theorem; and
\(\aleph_\lambda\defeq\sup\set{\aleph_\beta:\beta<\lambda}\) at a
limit ordinal \(\lambda\).  Under the Axiom of Choice, define
\(2^{\kappa}\defeq\abs{\Pow(\kappa)}\) for every cardinal \(\kappa\).
The \emph{continuum hypothesis} is the sentence
\[
  \mathsf{CH}\defeq(2^{\aleph_0}=\aleph_1),
\]
and the \emph{generalized continuum hypothesis} is
\[
  \mathsf{GCH}\defeq\forall\alpha\,(2^{\aleph_\alpha}=\aleph_{\alpha+1}),
\]
each read as the evident \(\Lang_{\in}\)-formalization of the displayed
cardinal equality.
\end{definition}

\begin{warning}[The Axiom of Choice hidden in cardinal notation]
The assertion that any two cardinalities are comparable, the assertion that
every set has an ordinal cardinal of the same size, and the unrestricted use of
cardinal arithmetic all use consequences of the Axiom of Choice.  A proof in \(\ZF\) must not import
those principles through notation.
\end{warning}

\begin{exercise}
\label{ex:cumulative-hierarchy-01}
Prove that \(\alpha\subseteq\beta\) whenever \(\alpha\leq\beta\) are
ordinals.  Identify where transitivity is used, and explain why no
transfinite induction is required.
\end{exercise}

\begin{exercise}
\label{ex:cumulative-hierarchy-02}
Compute \(\rank(\set{\varnothing,\set{\varnothing}})\) directly from the
recursive rank equation and verify the result from the hierarchy definition.
\end{exercise}

\begin{exercise}
\label{ex:cumulative-hierarchy-03}
Explain why the Mostowski collapse can fail to be injective when extensionality
of \(R\) is removed.  Supply a finite counterexample.
\end{exercise}

\chapter[Constructibility and Forcing]{Constructibility, Forcing, and Relative Consistency}
\label{ch:constructibility-forcing}

\section{The constructible hierarchy}
\label{sec:constructibility-forcing-01}

\begin{definition}[Definable subsets]
For a set \(A\), let \(\operatorname{Def}(A)\) be the collection of subsets of
\(A\) definable over the structure \((A,\in)\) by a first-order formula with
parameters from \(A\):
\[
  \operatorname{Def}(A)
  \defeq
  \set{
    X\subseteq A:
    \begin{gathered}
      \text{for some formula }\varphi(y,\bar z)
      \text{ and }\bar a\in A^{<\omega},\\
      X=\set{x\in A:(A,\in),s_{x,\bar a}\entails\varphi}
    \end{gathered}
  }.
\]
Here \(A^{<\omega}\) denotes the set of finite tuples from \(A\), and
\(s_{x,\bar a}\) maps \(y\) to \(x\) and each displayed variable in
\(\bar z\) to the corresponding entry of \(\bar a\); its values on other
variables are irrelevant.
The satisfaction relation in this definition is the set-sized satisfaction
relation constructed in the metatheory.  It is not a truth predicate internal
to \(\ZF\) for the entire universe.
\end{definition}

\begin{definition}[Constructible hierarchy]
Define by transfinite recursion
\[
  L_0=\varnothing,\qquad
  L_{\alpha+1}=\operatorname{Def}(L_\alpha),\qquad
  L_\lambda=\bigcup_{\beta<\lambda}L_\beta
\]
for limit \(\lambda\), and write
\[
  L=\bigcup_{\alpha\in\operatorname{Ord}}L_\alpha.
\]
A set is constructible iff it belongs to \(L\).
\end{definition}

The ordinary Tarskian satisfaction relation fixed earlier takes a set-sized
structure on its left and therefore cannot take the proper class \(L\) as an
argument.  To state the next result without adding class structures or a class
truth predicate, use the following formula-by-formula abbreviation.  For a
standard formula \(\varphi\) of \(\Lang_{\in}\), let
\(\varphi^L\) be its metalinguistic relativization to the displayed definable
class: atomic formulas are unchanged, the operation commutes with the
propositional connectives, and
\[
  (\exists x\,\psi)^L\defeq
    \exists x\,(x\in L\land\psi^L),
  \qquad
  (\forall x\,\psi)^L\defeq
    \forall x\,(x\in L\to\psi^L).
\]
Here \(x\in L\) is eliminated using the formula defining constructibility.
Accordingly, \(L\entails\sigma\) abbreviates the ambient assertion
\(\sigma^L\), and \(L\entails T\) means that this holds for every standard
sentence, including every displayed schema instance, represented by the
theory expression \(T\).  This is a formula-by-formula metalinguistic
convention.  It neither treats the proper class \(L\) as a set-sized structure
nor posits a single truth predicate for all formula codes.

\begin{theorem}[The constructible class L]
Assume \(\ZF\) in the ambient universe.  The class \(L\) is transitive and
contains every ordinal.
\end{theorem}

\paragraph{Proof sketch.}
The recursion defining \(L\) yields transitivity and containment of all
ordinals: transfinite induction shows that every level \(L_\alpha\) is
transitive and that \(\alpha\in L_{\alpha+1}\).  The omitted verification is
the level-by-level induction, supplied by \cite{Jech2003}.

\paragraph{Cited theorem (relativization scheme for the constructible
class).}
For every axiom \(\sigma\) of \(\ZFC+\mathsf{GCH}+(V=L)\), including every
instance of the separation and replacement schemas,
\[
  \ZF\derives\sigma^{L},
\]
where \((\cdot)^{L}\) is the relativization convention displayed above.
Each instance of the scheme is a theorem of \(\ZF\) about the definable
class \(L\); the proofs---reflection to set-sized levels, the canonical
definable well-order of \(L\), and condensation---are supplied by
\cite[Chapter~13, pp.~175--200]{Jech2003}.  No satisfaction relation and no
consistency implication is asserted here; the exact relative-consistency
statements are derived from this scheme in the final section of this
chapter.

\section{Forcing over a countable transitive set: a pedagogical device}
\label{sec:constructibility-forcing-02}

\begin{convention}
For the remainder of this section, fix an externally countable transitive set
\(M\) such that the membership structure \((M,\in)\) satisfies \(\ZFC\), and
fix
\[
  (\mathbb P,\leq,1_{\mathbb P})\in M,
\]
where \((\mathbb P,\leq)\) is a partially ordered set and
\(1_{\mathbb P}\) is its greatest member.  The order is oriented so that
\(q\leq p\) means that \(q\) is stronger than \(p\).  In this section
\(p,q,d\) range over members of \(\mathbb P\), while \(\mu,\nu\) range over
\(\mathbb P\)-names.  This is an unnamed local convention on an ordinary
poset, not the introduction of another kind of order.  The countability and
transitivity assumptions are explicit pedagogical hypotheses stronger than
the arithmetical assertion \(\mathsf{Con}(\ZFC)\).
\end{convention}

\begin{remark}[Why the fixed hypothesis exceeds the consistency assertion]
The formal consistency sentence \(\mathsf{Con}(\ZFC)\) of
\cref{def:formal-consistency-sentence} is a sentence of arithmetic, while
\(\ZFC\) is a theory in \(\Lang_{\in}\); comparing the convention's
hypothesis with that sentence requires a translation of arithmetic into the
membership language, and none has been fixed so far.  Fix the canonical
interpretation: numbers are the finite ordinals, zero is \(\varnothing\),
successor is \(S(x)=x\cup\set{x}\), and addition and multiplication are the
recursive operations on \(\omega\).  For an arithmetic sentence \(\sigma\),
write \(\sigma^{\in}\) for its \(\Lang_{\in}\)-translation under this
interpretation; the notation is a metalinguistic abbreviation whose scope is
this chapter.

Suppose the hypothesized \(M\) exists.  Then \((M,\in)\) is a set-sized
structure satisfying \(\ZFC\), so soundness (\cref{thm:fol-soundness})
gives metatheoretic consistency of \(\ZFC\).  Correctness of the fixed
proof predicate---the warning on external and internal consistency in the
second-incompleteness chapter---then makes \(\mathsf{Con}(\ZFC)\) true in
the standard structure \(\mathcal N\).  Absoluteness of arithmetic
statements over transitive models containing \(\omega\)
\cite[Chapter~12, pp.~155--174]{Jech2003} gives
\((M,\in)\entails\ZFC+\mathsf{Con}(\ZFC)^{\in}\), and soundness again gives
consistency of \(\ZFC+\mathsf{Con}(\ZFC)^{\in}\).  That consistency is not
provable in \(\ZFC+\mathsf{Con}(\ZFC)^{\in}\) itself.
\Cref{thm:godel-second} does not apply verbatim: it covers computably
enumerable theories \(T\supseteq\mathsf Q\) in the language of arithmetic,
and \(\ZFC+\mathsf{Con}(\ZFC)^{\in}\) is a membership-language theory.  It
applies transported across the canonical interpretation, under which the
fixed provability predicate still satisfies the three derivability
conditions.  So the fixed hypothesis strictly exceeds
\(\mathsf{Con}(\ZFC)\).  The omitted verifications are the translation
convention itself, the absoluteness computation, and the interpreted
derivability-condition bookkeeping.
\textit{Status: proof sketch.}
\end{remark}

\begin{definition}[Poset filter, dense set, and genericity]
\emph{Necessity.}
The qualifier \emph{poset} is required because the previously used
model-theoretic filter on \(\Pow(I)\) and the filters used here have different
kinds of members and closure clauses.

\emph{Definition.}
A poset filter on the fixed \((\mathbb P,\leq)\) is a nonempty set
\(G\subseteq\mathbb P\) satisfying
\[
\begin{aligned}
  &p\in G\text{ and }p\leq q \quad\Longrightarrow\quad q\in G,
  \\
  &p,q\in G \quad\Longrightarrow\quad
    \exists d\in G\,(d\leq p\land d\leq q).
\end{aligned}
\]
Thus the first clause closes \(G\) upward toward weaker members of
\(\mathbb P\), and the
second makes it downward directed toward common strengthenings.

\emph{Standing.}
The phrase \emph{poset filter} is a metalinguistic predicate of subsets \(G\)
of the fixed poset \(\mathbb P\); it is not an object-language
formula or a semantic structure.

\emph{Nonredundancy.}
A model-theoretic filter on \(\Pow(I)\) has subsets of \(I\) as its members,
is upward closed under inclusion, and is closed under finite intersections.
A poset filter is a subset of \(\mathbb P\) and uses upward closure in
\(\leq\) together with directed common strengthenings.  Neither
closure specification supplies the other without additional structure.

\emph{Use.}
The qualified term is used below to define \(M\)-genericity, to interpret
\(\mathbb P\)-names, and to prove monotonicity of the semantic forcing
relation.

A set \(D\subseteq\mathbb P\) is dense iff
\[
  \forall p\in\mathbb P\,\exists q\in D\,(q\leq p).
\]
A poset filter \(G\subseteq\mathbb P\) is \(M\)-generic iff
\[
  \forall D\in M\,
  \bigl(D\subseteq\mathbb P\text{ is dense}\to G\cap D\ne\varnothing\bigr).
\]
Genericity is evaluated in the ambient universe.  Because \(M\) is
transitive, a statement built from membership whose quantifiers are all
bounded by members of \(M\)---such as \(p,q\in\mathbb P\),
\(q\leq p\), or \(D\subseteq\mathbb P\)---holds in \((M,\in)\) iff it
holds in the ambient universe; this is verified by induction on such
bounded statements, and it is the only transfer between \(M\) and the
ambient universe that this section uses.
\end{definition}

\begin{lemma}[Existence of a generic filter under the fixed hypothesis]
For every \(p\in\mathbb P\), an \(M\)-generic poset filter containing \(p\)
exists in the ambient universe.
\end{lemma}

\begin{proof}
Externally enumerate the dense subsets of \(\mathbb P\) belonging to \(M\) as
\(D_0,D_1,\ldots\).  Choose \(p_0\leq p\), and recursively choose
\(p_{n+1}\leq p_n\) with \(p_{n+1}\in D_n\).  Let
\[
  G=\set{q\in\mathbb P:\exists n\,(p_n\leq q)}.
\]
Then \(p\in G\), \(G\) is upward closed, and any two members of \(G\) have a
common strengthening \(p_k\in G\) for a sufficiently large \(k\).  Moreover,
\(p_{n+1}\in G\cap D_n\) for every \(n\).  Hence \(G\) is an \(M\)-generic
poset filter containing \(p\).
\end{proof}

\begin{definition}[Names, interpretation, and extension]
Define the hierarchy of \(\mathbb P\)-names by ordinal recursion:
\[
  V^{\mathbb P}_\alpha
  \defeq
  \Pow\!\left(
    \left(\bigcup_{\beta<\alpha}V^{\mathbb P}_\beta\right)
    \times\mathbb P
  \right),
  \qquad
  V^{\mathbb P}\defeq
  \bigcup_{\alpha\in\operatorname{Ord}}V^{\mathbb P}_\alpha.
\]
A \(\mathbb P\)-name is a member of the class \(V^{\mathbb P}\).  Thus each
member of a name \(\nu\) is a pair \(\tuple{\mu,p}\), where \(\mu\) has lower
name rank and \(p\in\mathbb P\).  For a poset filter \(G\), rank recursion
defines
\[
  \nu^G
  \defeq
  \set{\mu^G:
        \exists p\in G\,(\tuple{\mu,p}\in\nu)}.
\]
For each \(x\in M\), the notation \(\check x\) is needed to exhibit a
\(\mathbb P\)-name in \(M\) whose interpretation is \(x\).  Define it by rank
recursion:
\[
  \check x\defeq
  \set{\tuple{\check y,1_{\mathbb P}}:y\in x}.
\]
Its standing is a set of name--condition pairs, hence a
\(\mathbb P\)-name, rather than the set \(x\) itself.  This distinction
is nonredundant because \(x\) need not have the pair structure required of a
name.  Because a nonempty poset filter contains \(1_{\mathbb P}\), induction on
the rank of \(x\) gives
\(\check x\in M\cap V^{\mathbb P}\) and \((\check x)^G=x\).  This calculation
is the exact link by which members of \(M\) enter \(M[G]\), and the
notation is used below in the name for the Cohen real.  It is not an appeal
to the cited forcing theorem below.
The associated extension is the set
\[
  M[G]
  \defeq
  \set{\nu^G:\nu\in M\cap V^{\mathbb P}}.
\]
Both recursions are well founded on ordinal rank; neither clause presupposes a
forcing relation.
\end{definition}

\begin{definition}[Semantic forcing under the fixed hypothesis]
\emph{Necessity.}
The qualifier \emph{semantic} is required to distinguish the exact external
relation defined here from the separate recursively defined forcing relation
whose construction is only cited below.

\emph{Definition.}
Let \(\varphi\) be a formula of the language of membership whose free variables
are among the pairwise distinct variables \(v_1,\ldots,v_n\), let
\(\nu_1,\ldots,\nu_n\in M\cap V^{\mathbb P}\), and let \(p\in\mathbb P\).
In this section the notation
\[
  p\forces^{\mathrm{sem}}_{M,\mathbb P}
    \varphi(\nu_1,\ldots,\nu_n)
\]
means exactly that, for every \(M\)-generic poset filter
\(G\subseteq\mathbb P\) with \(p\in G\) and every assignment \(s_G\) into
\(M[G]\) satisfying \(s_G(v_i)=\nu_i^G\) for \(1\leq i\leq n\),
\[
  (M[G],\in),s_G\entails\varphi.
\]
The preceding existence lemma prevents this universal condition from being
vacuous.

\emph{Standing.}
For the fixed \(M\) and \(\mathbb P\), semantic forcing is a metalinguistic
relation among a member \(p\) of \(\mathbb P\), a membership-language formula,
and a tuple of \(\mathbb P\)-names.

\emph{Nonredundancy.}
This relation is not ordinary satisfaction: it quantifies over every
\(M\)-generic poset filter containing \(p\).  It is also not the recursively
defined relation used inside \(M\) to prove definability and preservation;
no such recursive relation is defined in this manuscript.

\emph{Use.}
The notation is used immediately in the monotonicity lemma and later in the
source-pinned comparison with recursive forcing.
\end{definition}

\begin{lemma}[Monotonicity of semantic forcing]
Let the formula and names be as in the preceding definition.  If \(q\leq p\)
and
\[
  p\forces^{\mathrm{sem}}_{M,\mathbb P}
    \varphi(\nu_1,\ldots,\nu_n),
\]
then
\[
  q\forces^{\mathrm{sem}}_{M,\mathbb P}
    \varphi(\nu_1,\ldots,\nu_n).
\]
\end{lemma}

\begin{proof}
Let \(G\) be any \(M\)-generic poset filter with \(q\in G\), and let \(s_G\)
be any assignment having the values required in the semantic-forcing
definition.  Since \(q\leq p\), upward closure of the poset filter gives
\(p\in G\).  The hypothesis at \(p\) therefore yields
\((M[G],\in),s_G\entails\varphi\).  Since \(G\) and \(s_G\) were arbitrary,
the defining universal condition holds at \(q\).
\end{proof}

\paragraph{Cited theorem (forcing theorem).}
Fix the countable transitive \((M,\in)\), the poset \(\mathbb P\in M\), and
the name interpretation already displayed.  For every fixed standard formula
\(\varphi(v_1,\ldots,v_n)\) of \(\Lang_{\in}\), Jech's forcing development
\cite[pp.~201--224]{Jech2003} constructs, by simultaneous recursion on name
rank in the membership and equality cases and then by recursion on formula
complexity, the source's forcing assertion for triples
\((p,\varphi,\bar\nu)\).  For \(p\in\mathbb P\),
\(\bar\nu\in(M\cap V^{\mathbb P})^n\), and every \(M\)-generic poset filter
\(G\), it proves the four claims below.  In item~3,
\(\bar\nu^G\) abbreviates the interpreted tuple
\((\nu_1^G,\ldots,\nu_n^G)\), distinguishing it from the tuple of names and
supplying exactly the parameter values used there.
\begin{enumerate}
  \item for that fixed \(\varphi\), an \(\Lang_{\in}\)-formula, with the fixed
        poset as a parameter, represents the recursive assertion over
        \((M,\in)\);
  \item the recursive assertion at \(p\) holds iff
        \(p\forces^{\mathrm{sem}}_{M,\mathbb P}\varphi(\bar\nu)\);
  \item \((M[G],\in)\entails\varphi(\bar\nu^G)\) iff some
        \(p\in G\) satisfies the recursive assertion for
        \(\varphi(\bar\nu)\); and
  \item \((M[G],\in)\) satisfies \(\ZFC\) and has the same ordinals as
        \((M,\in)\).
\end{enumerate}
The first clause is formula-by-formula definability, not a single internal
truth predicate ranging over all standard formula codes.  The four conclusions
are imported as a cited theorem; this manuscript has defined only the external
semantic relation used in item~2.

\paragraph{Internal proof boundary.}
To turn that citation into a complete proof here would require printing the two
simultaneous atomic clauses, the recursive clauses for the active primitive
logical symbols
\(\bottom,\neg,\land,\lor,\to,\forall,\exists\), and the precise recursion
measure; proving existence, uniqueness, monotonicity, and formula-by-formula
definability; proving the density and truth lemmas in every constructor case;
and verifying every \(\ZFC\) axiom in the extension, with the Replacement
argument stated in full.  Those source-dependent steps are not reproduced, so
no local recursive relation or local completeness claim is introduced.

\begin{example}[One Cohen real]
Let
\[
  \mathbb P=\operatorname{Fn}(\omega,2,{<}\omega),
\]
the finite partial functions from \(\omega\) to \(2\), ordered by reverse
inclusion.  A name is necessary here because the generic function depends on
the external filter \(G\), while membership in \(M[G]\) must be witnessed by
an object already in \(M\).  Define
\[
  \dot c\defeq
  \set{
    \tuple{\check{\tuple{n,i}},p}:
    p\in\mathbb P,\ n\in\dom(p),\ p(n)=i
  }.
\]
The symbol \(\dot c\) stands for a \(\mathbb P\)-name in \(M\), not for the
function \(c\) in the extension; the dot records precisely this type
difference.  This is nonredundant with \(\check x\): check names interpret as
fixed members of \(M\), whereas \(\dot c\) interprets differently with the generic
filter.  The ground-set-name recursion and separation in \(M\) give
\(\dot c\in M\cap V^{\mathbb P}\).  Directly from name interpretation,
\[
\begin{aligned}
  (\dot c)^G
    &=\set{\tuple{n,i}:\exists p\in G\,(p(n)=i)}\\
    &=\bigcup G.
\end{aligned}
\]
Thus, if \(c\defeq\bigcup G\), then \(c=(\dot c)^G\in M[G]\); this immediate
use is the reason for introducing \(\dot c\).  Directedness
of \(G\) makes the union single-valued, and for each \(n\) the set
\[
  D_n=\set{p\in\mathbb P:n\in\dom(p)}
\]
is dense and belongs to \(M\), so genericity makes \(c\) a total function
\(\omega\to2\).  Write \(2^\omega\) for the set of total functions
\(\omega\to2\).  For \(b\in M\cap2^\omega\), the set
\[
  E_b=\set{p\in\mathbb P:
            \exists n\in\dom(p)\,(p(n)\ne b(n))}
\]
belongs to \(M\) and is dense: extend any condition at one fresh coordinate
with the value opposite to \(b\).  Since \(G\cap E_b\ne\varnothing\),
\(c\ne b\).  Hence \(c\notin M\), while the displayed name calculation proves
that the extension contains this new real.  It does not by itself prove
\(\neg\mathsf{CH}\);
that conclusion requires a poset adding enough reals together with
cardinal-preservation arguments.
\end{example}

\section{Exact relative-consistency status}
\label{sec:constructibility-forcing-03}

\begin{lemma}[Relativization of derivations]
\label{lem:relativization-derivations}
Let \(\varphi_1,\ldots,\varphi_k,\psi\) be formulas of \(\Lang_{\in}\) whose
free variables are among the pairwise distinct variables
\(v_1,\ldots,v_m\).  If \(\set{\varphi_1,\ldots,\varphi_k}\derives\psi\) in
the fixed calculus \(\ND_{\FOL}\), then
\[
  \ZF\derives
  \forall v_1\cdots\forall v_m\,
  \Bigl(
    \bigwedge_{j\leq m}v_j\in L
    \land
    \bigwedge_{i\leq k}\varphi_i^{L}
    \to\psi^{L}
  \Bigr).
\]
For \(m=0\) or \(k=0\) the corresponding conjunction is absent from the
displayed conclusion.
\end{lemma}

\begin{proof}
Strengthen the claim: for every derivation \(D\) with open assumptions
\(\Gamma\) and conclusion \(\chi\), and for every finite list
\(u_1,\ldots,u_r\) of pairwise distinct variables containing every variable
free in \(\Gamma\cup\set{\chi}\), the theory \(\ZF\) derives the displayed
guarded implication formed from \(\Gamma\), \(\chi\), and
\(u_1,\ldots,u_r\).  The lemma is the case in which \(D\) derives \(\psi\)
from \(\varphi_1,\ldots,\varphi_k\) and the list is \(v_1,\ldots,v_m\).
Induct on \(D\); throughout, argue inside \(\ZF\) under the quantifiers
\(\forall u_1\cdots\forall u_r\) and the guard hypotheses \(u_j\in L\).

Two bookkeeping devices recur.  First, \(\ZF\derives\exists u\,(u\in L)\):
by the comprehension reading of \(\operatorname{Def}\), separation proves
\(\set{x\in L_0:\ldots}=\varnothing\) whatever the satisfaction clause
returns over the empty domain, which the fixed structure convention
excludes from semantic evaluation; hence
\(\varnothing\in\operatorname{Def}(L_0)=L_1\), so \(\varnothing\in L\).
This derivation does not route through the proof-sketch theorem that \(L\)
contains every ordinal.  Whenever the induction hypothesis supplies a claim
whose variable list contains a variable \(u\) absent from the target node,
instantiate the universally quantified guarded implication at this witness;
the vanished variable and its guard disappear.  Second, where the
capture-avoiding substitution of the calculus renames bound variables, the
relativizations of the formulas concerned are alpha-variants, because
\((\cdot)^{L}\) inserts the same parameter-free guard under alpha-variant
binders; the explicit \(\alpha\) rule of \(\ND_{\FOL}\) bridges them.

\emph{Assumptions, propositional rules, and classical rules.}  An
assumption leaf \(\chi\) maps to the relativized hypothesis \(\chi^{L}\) of
the guarded implication.  The operation \((\cdot)^{L}\) commutes with the
propositional connectives and \(\bottom^{L}=\bottom\), so every
propositional and classical rule instance maps to the same rule instance on
the relativized premises, with discharged assumptions mapped to discharged
relativized assumptions.

\emph{Quantifier rules.}  Every \(\Lang_{\in}\)-term is a variable, since
the signature has no function or constant symbols.  A \(\forall E\) step
from \(\forall x\,\chi\) to \(\chi[u/x]\) maps to a \(\forall E\) step on
\((\forall x\,\chi)^{L}=\forall x\,(x\in L\to\chi^{L})\), followed by
implication elimination against the guard \(u\in L\) supplied by the
antecedent.  An \(\exists I\) step from \(\chi[u/x]\) to
\(\exists x\,\chi\) maps to conjunction introduction with the guard
\(u\in L\) followed by \(\exists I\), concluding
\(\exists x\,(x\in L\land\chi^{L})\).  In both steps, if \(x\) is not free
in \(\chi\), the instantiating variable is absent from the formulas; use
the witness supplied by \(\exists u\,(u\in L)\) in place of a guard from
the antecedent.  A \(\forall I\) step with generalized variable \(x\) maps
to implication introduction discharging the hypothesis \(x\in L\),
followed by \(\forall I\); if \(x\) is not free in the premise, first
conjoin the hypothesis \(x\in L\) to the relativized premise and project
the premise back out, so that the discharged assumption occurs in its
derivation.  An \(\exists E\) step with eigenvariable \(y\) maps to
\(\exists E\) on the relativized existential premise; the minor derivation
assumes \(y\in L\land\chi^{L}[y/x]\), and conjunction eliminations recover
the guard \(y\in L\) and the relativized assumption.  The eigenvariable
conditions are unchanged, because the inserted guard \(x\in L\) is
parameter-free and adds no free variables.

\emph{Equality rules.}  Atomic formulas are unchanged by \((\cdot)^{L}\),
so an \(=I\) step concludes the unchanged atom.  Substitution of variables
for variables commutes with \((\cdot)^{L}\) up to the alpha-variant
bookkeeping above, so an \(=E\) step on \(\varphi\) maps to the \(=E\) step
on \(\varphi^{L}\).

\emph{Alpha rule.}  An explicit \(\alpha\) step maps to an \(\alpha\) step
by the second bookkeeping device.

At the root, the open relativized assumptions are
\(\varphi_1^{L},\ldots,\varphi_k^{L}\) and the guarded variables are
\(v_1,\ldots,v_m\); implication introduction and universal introduction in
the ambient \(\ZF\) derivation yield the displayed conclusion.
\end{proof}

\begin{theorem}[Relative consistency of CH]
\label{thm:ch-relative-consistency}
If \(\ZF\) is consistent, then \(\ZFC+\mathsf{GCH}+(V=L)\) is consistent.
Hence if \(\ZFC\) is consistent, then \(\ZFC+\mathsf{CH}\) is consistent.
\end{theorem}

\begin{proof}
Suppose a derivation concludes \(\bottom\) from axioms
\(\sigma_1,\ldots,\sigma_k\) of \(\ZFC+\mathsf{GCH}+(V=L)\).  All the
\(\sigma_i\) and \(\bottom\) are sentences, so
\cref{lem:relativization-derivations} with \(m=0\) and \(\psi=\bottom\)
gives
\(\ZF\derives\sigma_1^{L}\land\cdots\land\sigma_k^{L}\to\bottom\), and the
relativization scheme cited in \cref{sec:constructibility-forcing-01}
gives \(\ZF\derives\sigma_i^{L}\) for each \(i\).  Conjunction
introduction and implication elimination yield \(\ZF\derives\bottom\),
contradicting the consistency of \(\ZF\).

For the second claim, every \(\ZF\)-derivation is a \(\ZFC\)-derivation,
so consistency of \(\ZFC\) gives consistency of \(\ZF\) and, by the first
claim, of \(\ZFC+\mathsf{GCH}+(V=L)\).  That theory derives
\(\mathsf{CH}\): universal elimination at the ordinal zero---formally,
existential elimination from the \(\ZFC\)-theorem that the ordinal zero
exists, followed by universal elimination and equality
substitutivity---produces the \(\alpha=0\) instance of the formalized
\(\mathsf{GCH}\) display, which is the formalized \(\mathsf{CH}\).  Now
suppose some derivation concludes \(\bottom\) with open assumptions in
\(\ZFC\cup\set{\mathsf{CH}}\).  At least one open occurrence of
\(\mathsf{CH}\) remains in it, since otherwise its open assumptions lie in
\(\ZFC\) and \(\ZFC\derives\bottom\), contrary to the assumed consistency
of \(\ZFC\).  Applying \(\neg I\) at the root, discharging every such
occurrence, gives \(\ZFC\derives\neg\mathsf{CH}\); monotonicity of the
derivability relation gives
\(\ZFC+\mathsf{GCH}+(V=L)\derives\neg\mathsf{CH}\); and negation
elimination against the derivation of \(\mathsf{CH}\) above gives
\(\ZFC+\mathsf{GCH}+(V=L)\derives\bottom\), contradicting the consistency
established in the first claim.  Hence \(\ZFC+\mathsf{CH}\) is
consistent.
\end{proof}

\paragraph{Cited theorem (relative consistency of the negation of CH).}
If \(\ZFC\) is consistent, then \(\ZFC+\neg\mathsf{CH}\) is consistent
\cite[Chapter~14, pp.~201--224]{Jech2003}.  The import is a family of
\(\ZFC\)-theorem schemes about the definable class \(V^{B}\) of
\(B\)-valued names, where \(B\) is the Boolean completion of the poset
\(\operatorname{Fn}(\aleph_2\times\omega,2,{<}\omega)\) of finite partial
functions from \(\aleph_2\times\omega\) to \(2\), and
\(\lVert\sigma\rVert_{B}\) denotes the defined \(\ZFC\)-term assigning to
a membership-language sentence \(\sigma\) with name parameters its Boolean
truth value in \(B\); both notations are scoped to this import.  The
cited schemes are:
\begin{enumerate}
  \item \(\lVert\sigma\rVert_{B}=1\) for each axiom \(\sigma\) of
        \(\ZFC\), including every schema instance;
  \item Boolean soundness: if
        \(\set{\sigma_1,\ldots,\sigma_k}\derives\psi\), then
        \(\ZFC\derives
          \lVert\sigma_1\rVert_{B}\wedge\cdots\wedge
          \lVert\sigma_k\rVert_{B}\leq\lVert\psi\rVert_{B}\),
        where \(\lVert\bottom\rVert_{B}=0\) and \(\ZFC\) proves \(B\)
        nontrivial, that is, \(0\ne1\) in \(B\);
  \item \(B\) satisfies the countable chain condition, and
        chain-condition forcing preserves cardinals:
        \(\lVert\check\aleph_1=\aleph_1\land
           \check\aleph_2=\aleph_2\rVert_{B}=1\); and
  \item the distinctness density argument of the one-Cohen-real example,
        computed per pair of \(\aleph_2\)-indexed sections of the adjoined
        function, gives
        \(\lVert 2^{\aleph_0}\geq\check\aleph_2\rVert_{B}=1\), hence
        \(\lVert\neg\mathsf{CH}\rVert_{B}=1\).
\end{enumerate}
A derivation of \(\bottom\) from \(\ZFC\cup\set{\neg\mathsf{CH}}\) would
therefore give \(\ZFC\derives 1=0\) in \(B\), contradicting
nontriviality.  This Boolean-valued route is the discharge demanded by
the warning on countable transitive structures below; it hypothesizes no
countable transitive model.  The source states these schemes for its own
calculus; consistency transfers to the fixed calculus \(\ND_{\FOL}\)
because soundness (\cref{thm:fol-soundness}) and strong completeness
(\cref{thm:fol-completeness}) make derivability calculus-invariant.

\begin{remark}[Arithmetized form]
The two derivation transformations above are primitive recursive, so the
implications
\(\mathsf{Con}(\ZFC)\to\mathsf{Con}(\ZFC+\mathsf{CH})\) and
\(\mathsf{Con}(\ZFC)\to\mathsf{Con}(\ZFC+\neg\mathsf{CH})\), read as
arithmetized statements under the fixed proof predicate, are theorems of
\(\PA\).  The omitted verification is the arithmetization, including a
primitive recursive translation between the fixed calculus \(\ND_{\FOL}\)
and the source calculus of the cited theorem above: the semantic
calculus-invariance argument there does not supply that translation; the
deduction-theorem translation does.
\textit{Status: proof sketch.}
\end{remark}

\begin{theorem}[Nonderivability]
\label{thm:ch-nonderivability}
If \(\ZFC\) is consistent, then \(\ZFC\notderives\mathsf{CH}\) and
\(\ZFC\notderives\neg\mathsf{CH}\).
\end{theorem}

\begin{proof}
Suppose \(\ZFC\derives\mathsf{CH}\).  A witnessing derivation has its
open assumptions in \(\ZFC\subseteq\ZFC\cup\set{\neg\mathsf{CH}}\), so
monotonicity of the derivability relation gives
\(\ZFC+\neg\mathsf{CH}\derives\mathsf{CH}\); the assumption rule gives
\(\ZFC+\neg\mathsf{CH}\derives\neg\mathsf{CH}\); and negation elimination
gives \(\ZFC+\neg\mathsf{CH}\derives\bottom\).  The contrapositive of the
cited relative-consistency theorem for the negation of \(\mathsf{CH}\)
then refutes the consistency of \(\ZFC\), contrary to hypothesis.
Symmetrically, if \(\ZFC\derives\neg\mathsf{CH}\), then monotonicity, the
assumption rule, and negation elimination give
\(\ZFC+\mathsf{CH}\derives\bottom\), contradicting the consistency of
\(\ZFC+\mathsf{CH}\) supplied by \cref{thm:ch-relative-consistency}.
\end{proof}

\begin{warning}[What the consistency theorem does not say]
Neither construction proves \(\mathsf{Con}(\ZFC)\) without a consistency
hypothesis.  The exact second incompleteness result in
\cref{thm:godel-second} applies to a consistent computably enumerable theory
\(T\supseteq\mathsf Q\) and the fixed provability predicate satisfying all
three derivability conditions; under those hypotheses, \(T\) does not derive
its displayed arithmetized consistency sentence.
\end{warning}

\begin{warning}[Countable transitive structures]
The elementary presentation ``take a countable transitive structure
\((M,\in)\) satisfying \(\ZFC\) and construct an \(M\)-generic filter''
assumes more than the bare
arithmetical statement \(\mathsf{Con}(\ZFC)\).  Countability permits an
external enumeration of the dense sets in \(M\), and transitivity makes the
name interpretation transparent.  A proof of the displayed
relative-consistency implication must discharge that pedagogical assumption
by a
syntactic, Boolean-valued, or suitable model-theoretic version of forcing.
The cited relative-consistency theorem for the negation of \(\mathsf{CH}\)
above supplies exactly the Boolean-valued discharge.
\end{warning}

\begin{exercise}
\label{ex:constructibility-forcing-01}
For Cohen forcing, prove that
\(D_n=\set{p\in\mathbb P:n\in\dom(p)}\) is dense for every \(n\in\omega\).
Then prove that \(\bigcup G\) is a function.
\end{exercise}

\begin{exercise}
\label{ex:constructibility-forcing-02}
Distinguish the following three assertions: \(L\entails\mathsf{CH}\),
\(\ZFC\derives V=L\to\mathsf{CH}\), and
\(\mathsf{Con}(\ZFC)\to\mathsf{Con}(\ZFC+\mathsf{CH})\).
State the language level and required theorem for each implication between
them.
\end{exercise}

\part{Proof Theory}
\label{part:proof-theory}

\chapter[Classical and Intuitionistic Calculi]{Classical and Intuitionistic Sequent Calculi}
\label{ch:lk-lj}

\begin{convention}[Proof-theory formula metavariables]
Throughout this part, \(A,B,C\) range over formulas unless a section explicitly
declares a narrower local role.  The symbols \(\Gamma,\Delta,\Sigma,\Pi\)
retain their normative standing as formula contexts.  The simply typed lambda
calculus section explicitly redeclares \(A,B\) as type metavariables and
\(M,N\) as term metavariables only within that section.
\end{convention}

The sequent calculi \(\LK\) and \(\LJ\), and the cut-elimination
theorem of \cref{ch:cut-elimination}, originate with Gentzen
\cite{Gentzen1934a,Gentzen1935b}.

Structural proof theory studies derivations as formal objects.  The
active calculus, its context discipline, and its conversion relation must be
fixed before a claim about cut elimination or normalization is meaningful.
The presentation below follows the standard separation between sequent
calculus and natural deduction in
\cite{TroelstraSchwichtenberg2000,NegriVonPlato2001}.

\section{Sequents and contexts}
\label{sec:lk-lj-01}

\begin{convention}[Contexts]
In this chapter a context is a finite multiset of formulas.  Commas denote
multiset union.  Exchange is therefore built into the representation of a
context; weakening and contraction are not.  If \(\Gamma\) is a context, then
\(\Gamma,A\) adds one occurrence of \(A\).
\end{convention}

\begin{definition}[Classical sequent]
A classical sequent is an expression
\[
  \Gamma\Rightarrow\Delta,
\]
where \(\Gamma\) and \(\Delta\) are finite multisets of formulas.  Its semantic
reading is
\[
  \bigwedge\Gamma\to\bigvee\Delta,
\]
with \(\bigwedge\varnothing=\topv\) and
\(\bigvee\varnothing=\bottom\).  This reading states the intended semantics; a
sequent is itself a syntactic object, not an object-language conditional.
\end{definition}

\begin{definition}[Intuitionistic sequent]
An intuitionistic sequent has the form
\[
  \Gamma\Rightarrow A,
\]
with exactly one succedent formula.  The restriction to one succedent is part
of the definition of the calculus \(\LJ\), not a typographic abbreviation for
a classical sequent.
\end{definition}

\section[The propositional rules of LK]{The propositional rules of \(\LK\)}
\label{sec:lk-lj-02}

The initial and structural rules are
\[
\frac{}{A\Rightarrow A}\;\textsc{Id}
\qquad
\frac{}{\bottom,\Gamma\Rightarrow\Delta}\;\textsc{\(\bottom L\)},
\]
\[
\frac{\Gamma\Rightarrow\Delta}
     {A,\Gamma\Rightarrow\Delta}\;\textsc{WL}
\qquad
\frac{\Gamma\Rightarrow\Delta}
     {\Gamma\Rightarrow\Delta,A}\;\textsc{WR},
\]
\[
\frac{A,A,\Gamma\Rightarrow\Delta}
     {A,\Gamma\Rightarrow\Delta}\;\textsc{CL}
\qquad
\frac{\Gamma\Rightarrow\Delta,A,A}
     {\Gamma\Rightarrow\Delta,A}\;\textsc{CR},
\]
and
\[
\frac{\Gamma\Rightarrow\Delta,A
      \qquad
      A,\Pi\Rightarrow\Sigma}
     {\Gamma,\Pi\Rightarrow\Delta,\Sigma}\;\textsc{Cut}.
\]
The displayed occurrence of \(A\) in cut is the \emph{cut formula}.

The logical rules for conjunction and disjunction are
\[
\frac{A,B,\Gamma\Rightarrow\Delta}
     {A\land B,\Gamma\Rightarrow\Delta}\;\textsc{\(\land L\)}
\qquad
\frac{\Gamma\Rightarrow\Delta,A
      \qquad
      \Gamma\Rightarrow\Delta,B}
     {\Gamma\Rightarrow\Delta,A\land B}\;\textsc{\(\land R\)},
\]
\[
\frac{A,\Gamma\Rightarrow\Delta
      \qquad
      B,\Gamma\Rightarrow\Delta}
     {A\lor B,\Gamma\Rightarrow\Delta}\;\textsc{\(\lor L\)}
\qquad
\frac{\Gamma\Rightarrow\Delta,A,B}
     {\Gamma\Rightarrow\Delta,A\lor B}\;\textsc{\(\lor R\)}.
\]
The implication rules are
\[
\frac{\Gamma\Rightarrow\Delta,A
      \qquad
      B,\Pi\Rightarrow\Sigma}
     {A\to B,\Gamma,\Pi\Rightarrow\Delta,\Sigma}
     \;\textsc{\(\to L\)}
\qquad
\frac{A,\Gamma\Rightarrow\Delta,B}
     {\Gamma\Rightarrow\Delta,A\to B}
     \;\textsc{\(\to R\)}.
\]
Negation is the abbreviation \(\neg A\defeq A\to\bottom\).  If negation is
instead made primitive, its rules must be added and the cut-elimination proof
must include their cases.

\section{Quantifier rules}
\label{sec:lk-lj-03}

For first-order formulas the quantifier rules are
\[
\frac{A[t/x],\Gamma\Rightarrow\Delta}
     {\forall x\,A,\Gamma\Rightarrow\Delta}\;\textsc{\(\forall L\)}
\qquad
\frac{\Gamma\Rightarrow\Delta,A[y/x]}
     {\Gamma\Rightarrow\Delta,\forall x\,A}\;\textsc{\(\forall R\)},
\]
\[
\frac{A[y/x],\Gamma\Rightarrow\Delta}
     {\exists x\,A,\Gamma\Rightarrow\Delta}\;\textsc{\(\exists L\)}
\qquad
\frac{\Gamma\Rightarrow\Delta,A[t/x]}
     {\Gamma\Rightarrow\Delta,\exists x\,A}\;\textsc{\(\exists R\)}.
\]
Every displayed instance \(A[t/x]\) uses the total capture-avoiding
substitution fixed in the first-order syntax chapter; bound variables are
renamed when required, so there is no additional definedness side condition.
In \(\forall R\) and \(\exists L\), the eigenvariable \(y\) must not occur
free in the conclusion sequent.  This freshness requirement is a rule
condition, not a convention that may be repaired after the inference.

\begin{definition}[\(\LK\)-derivation]
An \(\LK\)-derivation is a finite tree of classical sequents in which
every leaf is the conclusion of an instance of a displayed zero-premise
rule, and every nonleaf is the conclusion of exactly one instance of one of
the remaining displayed \(\LK\) rules, whose premises are its immediate
predecessors.  A sequent
is \(\LK\)-derivable, written
\[
  \LK\derives\Gamma\Rightarrow\Delta,
\]
when it is the root of an \(\LK\)-derivation.  A derivation is cut-free when it
contains no cut instance.
\end{definition}

\section[The calculus LJ]{The calculus \(\LJ\)}
\label{sec:lk-lj-04}

The calculus \(\LJ\) uses the same finite-multiset antecedent contexts and the
single-succedent sequent form \(\Gamma\Rightarrow A\).  Its complete primitive
rule surface is as follows.  The zero-premise and left-structural rules are
\[
\frac{}{A\Rightarrow A}\;\textsc{Id}
\qquad
\frac{}{\bottom,\Gamma\Rightarrow C}\;\textsc{\(\bottom L\)},
\]
\[
\frac{\Gamma\Rightarrow C}
     {A,\Gamma\Rightarrow C}\;\textsc{WL}
\qquad
\frac{A,A,\Gamma\Rightarrow C}
     {A,\Gamma\Rightarrow C}\;\textsc{CL}.
\]
Exchange is built into finite-multiset contexts.  The conjunction and
disjunction rules are
\[
\frac{A,B,\Gamma\Rightarrow C}
     {A\land B,\Gamma\Rightarrow C}\;\textsc{\(\land L\)}
\qquad
\frac{\Gamma\Rightarrow A\qquad\Gamma\Rightarrow B}
     {\Gamma\Rightarrow A\land B}\;\textsc{\(\land R\)},
\]
\[
\frac{A,\Gamma\Rightarrow C\qquad B,\Gamma\Rightarrow C}
     {A\lor B,\Gamma\Rightarrow C}\;\textsc{\(\lor L\)}.
\]
The two disjunction-right rules are
\[
\frac{\Gamma\Rightarrow A}
     {\Gamma\Rightarrow A\lor B}\;\textsc{\(\lor R_1\)}
\qquad
\frac{\Gamma\Rightarrow B}
     {\Gamma\Rightarrow A\lor B}\;\textsc{\(\lor R_2\)}.
\]
The implication rules are
\[
\frac{\Gamma\Rightarrow A
      \qquad
      B,\Pi\Rightarrow C}
     {A\to B,\Gamma,\Pi\Rightarrow C}\;\textsc{\(\to L\)}
\qquad
\frac{A,\Gamma\Rightarrow B}
     {\Gamma\Rightarrow A\to B}\;\textsc{\(\to R\)}.
\]
The quantifier rules are
\[
\begin{gathered}
\frac{A[t/x],\Gamma\Rightarrow C}
     {\forall x\,A,\Gamma\Rightarrow C}\;\textsc{\(\forall L\)}
\\[1ex]
\frac{\Gamma\Rightarrow A[y/x]}
     {\Gamma\Rightarrow\forall x\,A}\;\textsc{\(\forall R\)}
\\[1ex]
\frac{A[y/x],\Gamma\Rightarrow C}
     {\exists x\,A,\Gamma\Rightarrow C}\;\textsc{\(\exists L\)}
\\[1ex]
\frac{\Gamma\Rightarrow A[t/x]}
     {\Gamma\Rightarrow\exists x\,A}\;\textsc{\(\exists R\)}.
\end{gathered}
\]
As above, substitution is total and capture avoiding, and \(y\) must not occur
free in the conclusion of \(\forall R\) or \(\exists L\).  There is no right
weakening or right contraction because there is exactly one succedent formula.
Cut is
\[
\frac{\Gamma\Rightarrow A
      \qquad
      A,\Pi\Rightarrow C}
     {\Gamma,\Pi\Rightarrow C}\;\textsc{Cut}.
\]
Negation remains the defined connective
\(\neg A\defeq A\to\bottom\); there are no additional primitive negation
rules.

\begin{definition}[\(\LJ\)-derivation]
An \(\LJ\)-derivation is a finite rooted tree of intuitionistic sequents.  A
leaf is the conclusion of an instance of \textsc{Id} or
\textsc{\(\bottom L\)}, and every nonleaf is the conclusion of exactly one
instance of one of the remaining displayed \(\LJ\) rules, whose premises are
its immediate predecessors.  A sequent is \(\LJ\)-derivable, written
\[
  \LJ\derives\Gamma\Rightarrow A,
\]
iff it is the root of such a tree.  The derivation is cut-free iff
none of its nodes is obtained by \textsc{Cut}.
\end{definition}

\begin{theorem}[Soundness of \(\LK\)]
Let \(\Gamma\) and \(\Delta\) be finite contexts.
If \(\LK\derives\Gamma\Rightarrow\Delta\), then every classical structure
and assignment satisfying every formula in \(\Gamma\) satisfies at least one
formula in \(\Delta\).
\end{theorem}

\paragraph{Proof sketch.}
Induct on the derivation.  Initial and structural cases follow directly from
the semantic reading of a sequent.  Each logical case unfolds the semantic
clause for its principal connective.  The quantifier cases use the substitution
lemma; the eigenvariable cases use the fact that the eigenvariable is absent
from the conclusion.  The individual case calculations are omitted.  Semantic
soundness of \(\LJ\) is stated only after its
intuitionistic semantics is defined in \cref{ch:intuitionistic-logic}.

\begin{example}[A cut-free derivation]
The commutativity sequent \(A\land B\Rightarrow B\land A\) has the cut-free
derivation
\[
\frac{
  \dfrac{B\Rightarrow B}{A,B\Rightarrow B}\;\textsc{WL}
  \qquad
  \dfrac{A\Rightarrow A}{A,B\Rightarrow A}\;\textsc{WL}
}{
  \dfrac{A,B\Rightarrow B\land A}{A\land B\Rightarrow B\land A}
  \;\textsc{\(\land L\)}
}\;\textsc{\(\land R\)}.
\]
The formulas above the last inference are subformulas of the endsequent.  That
observation becomes a theorem only after cut elimination.
\end{example}

\section{Primitive, derived, and admissible rules}
\label{sec:lk-lj-05}

\begin{definition}[Primitive rule]
A primitive rule is a rule occurring in the definition of the active calculus.
Its status is syntactic: it is not inferred from soundness or semantic
validity.
\end{definition}

\begin{definition}[Derived rule]
A rule is derived in a calculus when there is a uniform derivation of its
conclusion from its premises using the primitive rules.  Replacing an
application of a derived rule by that derivation preserves the open leaves.
\end{definition}

\begin{definition}[Admissible rule]
A rule is admissible in a calculus when every substitution instance whose
premises are derivable has a derivable conclusion.  Adding an admissible rule
does not enlarge the set of derivable endsequents.
\end{definition}

Every derived rule is admissible.  The converse need not hold.  Admissibility
is a closure property of the derivability relation; a derived rule supplies
a uniform internal derivation of its conclusion from its premises as open
leaves.

\begin{example}[Cut after cut elimination]
Cut is primitive in the displayed formulations of \(\LK\) and \(\LJ\).  The
cut-elimination theorem shows that cut is admissible in the cut-free fragments.
This does not make cut a derived rule of the cut-free fragment: a single
uniform cut-free derivation with the two arbitrary cut premises as open leaves
would be a stronger assertion than admissibility.
\end{example}

\begin{warning}[A sequent is not a consequence assertion]
The sequent \(\Gamma\Rightarrow\Delta\), the derivability assertion
\(\LK\derives\Gamma\Rightarrow\Delta\), and its semantic validity are three
different objects.  Soundness and completeness connect them; typography does
not identify them.
\end{warning}

\begin{definition}[Principal formula]
In an instance of a displayed logical rule, the \emph{principal formula} is
the compound formula occurrence displayed in the conclusion and introduced
by that instance; the displayed occurrences in the premises from which it
is formed are its \emph{active formulas}, and all other formula occurrences
are context.  In an instance of weakening or contraction, the principal
formula is the displayed weakened or contracted occurrence.  A zero-premise
rule instance has no principal formula.
\end{definition}

\begin{exercise}
\label{ex:lk-lj-01}
Give a cut-free \(\LJ\)-derivation of
\(A\land(B\lor C)\Rightarrow(A\land B)\lor(A\land C)\).  Mark the principal
formula of every logical inference.
\end{exercise}

\begin{exercise}
\label{ex:lk-lj-02}
Show that deleting the eigenvariable condition from \(\forall R\) makes the
calculus unsound by deriving a sequent and giving a first-order structure and
assignment that satisfy every antecedent formula and falsify every succedent
formula.
\end{exercise}

\chapter{Cut Elimination and Analytic Derivations}
\label{ch:cut-elimination}

\section{Measures on cuts}
\label{sec:cut-elimination-01}

\begin{definition}[Formula rank and derivation height]
The rank \(r(A)\) of a formula is the number of logical symbols in it.  The
height \(h(D)\) of a derivation is the greatest number of inference steps on a
branch from a leaf to its root.  These measures are different: formula rank
measures syntax, whereas derivation height measures a proof tree.
\end{definition}

The reduction below is run for the following generalization of cut, the
\emph{mix} (multicut) rule.  Here \(A^{k}\) abbreviates \(k\)
comma-separated occurrences of \(A\).  For \(\LK\), with \(m,n\geq1\),
and with \(\Delta\) and \(\Sigma\) permitted to contain further
occurrences of \(A\):
\[
\frac{\Gamma\Rightarrow\Delta,A^{m}
      \qquad
      A^{n},\Pi\Rightarrow\Sigma}
     {\Gamma,\Pi\Rightarrow\Delta,\Sigma}\;\textsc{Mix}.
\]
For \(\LJ\) the succedent is single, so the left premise's succedent
multiplicity \(m\) is fixed at one: from \(\Gamma\Rightarrow A\) and
\(A^{n},\Pi\Rightarrow C\) with \(n\geq1\), \textsc{Mix} concludes
\(\Gamma,\Pi\Rightarrow C\).  The displayed occurrences of \(A\) are the
\emph{mix formula} of the instance.

\begin{lemma}[Equipollence of cut and mix]
\label{lem:cut-mix-equipollence}
The calculus obtained from \(\LK\) by replacing \textup{\textsc{Cut}} by
\textup{\textsc{Mix}} derives exactly the \(\LK\)-derivable sequents, and likewise
for \(\LJ\).  Consequently, a derivation in the mix formulation with no
mix instance is a cut-free derivation, and eliminating mix eliminates
cut.
\textit{Status: complete proof.}
\end{lemma}
\begin{proof}
Cut is the mix instance \(m=n=1\), so reading each cut as a mix converts
a derivation with cuts into a derivation in the mix formulation with the
same endsequent.  Conversely, an instance of \textsc{Mix} with
multiplicities \(m\) and \(n\) is derived from the displayed primitive
rules: \(m-1\) applications of \textsc{CR} to the left premise contract
the succedent occurrences \(A^{m}\) to a single \(A\), \(n-1\)
applications of \textsc{CL} to the right premise contract \(A^{n}\) to a
single \(A\), and one \textsc{Cut} concludes
\(\Gamma,\Pi\Rightarrow\Delta,\Sigma\).  Replacing every mix instance,
uppermost first, by this derivation yields a derivation in the original
formulation with the same endsequent.  In \(\LJ\) the left premise
already has the single succedent \(A\), so only \textsc{CL} is used.
Hence the cut and mix formulations derive the same sequents, and a
mix-elimination procedure eliminates cut: first read every cut as the mix
instance \(m=n=1\), then eliminate the mixes.
\end{proof}

For an instance of mix with mix formula \(A\) and immediate
subderivations \(D_1\) and \(D_2\), its rank is \(r(A)\) and its height
is \(h(D_1)+h(D_2)\), the sum of the heights of those subderivations.
The reduction is controlled by the lexicographic order on
\[
  \tuple{r(A),h(D_1)+h(D_2)}.
\]
For a cut, the mix instance \(m=n=1\), these measures are its \emph{cut
rank} and \emph{cut height}.

\begin{theorem}[Cut elimination for \(\LK\) and \(\LJ\)]
Every \(\LK\)-derivation can be transformed into a cut-free
\(\LK\)-derivation of the same endsequent.  Every \(\LJ\)-derivation can be
transformed into a cut-free \(\LJ\)-derivation of the same endsequent.
\textit{Status: proof sketch.}
\end{theorem}

\paragraph{Proof sketch.}
By \cref{lem:cut-mix-equipollence}, first replace every cut in the given
derivation by the mix instance \(m=n=1\); the reduction below eliminates
mixes, and steps 1, 2, 5, and 6 apply verbatim with ``mix'' for ``cut''
and the mix formula for the cut formula.
Choose an uppermost cut of maximal cut rank.
\begin{enumerate}
  \item If the cut formula is not principal in the final inference of one
        premise, and that inference is not a contraction of the cut formula,
        commute the cut above that inference.  The cut rank is
        unchanged and the cut height decreases.
  \item If the cut formula is introduced by weakening in one premise, delete
        the cut and restore the conclusion by weakenings and contractions.
  \item If a premise of a mix ends in a contraction of the mix formula
        (\textsc{CL} above the right premise, \textsc{CR} above the
        left), replace that premise by the contraction's premise and
        raise the corresponding multiplicity by one: the rank is
        unchanged and the height decreases by one.
  \item If the cut formula is principal on both sides, replace the cut by cuts
        on proper subformulas.  For example, a cut between \(\land R\) and
        \(\land L\) is replaced by cuts on the two conjuncts.  Every new cut
        has strictly smaller cut rank, so the number of cuts of maximal cut
        rank decreases.  When the mix formula is principal on both sides
        with multiplicities \((m,n)\), the principal reduction yields
        mixes on proper subformulas, of smaller rank, together with
        residual mixes on the remaining occurrences, of equal rank but
        smaller height.
  \item If an initial sequent is a premise of the cut, delete the cut directly.
  \item For quantifiers, rename eigenvariables before permutation and use the
        substitution lemma.  This prevents capture and preserves the
        eigenvariable conditions.
  \item Repeat according to the well-founded lexicographic measure.  The
        resulting derivation has no cuts.
\end{enumerate}
The calculus-specific table of commuting and principal reductions is the
omitted verification; it occupies the
complete proof in
\cite{TroelstraSchwichtenberg2000,NegriVonPlato2001}.

\begin{definition}[Subformula property]
A propositional derivation has the subformula property iff every
formula occurrence in it is a subformula of a formula in its endsequent.
\end{definition}

\begin{lemma}[Subformula property of cut-free derivations]
\label{lem:cutfree-subformula}
Every cut-free propositional \(\LK\)- or \(\LJ\)-derivation has the
subformula property.
\textit{Status: complete proof.}
\end{lemma}

\begin{proof}
In a cut-free derivation, inspect each rule upward from the endsequent.  Every
premise formula is either copied from the conclusion or is an immediate
subformula of the principal formula.  Induction on derivation height completes
the argument.
\end{proof}

\begin{corollary}[Propositional subformula property]
Every derivable propositional \(\LK\)- or \(\LJ\)-sequent has a cut-free
derivation with the subformula property.
\textit{Status: proof sketch.}
\end{corollary}

\begin{proofplan}
The cut-elimination theorem transforms a given derivation of the sequent
into a cut-free derivation of the same endsequent, and
\cref{lem:cutfree-subformula} shows that every such derivation has the
subformula property.
This derivation uses the preceding proof sketch of the cut-elimination
theorem; a complete proof supplies its reduction case table.
\end{proofplan}

\begin{remark}[First-order qualification]
For first-order cut-free derivations, quantifier rules introduce substitution
instances such as \(A[t/x]\).  The literal propositional subformula property
therefore fails.  The correct statement is a subformula property modulo term
substitution and eigenvariable renaming.  State the selected formulation
before using it in proof search.
\end{remark}

\begin{corollary}[Underivability of the empty sequent]
There is no \(\LK\)-derivation of the empty sequent
\(\varnothing\Rightarrow\varnothing\).
\textit{Status: proof sketch.}
\end{corollary}

\begin{proofplan}
If the empty sequent were derivable, cut elimination would give a cut-free
derivation of it.  The empty sequent is not an instance of a zero-premise
rule, and the conclusion of every other cut-free rule contains at least one
displayed formula occurrence.  Hence no cut-free derivation has the empty
sequent at its root, a contradiction.
This derivation uses the preceding proof sketch of the cut-elimination
theorem; a complete proof supplies its reduction case table.
\end{proofplan}

\section{What cut elimination establishes}
\label{sec:cut-elimination-02}

Cut elimination establishes admissibility of cut and yields analytic
derivations.  Depending on the calculus, it supports consistency,
interpolation, conservativity, and proof-search results.  None of those
consequences follows merely from the phrase ``cut-free'': each requires a
separate argument using the exact rule set.

Cut elimination does not imply that removing cut makes proof search efficient.
Elimination can cause a non-elementary increase in derivation size, and a
cut-free search space may still be infinite because of quantifier terms.

\begin{warning}[Calculus dependence]
Adding equality rules, induction, fixed points, modal rules, or infinitary rules
changes the cut reductions.  A theorem for the displayed \(\LK\) or \(\LJ\)
does not automatically apply to the extension.
\end{warning}

\begin{exercise}
\label{ex:cut-elimination-01}
Write the principal cut reduction for a cut formula \(A\land B\) when the left
premise ends in \(\land R\) and the right premise ends in \(\land L\).  Show
that each new cut has smaller cut rank.
\end{exercise}

\begin{exercise}
\label{ex:cut-elimination-02}
Explain why the subformula property gives a decision procedure for
propositional \(\LJ\), but not directly for first-order \(\LJ\).
\end{exercise}

\section{Interpolation and definability}
\label{sec:cut-elimination-interpolation}

Interpolation is one of the consequences of cut elimination named in the
preceding section; its statement, and the definability theorem derived from
it, are recorded here at cited-theorem level.

\begin{theorem}[Craig interpolation]
\label{thm:craig-interpolation}
Let \(\varphi\) and \(\psi\) be first-order sentences with
\(\entails\varphi\to\psi\).  Then there is a sentence \(\theta\) such
that \(\entails\varphi\to\theta\), \(\entails\theta\to\psi\), and
every nonlogical symbol of \(\theta\) occurs in both \(\varphi\) and
\(\psi\); equality and \(\bottom\) may occur in \(\theta\)
regardless.  The degenerate cases are covered because \(\bottom\) and
\(\neg\bottom\) are available as interpolants when \(\varphi\) is
unsatisfiable or \(\psi\) is valid.
\textit{Status: cited theorem.}
\end{theorem}

See \cite{TroelstraSchwichtenberg2000} for the proof-theoretic route
through Maehara's lemma over a cut-free sequent calculus, and
\cite{Hodges1993} for the model-theoretic route through compactness and
elementary diagrams.

\begin{theorem}[Beth definability]
\label{thm:beth}
Let \(T\) be a first-order theory in a signature containing a
distinguished relation symbol \(R\), and suppose \(T\) defines \(R\)
implicitly: any two models of \(T\) with the same reduct to the signature
without \(R\) interpret \(R\) identically.  Then \(T\) defines \(R\)
explicitly: some formula \(\theta(\bar x)\) not containing \(R\)
satisfies
\(T\derives\forall\bar x\,(R(\bar x)\leftrightarrow\theta(\bar x))\).
\textit{Status: cited theorem.}
\end{theorem}

Beth definability follows from Craig interpolation; see \cite{Hodges1993}.
The Robinson joint-consistency theorem is a further standard equivalent of
interpolation and is stated and proved in the same source.  The theorem is
converse to the conservativity of definitional extensions
(\cref{thm:definitional-conservativity}): adding a defining axiom for
\(R\) produces a theory that defines \(R\) explicitly, hence implicitly,
while proving no new sentences of the smaller signature; Beth definability
states that implicit definability already guarantees an explicit
definition.

\chapter[Normalization and Ordinal Analysis]{Normalization, Proof Terms, and Ordinal Analysis}
\label{ch:normalization-curry-howard}

\section{Natural-deduction normalization}
\label{sec:normalization-curry-howard-01}

The normalization theory of natural deduction developed in this chapter
originates with Prawitz \cite{Prawitz1965}.

\begin{definition}[Detour]
A detour is a segment of a natural-deduction derivation in which a formula is
introduced by an introduction rule and is immediately used as the major
premise of the matching elimination rule; the \emph{major premise} of an
elimination rule is the premise exhibiting the connective or quantifier
that the rule eliminates.  A derivation is normal when it has
no detour and satisfies the required permutation conditions for disjunction
and existential elimination.
\end{definition}

For implication, the principal reduction has the schematic form
\[
  \dfrac{
    \dfrac{\begin{matrix}{[A]}^i\\ \vdots\,D\\ B\end{matrix}}
          {A\to B}\;\to I_i
    \qquad
    \begin{matrix}\vdots\,E\\ A\end{matrix}
  }{B}\;\to E
  \quad\leadsto\quad
  \begin{matrix}\vdots\,E\\ A\\ \vdots\,D\\ B\end{matrix}\;,
\]
where \(D[E/A^i]\) substitutes the derivation \(E\) for every discharged
occurrence carrying label \(i\).  This is substitution of derivations, not term
substitution inside formulas.

For the theorem below, the calculus is exactly the finite-tree
natural-deduction presentation of \cref{ch:prop-deduction,ch:fol-deduction}
with the rules
\(\land I,\land E_1,\land E_2,\lor I_1,\lor I_2,\lor E,\to I,\to E,
\bottom E,\forall I,\forall E,\exists I,\exists E\).  It omits
\(\mathsf{RAA}\), the primitive-negation rules, and the equality rules, and
uses \(\neg A\defeq A\to\bottom\).  Its quantifier substitutions are the total
capture-avoiding substitutions fixed earlier, and its discharge and
eigenvariable conditions are literally those attached to the cited rules.
This declaration fixes the derivation objects to which ``reduces'' and
``normal'' refer; it does not posit an unnamed standard calculus.

\begin{theorem}[Normalization for intuitionistic natural deduction]
Every derivation in the calculus just fixed reduces to a normal derivation with
the same conclusion and no undischarged assumption that was not already
undischarged in the original derivation.  Equivalently, after identifying
alpha-equivalent formulas, the resulting set of undischarged assumption
formulas is a subset of the original set.
\end{theorem}

\paragraph{Proof sketch.}
Define maximal formulas and maximal segments.  Principal reductions remove an
introduction--elimination detour; permutation conversions move eliminations
above disjunction or existential eliminations.  A naive count of formula
occurrences is insufficient because a reduction can duplicate a
subderivation.  The standard proof uses a lexicographic measure on the rank
of maximal formulas and the lengths of maximal segments, or translates to a
typed term calculus and proves strong normalization there.  The full
calculus-specific proof is given in \cite{TroelstraSchwichtenberg2000}.

Normalization and cut elimination express related proof transformations, but
they are not one theorem.  Normalization concerns a reduction relation on
natural-deduction derivations; cut elimination concerns removal of a sequent
rule.

\section{The simply typed lambda calculus}
\label{sec:normalization-curry-howard-02}

Fix a set \(\mathsf{Base}\) of base types.  Types and terms are generated by
the following scoped metavariables: in this section only, \(A,B\) range over
types, \(M,N\) range over terms, and \(\alpha\) ranges over members of
\(\mathsf{Base}\).  The symbol \(\Theta_{\mathrm{ty}}\) ranges over finite
lists assigning distinct term variables to types; these lists are the typing
contexts of this section.  These are metalinguistic roles, not additional
object-language variables.
\[
  A,B::=\alpha\mid A\to B,
  \qquad \alpha\in\mathsf{Base},
\]
and
\[
  M,N::=x\mid \lambda x{:}A.M\mid MN.
\]
The typing judgment \(\Theta_{\mathrm{ty}}\derives M:A\) is generated by
\[
\frac{x{:}A\in\Theta_{\mathrm{ty}}}
     {\Theta_{\mathrm{ty}}\derives x:A}\;\textsc{Var},
\qquad
\frac{\Theta_{\mathrm{ty}},x{:}A\derives M:B}
     {\Theta_{\mathrm{ty}}\derives\lambda x{:}A.M:A\to B}
     \;\textsc{\(\to I\)},
\]
In \textsc{\(\to I\)}, the variable \(x\) does not occur in
\(\Theta_{\mathrm{ty}}\); this can always be arranged by renaming the
bound variable of the abstraction, and typing judgments are stable under
such renaming of the subject term.
\[
\frac{\Theta_{\mathrm{ty}}\derives M:A\to B
      \qquad
      \Theta_{\mathrm{ty}}\derives N:A}
     {\Theta_{\mathrm{ty}}\derives MN:B}\;\textsc{\(\to E\)}.
\]
The principal computation is beta reduction:
\[
  (\lambda x{:}A.M)N\longrightarrow_\beta M[N/x],
\]
where term substitution is capture-avoiding.

\begin{theorem}[Subject reduction]
If \(\Theta_{\mathrm{ty}}\derives M:A\) and
\(M\longrightarrow_\beta N\), then
\(\Theta_{\mathrm{ty}}\derives N:A\).
\end{theorem}

\paragraph{Proof sketch.}
Prove the typing substitution lemma: if
\(\Theta_{\mathrm{ty}},x{:}B\derives M:A\) and
\(\Theta_{\mathrm{ty}}\derives N:B\), then
\(\Theta_{\mathrm{ty}}\derives M[N/x]:A\).  Apply it to the contracted beta
redex and use induction on the compatible closure of beta reduction for
reductions inside a term.

\begin{theorem}[Strong normalization]
Every well-typed term in the simply typed lambda calculus has no infinite
beta-reduction sequence.
\end{theorem}

\paragraph{Proof sketch.}
Assign to each type a reducibility candidate.  A term belongs to the candidate
for a base type when it is strongly normalizing; it belongs to the candidate
for \(A\to B\) when application to every member of the candidate for \(A\)
produces a member of the candidate for \(B\).  The fundamental lemma proves by
induction on typing derivations that every well-typed term belongs to its
type's candidate.  Strong normalization follows.  This proof sketch gives the
logical-relations architecture; closure lemmas for the candidates are omitted.

\section{Curry--Howard correspondence}
\label{sec:normalization-curry-howard-03}

Under the Curry--Howard correspondence, a proposition is represented by a
type, a derivation by a typed term, implication introduction by lambda
abstraction, implication elimination by application, and normalization by
computation.  The correspondence is structural: it does not identify every
formal derivation with a program without first fixing a calculus and a
term assignment.

\begin{example}[Composition]
The term
\[
  \lambda f{:}(B\to C).\,
  \lambda g{:}(A\to B).\,
  \lambda x{:}A.\,f(gx)
\]
has type
\[
  (B\to C)\to(A\to B)\to A\to C.
\]
Reading types as propositions gives a proof of the corresponding intuitionistic
formula.  Each abstraction discharges one typed assumption, and each
application is an implication-elimination step.
\end{example}

Products and sums correspond to conjunction and disjunction after their term
constructors, eliminators, and reduction rules are defined.  The empty type
corresponds to \(\bottom\).  Classical principles require additional control
operators, translations, or axioms; they are not terms of the pure simply
typed lambda calculus.

\section{Scope of ordinal analysis}
\label{sec:normalization-curry-howard-04}

The landscape of ordinal-informative consistency analyses is surveyed in
\cite{Buss1998}.

\paragraph{Cited theorem.}
Gentzen's consistency analysis of arithmetic, in the sequent-reduction
presentation surveyed in
\cite[Chapter~10, especially Sections~10.1--10.4, pp.~318--337]
     {TroelstraSchwichtenberg2000}, assigns notations below
\(\varepsilon_0\) to arithmetic derivations so that each critical reduction
strictly lowers the assigned notation.  In a metatheory that accepts
transfinite induction for the predicates and primitive-recursive ordering used
in that particular reduction analysis, termination yields
\(\mathsf{Con}(\PA)\).  This is an analysis-relative imported result;
it is not a definition of a canonical invariant and it does not assert that an
unspecified ordinal ``measures'' a theory.

\paragraph{Formal objects needed to read the citation.}
The symbol \(\varepsilon_0\) denotes the least nonzero ordinal
\(\alpha\) satisfying \(\omega^\alpha=\alpha\).  A standard external way to
write ordinals below it uses finite Cantor-normal-form terms
\[
  a=\omega^{a_1}\cdot m_1+\cdots+
    \omega^{a_k}\cdot m_k,
\]
where \(k\geq1\), each \(m_i\) is a positive natural number, the exponent
terms satisfy \(a_1>\cdots>a_k\), and \(0\) is also a term.  Exponents are
formed recursively by the same clause.  Interpreting the symbols by ordinal
addition and exponentiation assigns each such finite term an ordinal below
\(\varepsilon_0\).  This display supplies notation for the following worked
calculation; it is not asserted to be the source's complete primitive-recursive
coding.

For example, \(\lambda=\omega^\omega\) is a limit notation.  The sequence
\[
  \lambda[n]=\omega^n\qquad(n\geq1)
\]
satisfies \(\lambda[n]<\lambda\) for every \(n\) and
\(\sup_{n\geq1}\lambda[n]=\lambda\).  Likewise, for
\(\mu=\omega^3\), define \(\mu[n]=\omega^2\cdot n\) for \(n\geq1\).  Then
\[
  \mu[4]=\omega^2\cdot4<\omega^3=\mu.
\]
This is the elementary shape of a fundamental-sequence descent: a limit
notation is replaced by a selected smaller approximation.  It does not show
that any particular proof reduction receives \(\mu\), because that claim
depends on the source's derivation coding and ordinal assignment.

The cited result therefore has three separately typed ingredients: finite
derivation codes for the selected sequent presentation of first-order
arithmetic; ordinal-notation codes with a primitive-recursive comparison and
fundamental-sequence operation; and a primitive-recursive reduction on
derivation codes together with an assignment from derivations to notation
codes.  Its decisive lemma says, case by case for that reduction, that a
critical reduction strictly decreases the assigned notation.  Transfinite
induction for the predicates used in the reduction then excludes an infinite
reduction sequence, and the source extracts the consistency statement from
the resulting normal form.

\paragraph{Internal proof boundary.}
A complete proof in this manuscript would still have to reproduce the exact
arithmetic rule surface and derivation coding, the complete notation grammar
and algorithms, every reduction clause, the assignment, every case of the
decrease lemma, the precise transfinite-induction schema, and the final
consistency extraction.  A later use of the technical phrase
\emph{proof-theoretic ordinal} for an arbitrary theory would additionally need
a stated reduction notion, a leastness or supremum clause, and a comparison
showing what is invariant under alternative presentations.  No such general
invariant is defined here.

The imported result is relative.  It does not give a finitary proof of
\(\mathsf{Con}(\PA)\) formalizable in \(\PA\) itself.  If \(\PA\) is
consistent, the second incompleteness theorem excludes such an internal proof
for the standard arithmetized consistency statement.  Stronger theories can
require notation systems far beyond \(\varepsilon_0\), and some strength
comparisons are not captured by a single canonical ordinal.

\begin{warning}[Unqualified strength claims]
Do not infer that a larger displayed ordinal alone makes one theory stronger
than another.  Compare the theories, notation systems, base metatheories, and
reflection principles used in the analyses.
\end{warning}

\begin{exercise}
\label{ex:normalization-curry-howard-01}
Type the term \(\lambda x{:}A.\lambda y{:}B.x\), and give the corresponding
natural-deduction derivation.  Identify the structural weakening implicit in
the unused variable \(y\).
\end{exercise}

\begin{exercise}
\label{ex:normalization-curry-howard-02}
Prove subject reduction for a single outermost beta redex.  State the typing
substitution lemma exactly and mark where capture avoidance is used.
\end{exercise}

\begin{definition}[Abstract reduction relation and normalization]
An \emph{abstract reduction relation} is a binary relation \(\to\) on a
set.  An element is \emph{strongly normalizing} iff no infinite
\(\to\)-sequence starts from it; the relation is strongly normalizing iff
every element is.  An element is \emph{weakly normalizing} iff some finite
\(\to\)-sequence from it ends in an element with no \(\to\)-successor;
the relation is weakly normalizing iff every element is.  Beta reduction on
well-typed terms is an abstract reduction relation, and the
strong-normalization theorem above states that it is strongly normalizing.
\end{definition}

\begin{exercise}
\label{ex:normalization-curry-howard-03}
Explain why a strongly normalizing abstract reduction relation is weakly
normalizing, and exhibit an abstract reduction relation that is weakly but
not strongly normalizing.
\end{exercise}

\part{Second-Order and Higher-Order Logic}
\label{part:higher-order}

\chapter{Second-Order Syntax and Two Semantics}
\label{ch:second-order-logic}

Second-order syntax permits quantification over relations on the first-order
domain.  That syntactic extension admits at least two distinct semantic regimes.  Full semantics quantifies over every relation of the
required arity.  Henkin semantics quantifies over a specified collection of
admissible relations.  The distinction controls categoricity, compactness,
L\"owenheim--Skolem properties, and completeness
\cite{Henkin1950,BarwiseFeferman1985}.

\section{Second-order formulas}
\label{sec:second-order-logic-01}

\begin{definition}[Second-order signature]
A second-order signature \(\Lang_2\) contains the constant, function, and
relation symbols of a first-order signature.  Its logical alphabet also
contains, for every \(n\geq1\), a countably infinite set
\(\mathsf{RelVar}_n\) of \(n\)-ary relation variables.  Relation variables are
logical variables; nonlogical relation symbols remain members of \(\Lang_2\).
\end{definition}

First-order terms are generated as before.  If \(X\in\mathsf{RelVar}_n\) and
\(t_1,\ldots,t_n\) are terms, then \(X(t_1,\ldots,t_n)\) is atomic.  Formulas
are closed under the propositional connectives, first-order quantifiers
\(\forall x,\exists x\), and second-order quantifiers
\(\forall X^n,\exists X^n\).  The superscript records arity and may be omitted
only when the arity is fixed by context.

\begin{definition}[Second-order sentence]
A second-order sentence is a second-order formula having no free first-order
variable occurrence and no free relation-variable occurrence.
\end{definition}

Term substitution for first-order variables and relation-formula substitution
for relation variables are different operations.  For example, substituting a
formula \(\varphi(x_1,\ldots,x_n)\) for a free occurrence of \(X^n\) requires
renaming every binder that could capture a variable free in \(\varphi\).

\section{Full semantics}
\label{sec:second-order-logic-02}

\begin{definition}[Full second-order assignment]
Let \(\mathcal M\) be a \(\Lang_2\)-structure with domain \(M\).  A full
second-order assignment \(s\) maps each first-order variable to an element of
\(M\) and each \(n\)-ary relation variable to an arbitrary member of
\(\Pow(M^n)\).
\end{definition}

The first-order atomic and connective clauses are unchanged.  The new atomic
and quantifier clauses are
\[
  \mathcal M,s\entails X(t_1,\ldots,t_n)
  \quad\Longleftrightarrow\quad
  \tuple{\sem{t_1}^{\mathcal M}_s,\ldots,
         \sem{t_n}^{\mathcal M}_s}\in s(X),
\]
and
\[
\begin{aligned}
  \mathcal M,s\entails\forall X^n\,\varphi
  \quad\Longleftrightarrow\quad
  &\text{for every }R\subseteq M^n,
    \ \mathcal M,s[X\mapsto R]\entails\varphi,\\
  \mathcal M,s\entails\exists X^n\,\varphi
  \quad\Longleftrightarrow\quad
  &\text{for some }R\subseteq M^n,
    \ \mathcal M,s[X\mapsto R]\entails\varphi.
\end{aligned}
\]

\begin{definition}[Full semantic consequence]
For a set \(\Gamma\) of second-order formulas,
\[
  \Gamma\entails_{\mathrm{full}}\varphi
\]
iff every structure under every full second-order assignment that
satisfies every member of \(\Gamma\) also satisfies \(\varphi\).
\end{definition}

The word \emph{full} modifies the range of second-order quantification.  It
does not mean that the first-order domain contains every set or that the
structure is a set-theoretic universe.

\begin{theorem}[Full comprehension validity]
Let \(\varphi(\bar x)\) be a formula in which \(X^n\) is not free.  Under full
semantics the universal closure of
\[
  \exists X^n\,\forall x_1\cdots\forall x_n\,
  \bigl(X(\bar x)\leftrightarrow\varphi(\bar x)\bigr)
\]
is valid.
\end{theorem}

\begin{proof}
Fix \(\mathcal M\) and a full second-order assignment \(s\) of all
parameters.  Let
\[
  R=\set{\bar a\in M^n:
          \mathcal M,s[\bar x\mapsto\bar a]\entails\varphi(\bar x)}.
\]
Full semantics permits \(X\) to range over this particular subset of \(M^n\).
With \(s(X)=R\), the biconditional follows from the definition of \(R\).
\end{proof}

\section{Full second-order categoricity}
\label{sec:second-order-logic-03}

Let \(\Lang_{\mathsf N}=\set{0,S}\), with a constant symbol \(0\) and a unary
function symbol \(S\).  The full second-order Peano axioms are
\[
  \forall x\,(Sx\ne0),
\]
\[
  \forall x\,\forall y\,(Sx=Sy\to x=y),
\]
and the second-order induction sentence
\[
  \forall X\,
  \left[
    \left(X(0)\land
      \forall x\,(X(x)\to X(Sx))\right)
    \to
    \forall x\,X(x)
  \right].
\]

\begin{theorem}[Categoricity of full second-order arithmetic]
Every two structures satisfying the displayed second-order Peano axioms under
full semantics are isomorphic.  In particular, every such structure is
isomorphic to \((\N,0,n\mapsto n+1)\).
\end{theorem}

\begin{proof}
Let \(\mathcal M=(M,0^{\mathcal M},S^{\mathcal M})\) satisfy the axioms.  Define
\(f:\N\to M\) recursively by
\[
  f(0)=0^{\mathcal M},
  \qquad
  f(n+1)=S^{\mathcal M}(f(n)).
\]
The map \(f\) is injective.  The proof is a metatheoretic induction on
\(n\in\N\) of the claim: for every \(m\in\N\), if \(f(n)=f(m)\), then
\(n=m\).  Base case \(n=0\): suppose \(f(0)=f(m)\) with \(m\ne0\).  Then
\(m=k+1\) for some \(k\in\N\), so the supposition reads
\(0^{\mathcal M}=S^{\mathcal M}(f(k))\), contradicting the first axiom;
hence \(m=0\).  Induction step: assume the claim for \(n\) and suppose
\(f(n+1)=f(m)\).  If \(m=0\), then
\(S^{\mathcal M}(f(n))=0^{\mathcal M}\), and the first axiom excludes this;
so \(m=k+1\) for some \(k\in\N\), and the supposition reads
\(S^{\mathcal M}(f(n))=S^{\mathcal M}(f(k))\).  The second axiom gives
\(f(n)=f(k)\), the induction hypothesis gives \(n=k\), and therefore
\(n+1=k+1=m\).  The range \(R\subseteq M\) of \(f\) contains
\(0^{\mathcal M}\) and is closed under \(S^{\mathcal M}\).  Under full
semantics, the induction quantifier ranges over this actual subset \(R\).
Induction therefore gives \(R=M\), so \(f\) is surjective.  It preserves the
constant and successor by construction and is therefore an isomorphism.
\end{proof}

If addition and multiplication are included with their recursive axioms, the
same isomorphism preserves them by induction on \(\N\).  The categoricity claim
depends on full semantics; it is not a theorem about arbitrary Henkin
structures.

\section{Failure of compactness and effective completeness}
\label{sec:second-order-logic-04}

\begin{theorem}[Failure of compactness under full semantics]
Full second-order semantic consequence is not compact.
\end{theorem}

\begin{proof}
Expand \(\Lang_{\mathsf N}\) by a constant \(c\).  Let \(T\) contain the full
second-order Peano axioms and, for every \(n\in\N\), the sentence
\[
  c\ne S^n0.
\]
Every finite subset of \(T\) is satisfied under full semantics by a suitable
structure: use the standard natural numbers and interpret \(c\) as a number
larger than every index occurring in that finite subset.  If one structure
satisfied all of \(T\) under full semantics, categoricity would identify its
domain with \(\N\), while the displayed sentences would require \(c\) to
differ from every natural number.  Hence no such structure exists.
\end{proof}

\begin{theorem}[No effective complete calculus for full validity]
Here an effective calculus means one whose finite derivations are
computably checkable, so its theorems form a computably enumerable set.  There
is no such calculus \(\mathsf D\) for which, for every second-order sentence
\(\sigma\),
\[
  \mathsf D\derives\sigma
  \quad\Longleftrightarrow\quad
  \entails_{\mathrm{full}}\sigma.
\]
Indeed, the set of full second-order validities is not computably enumerable.
\end{theorem}

\paragraph{Proof sketch.}
Use the finite full second-order axiomatization of arithmetic in the signature
with \(0,S,+,\times\).  For each first-order arithmetic sentence \(\sigma\),
effectively form the second-order sentence
\(\mathsf{PA}_2\to\sigma\).  Full categoricity gives
\[
  \entails_{\mathrm{full}}(\mathsf{PA}_2\to\sigma)
  \quad\Longleftrightarrow\quad
  (\N,0,S,+,\times)\entails\sigma.
\]
If full second-order validities were computably enumerable, true
first-order arithmetic would therefore be computably enumerable,
contradicting the corollary of Tarski's undefinability theorem
(\cref{thm:tarski-undefinability}) that true arithmetic is not computably
enumerable.  The reduction and its effectiveness are the omitted
steps controlled by \cite{BarwiseFeferman1985,Shapiro1991}; the computability
proof itself is not reproduced.

This theorem excludes an effective, sound, and complete calculus for full
validity.  It does not exclude an ineffective presentation that takes all full
validities as axioms.  The effectiveness qualification is therefore part of
the theorem, not an optional refinement; see
\cite{BarwiseFeferman1985,Shapiro1991}.

\begin{theorem}[Categoricity of the full second-order reals]
\label{thm:second-order-reals-categorical}
Let \(\Lang_{\mathsf R}=\set{0,1,+,\times,<}\).  Take the finitely many
first-order axioms stating that \((0,1,+,\times,<)\) is an ordered field,
together with the second-order completeness sentence
\[
  \forall X\,
  \Bigl[
    \exists x\,X(x)
    \land
    \exists y\,\forall x\,(X(x)\to x\le y)
    \;\to\;
    \exists z\,
    \bigl(
      \forall x\,(X(x)\to x\le z)
      \land
      \forall y\,(\forall x\,(X(x)\to x\le y)\to z\le y)
    \bigr)
  \Bigr],
\]
where \(u\le v\) abbreviates \(u<v\lor u=v\) in this scope: every nonempty
subset with an upper bound has a least upper bound.  Any two
\(\Lang_{\mathsf R}\)-structures satisfying these axioms under full
semantics are isomorphic.  In particular, every such structure is isomorphic
to the ordered field \(\R\), and its domain is uncountable.
\textit{Status: cited theorem.}
\end{theorem}

See \cite{Shapiro1991} for the categoricity proof, which sends each element
of one model to the element of the other determined by the same cut of the
common rationals and verifies that this map is an ordered-field isomorphism;
the ordered-field axioms abbreviated above are listed there.  The
uncountability of \(\R\) is supplied by \cite{Jech2003}.

Full second-order logic also lacks the first-order L\"owenheim--Skolem
properties.  The conjunction of the displayed second-order Peano axioms is
satisfied by no uncountable structure
under full semantics, so upward L\"owenheim--Skolem fails.  In the other
direction, the axioms of \cref{thm:second-order-reals-categorical} are
satisfied by the infinite structure \(\R\), while
\cref{thm:second-order-reals-categorical} makes every structure satisfying
them under full semantics isomorphic to \(\R\) and therefore uncountable.
These axioms are thus satisfied by an infinite structure but by no countable
structure; downward L\"owenheim--Skolem fails.

\begin{theorem}[Lindstr\"om]
\label{thm:lindstrom}
Let \(\mathcal L\) be a regular abstract logic extending first-order
logic: one closed under the Boolean connectives, quantification, and
relativization, with satisfaction invariant under isomorphism.  If
\(\mathcal L\) satisfies countable compactness and the downward
L\"owenheim--Skolem property to countable structures, then
\(\mathcal L\) has exactly the expressive power of first-order logic.
\textit{Status: cited theorem.}
\end{theorem}

See \cite{Lindstrom1969} for the original proof and
\cite[Chapter~XIII]{EbbinghausFlumThomas1994} for a complete textbook
treatment, including the exact closure conditions abbreviated here as
regularity.  The theorem explains the pattern of this section: full
second-order logic strictly exceeds first-order expressive power, so it
must lose countable compactness or the downward L\"owenheim--Skolem
property, and this section shows that it loses both.

\section{Henkin semantics}
\label{sec:second-order-logic-05}

\begin{definition}[Henkin structure]
A Henkin structure for the relational second-order language is a structure
\[
  \mathcal H=(M,(\mathcal R_n)_{n\geq1},\ldots),
\]
where \(M\ne\varnothing\), each
\(\mathcal R_n\) is a nonempty subset of \(\Pow(M^n)\), and the remaining
components interpret the nonlogical symbols.  An \(n\)-ary relation variable
ranges over \(\mathcal R_n\), not over all of \(\Pow(M^n)\).
\end{definition}

A systematic treatment of Henkin semantics and of the translation between
second-order logic and many-sorted first-order logic is
\cite{Manzano1996}.

When the calculus contains a comprehension schema, the semantic class
must be restricted accordingly: every relation definable by an eligible
formula with parameters from the Henkin domains must belong to the relevant
\(\mathcal R_n\).  This closure is a hypothesis on the semantic class.  It is
not automatic from \(\mathcal R_n\subseteq\Pow(M^n)\).

\begin{definition}[Henkin semantic consequence]
Write
\[
  \Gamma\entails_{\mathrm H}\varphi
\]
when every Henkin structure under every sorted assignment satisfying
\(\Gamma\) satisfies \(\varphi\).  A full second-order structure is the
special case
\(\mathcal R_n=\Pow(M^n)\) for every \(n\).
\end{definition}

\paragraph{Cited theorem.}
For the typed formula system and general structures fixed in
\cite[pp.~81--91]{Henkin1950}, a formula is a theorem of that source's system
iff it is valid in that source's general structures.  The hypotheses
include the source's typed axiom and rule surface and its corresponding closure
conditions; the conclusion is completeness for that paired syntax and
semantics, not for full second-order validity.  This is a cited theorem.  It
does not define derivability for this chapter, whose relational language has
not imported that rule surface or selected a comprehension schema.

\paragraph{Worked many-sorted translation.}
The semantic reduction underlying Henkin completeness can nevertheless be
stated exactly without choosing a proof calculus.  Use a many-sorted
first-order language with an object sort \(D\), a relation sort \(S_n\) for
each \(n\geq1\), and an \((n+1)\)-ary predicate symbol \(E_n\) whose argument
sorts are \((S_n,D,\ldots,D)\).  All original constants, function symbols, and
relation symbols keep their object-sort interpretations.

To distinguish a second-order formula from its many-sorted image, define the
metalinguistic translation \(\varphi\mapsto\varphi^\sharp\) recursively.  It
leaves object terms, equality, and nonlogical relation atoms unchanged,
commutes with every propositional connective, and has the clauses
\[
\begin{aligned}
  \bigl(X^n(t_1,\ldots,t_n)\bigr)^\sharp
    &\defeq E_n(X,t_1,\ldots,t_n),\\
  (\forall x\,\psi)^\sharp
    &\defeq\forall x{:}D\,\psi^\sharp,
  & (\exists x\,\psi)^\sharp
    &\defeq\exists x{:}D\,\psi^\sharp,\\
  (\forall X^n\,\psi)^\sharp
    &\defeq\forall X{:}S_n\,\psi^\sharp,
  & (\exists X^n\,\psi)^\sharp
    &\defeq\exists X{:}S_n\,\psi^\sharp.
\end{aligned}
\]
This translation is not mere renaming: it replaces relation application by an
incidence predicate and makes the two quantifier ranges into sorts.  Its
standing is a recursively eliminable metalinguistic operation, and it is used
in the satisfaction equivalence below.

From a Henkin structure
\(\mathcal H=(M,(\mathcal R_n)_{n\geq1},\ldots)\), form the many-sorted
structure \(\mathcal H^\sharp\) by taking
\[
  D^{\mathcal H^\sharp}=M,
  \qquad
  S_n^{\mathcal H^\sharp}=\mathcal R_n,
  \qquad
  E_n^{\mathcal H^\sharp}(R,\bar a)
  \Longleftrightarrow \bar a\in R.
\]
This structure is needed because \(\varphi^\sharp\) contains sorted variables
and the incidence symbols \(E_n\), which the original second-order structure
does not interpret as first-order symbols.  Its standing is a many-sorted
first-order structure.  It is nonredundant with \(\mathcal H\): relations in
\(\mathcal H\) are values of second-order variables, whereas the same
relations are first-order objects of the sorts \(S_n\) in
\(\mathcal H^\sharp\).  The structure is used immediately in the following
equivalence.
For every second-order formula \(\varphi\) and every matching sorted assignment
\(s\),
\[
  \mathcal H,s\entails\varphi
  \quad\Longleftrightarrow\quad
  \mathcal H^\sharp,s\entails\varphi^\sharp.
\]
Indeed, induction on \(\varphi\) handles an atom \(X^n(\bar t)\) by the
definition of \(E_n\); first-order atoms because the object interpretations are
unchanged; connectives by the induction hypothesis; object quantifiers because
both sides range over \(M\); and relation quantifiers because both sides range
over \(\mathcal R_n\).  These cases exhaust the formula grammar.

If a selected calculus contains a comprehension instance for an eligible
formula \(\theta(x_1,\ldots,x_n,\bar z,\bar Y)\), where \(X^n\) is not free
in \(\theta\), its semantic counterpart is the many-sorted sentence
\[
  \forall\bar z\,\forall\bar Y\,
  \exists X{:}S_n\,\forall\bar x{:}D\,
  \bigl(E_n(X,\bar x)\leftrightarrow
        \theta^\sharp(\bar x,\bar z,\bar Y)\bigr),
\]
where \(\bar x=(x_1,\ldots,x_n)\), and every displayed parameter is quantified
in its declared sort.  Thus the corresponding closure condition says exactly
that the \(n\)-ary relation defined by each eligible instance belongs to
\(\mathcal R_n\); it is not automatic for an arbitrary family of relation
domains.

\paragraph{Internal proof boundary.}
The translation proves a semantic equivalence, not a local completeness
theorem.  Before this chapter may assert Henkin completeness, compactness, or a
downward L\"owenheim--Skolem result for its own calculus, it must still list
every axiom, rule, structural convention, equality and extensionality
principle, and eligible comprehension instance; define finite derivations;
match those instances with semantic closure; and carry the chosen many-sorted
first-order completeness construction back to exactly that rule surface.  No
calculus name or derivability--validity biconditional is introduced here.

Henkin completeness and full categoricity concern different consequence
relations.  Combining them into the assertion ``second-order logic is complete
and categorical'' suppresses the semantic parameter and is incorrect.

\begin{example}[Reachability]
In a directed graph with edge relation \(E\), the formula
\[
\begin{aligned}
  \operatorname{Reach}(a,b)\defeq
  \forall X\,\bigl[&
    \bigl(X(a)\land
      \forall u\,\forall v\,
      ((X(u)\land E(u,v))\to X(v))\bigr)\\
    &\to X(b)
  \bigr]
\end{aligned}
\]
holds under full semantics iff a finite directed path leads from
\(a\) to \(b\).  If a path exists, every edge-closed set containing \(a\)
contains its successive vertices.  Conversely, the actual set of vertices
reachable from \(a\) is edge-closed; full quantification permits it as a value
of \(X\).  The converse can fail in a Henkin structure if that reachable set
is not among the admissible unary relations.
\end{example}

\begin{warning}[Semantic parameter omitted]
Never write \(\Gamma\entails\varphi\) for a second-order argument before
declaring full or Henkin semantics.  The same sentence can have different
status under the two consequence relations.
\end{warning}

\begin{exercise}
\label{ex:second-order-logic-01}
Prove that the full second-order induction axiom is satisfied by the standard
natural numbers.  Your proof must quantify over an arbitrary subset of
\(\N\), not merely a definable subset.
\end{exercise}

\begin{exercise}
\label{ex:second-order-logic-02}
Explain precisely which step of the categoricity proof fails in a Henkin
structure whose unary-relation domain does not contain the range of \(f\).
\end{exercise}

\begin{exercise}
\label{ex:second-order-logic-03}
Use the compactness counterexample above to prove that no finitary sound and
strongly complete calculus exists for full second-order semantic consequence.
Distinguish this argument from the stronger non-enumerability theorem.
\end{exercise}

\chapter[Simple Type Theory]{Simple Type Theory and Higher-Order Logic}
\label{ch:simple-type-theory}

The term \emph{higher-order logic} does not determine one calculus.  A formal
presentation must specify its types, term constructors, equality, logical
constants, conversion rules, axioms, and semantic regime.  This chapter fixes
a representative classical simple-type syntax and two semantic classes.  It
does not select a proof calculus.

\section{Types, terms, and formulas}
\label{sec:simple-type-theory-01}

Within this chapter, \(\alpha\) and \(\beta\) are metalinguistic variables
ranging over simple types, while \(F,A,B\) range over well-typed terms; a term
of type \(o\) is a formula.  These are scoped metalinguistic roles, not
object-language variables.  The type metavariables do not retain the earlier
ordinal role and do not reuse the normative sentence metavariables
\(\sigma,\tau\) or substitution metavariable \(\rho\).  The symbols \(\iota\)
and \(o\) are object-language type constants.

\begin{definition}[Simple types]
Let \(\iota\) be the type of individuals and \(o\) the type of truth values.
The simple types are generated by
\[
  \alpha,\beta::=\iota\mid o\mid\alpha\to\beta.
\]
This is a free grammar: \(\iota\ne o\), neither base type is an arrow type,
and
\[
  (\alpha\to\beta)=(\alpha'\to\beta')
  \quad\Longleftrightarrow\quad
  (\alpha=\alpha'\text{ and }\beta=\beta').
\]
In particular, no finite simple type \(\alpha\) satisfies
\(\alpha=\alpha\to o\).
\end{definition}

For each type, fix variables and a possibly empty collection of nonlogical
constants, each annotated with that type.  The well-typed terms are the least
family containing those variables and constants and closed under exactly these
two formation clauses:
\[
  \frac{F:\alpha\to\beta\qquad A:\alpha}{FA:\beta},
  \qquad
  \frac{x:\alpha\qquad B:\beta}
       {\lambda x{:}\alpha.B:\alpha\to\beta}.
\]
Formulas are exactly the terms of type \(o\).
Equality at type \(\alpha\) has type
\[
  =_\alpha:\alpha\to\alpha\to o,
\]
and the universal quantifier at type \(\alpha\) has type
\[
  \forall_\alpha:(\alpha\to o)\to o.
\]
Thus \(\forall x{:}\alpha.\,A\) abbreviates
\(\forall_\alpha(\lambda x{:}\alpha.A)\).

Alpha-equivalence, capture-avoiding substitution, beta reduction, and eta
conversion are typed operations.  A purported application with mismatched
types is not a false formula; it is not a term.

\section{Standard and Henkin semantics}
\label{sec:simple-type-theory-02}

\begin{definition}[Standard type structure]
A standard type structure \(\mathcal M\) assigns nonempty sets \(D_\alpha\) to types such
that
\[
  D_o=\set{0,1},
  \qquad
  D_{\alpha\to\beta}=D_\beta^{D_\alpha},
\]
the full set of functions from \(D_\alpha\) to \(D_\beta\).  Application is
function application, lambda abstraction denotes the corresponding function,
equality denotes identity at each type, and the logical constants have their
classical meanings.  Every admitted nonlogical constant \(c{:}\alpha\) is
assigned an element \(c^{\mathcal M}\in D_\alpha\).
\end{definition}

\begin{definition}[Henkin type structure]
A Henkin type structure \(\mathcal M\) assigns a nonempty set
\(D_\alpha\) to every type, with \(D_o=\set{0,1}\), and
a selected set
\[
  D_{\alpha\to\beta}\subseteq D_\beta^{D_\alpha}
\]
to every function type.  The selected function sets must contain the
denotation, under the recursion below, of every well-typed lambda term under
every assignment.  Application and the logical constants obey the same typed
clauses as in a standard type structure, and every admitted nonlogical constant
\(c{:}\alpha\) is assigned an element
\(c^{\mathcal M}\in D_\alpha\).
\end{definition}

In either kind of structure, a typed assignment \(s\) maps every variable
\(x{:}\alpha\) to a member of \(D_\alpha\).  Denotation is the recursion
\[
\begin{aligned}
  \sem{x}^{\mathcal M}_s&=s(x),
  &\sem{c}^{\mathcal M}_s&=c^{\mathcal M},\\
  \sem{FA}^{\mathcal M}_s
    &=\sem{F}^{\mathcal M}_s\bigl(\sem{A}^{\mathcal M}_s\bigr),
  &\sem{\lambda x{:}\alpha.B}^{\mathcal M}_s
    &=\bigl(a\mapsto
       \sem{B}^{\mathcal M}_{s[x\mapsto a]}\bigr).
\end{aligned}
\]
For a Henkin type structure, its closure requirement says precisely that the
function in the last clause belongs to the selected
\(D_{\alpha\to\beta}\) for every well-typed abstraction and matching
assignment.  Equality at type \(\alpha\) denotes the characteristic function
of identity on \(D_\alpha\), and \(\forall_\alpha\) returns \(1\) exactly on
those predicates in \(D_{\alpha\to o}\) that take value \(1\) at every member
of \(D_\alpha\).

Every standard type structure is a Henkin type structure.  The converse fails
when a selected function space omits set-theoretic functions.

\paragraph{Cited theorem.}
For the formal theory of types and the general structures fixed in
\cite[pp.~81--91]{Henkin1950}, soundness and completeness hold for that exact
paired syntax and semantics.  The hypotheses include the source's primitive
typed constants, axioms, rules, conversion principles, and closure conditions.
In its general structures a function domain need contain the functions denoted
by the source's terms under assignments, rather than every set-theoretic
function of the appropriate type.  This is a cited theorem.  It neither makes
full-function semantics complete nor selects a derivability relation for the
syntax-only presentation in this chapter.

\paragraph{Worked semantic calculation.}
The closure clause in a Henkin type structure is required already to interpret
lambda abstraction.  Let \(x{:}\alpha\), let \(B{:}\beta\), let
\(A{:}\alpha\), let \(\mathcal M\) be a standard or Henkin type structure,
and let \(s\) be a typed assignment.  Capture-avoiding typed substitution
satisfies the semantic substitution identity
\[
  \sem{B[A/x]}^{\mathcal M}_s
  =
  \sem{B}^{\mathcal M}_{
    s[x\mapsto\sem{A}^{\mathcal M}_s]}.
\]
It is proved by induction on the typed construction of \(B\): the variable
case reads the updated assignment, application uses the two induction
hypotheses, and abstraction first renames a bound variable away from the free
variables of \(A\) and then applies the induction hypothesis under the updated
assignment.  Constants are unchanged.  Consequently a beta redex has the same
denotation as its contractum in both standard and Henkin structures:
\[
\begin{aligned}
  \sem{(\lambda x{:}\alpha.B)A}^{\mathcal M}_s
    &=\sem{\lambda x{:}\alpha.B}^{\mathcal M}_s
       \bigl(\sem{A}^{\mathcal M}_s\bigr)\\
    &=\sem{B}^{\mathcal M}_{
       s[x\mapsto\sem{A}^{\mathcal M}_s]}\\
    &=\sem{B[A/x]}^{\mathcal M}_s.
\end{aligned}
\]
The middle value is available because the denotation of the abstraction is a
member of the selected domain \(D_{\alpha\to\beta}\).  Thus term closure is a
definedness requirement, not an optional strengthening added only for a
completeness proof.

Quantification makes the semantic difference visible.  In
\(\forall f{:}(\alpha\to\beta).\,C\), a standard structure varies \(f\) over
all of \(D_\beta^{D_\alpha}\), whereas a Henkin structure varies it only over
the declared \(D_{\alpha\to\beta}\).  The same typed formula can therefore
have different validity status in the two semantic classes even though the
beta calculation remains sound in both.

\paragraph{Internal proof boundary.}
An internal completeness theorem here would require a co-located list of every
primitive typed constant, conversion rule, logical axiom, equality and
extensionality principle, inference rule, and structural convention; an exact
definition of finite derivations; semantic closure matching each syntactic
principle; and a Henkin construction and truth lemma for precisely that
package.  These data are source-dependent and are not reproduced.  The chapter
therefore introduces no calculus symbol and asserts no local
derivability--validity biconditional.

\begin{example}[A polymorphism that simple types do not express]
At each fixed type \(\alpha\), the identity term
\[
  \lambda x{:}\alpha.x
\]
has type \(\alpha\to\alpha\).  The displayed \(\alpha\) is a metavariable in a
schema of typed terms.  It is not a type variable in the object language of
simple type theory.  Object-language quantification over types would require a
different, explicitly polymorphic type system.
\end{example}

\begin{example}[Typed representation of a binary relation]
A binary relation on individuals has type
\[
  \iota\to\iota\to o,
\]
with application associated to the left.  If
\(R:\iota\to\iota\to o\) and \(a,b:\iota\), then \(Rab:o\) is a formula.
The terms \(Ra\) and \(R\) are not formulas: their types are
\(\iota\to o\) and \(\iota\to\iota\to o\), respectively.
\end{example}

\section{Type discipline and foundational scope}
\label{sec:simple-type-theory-03}

Simple types block unrestricted self-application.  If \(x:\alpha\), then
\(xx\) would require simultaneously \(x:\alpha\to\beta\) and \(x:\alpha\),
which the simple grammar cannot satisfy.  This syntactic fact blocks the form
of self-membership used in Russell's paradox.  It does not prove that every
extension of type theory is consistent; consistency remains a metatheorem
about a specified calculus.

Extensionality and the Axiom of Choice at higher types are separate principles.  A system
may take functional extensionality as an axiom,
\[
  \forall f,g{:}(\alpha\to\beta).\,
  \bigl(\forall x{:}\alpha.\,fx=gx\bigr)\to f=g,
\]
while another builds it into equality semantics.  A higher-order Axiom of Choice
operator introduces additional axioms.  Neither principle follows merely from
the type grammar.

\begin{warning}[Proof-assistant metalogics]
A proof assistant based on a dependent type theory, a simple type theory, or a
logical framework does not thereby prove a theorem about an encoded object
logic.  An encoding requires a theorem connecting encoded derivability or
satisfaction to the intended external judgment.
\end{warning}

\begin{exercise}
\label{ex:simple-type-theory-01}
Assign types to every subterm of
\(\lambda R{:}(\iota\to\iota\to o).\lambda x{:}\iota.Rxx\).
\end{exercise}

\begin{exercise}
\label{ex:simple-type-theory-02}
State the semantic clause for \(\forall f{:}(\alpha\to\beta).\,A\) under
standard semantics and under Henkin semantics.  Identify the only difference.
\end{exercise}

\begin{exercise}
\label{ex:simple-type-theory-03}
Explain why the assertion ``higher-order logic is complete'' is incomplete.
Give one precise true statement and one precise false statement obtained by
supplying a calculus and a semantic class.
\end{exercise}

\part{Nonclassical Logic}
\label{part:nonclassical}

\chapter{Intuitionistic Logic}
\label{ch:intuitionistic-logic}

\begin{convention}[Nonclassical formula metavariables]
Throughout this part, \(A,B,C\) range over formulas and \(p,q\) range over
propositional variables.  The symbols \(\Gamma,\Delta,\Sigma,\Pi\) retain their
normative standing as formula contexts.  Any narrower local role is declared
before use.
\end{convention}

Intuitionistic logic is specified positively by its calculus and semantics.  It
is not defined as classical logic with selected theorems deleted.  This chapter
uses \(\LJ\) from \cref{ch:lk-lj} as the calculus and Kripke
structures for its semantics; see
\cite{Kripke1965,TroelstraSchwichtenberg2000}.

\section{Language, calculus, and Kripke structures}
\label{sec:intuitionistic-logic-01}

The propositional language contains
\(\land,\lor,\to,\bottom\), and
\[
  \neg A\defeq A\to\bottom.
\]
Because \(\LJ\)-antecedents are finite multisets whereas premise collections
may be arbitrary sets, derivability is defined by finite support:
\[
  \Gamma\derives_{\IPL}A
  \quad:\Longleftrightarrow\quad
  \begin{gathered}
    \text{there is a finite multiset context \(\Gamma_0\), every formula}\\
    \text{occurring in which belongs to \(\Gamma\), such that}
    \quad \LJ\derives\Gamma_0\Rightarrow A.
  \end{gathered}
\]
This defines the left side for every formula set \(\Gamma\) while using only
contexts admitted by \(\LJ\).  The unlifted sequent notation remains available
for a particular finite multiset \(\Gamma_0\).

\begin{definition}[Intuitionistic Kripke frame and structure]
An intuitionistic Kripke frame is a partially ordered set \((W,\leq)\).  A
Kripke structure is \(\mathcal M=(W,\leq,V)\), where
\(V:\Prop\to\Pow(W)\) is persistent:
\[
  w\in V(p)\text{ and }w\leq v
  \quad\Longrightarrow\quad
  v\in V(p).
\]
The order represents growth of information.  This interpretation motivates
the semantics but is not part of its formal definition.
\end{definition}

The forcing relation \(\mathcal M,w\forces A\) is defined recursively:
\[
\begin{aligned}
  \mathcal M,w\forces p
    &\quad\Longleftrightarrow\quad w\in V(p),\\
  \mathcal M,w\not\forces\bottom,
    &\\
  \mathcal M,w\forces A\land B
    &\quad\Longleftrightarrow\quad
      \mathcal M,w\forces A\text{ and }\mathcal M,w\forces B,\\
  \mathcal M,w\forces A\lor B
    &\quad\Longleftrightarrow\quad
      \mathcal M,w\forces A\text{ or }\mathcal M,w\forces B,\\
  \mathcal M,w\forces A\to B
    &\quad\Longleftrightarrow\quad
      \forall v\geq w\,
      (\mathcal M,v\forces A\to\mathcal M,v\forces B).
\end{aligned}
\]
The arrow inside the final parenthesis is metalinguistic implication.

\begin{lemma}[Persistence]
For every formula \(A\), if \(w\leq v\) and \(\mathcal M,w\forces A\), then
\(\mathcal M,v\forces A\).
\end{lemma}

\begin{proof}
Proceed by structural induction on \(A\).  The atomic case is persistence of
\(V\); the conjunction and disjunction cases use the induction hypotheses.
For implication, suppose \(w\forces A\to B\) and \(w\leq v\).  If
\(v\leq u\) and \(u\forces A\), then \(w\leq u\), so \(u\forces B\).
Thus \(v\forces A\to B\).  The bottom case is vacuous.
\end{proof}

\begin{definition}[Intuitionistic semantic consequence]
\[
  \Gamma\entails_{\IPL}A
\]
iff, for every intuitionistic Kripke structure \(\mathcal M\) and every
world \(w\), forcing all members of \(\Gamma\) at \(w\) implies forcing \(A\)
at \(w\).
\end{definition}

\begin{theorem}[Soundness and strong completeness of \(\IPL\)]
For every formula set \(\Gamma\) and formula \(A\),
\[
  \Gamma\derives_{\IPL}A
  \quad\Longleftrightarrow\quad
  \Gamma\entails_{\IPL}A.
\]
Moreover, if \(\Gamma\notderives_{\IPL}A\) for finite \(\Gamma\), there is a
finite rooted Kripke structure with a world forcing every member of
\(\Gamma\) and not forcing \(A\).
\end{theorem}

\paragraph{Proof sketch.}
Soundness is induction on an \(\LJ\)-derivation, using persistence in the
implication rules.  For completeness, extend a nonderiving premise set to
prime theories---sets \(T\) that are consistent, are deductively closed
under \(\derives_{\IPL}\), and contain, for every disjunction \(A\lor B\)
with \(T\derives_{\IPL}A\lor B\), the disjunct \(A\) or the disjunct
\(B\)---and order them by inclusion.  Deductive closure makes
\(T\derives_{\IPL}A\lor B\) equivalent to \(A\lor B\in T\); the defining
condition is derivability from \(T\), not derivability from the empty set.
Force a propositional variable at a prime theory
iff the variable belongs to it.  The truth lemma
\(w\forces B\Longleftrightarrow B\in w\) is proved by structural induction;
the implication case uses an extension lemma for failed implications.  A
finite filtration---a quotient identifying worlds that agree on the
finitely many relevant subformulas---through the subformulas of
\(\Gamma\cup\set A\) yields the
stated finite rooted Kripke structure and world.  The prime-extension lemma,
failed-implication extension lemma, and filtration verification are omitted
and are supplied by the controlling sources.

\begin{example}[Counterexample structure and world]
Let \(W=\set{w_0,w_1}\) with \(w_0<w_1\), and let
\(V(p)=\set{w_1}\).  At \(w_0\), \(p\) is not forced.  The formula
\(\neg p\) is also not forced, because the extension \(w_1\) forces \(p\) and
does not force \(\bottom\).  Therefore
\[
  \mathcal M,w_0\not\forces p\lor\neg p.
\]
No world above \(w_0\) forces \(\neg p\), so
\(w_0\forces\neg\neg p\).  Since \(w_0\not\forces p\), the same structure
refutes \(\neg\neg p\to p\).
\end{example}

\begin{example}[A constructive derivation]
The formula \(A\land B\to B\land A\) is derivable in \(\IPL\).  In
\(\LJ\), from the initial sequents \(A,B\Rightarrow B\) and
\(A,B\Rightarrow A\), apply \(\land R\) to obtain
\(A,B\Rightarrow B\land A\), then \(\land L\) to obtain
\(A\land B\Rightarrow B\land A\), and \(\to R\) to conclude
\(\Rightarrow A\land B\to B\land A\).  No classical principle is used.
\end{example}

\begin{warning}[Disjunction at a world]
The clause for \(A\lor B\) requires one disjunct to be forced at the current
world.  It is not enough that every future world eventually forces some
possibly varying disjunct.
\end{warning}

\begin{exercise}
\label{ex:intuitionistic-logic-01}
Give a Kripke structure and a world in it that falsify Peirce's formula
\(((A\to B)\to A)\to A\).
\end{exercise}

\begin{exercise}
\label{ex:intuitionistic-logic-02}
Prove semantically that
\(\neg(A\lor B)\entails_{\IPL}\neg A\land\neg B\).  Then determine whether the
converse holds.
\end{exercise}

\section{Relations to classical derivability}
\label{sec:intuitionistic-classical}

\begin{theorem}[Glivenko]
\label{thm:glivenko}
For every propositional formula \(A\) and finite multiset \(\Gamma\) of
propositional formulas, with both calculi restricted to their propositional
rules:
\(\LK\derives\Gamma\Rightarrow A\) iff
\(\LJ\derives\Gamma\Rightarrow\neg\neg A\).
\textit{Status: cited theorem.}
\end{theorem}

See \cite{TroelstraVanDalen1988,TroelstraSchwichtenberg2000}.  Glivenko's
theorem fails for the quantifier rules; the negative translation repairs
the failure symbol by symbol.

\begin{theorem}[G\"odel--Gentzen negative translation]
\label{thm:negative-translation}
Assign to every first-order formula \(A\) a formula \(A^N\) by recursion:
\(\bottom^N\defeq\bottom\), and this clause takes priority in the sense
that \(\bottom\) is formed by its own grammar clause, is not an atomic
formula (\cref{sec:fol-syntax-semantics-02}), and is not treated by the next
clause; \(A^N\defeq\neg\neg A\) for atomic \(A\); the translation commutes
with \(\land\), \(\to\), and
\(\forall\); \((A\lor B)^N\defeq\neg(\neg A^N\land\neg B^N)\); and
\((\exists x\,A)^N\defeq\neg\forall x\,\neg A^N\).  Here \(\neg\) is the
abbreviation \(\neg A\defeq A\to\bottom\) of \cref{sec:lk-lj-02}, so a
negated formula is treated by the \(\to\) clause; by unique readability
(\cref{thm:fol-unique-readability}), exactly one clause applies to each
formula, and the recursion is well defined.  Then
\(\LK\derives\Gamma\Rightarrow A\) iff
\(\LJ\derives\Gamma^N\Rightarrow A^N\), where \(\Gamma^N\) translates
every member.
\textit{Status: cited theorem.}
\end{theorem}

See \cite{TroelstraVanDalen1988}.

\begin{theorem}[Disjunction property]
\label{thm:disjunction-property}
If \(\LJ\derives\;\Rightarrow A\lor B\), then
\(\LJ\derives\;\Rightarrow A\) or \(\LJ\derives\;\Rightarrow B\).
\textit{Status: complete proof.}
\end{theorem}

\begin{proof}
By cut elimination for \(\LJ\) (\cref{ch:cut-elimination}), fix a
cut-free \(\LJ\)-derivation of \(\Rightarrow A\lor B\).  Its final
inference is not \textsc{Id} or \(\bottom L\), whose conclusions have
nonempty antecedents, and not a left logical or left structural rule, for
the same reason.  The calculus has no right weakening or right contraction.
Among the remaining right rules, only \(\lor R_1\) and \(\lor R_2\)
conclude a sequent whose formula is a disjunction.  Hence the final
inference has premise \(\Rightarrow A\) or \(\Rightarrow B\).
\end{proof}

Classical derivability lacks this property: \(P\lor\neg P\) is
classically derivable while neither disjunct is.

\chapter{Modal, Temporal, and Epistemic Logic}
\label{ch:modal-temporal-epistemic}

\section{Normal propositional modal logic}
\label{sec:modal-temporal-epistemic-01}

Modal formulas add \(\Box A\) to the classical propositional grammar.  Define
\[
  \Diamond A\defeq\neg\Box\neg A.
\]
The symbols have no fixed application until a class of frames and a
satisfaction relation are specified.

\begin{definition}[Kripke frame and Kripke structure]
A modal Kripke frame is a pair \(\mathcal F=(W,R)\), where \(W\ne\varnothing\)
and \(R\subseteq W\times W\).  A modal Kripke structure is
\(\mathcal M=(W,R,V)\), where \(V:\Prop\to\Pow(W)\).  Satisfaction is
classical for propositional connectives and
\[
  \mathcal M,w\entails\Box A
  \quad\Longleftrightarrow\quad
  \forall v\,(wRv\to\mathcal M,v\entails A).
\]
Consequently, \(\mathcal M,w\entails\Diamond A\) iff some
\(R\)-successor of \(w\) satisfies \(A\).
\end{definition}

Relational semantics for normal modal logic originates with
\cite{Kripke1963}.

\begin{definition}[Structure validity and frame validity]
A formula is valid in a Kripke structure when it is satisfied at every world
of that structure.  It is valid on a frame when it is valid in every Kripke
structure based on that frame.  For a class \(\mathcal C\) of frames, write
\(\entails_{\mathcal C}A\) when \(A\) is valid on every frame in
\(\mathcal C\).
\end{definition}

\begin{definition}[Local and global modal consequence]
Local consequence is
\[
  \Gamma\entails^{\ell}_{\mathcal C}A
\]
iff, in every Kripke structure on a frame in \(\mathcal C\) and at every world
\(w\), satisfaction of \(\Gamma\) at \(w\) implies satisfaction of \(A\) at
\(w\).  Global consequence is
\[
  \Gamma\entails^{g}_{\mathcal C}A
\]
iff, in every such structure, satisfaction of every member of \(\Gamma\)
at every world implies satisfaction of \(A\) at every world.  These relations
can differ.
\end{definition}

The Hilbert calculus \(\K\) consists of all uniform substitution
instances of classically valid propositional formulas, modus ponens, the
distribution schema
\[
  \Box(A\to B)\to(\Box A\to\Box B),
  \tag{K}
\]
and necessitation
\[
  \frac{\derives_{\K}A}{\derives_{\K}\Box A}\;\textsc{Nec}.
\]
Necessitation applies to theorems, not to a formula depending on an
undischarged premise.  Define
\[
\begin{array}{rcl}
  \mathsf T&=&\K+(\Box A\to A),\\
  \Sfour&=&\mathsf T+(\Box A\to\Box\Box A),\\
  \Sfive&=&\mathsf T+(\Diamond A\to\Box\Diamond A).
\end{array}
\]
For \(\mathsf L\in\{\K,\mathsf T,\Sfour,\Sfive\}\) and a premise set
\(\Gamma\), an \(\mathsf L\)-derivation from \(\Gamma\) is a finite sequence
whose lines are members of \(\Gamma\), instances of the axiom schemas of
\(\mathsf L\), consequences of earlier lines by modus ponens, or formulas
\(\Box A\) for which \(A\) has a separate no-premise
\(\mathsf L\)-derivation.  Write \(\Gamma\derives_{\mathsf L}A\) when such a
sequence ends in \(A\).  Thus necessitation remains theorem-only even in a
derivation from premises.  This definition supplies the premise-sensitive
standing needed by strong completeness.

\begin{theorem}[Frame correspondence and completeness]
In each row, let \(\mathsf L\) be the displayed calculus and let
\(\mathcal C_{\mathsf L}\) be the displayed class of frames:
\[
\begin{array}{c|c}
  \text{calculus}&\text{frame class}\\ \hline
  \K&\text{all frames}\\
  \mathsf T&\text{reflexive frames}\\
  \Sfour&\text{reflexive transitive frames}\\
  \Sfive&\text{frames whose accessibility relation is an equivalence relation.}
\end{array}
\]
Then, for every premise set \(\Gamma\) and formula \(A\),
\[
  \Gamma\derives_{\mathsf L}A
  \quad\Longleftrightarrow\quad
  \Gamma\entails^{\ell}_{\mathcal C_{\mathsf L}}A.
\]
Thus each calculus is sound and strongly complete for local consequence on its
paired frame class.
\end{theorem}

\paragraph{Proof sketch.}
Soundness checks each axiom against its frame condition and verifies modus
ponens and necessitation.  Completeness constructs the canonical structure of
maximally consistent sets.  Define \(wRv\) when every \(\Box A\in w\) has
\(A\in v\), and prove the truth lemma by induction.  The axioms T, 4, and 5
force the canonical relation to be reflexive, transitive, and Euclidean;
reflexivity plus Euclideanness yields an equivalence relation.  See
\cite{BlackburnDeRijkeVenema2001} for the detailed canonical-structure
arguments.

\begin{example}[Three counterexample structures and worlds]
\begin{enumerate}
  \item Let \(W=\set w\), let \(R=\varnothing\), and make \(p\) false at
        \(w\) in the resulting structure \(\mathcal M\).  Then
        \(\mathcal M,w\entails\Box p\) and
        \(\mathcal M,w\notentails p\), so the structure and world falsify
        \(\Box p\to p\).  This refutes T on a nonreflexive frame.
  \item Let \(wRv\) and \(vRu\), but not \(wRu\).  Make \(p\) true at \(v\)
        and false at \(u\).  If \(v\) is the only successor of \(w\), then
        the resulting structure \(\mathcal M\) has
        \(\mathcal M,w\entails\Box p\) but
        \(\mathcal M,w\notentails\Box\Box p\).  Thus the structure and world
        falsify \(\Box p\to\Box\Box p\), refuting 4 on a nontransitive frame.
  \item Let \(wRv\) and \(wRu\), let \(v\) have no successors, and make \(p\)
        true only at \(u\) in the resulting structure \(\mathcal M\).  Then
        \(\mathcal M,w\entails\Diamond p\), while
        \(\mathcal M,v\notentails\Diamond p\); hence
        \(\mathcal M,w\notentails\Box\Diamond p\).  Thus the structure and
        world \(w\) falsify
        \(\Diamond p\to\Box\Diamond p\), refuting 5 on a non-Euclidean frame.
\end{enumerate}
\end{example}

\section{Linear temporal logic}
\label{sec:modal-temporal-epistemic-02}

Linear temporal logic over discrete infinite time adds the operators
\(\mathsf X\) (next) and \(\mathbin{\mathsf U}\) (until):
\[
  A::=p\mid\bottom\mid\neg A\mid(A\land A)\mid\mathsf X A\mid
       (A\mathbin{\mathsf U}A).
\]
Because this grammar retains only \(\neg\) and \(\land\) as primitive Boolean
connectives, the following local eliminable abbreviations are necessary for the
calculus below:
\[
\begin{aligned}
  A\lor B&\defeq\neg(\neg A\land\neg B),
  & A\to B&\defeq\neg(A\land\neg B),\\
  A\leftrightarrow B&\defeq(A\to B)\land(B\to A),
  & \topv&\defeq\neg\bottom.
\end{aligned}
\]
Their standing is metalinguistic abbreviation, not additional temporal
constructors.  Each expands to a formula of the displayed grammar, differs
from a primitive constructor precisely by being eliminable, and is used in the
axiom schemas, rules, and definitions of \(\mathsf F\) and \(\mathsf G\) below.
A temporal structure is an infinite word
\(\mathcal M=(\N,V)\), with
\(V:\N\to\Pow(\Prop)\).  Atomic and Boolean satisfaction is defined
pointwise by
\[
\begin{aligned}
  \mathcal M,n\entails p
    &\quad\Longleftrightarrow\quad p\in V(n),\\
  \mathcal M,n&\notentails\bottom,\\
  \mathcal M,n\entails\neg A
    &\quad\Longleftrightarrow\quad\mathcal M,n\notentails A,\\
  \mathcal M,n\entails A\land B
    &\quad\Longleftrightarrow\quad
      (\mathcal M,n\entails A\text{ and }\mathcal M,n\entails B).
\end{aligned}
\]
The temporal clauses are
\[
  \mathcal M,n\entails\mathsf X A
  \quad\Longleftrightarrow\quad
  \mathcal M,n+1\entails A,
\]
and
\[
\begin{aligned}
  \mathcal M,n\entails A\mathbin{\mathsf U}B
  \quad\Longleftrightarrow\quad
  \exists m\geq n\,[&\mathcal M,m\entails B\ \land{}\\
   &\forall k\,(n\leq k<m\to\mathcal M,k\entails A)].
\end{aligned}
\]
Define
\[
  \mathsf F A\defeq\topv\mathbin{\mathsf U}A,
  \qquad
  \mathsf G A\defeq\neg\mathsf F\neg A.
\]
The temporal validity condition used below is
\[
  \mathcal M,n\entails A
  \quad\text{for every temporal structure \(\mathcal M\) and every
  \(n\in\N\)}.
\]
It is a property of one formula, not an undefined premise-sensitive consequence
relation.

Fix the following Hilbert calculus for theoremhood in this language.  A
derivation is a finite sequence of formulas in which each line is a uniform
substitution instance of a classically valid propositional formula, an
instance of
one of the three schemas
\[
\begin{aligned}
  \mathsf X(A\to B)&\to(\mathsf XA\to\mathsf XB),\\
  \mathsf X\neg A&\leftrightarrow\neg\mathsf XA,\\
  A\mathbin{\mathsf U}B
    &\leftrightarrow
      B\lor(A\land\mathsf X(A\mathbin{\mathsf U}B)),
\end{aligned}
\]
or follows from earlier lines by one of the rules
\[
  \frac{A\qquad A\to B}{B}\;\textsc{MP},
  \qquad
  \frac{A}{\mathsf XA}\;\textsc{XNec},
  \qquad
  \frac{B\lor(A\land\mathsf X C)\to C}
       {A\mathbin{\mathsf U}B\to C}\;\textsc{UInd}.
\]
All metavariables in these schemas and rules range over temporal formulas.
Write \(\derives_{\mathsf{LTL}}A\) when such a no-premise derivation ends in
\(A\).  This is a syntactic theoremhood relation, distinct from temporal
validity, and it is introduced because the next theorem compares those two
standings.  Since derivations have no open premises, both necessitation and the
induction rule apply only to theorems.

\begin{theorem}[Soundness and completeness of linear temporal logic]
For every temporal formula \(A\),
\[
  \derives_{\mathsf{LTL}}A
  \quad\Longleftrightarrow\quad
  \mathcal M,n\entails A
  \text{ for every temporal structure \(\mathcal M\) and every \(n\in\N\)}.
\]
This is weak completeness for validity, not a strong-completeness assertion
for arbitrary premise sets.  It concerns discrete infinite linear time;
finite traces, dense time, and branching time require different semantics and
calculi.
\end{theorem}

\paragraph{Proof sketch.}
Soundness of the first two temporal schemas and of
\textsc{XNec} follows directly from the unique successor.  The unfolding
schema follows by separating the cases where the witness for until is the
current position or a later position.  For \textsc{UInd}, assume its premise is
valid and reason backward from the least witness for \(B\) to show \(C\) at
each preceding position at which \(A\) holds.  Completeness closes a formula
under subformulas, negation, and the temporal obligations generated by until;
it constructs maximally consistent closure sets linked by the \(\mathsf X\)
requirements, proves that every pending until obligation is eventually
realized, and unravels the resulting path into a temporal structure.  The
truth
lemma then yields a falsifying temporal structure for every underivable
formula.  The consistency, eventuality-realization, and truth-lemma case
calculations are omitted---this is why the theorem is presented with a
proof sketch rather than a complete proof---and are controlled by
\cite{GabbayPnueliShelahStavi1980}.

\begin{example}[Eventuality is not recurrence]
Let \(p\) hold only at time \(1\).  At time \(0\), \(\mathsf Fp\) is true, but
\(\mathsf G\mathsf Fp\) is false because no occurrence of \(p\) remains after
time \(1\).  Thus \(\mathsf Fp\to\mathsf G\mathsf Fp\) is not valid:
the displayed temporal structure and position \(0\) falsify it.
\end{example}

\section{Epistemic logic}
\label{sec:modal-temporal-epistemic-03}

Fix a nonempty set \(I\) of agents.  The language adds \(K_iA\) for each
\(i\in I\).  An epistemic frame is
\[
  (W,(\sim_i)_{i\in I}),
\]
where \(W\ne\varnothing\) and each
\(\sim_i\subseteq W\times W\).  An epistemic structure is
\(\mathcal M=(W,(\sim_i)_{i\in I},V)\), where
\(V:\Prop\to\Pow(W)\), with the classical connective clauses and
\[
  \mathcal M,w\entails K_iA
  \quad\Longleftrightarrow\quad
  \forall v\,(w\sim_i v\to\mathcal M,v\entails A).
\]
For knowledge, each \(\sim_i\) is normally required to be an equivalence
relation.  The multi-agent calculus used below has all uniform substitution
instances of classically valid propositional formulas and, for each
\(i\in I\), the
schemas
\[
\begin{aligned}
  K_i(A\to B)&\to(K_iA\to K_iB),&
  K_iA&\to A,\\
  K_iA&\to K_iK_iA,&
  \neg K_iA&\to K_i\neg K_iA.
\end{aligned}
\]
For a premise set \(\Gamma\), a derivation is a finite sequence in which each
line is a member of \(\Gamma\), an instance of one of the displayed schemas,
follows by modus ponens, or has the form \(K_iA\) where \(A\) has a separate
no-premise derivation.  Write \(\Gamma\derives_{\Sfive_I}A\) when such a
derivation ends in \(A\).  This theorem-only restriction on necessitation is
part of the calculus.  The subscript \(I\) records that the construction is the
componentwise multi-agent extension of the previously defined \(\Sfive\)
calculus; it is needed to distinguish the agent-indexed language and adds no
interaction axioms between distinct agents.

\begin{theorem}[Multi-agent \(\Sfive\)]
For every premise set \(\Gamma\) and formula \(A\),
\(\Gamma\derives_{\Sfive_I}A\) iff \(A\) is a local consequence of
\(\Gamma\) over epistemic frames in which every \(\sim_i\) is an equivalence
relation.
\end{theorem}

\paragraph{Proof sketch.}
Soundness checks each agent-indexed axiom against
reflexivity, transitivity, and Euclideanness of \(\sim_i\); necessitation is
theorem-only.  For completeness, take maximally consistent sets as worlds and
define \(w\sim_i v\) iff every \(K_iA\in w\) has \(A\in v\).  The
truth lemma is a structural induction.  The three non-distribution schemas
make each canonical relation reflexive, transitive, and Euclidean, hence an
equivalence relation; no step relates \(\sim_i\) to \(\sim_j\) for
\(i\ne j\).  Lindenbaum extension, the modal existence lemma, and the full
truth-lemma calculation are omitted; the one-agent instances are controlled by
\cite[Chapter 4]{BlackburnDeRijkeVenema2001}.

For a finite group \(G\subseteq I\), the language further adds the
operators \(E_G\) and \(C_G\), with \(E_GA\) and \(C_GA\) formulas
whenever \(A\) is a formula.  Everyone-knows is definable by
\[
  E_GA\defeq\bigwedge_{i\in G}K_iA.
\]
Common knowledge \(C_GA\) receives a satisfaction clause of its own:
\(C_GA\) is true at \(w\) iff \(A\) is true at every world reachable
from \(w\) by a path of length \(n\geq0\) whose edges belong to
\(\bigcup_{i\in G}\sim_i\); the length-zero path makes \(w\) itself
reachable, so the clause is reflexive.  Common knowledge is therefore a
reflexive-transitive closure modality; it is not a finite iteration of
\(E_G\) bounded in advance.

\begin{example}[Truth without knowledge]
Let \(w\sim_i v\), let \(p\) be true at \(w\), and let \(p\) be false at
\(v\).  Then \(p\) is true at \(w\) but \(K_ip\) is false there.  The structure
does not refute factivity \(K_ip\to p\): equivalence makes \(w\sim_i w\), so
knowledge still implies truth at the actual world.
\end{example}

\begin{warning}[Changing the accessibility condition]
If an epistemic relation is not reflexive, \(K_iA\to A\) can fail.  Such a
modality may represent belief or information rather than factive knowledge, but
the change must be made in both terminology and frame class.
\end{warning}

\begin{exercise}
\label{ex:modal-temporal-epistemic-01}
Give a two-world Kripke structure and a world that falsify
\(\Diamond p\to\Box\Diamond p\) on a reflexive but non-Euclidean frame.
\end{exercise}

\begin{exercise}
\label{ex:modal-temporal-epistemic-02}
Evaluate \(p\mathbin{\mathsf U}q\) at every position of a five-step prefix in
which \(p\) holds at positions \(0,1,2\) and \(q\) first holds at position
\(3\).  State what must be specified beyond the prefix.
\end{exercise}

\begin{exercise}
\label{ex:modal-temporal-epistemic-03}
Give an epistemic structure in which every agent in \(G\) knows \(p\) at \(w\), but
\(C_Gp\) fails at \(w\).
\end{exercise}

\chapter{Many-Valued and Paraconsistent Logic}
\label{ch:many-valued-paraconsistent}

\section{Logical matrices}
\label{sec:many-valued-paraconsistent-01}

\begin{definition}[Logical matrix]
A finite logical matrix for a propositional signature is the structure
\[
  \mathcal M=(\mathsf V,D,(f_c)_{c\in\Lang}),
\]
where \(\mathsf V\) is a finite nonempty set of values,
\(D\subseteq\mathsf V\) is the set of
designated values, and each \(n\)-ary connective \(c\) is interpreted by a
function \(f_c:\mathsf V^n\to\mathsf V\).  A valuation maps propositional
variables to \(\mathsf V\) and extends homomorphically to every formula.
\end{definition}

\begin{definition}[Matrix consequence]
For a formula set \(\Gamma\),
\[
  \Gamma\entails_{\mathcal M}A
\]
iff every valuation assigning a designated value to every member of
\(\Gamma\) also assigns a designated value to \(A\).
\end{definition}

The number of values does not determine a consequence relation.  The
operations and the designated set are equally essential.

\section{A value-labelled tableau calculus}
\label{sec:many-valued-paraconsistent-02}

For a fixed finite matrix \(\mathcal M\), a labelled formula \(A{:}a\) asserts
that \(A\) has exact value \(a\in\mathsf V\).  Expand
\[
  c(A_1,\ldots,A_n){:}a
\]
into one branch for every tuple \(\tuple{a_1,\ldots,a_n}\) satisfying
\(f_c(a_1,\ldots,a_n)=a\), placing
\(A_1{:}a_1,\ldots,A_n{:}a_n\) on that branch.  A branch closes when it
contains \(p{:}a\) and \(p{:}b\) for the same variable \(p\) and distinct
values \(a\ne b\), or when it contains a labelled formula
\(c(A_1,\ldots,A_n){:}a\) for which no tuple satisfies
\(f_c(a_1,\ldots,a_n)=a\), so that its expansion produces no branches.

To test \(\Gamma\entails_{\mathcal M}B\), start one tableau for each assignment
of a designated label to every member of \(\Gamma\) and a nondesignated label
to \(B\).  The tableau succeeds when every branch of every such root closes.

\begin{definition}[Saturated branch]
A branch is \emph{saturated} iff, for every labelled formula
\(c(A_1,\ldots,A_n){:}a\) on it, there is at least one tuple
\(\tuple{a_1,\ldots,a_n}\) with \(f_c(a_1,\ldots,a_n)=a\) such that
\(A_1{:}a_1,\ldots,A_n{:}a_n\) all lie on the branch.
\end{definition}

\begin{theorem}[Finite-matrix tableau soundness and completeness]
For every finite matrix \(\mathcal M\), finite formula set \(\Gamma\), and
formula \(B\), the value-labelled tableau for \(\Gamma\) and \(B\) closes
iff \(\Gamma\entails_{\mathcal M}B\).
\textit{Status: complete proof.}
\end{theorem}

\begin{proof}
Call a valuation \(v\) \emph{compatible} with a branch iff \(v(A)=a\) for
every labelled formula \(A{:}a\) on the branch.  Full expansion terminates:
expanding a labelled occurrence removes it from the unexpanded stock of its
branch and adds only labelled occurrences of its immediate subformulas, so
the sum of the complexities of the unexpanded labelled occurrences on a
branch strictly decreases at each step; each expansion produces at most
\(\abs{\mathsf V}^{n}\) successor branches, so the fully expanded tree is
finite.

Suppose first that \(\Gamma\notentails_{\mathcal M}B\).  Fix a valuation
\(v\) assigning a designated value to every member of \(\Gamma\) and a
nondesignated value to \(B\).  Among the started roots is the one labelling
each \(C\in\Gamma\) with \(v(C)\) and \(B\) with \(v(B)\), and \(v\) is
compatible with it.  If \(v\) is compatible with a branch and the labelled
formula \(c(A_1,\ldots,A_n){:}a\) on it is expanded, then
\[
  f_c(v(A_1),\ldots,v(A_n))=v(c(A_1,\ldots,A_n))=a,
\]
so \(\tuple{v(A_1),\ldots,v(A_n)}\) is among the expansion tuples---in
particular the expansion is not empty---and the successor branch it produces
is again compatible with \(v\).  Induction on the number of expansion steps
therefore keeps one branch compatible with \(v\) at every stage of this
root.  A branch compatible with \(v\) satisfies neither closure condition:
labels \(p{:}a\) and \(p{:}b\) with \(a\ne b\) would give \(a=v(p)=b\), and
the displayed calculation supplies a tuple for every labelled compound
formula.  Hence this root never closes, and the tableau does not close.

Conversely, suppose the tableau does not close: some root, fully expanded,
retains a branch that does not close.  That branch is saturated: every
labelled compound formula on it has been expanded, the second closure
condition excludes an empty expansion, and the branch passes through one of
the produced successor branches, which placed the components of a witnessing
tuple on it.  Define a valuation \(v\) as follows: a variable \(p\) labelled
on the branch carries exactly one label \(p{:}a\), because the first closure
condition fails, and \(v(p)\defeq a\); every variable not labelled on the
branch receives one fixed value of the nonempty set \(\mathsf V\).  By
induction on formula complexity, \(v(A)=a\) for every labelled formula
\(A{:}a\) on the branch.  If \(A\) is a variable, this is the definition of
\(v\).  If \(A=c(A_1,\ldots,A_n)\), saturation supplies a tuple
\(\tuple{a_1,\ldots,a_n}\) with \(f_c(a_1,\ldots,a_n)=a\) and every
\(A_i{:}a_i\) on the branch; the induction hypothesis for the immediate
subformulas \(A_i\) gives \(v(A_i)=a_i\), and the homomorphic extension of
\(v\) gives
\[
  v(A)=f_c(v(A_1),\ldots,v(A_n))=f_c(a_1,\ldots,a_n)=a.
\]
Applied to the root labels, \(v\) assigns every member of \(\Gamma\) its
designated root label and \(B\) its nondesignated root label, so
\(\Gamma\notentails_{\mathcal M}B\).
\end{proof}

\section{Kleene's three-valued matrix}
\label{sec:many-valued-paraconsistent-03}

Let
\[
  \mathsf V=\set{0,\tfrac12,1},
  \qquad D=\set1,
\]
and define
\[
  \neg a=1-a,
  \qquad
  a\land b=\min(a,b),
  \qquad
  a\lor b=\max(a,b).
\]
This matrix determines the strong Kleene consequence relation for the
displayed language; the matrix is denoted \(\mathbf K_3\), and
\(\entails_{\mathbf K_3}\) is its matrix consequence relation.

\begin{example}[Excluded middle is not designated]
If \(v(p)=\tfrac12\), then
\[
  v(p\lor\neg p)
  =\max(\tfrac12,\tfrac12)
  =\tfrac12\notin D.
\]
Therefore \(\notentails_{\mathbf K_3}p\lor\neg p\).  The calculation is a
countervaluation, not a derivation.
\end{example}

\section{The Logic of Paradox}
\label{sec:many-valued-paraconsistent-04}

Use the same value set and operations, but designate
\[
  D_{\mathsf{LP}}=\set{\tfrac12,1}.
\]
Interpret \(\tfrac12\) as carrying both positive and negative support.  The
resulting matrix defines the Logic of Paradox \(\mathsf{LP}\).  Its
value-labelled tableau is sound and complete by the finite-matrix theorem.

\begin{definition}[Paraconsistent consequence relation]
A consequence relation is paraconsistent iff there are formulas
\(A,B\) such that
\[
  A,\neg A\notentails B.
\]
This definition denies unrestricted explosion.  It does not assert that any
particular contradiction is derivable or true.
\end{definition}

\begin{example}[Counterexample valuation to explosion in \(\mathsf{LP}\)]
Set
\[
  v(p)=\tfrac12,
  \qquad
  v(q)=0.
\]
Both \(p\) and \(\neg p\) receive the designated value \(\tfrac12\), while
\(q\) is nondesignated.  Hence
\[
  p,\neg p\notentails_{\mathsf{LP}}q.
\]
At the same valuation, \(p\lor\neg p\) is designated.  Paraconsistency and the
failure of excluded middle are therefore independent properties; compare the
two choices of designated values on the same algebra.
\end{example}

\begin{warning}[Truth values are formal data]
Labels such as ``both'' or ``indeterminate'' are interpretations of matrix
values.  The matrix consequence relation is determined by the truth-value set,
operations, and designated set.  An informal label cannot replace those data.
\end{warning}

\begin{exercise}
\label{ex:many-valued-paraconsistent-01}
Construct the complete value-labelled tableau for
\(p\land q\entails_{\mathbf K_3}p\).  Explain why the same tableau shape is
sound for \(\mathsf{LP}\).
\end{exercise}

\begin{exercise}
\label{ex:many-valued-paraconsistent-02}
Find a formula that is valid in \(\mathsf{LP}\) but not valid in
\(\mathbf K_3\), and verify both claims by the three possible values of its
propositional variable.
\end{exercise}

\chapter[Relevance and Linear Logic]{Relevance, Substructural Systems, and Linear Logic}
\label{ch:relevance-linear}

\section{Structural rules as logical data}
\label{sec:relevance-linear-01}

In a sequent calculus, structural rules determine how formula occurrences may
be used:
\[
\frac{\Gamma\Rightarrow B}{\Gamma,A\Rightarrow B}\;\textsc{Weakening},
\qquad
\frac{\Gamma,A,A\Rightarrow B}{\Gamma,A\Rightarrow B}\;\textsc{Contraction},
\]
and exchange permutes adjacent occurrences.  A substructural logic omits or
controls at least one structural rule.  The context representation must
therefore remain visible: sets build exchange and contraction into syntax,
multisets build in exchange only, and sequences build in neither.

\section{Routley--Meyer semantics for relevance}
\label{sec:relevance-linear-02}

The formula language used in this section is exactly
\[
  A::=p\mid\neg A\mid(A\land A)\mid(A\lor A)\mid(A\to A),
  \qquad p\in\Prop.
\]
No consequence relation or proof calculus is determined by this formation
grammar alone.

\begin{definition}[Routley--Meyer frame and structure]
A Routley--Meyer frame is a tuple
\[
  \mathcal F=(W,N,R,{*}),
\]
where \(W\ne\varnothing\), \(N\subseteq W\) is a nonempty set of normal worlds,
\(R\subseteq W^3\), and \({*}:W\to W\) is an involution.  The base normality
condition used here is
\[
  n\in N
  \quad\Longrightarrow\quad
  (R(n,u,v)\Longleftrightarrow u=v).
\]
A Routley--Meyer structure is
\(\mathcal M=(W,N,R,{*},V)\), where the frame components satisfy the displayed
conditions and \(V:\Prop\to\Pow(W)\) is a valuation.
\end{definition}

Atomic satisfaction is
\[
  \mathcal M,w\entails p
  \quad\Longleftrightarrow\quad
  w\in V(p).
\]
Satisfaction for conjunction and disjunction is local and classical.  The
relevance-sensitive clauses are
\[
  \mathcal M,w\entails\neg A
  \quad\Longleftrightarrow\quad
  \mathcal M,w^*\notentails A,
\]
and
\[
  \mathcal M,w\entails A\to B
  \quad\Longleftrightarrow\quad
  \forall u,v\,
  \bigl((R(w,u,v)\land\mathcal M,u\entails A)
        \to\mathcal M,v\entails B\bigr).
\]
Validity requires satisfaction at every normal world.  Different relevance
logics impose additional first-order conditions on \(R\), \({*}\), and \(N\),
and use matching Hilbert, display, or labelled calculi.  No such calculus is
selected by the semantic clauses alone.

\paragraph{Worked semantic calculation.}
At a normal world \(n\in N\), the displayed base condition reduces implication
to global inclusion of truth sets:
\[
\begin{aligned}
  \mathcal M,n\entails A\to B
  &\Longleftrightarrow
    \forall u,v\,
      ((u=v\land\mathcal M,u\entails A)
        \Longrightarrow\mathcal M,v\entails B)\\
  &\Longleftrightarrow
    \forall u\in W\,
      (\mathcal M,u\entails A
       \Longrightarrow\mathcal M,u\entails B).
\end{aligned}
\]
Thus implication at \(n\) is not evaluated merely from the truth values of
\(A\) and \(B\) at \(n\).  At a nonnormal world, the ternary relation can
instead test an \(A\)-world and a possibly different \(B\)-world.  The
four-world countermodel below exploits exactly that difference.

For negation, the star operation supplies the comparison point.  Let
\(W=\set{w,w^*}\), let \(N=\set w\), let star swap the two worlds, put
\[
  R=\set{\tuple{w,x,x}:x\in W},
  \qquad
  V(p)=\set w,
\]
and let \(\mathcal M=(W,N,R,{*},V)\), with the values of other
propositional variables arbitrary.
Then \(\mathcal M,w\entails\neg p\) holds while
\(\mathcal M,w^*\notentails\neg p\).  This is a direct calculation from the
semantic clause, not an assertion that every relevance logic validates the
same additional negation principles.

\paragraph{Cited theorem.}
Routley and Meyer \cite[pp.~53--73]{RoutleyMeyer1972} pair the entailment
systems specified in that source with ternary relational semantics and prove
the corresponding source-scoped semantic results.  The controlling hypotheses
are the selected source calculus and its matching postulates on the ternary
relation, normal points, and star operation; changing either side changes the
result.  This citation supports the displayed semantic architecture and the
need for a matched pair.  It is not a completeness theorem for the deliberately
broad base frames defined in this chapter.

\paragraph{Internal proof boundary.}
A local completeness theorem would require selecting one existing system and
reproducing its exact formula language; every axiom and inference rule; the
context, label, and structural policies; all side conditions; finite derivation
objects; and the precise corresponding frame postulates.  A complete proof
would then need rule-by-rule soundness and a canonical-model or
saturated-branch construction with an existence lemma and truth lemma for
every constructor.  Those source-dependent data are not supplied.  The chapter
therefore uses only the explicitly displayed satisfaction clauses and
normal-world validity, and introduces no calculus name or
derivability--validity biconditional.

\begin{example}[Failure of irrelevant weakening]
Let \(W=\set{n,a,u,v}\) consist of four distinct worlds, let \(N=\set n\),
let \({*}\) be the identity, and put
\[
  R=\set{\tuple{n,x,x}:x\in W}\cup\set{\tuple{a,u,v}}.
\]
Let \(V(p)=\set a\), \(V(q)=\set u\), and let
\(\mathcal M=(W,N,R,{*},V)\).  Then \(q\to p\) is false at \(a\), since
\(R(a,u,v)\), \(q\) is true at \(u\), and \(p\) is false at \(v\).  Since
\(R(n,a,a)\) and \(p\) is true at \(a\),
\[
  \mathcal M,n\notentails p\to(q\to p).
\]
Thus the classical weakening principle need not be relevantly valid in this
structure.
\end{example}

The example establishes failure in the displayed base structure.  It does not
by itself classify the formula in every system called a relevance logic.

\section{Intuitionistic multiplicative linear logic}
\label{sec:relevance-linear-03}

The formulas of the multiplicative fragment are
\[
  A::=p\mid\mathbf 1\mid A\otimes A\mid A\multimap A.
\]
Sequents have one succedent and multiset antecedents.  The calculus contains
identity and cut together with
\[
\frac{\Gamma,A,B\Rightarrow C}
     {\Gamma,A\otimes B\Rightarrow C}\;\textsc{\(\otimes L\)},
\qquad
\frac{\Gamma\Rightarrow A\qquad\Delta\Rightarrow B}
     {\Gamma,\Delta\Rightarrow A\otimes B}\;\textsc{\(\otimes R\)},
\]
\[
\frac{\Gamma,A\Rightarrow B}
     {\Gamma\Rightarrow A\multimap B}\;\textsc{\(\multimap R\)},
\qquad
\frac{\Gamma\Rightarrow A\qquad\Delta,B\Rightarrow C}
     {\Gamma,\Delta,A\multimap B\Rightarrow C}\;\textsc{\(\multimap L\)},
\]
and
\[
\frac{}{\Rightarrow\mathbf1}\;\textsc{\(\mathbf1R\)},
\qquad
\frac{\Gamma\Rightarrow C}
     {\Gamma,\mathbf1\Rightarrow C}\;\textsc{\(\mathbf1L\)}.
\]
Weakening and contraction are absent.  Exchange is present because contexts
are multisets.

\begin{definition}[Commutative residuated partially ordered monoid]
A commutative residuated partially ordered monoid is
\[
  \mathcal A=(A,\leq,\otimes,1,\multimap)
\]
such that \(\leq\) is a partial order on \(A\), \((A,\otimes,1)\) is a
commutative monoid, \(\otimes\) is monotone in each argument with respect
to \(\leq\), and
\[
  a\otimes b\leq c
  \quad\Longleftrightarrow\quad
  b\leq a\multimap c.
\]
\end{definition}

A valuation extends by
\[
  \sem{A\otimes B}=\sem A\otimes\sem B,
  \qquad
  \sem{A\multimap B}=\sem A\multimap\sem B,
  \qquad
  \sem{\mathbf1}=1.
\]
The sequent \(A_1,\ldots,A_n\Rightarrow B\) is valid when
\[
  \sem{A_1}\otimes\cdots\otimes\sem{A_n}\leq\sem B
\]
under every valuation in every such monoid.

\begin{theorem}[Soundness, completeness, and cut elimination]
The displayed multiplicative intuitionistic linear calculus is sound and
complete for commutative residuated partially ordered monoids.  Cut can be
eliminated: every derivable sequent has a cut-free derivation;
equivalently, cut is admissible in the cut-free fragment of the calculus.
\end{theorem}

\paragraph{Proof sketch.}
Soundness is a rule-by-rule use of monoid laws and residuation.  Algebraic
completeness orders formulas by derivability and forms the Lindenbaum quotient;
tensor supplies the monoid operation and linear implication supplies its
residual.  Cut elimination uses principal and commuting reductions adapted to
context splitting.  The broader classical linear calculus and its
proof-theoretic semantics originate with \cite{Girard1987}.

\begin{example}[Counterexample algebra and valuations]
Let \(A=\Z\), let the order be equality, define
\[
  a\otimes b=a+b,
  \qquad
  1^{\mathcal A}=0,
  \qquad
  a\multimap b=b-a.
\]
This is a commutative residuated partially ordered monoid because
\[
  a+b=c\quad\Longleftrightarrow\quad b=c-a.
\]
If \(v(p)=v(q)=1\), then \(p\Rightarrow p\otimes p\) would require
\(1=2\), so it is invalid; unrestricted contraction would derive it from
\(p,p\Rightarrow p\otimes p\).  Similarly, \(p,q\Rightarrow p\) would require
\(2=1\), so unrestricted left weakening is invalid.
\end{example}

The exponential \(!A\) marks formulas for which controlled structural rules
are available.  A standard extension includes dereliction, weakening,
contraction, and promotion in the forms
\[
\frac{\Gamma,A\Rightarrow B}{\Gamma,!A\Rightarrow B},
\qquad
\frac{\Gamma\Rightarrow B}{\Gamma,!A\Rightarrow B},
\qquad
\frac{\Gamma,!A,!A\Rightarrow B}{\Gamma,!A\Rightarrow B},
\]
\[
  \frac{!\Gamma\Rightarrow A}{!\Gamma\Rightarrow !A}.
\]
The side condition in promotion requires every antecedent formula to be
exponential.  Adding \(!\) changes both the syntax and the cut-elimination
proof.

\section{Relations among the systems}
\label{sec:relevance-linear-04}

Relevance logic and linear logic are both substructural, but they are not one
logic.  Relevance systems constrain implication so that premises bear an
appropriate relation to conclusions; many retain contraction while rejecting
unrestricted weakening.  Linear logic tracks multiplicity and normally rejects
both weakening and contraction outside exponentials.  A noncommutative
substructural calculus also rejects exchange and uses sequence contexts, as in
Lambek calculi.

\begin{warning}[Resource metaphors]
Calling a formula a ``resource'' can motivate a rule, but the formal content is
the context representation and the presence or absence of structural rules.
The metaphor cannot establish soundness, completeness, or cut elimination.
\end{warning}

\begin{exercise}
\label{ex:relevance-linear-01}
Verify soundness of \(\multimap R\) and \(\multimap L\) directly from the
residuation law.
\end{exercise}

\begin{exercise}
\label{ex:relevance-linear-02}
Give a cut-free derivation of
\(A\multimap B,A\Rightarrow B\).  Explain why no weakening is required.
\end{exercise}

\begin{exercise}
\label{ex:relevance-linear-03}
For each of exchange, weakening, and contraction, state how representing a
context as a sequence, multiset, or set affects whether the rule is built into
syntax.
\end{exercise}

\begin{exercise}
\label{ex:relevance-linear-04}
The formula \(p\to(q\to p)\) is classically valid and fails in the displayed
Routley--Meyer structure.  Identify the exact semantic clause that makes the
counterexample possible and explain why an ordinary binary Kripke relation would
not encode the same conditional clause.
\end{exercise}
\part{Formal Practice, Automation, and Verification}

\chapter{Constructing Derivations, Models, and Countermodels}
\label{ch:formal-practice}

This chapter collects methods that apply across the logical systems already
defined.  A method is not an inference rule.  It guides the search for a
derivation or semantic object; the completed result must still be checked by
the grammar, calculus, and semantics of the active logical system.

\section{A derivation-search method}
\label{sec:formal-practice-01}

\begin{method}[Backward and forward search]
Given premises \(\Gamma\), conclusion \(\varphi\), and calculus
\(\mathcal C\):
\begin{enumerate}
  \item Copy the exact judgment \(\Gamma\derives_{\mathcal C}\varphi\).
  \item Identify the principal symbol of \(\varphi\), its main
        connective in the propositional case.
  \item Work backward through introduction rules whose conclusion has that
        principal symbol.
  \item Work forward through elimination rules whose major premise is already
        available.
  \item Record every assumption to be discharged and the rule that will
        discharge it.
  \item For a quantifier rule, verify the displayed total capture-avoiding
        substitution and, when present, write the eigenvariable condition
        before applying the rule.
  \item When the branches meet, rewrite the result as a derivation tree and
        check every leaf, rule instance, and discharge mark.
\end{enumerate}
\end{method}

\begin{example}[Mixed search]
To derive \((P\to R)\land(Q\to R)\) from \(P\lor Q\) and
\((P\lor Q)\to R\), backward search on the conjunction creates the two targets
\(P\to R\) and \(Q\to R\).  Each target creates an assumption, to be
discharged by
implication introduction.  Forward search applies the second premise to the
first premise and yields \(R\), independently of either assumption.
Because a discharging application must discharge at least one open
occurrence of its assumption, each implication is then concluded through
a conjunction detour: assume \(P\), conjoin it with the derived \(R\) by
\(\land I\), recover \(R\) by \(\land E_2\), and discharge the occurrence
of \(P\) used by \(\land I\); the case for \(Q\) is symmetric.  The
derivation is longer than
necessary, and that fact is itself useful: the stronger intermediate result
\(R\) should be derived first, after which both implications follow by
the detours just described.
\end{example}

\section{A propositional countervaluation method}
\label{sec:formal-practice-02}

To refute \(\Gamma\entails\varphi\), the target conditions are
\[
  \widehat v(\gamma)=\True\quad(\gamma\in\Gamma),
  \qquad
  \widehat v(\varphi)=\False.
\]
Work inward from these target values.  A false implication requires a true
antecedent and false consequent; a true conjunction requires two true
conjuncts; a false disjunction requires two false disjuncts.  When a branch
assigns both truth values to the same propositional variable, close that branch.
A surviving total valuation is a countervaluation.

\begin{example}
For \((P\to Q)\lor(Q\to P)\), an attempted false valuation requires both
implications false.  The first requires \(P\) true and \(Q\) false; the second
requires \(Q\) true and \(P\) false.  The requirements conflict, so no
countervaluation exists.  This search establishes validity only because every
way for the outer disjunction to be false was exhausted.
\end{example}

\section{A first-order semantic-refutation method}
\label{sec:formal-practice-03}

\begin{method}[Small-structure search]
To refute a first-order semantic consequence:
\begin{enumerate}
  \item Negate the conclusion at the metalinguistic level by requiring it not to
        be satisfied.
  \item Expand the semantic clauses for the premises and failed conclusion.
  \item Start with a one-element domain unless a premise or failed conclusion
        requires distinct elements.
  \item Introduce witnesses for existential requirements and counterexamples
        for failed universal requirements.
  \item Interpret each relation and function symbol only as far as the expanded
        conditions require.
  \item Complete every partially specified function interpretation to a
        total function on the domain; the remaining values are
        unconstrained, and any choice may be recorded.
  \item Check every premise from the original formulas, not merely the expanded
        notes.
  \item Check that the conclusion fails under the required variable assignment,
        or in the structure when the conclusion is a sentence.
\end{enumerate}
\end{method}

\begin{example}
The sentence \(\forall x\exists y\,R(x,y)\) does not semantically entail
\(\exists y\forall x\,R(x,y)\).  Let the domain be \(\set{0,1}\), and
interpret \(R\) as \(\set{\tuple{0,0},\tuple{1,1}}\).  Every element is
related to itself, so the premise is true.  No single element is related to
both elements, so the conclusion is false.
\end{example}

\section{Writing a metatheoretic proof}
\label{sec:formal-practice-04}

A metatheoretic proof should expose its induction object and invariant.

\begin{method}[Metatheorem template]
\begin{enumerate}
  \item State the theorem with the language, calculus, semantic class, and all
        side hypotheses.
  \item Expand the target technical term into its definition.
  \item Identify the construction or induction object: formula, term,
        derivation, finite set, enumeration stage, or computation.
  \item State the invariant in a quantified sentence.
  \item Prove every base case and constructor case.
  \item Isolate uses of compactness, the Axiom of Choice, classical reasoning, or an imported
        theorem.
  \item Restate exactly the conclusion obtained; do not strengthen it in the
        final sentence.
\end{enumerate}
\end{method}

\begin{warning}
An induction on formula complexity is not complete when it proves only the
connective cases and omits the quantifier cases.  An induction on a derivation
is not complete when it checks only logical rules and omits structural rules or
initial sequents present in the calculus.
\end{warning}

\section{Failure analysis}
\label{sec:formal-practice-05}

When a proposed proof fails, classify the failure before changing the theorem.
\begin{longtable}{@{}p{0.24\textwidth}p{0.66\textwidth}@{}}
\toprule
Failure & Meaning \\
\midrule
\endhead
No derivation found & The search may be incomplete; this alone gives no
semantic conclusion. \\
Semantic refutation constructed & Semantic consequence fails; by soundness,
derivability in the sound calculus also fails. \\
Rule side condition fails & That rule instance is unavailable; a different
derivation may exist. \\
Construction invariant fails & The proposed construction does not prove the
theorem; a corrected invariant or construction is required. \\
Hypothesis cannot be verified & The cited theorem cannot yet be applied. \\
Proof assistant rejects a term & The term does not type-check in the current
environment; the informal theorem may still be correct. \\
Solver returns unknown & No satisfiability or validity conclusion follows from
that response. \\
\bottomrule
\end{longtable}

\begin{exercise}
\label{ex:formal-practice-01}
Develop a countermodel to
\(\forall x(P(x)\lor Q(x))\entails
(\forall xP(x))\lor(\forall xQ(x))\).  Prove that a one-element domain cannot
serve and that a two-element domain can.
\end{exercise}

\begin{exercise}
\label{ex:formal-practice-02}
Write a proof plan for the substitution lemma.  Identify the induction object,
the strengthened induction statement, and the quantifier case in which bound
variable renaming is required.
\end{exercise}

\chapter{Automated Reasoning}
\label{ch:automated-reasoning}

Automated reasoning separates a calculus from a search procedure.  A calculus
determines which inferences are licensed.  A search procedure decides which
licensed inference to try next.  Soundness is a property of the inferences;
termination and efficiency are properties of the search procedure.  The
handbook source for the calculi and procedures of this chapter is
\cite{RobinsonVoronkov2001}.

\section{Propositional satisfiability}
\label{sec:automated-reasoning-01}

The terms \emph{propositional literal}, \emph{clause}, and
\emph{conjunctive normal form} retain their definitions from the classical
propositional chapter.  In particular, a clause \(C\) is a finite set of
propositional literals, not a formula, and
\(\operatorname{form}_{\CPL}(C)\) is its exact formula interpretation.  Let
\(S\) be a finite set of clauses.  A valuation satisfies \(S\) iff it
satisfies \(\operatorname{form}_{\CPL}(C)\) for every \(C\in S\).  The empty
set of clauses is satisfied by every valuation, whereas a set containing the
empty clause is satisfied by none.

\begin{definition}[Unit propagation]
Let \(S\) be a finite set of clauses.  If \(S\) contains a unit clause
\(\set L\), \emph{unit propagation} assigns \(L\) true, deletes every clause
containing \(L\), and deletes the complementary literal from every remaining
clause.
\end{definition}

\begin{definition}[Pure-literal elimination]
If a propositional variable occurs in a finite clause set with only one
polarity, \emph{pure-literal elimination} assigns the truth value that
satisfies all of those occurrences and deletes the clauses containing them.
\end{definition}

The Davis--Putnam--Logemann--Loveland (DPLL) procedure recursively applies
unit
propagation, optionally applies pure-literal elimination, and otherwise branches
on an unassigned variable \cite{DavisPutnam1960,DavisLogemannLoveland1962}.
The following pseudocode specifies the logical core; it does not specify a
branching heuristic.  The set \(V\) of variables occurring in the original
input is retained unchanged by every recursive call.

\begin{lstlisting}[caption={Logical core of a DPLL satisfiability procedure}]
DPLL(S, valuation, V):
  repeatedly unit-propagate S, extending valuation by each forced value
  if S has no clauses:
    extend valuation to every unassigned variable in V with value false
    return SAT with the resulting total valuation
  if S contains the empty clause: return UNSAT
  choose an unassigned variable p
  r := DPLL(S with p=true, valuation extended by p=true, V)
  if r is SAT: return r
  return DPLL(S with p=false, valuation extended by p=false, V)
\end{lstlisting}

\begin{theorem}[Correctness of the displayed DPLL procedure]
For every finite set \(S\) of clauses, the displayed
procedure terminates.  It returns \(\SAT\) with a satisfying valuation iff
\(S\) is satisfiable, and it returns \(\UNSAT\) iff \(S\) is
unsatisfiable.
\end{theorem}
\begin{proof}
Maintain the invariant that every deleted original clause already contains a
literal made true by the current partial valuation, and every undeleted clause
is its residual after removing literals made false.  Unit propagation preserves
this invariant: every satisfying extension of a unit clause makes its literal
true.  If no residual clauses remain, every original clause is already true;
assigning arbitrary values to the remaining input variables cannot falsify a
clause that already contains a true literal.  Thus the displayed completion by
\(\False\) returns a total satisfying valuation on \(V\).

At a branch on \(p\), every total valuation makes either \(p\) true or false,
so the current residual set is satisfiable iff at least one branch is
satisfiable.  Each branch assigns a previously unassigned variable; there are
finitely many variables, so every path terminates.  Induction on the number of
unassigned variables now establishes both return specifications.
\end{proof}

\section{Equivalence and equisatisfiability}
\label{sec:automated-reasoning-02}

\begin{definition}[Equisatisfiability]
Formulas \(\varphi\) and \(\psi\), possibly in different signatures, are
\emph{equisatisfiable} iff the two formulas are either both satisfiable
or both unsatisfiable.
\end{definition}

Logical equivalence is stronger: it compares truth under corresponding
valuations, whereas equisatisfiability compares only existence of a satisfying
valuation.  Tseitin clausification \cite{Tseitin1968} introduces a fresh
variable for each
selected subformula and constrains that variable to have the same truth value
as the subformula.  If \(p_\varphi\) names the complete input formula
\(\varphi\), the output also contains the unit clause \(\set{p_\varphi}\),
which requires the root variable to be true.  Call the resulting finite clause
set \(S_\varphi\).  The conjunction of the
\(\operatorname{form}_{\CPL}\)-images in any enumeration of
\(S_\varphi\) is equisatisfiable with \(\varphi\) in the expanded signature,
and
\(S_\varphi\) can have size linear in the size of \(\varphi\).  The output is
not literally the same formula and need not be logically equivalent in the
original signature.

\begin{example}[Tseitin constraint]
For a subformula \(\theta=(P\land Q)\), introduce fresh \(p_\theta\) and add
the finite set of clauses
\[
  S_\theta=
  \set{
    \set{\neg p_\theta,P},
    \set{\neg p_\theta,Q},
    \set{p_\theta,\neg P,\neg Q}
  }.
\]
Conjoining the three \(\operatorname{form}_{\CPL}\)-images in the displayed
clause order gives a formula logically equivalent to
\[
  (\neg p_\theta\lor P)\land
  (\neg p_\theta\lor Q)\land
  (p_\theta\lor\neg P\lor\neg Q),
\]
which expresses \(p_\theta\leftrightarrow(P\land Q)\).
Every valuation of \(P,Q\) extends uniquely to \(p_\theta\) while satisfying
\(S_\theta\).  If \(\theta\) is the complete input, the Tseitin output is
\(S_\theta\cup\set{\set{p_\theta}}\); the added unit clause requires the
extension to make \(\theta\) true.
\end{example}

\section{First-order unification}
\label{sec:automated-reasoning-03}

The definitions of finite-support simultaneous term substitution, composition,
the instance preorder, unifier, and most general unifier are exactly those in
the first-order resolution chapter.  In this chapter only, substitutions are
written \(\theta_{\mathrm u},\rho_{\mathrm u},\eta_{\mathrm u}\).  The
subscript is necessary because \(\theta\) already denotes a formula in the
preceding Tseitin example; it changes no definition.  Composition is oriented
by
\[
  (\theta_{\mathrm u}\eta_{\mathrm u})(x)
  =(\theta_{\mathrm u}(x))\eta_{\mathrm u}.
\]
Thus \(\theta_{\mathrm u}\) is a most general unifier of atoms \(A\) and
\(B\) iff \(A\theta_{\mathrm u}=B\theta_{\mathrm u}\) literally and every
unifier \(\rho_{\mathrm u}\) has the form
\(\theta_{\mathrm u}\eta_{\mathrm u}\) for some term substitution
\(\eta_{\mathrm u}\).

The unification procedure performs the occurs check when it considers an
equation \(x\doteq t\).  If \(x\) occurs properly in \(t\), the procedure
fails.  Thus the atoms \(P(x)\) and \(P(f(x))\) have no unifier: their
argument equation is \(x\doteq f(x)\).  Assigning \(x\mapsto f(x)\) does not
unify the atoms because their simultaneously substituted arguments are
\(f(x)\) and \(f(f(x))\).

\begin{example}
The atoms \(P(f(x,g(a)))\) and \(P(f(g(y),g(y)))\) have the most general
unifier
\(\theta_{\mathrm u}=\set{x\mapsto g(a),y\mapsto a}\).  Applying it to
either atom yields \(P(f(g(a),g(a)))\).
\end{example}

\section{First-order resolution}
\label{sec:automated-reasoning-04}

The term \emph{clause} retains its standing as a finite set of first-order
literals.  Its formula interpretation is
\(\operatorname{ucl}_{\Lang}(C)\), the universal closure of
\(\operatorname{form}_{\Lang}(C)\).  Before each inference, rename the two
parent clauses apart.  If \(A\in E\), \(\neg B\in F\), and
\(\theta_{\mathrm u}\) is a most general unifier of the atoms \(A\) and
\(B\), the active resolution rule is
\[
  \Rule[Resolution]{E\qquad F}
  {\bigl((E\setminus\set{A})\cup(F\setminus\set{\neg B})\bigr)
   \theta_{\mathrm u}}.
\]
The conclusion is the resolvent.  Factoring is also part of the active
calculus: if two same-polarity literals in \(E\) have unifiable underlying
atoms and \(\theta_{\mathrm u}\) is a most general unifier, infer
\(E\theta_{\mathrm u}\).

\begin{theorem}[Soundness of resolution]
Every structure satisfying the parent clauses of a resolution or factoring
inference satisfies its conclusion.
\end{theorem}
\begin{proof}
Renaming apart preserves the universally closed formula interpretations.
Substitution preserves satisfaction of each parent clause.  If an assignment
falsified every literal in the displayed resolvent, satisfaction of the first
substituted parent would force \(A\theta_{\mathrm u}\), while satisfaction of
the second would force \(\neg(B\theta_{\mathrm u})\).  The two atoms are
literally identical, a
contradiction.  Factoring is sound because simultaneous substitution preserves
clause satisfaction.
\end{proof}

For a finite equality-free clause set, resolution with factoring is
refutationally complete: after adding a constant if the signature has no
ground term, the set is unsatisfiable iff the empty clause has a resolution
derivation (\cref{thm:fol-resolution}; \cite{Robinson1965}).  Bare
resolution is not complete for
logical equality.  Modern saturation calculi add ordering,
redundancy, and equality reasoning; those additions change the calculus and
must be stated before their completeness results are used
\cite{BachmairGanzinger2001}.

\section{Equality and superposition}
\label{sec:automated-reasoning-05}

Replacing equals by equals is semantically valid, but unrestricted replacement
creates an impractical search space.  A superposition calculus integrates
oriented equality reasoning with resolution by selecting eligible terms and
using a reduction ordering.  The ordering restrictions and redundancy
criterion are part of the calculus, not optional implementation details.

\begin{warning}
Treating equality as an uninterpreted binary relation without the identity
axioms changes the semantics.  Conversely, passing a built-in equality symbol
to a prover whose completeness theorem assumes a superposition calculus
requires the prover's actual equality rules, not ordinary resolution alone.
\end{warning}

\section{SAT, SMT, and first-order proving}
\label{sec:automated-reasoning-06}

\begin{definition}[Satisfiability modulo theories]
Given a specified background theory \(T\), the \emph{satisfiability
modulo theories problem} asks whether some model of \(T\) satisfies the
input formula under some variable assignment; equivalently, whether
\(T\) together with the existential closure of the input formula is
satisfiable.
\end{definition}

An SMT solver commonly combines Boolean search with decision procedures for
fragments such as equality with uninterpreted functions, linear arithmetic,
arrays, or bit-vectors.  The exact input grammar, sort system, theory symbols,
commands, and responses are standardized by SMT-LIB.  Version 2.7 specifies a
many-sorted language, operational commands, and logical semantics
\cite{SMTLIB2025}.  A response of \texttt{sat} should be accompanied by a model
when one is needed; a response of \texttt{unsat} should be accompanied by a
checkable proof or trusted-solver assumption when the conclusion is used in a
high-assurance development.

The TPTP language supplies standardized formula roles, names, syntax, and
annotations for automated theorem proving \cite{TPTP}.  A parser-accepted file
is not thereby a correctly formalized theory.  The declarations, formula
roles, conjecture polarity, and imported axioms must still be audited.

\section{Complexity and practical status}
\label{sec:automated-reasoning-07}

A \emph{decision problem} is, under the fixed coding, a set of natural
numbers.  It is \emph{polynomial-time decidable} iff some deterministic
Turing machine decides it within \(p(n)\) steps on every input of length
\(n\), for some polynomial \(p\).  The class \(\NP\) consists of the
decision problems \(A\) for which there are a polynomial-time decidable
binary relation \(R\) and a polynomial \(q\) such that \(x\in A\) iff
some witness \(w\) whose length is at most \(q(n)\), where \(n\) is
the length of \(x\), satisfies \(R(x,w)\).  A problem is \(\NP\)-\emph{complete} iff
it belongs to \(\NP\) and every member of \(\NP\) reduces to it by a
polynomial-time computable many-one reduction.  Propositional
satisfiability is \(\NP\)-complete \cite{Cook1971}.  Whether
propositional satisfiability is decidable in worst-case polynomial time is
the open \(\mathsf P\) versus \(\NP\) question; no worst-case
polynomial-time decision procedure is known.  This complexity
result does not imply that every large instance is difficult; modern solvers
exploit structure.  It also does not justify replacing a proof with a solver
response whose trust path is unstated.

First-order validity is computably enumerable and undecidable.  A complete
first-order prover may therefore search forever on some nonvalid inputs.  A
timeout means only that the configured run did not finish.

\begin{method}[Automated-result audit]
Record:
\begin{enumerate}
  \item the exact input language and version;
  \item the intended logical semantics;
  \item every declaration and nonlogical axiom;
  \item transformations, including whether they preserve equivalence or only
        satisfiability;
  \item the calculus and search program;
  \item the result status and resource limit;
  \item the proof certificate or model, when available; and
  \item the independent checker and its trusted base.
\end{enumerate}
\end{method}

\begin{exercise}
\label{ex:automated-reasoning-01}
Run the logical core of DPLL by hand on
\[
  \set{
    \set{P,Q},
    \set{\neg P,R},
    \set{\neg Q,R},
    \set{\neg R}
  }.
\]
Give the sequence of unit propagations and the final result.
\end{exercise}

\begin{exercise}
\label{ex:automated-reasoning-02}
Show that the atoms \(P(f(x,x))\) and \(P(f(g(y),y))\) have no unifier.
Identify the equation that triggers the occurs check.
\end{exercise}

\begin{exercise}
\label{ex:automated-reasoning-03}
Transform the negation of
\(\forall x(P(x)\to\exists y\,R(x,y))\) into a finite set of first-order
clauses whose universal-closure interpretations are equisatisfiable with the
negated sentence in the expanded signature.  At every step state whether the
transformation preserves logical equivalence or only equisatisfiability.
\end{exercise}

\chapter{Mechanized Proof Checking}
\label{ch:mechanization}

\section{Kernel, elaboration, and automation}
\label{sec:mechanization-01}

\begin{definition}[Proof term]
Relative to a type-theoretic proof assistant, a \emph{proof term} is a term
whose type is the formal statement proved, according to the assistant's kernel
type theory.
\end{definition}

\begin{definition}[Kernel]
A \emph{kernel} is the component that checks declarations and proof terms
against the primitive typing and conversion rules of the implemented
foundation.
\end{definition}

\begin{definition}[Elaboration]
\emph{Elaboration} is the transformation of surface-language input containing
notation, omitted arguments, coercions, and tactic output into a kernel-level
term with all information required for checking.
\end{definition}

\begin{definition}[Tactic]
A \emph{tactic} is a program that transforms a proof state or constructs a
candidate proof term.  A tactic is not an additional inference rule when its
output is checked by the kernel.
\end{definition}

Lean 4, Rocq, and Isabelle expose different foundations and user languages.
Their official manuals distinguish surface syntax from kernel or framework
logic \cite{LeanManual2026,LeanTPIL,DeMouraUllrich2021,RocqManual2026,Isabelle2025}.
A theorem
checked in one assistant is relative to the definitions, axioms, imports, and
kernel of that development.  For a survey of the mechanization of
mathematics, see \cite{Avigad2018}.

\section{Propositions as types}
\label{sec:mechanization-02}

In a propositions-as-types foundation, a proposition is represented by a type
and a proof by a term of that type.  Implication corresponds to a function type:
from a term of type \(A\to B\) and a term of type \(A\), function application
produces a term of type \(B\).  Universal quantification corresponds to a
dependent function type when the quantified object may occur in the result.

This correspondence does not identify every object-language formula in this
reference with a proposition of the assistant without an encoding.  A deep
embedding represents object-language syntax by a datatype and defines an
object-level derivability or satisfaction relation.  A shallow embedding maps
object-language operations directly to operations in the assistant's
foundation.  The latter can be concise, but it may make adequacy and the exact
object logic harder to isolate.

\begin{method}[Encoding audit]
For an encoded logical system, require:
\begin{enumerate}
  \item a representation of signatures, terms, formulas, and variable binding;
  \item a definition of substitution with a capture argument;
  \item a representation of structures and satisfaction or of derivations;
  \item a theorem connecting the representation to the intended paper
        definition;
  \item explicit treatment of quotienting, alpha-equivalence, or de Bruijn
        indices;
  \item soundness of any reflected decision procedure; and
  \item a list of axioms and unsafe or opaque extensions used.
\end{enumerate}
\end{method}

\section{Trusted base}
\label{sec:mechanization-03}

\begin{definition}[Trusted base]
For a particular checked result, the \emph{trusted base} is the set of
components and assumptions whose correctness is required but is not itself
established by the checked proof term.
\end{definition}

The trusted base normally includes the kernel executable and hardware.  It may
also include axioms, foreign code, an external solver, generated proof
certificates, or an unverified parser, depending on the workflow.  The trusted
base is result-specific; naming the proof assistant is not a complete trust
report.

\begin{method}[Mechanized-result report]
Report separately:
\begin{enumerate}
  \item the natural-language theorem intended;
  \item the exact formal theorem statement;
  \item the correspondence argument between them;
  \item the assistant and exact version;
  \item the imported modules;
  \item the axioms printed or otherwise inspected;
  \item the kernel-check result;
  \item external tools and certificate checking; and
  \item remaining human review of definitions and fidelity.
\end{enumerate}
\end{method}

\section{A small propositions-as-types example}
\label{sec:mechanization-04}

The term
\[
  \lambda f.\,\lambda g.\,\lambda x.\,f(g(x))
\]
has type
\[
  (B\to C)\to(A\to B)\to A\to C.
\]
Under propositions as types it is a proof of the corresponding implicational
formula.  Each lambda abstraction introduces an assumption, and each
application eliminates an implication.  The type checker verifies the local
rule structure; it does not decide whether \(A,B,C\) faithfully encode the
informal concepts intended by an application.

\begin{warning}
The phrase ``formally verified'' is incomplete without the formal statement and
trusted base.  A proof of a weakened translation, a proof containing an
unintended axiom, or a proof about an incorrect definition may be kernel-correct
and still fail to establish the intended theorem.
\end{warning}

\begin{exercise}
\label{ex:mechanization-01}
Give proof terms for \(A\to A\),
\((A\land B)\to A\), and
\((A\to B)\to(B\to C)\to(A\to C)\) in a simply typed lambda calculus with
product types.
\end{exercise}

\begin{exercise}
\label{ex:mechanization-02}
Compare a deep and shallow embedding of propositional conjunction.  State one
adequacy theorem required by each representation.
\end{exercise}

\chapter{Writing and Typesetting Formal Logic}
\label{ch:latex-style}

\section{Academic prose carries the argument}
\label{sec:latex-style-01}

Every display belongs to a grammatical sentence.  Use prose to state why an
equation or derivation appears, which hypotheses are active, and what the
display establishes.  Symbols should compress repeated structure after it has
been defined; they should not replace the definition.

\begin{method}[Paragraph structure]
A proof paragraph should normally have four components:
\begin{enumerate}
  \item the local target;
  \item the definition or prior result used;
  \item the derivation, construction, or calculation; and
  \item the exact consequence for the theorem.
\end{enumerate}
When the calculation has several lines, align relation symbols and justify each
nontrivial transition.
\end{method}

The following rules govern academic style in this reference.
\begin{enumerate}
  \item State the language, calculus, semantic class, and ambient metatheory
        before asserting derivability, validity, completeness, or consistency.
  \item Quantify every object on which a claim depends.  Replace ``for suitable
        objects'' by the exact domain and hypotheses.
  \item Use \emph{definition}, \emph{calculation}, \emph{derivation}, and
        \emph{countermodel} only for the corresponding evidence.  A diagram or
        suggestive example is not a derivation.
  \item Use only the proof-status descriptions \emph{complete proof},
        \emph{proof sketch}, and \emph{cited theorem}.  A proof sketch must
        state its omitted obligations; a cited theorem must reproduce the
        hypotheses used and identify its source.
  \item Replace evaluative adverbs such as ``clearly'' and ``obviously'' by the
        definition, rule instance, or finite calculation that licenses the
        sentence.
  \item Keep object-language formulas distinct from metalinguistic assertions.
        In particular, \(\varphi\leftrightarrow\psi\) is a formula, whereas
        \(\varphi\equiv\psi\) is a declared metalinguistic relation.
  \item State the induction object, induction hypothesis, and every constructor
        case.  If a simultaneous or strengthened induction is required, state
        it before the base case.
  \item End a proof paragraph with the exact clause established.  Do not leave
        the reader to infer whether the argument proved existence, uniqueness,
        equivalence, or only one implication.
\end{enumerate}

\section{One symbol, one role within a declared scope}
\label{sec:latex-style-02}

Maintain a notation register during drafting.  For every symbol record its
standing, type, first definition, scope, and conflicts.  A glyph is not globally
reserved merely because one chapter assigns it a role.  For example,
\(\Sigma\) is a formula context in a displayed calculus and is also part of
\(\Sigma^0_n\) in the arithmetical hierarchy.  These roles are legitimate
because each scope is declared before use.  International conventions for
mathematical symbols are catalogued in \cite{ISO80000}; where the
registered notation of this document departs from them, the register
governs.
Within one declared scope, however, a glyph must not change standing silently.

Overloading \(\entails\) is controlled by argument shape:
\[
  \mathcal M,s\entails\varphi,
  \qquad
  \mathcal M\entails\sigma,
  \qquad
  \Gamma\entails\varphi.
\]
These respectively denote satisfaction, truth of a sentence in a structure,
and semantic consequence.  The propositional, modal, and matrix chapters
register the analogous shapes \(v\entails\varphi\),
\(\mathcal M\entails T\), \(\mathcal M,w\entails A\), and
\(\entails_{\mathbf K_3}\); no meaning outside these declared shapes is
assigned to the glyph.

The unification section of the automated-reasoning chapter
(\cref{sec:automated-reasoning-03}) uses
\(\theta_{\mathrm u},\rho_{\mathrm u},\eta_{\mathrm u}\) for substitutions.
The visible subscript distinguishes those maps from the formula \(\theta\) in
the Tseitin example while preserving the substitution definitions and
composition convention of the first-order resolution chapter.

\section{A complete before-and-after formalization}
\label{sec:latex-style-before-after}

Consider the sentence ``Clearly, if every \(P\) is a \(Q\) and something is a
\(P\), then something is a \(Q\).''  As written, it does not identify a
signature, a calculus, the status of the letters, or the inference that uses the
existential premise.  It is therefore not yet a formal-logical claim.

Fix a first-order signature with unary relation symbols \(P,Q\), use the
natural-deduction calculus \(\ND_{\FOL}\), and let \(x,y\) be object-language
variables.  The corrected claim is
\[
  \forall x(P(x)\to Q(x)),\ \exists xP(x)
  \derives_{\ND_{\FOL}} \exists xQ(x).
\]
Here is a complete derivation.  The superscript \(i\) marks the assumption
discharged by \(\exists E\):
\[
\begin{array}{rll}
1.&\forall x(P(x)\to Q(x))&\text{premise},\\
2.&\exists xP(x)&\text{premise},\\
3.&[P(y)]^i&\text{assumption},\\
4.&P(y)\to Q(y)&1,\ \forall E,\\
5.&Q(y)&3,4,\ \to E,\\
6.&\exists xQ(x)&5,\ \exists I,\\
7.&\exists xQ(x)&2,\ 3\text{--}6,\ \exists E\ (i).
\end{array}
\]
Choose \(y\) absent from the conclusion, from every other open assumption, and
from the existential premise.  Those are exactly the active
\(\exists E\) side conditions.  Lines 4--6 show that an arbitrary witness for
the premise supplies a witness for the conclusion; line 7 discharges the only
open assumption.  The revised passage therefore fixes the language,
calculus, metavariable standing, rule applications, discharge, and freshness
condition that the original sentence omitted.

\section{LaTeX structure}
\label{sec:latex-style-03}

Use semantic markup in the source.  A macro name should record the formal-logical
role, not the desired visual spacing.  For example:
\begin{lstlisting}[language={[LaTeX]TeX},caption={Semantic macros for logic}]
\newcommand{\Fm}{\mathsf{Fm}}
\newcommand{\Sent}{\mathsf{Sent}}
\newcommand{\FV}{\operatorname{FV}}
\newcommand{\derives}{\vdash}
\newcommand{\entails}{\models}
\end{lstlisting}

Use \texttt{amsmath} alignment environments; do not use raw double-dollar
displays.  Use \texttt{amsthm} environments to distinguish definitions,
theorems, examples, and remarks.  Label every referenced theorem or equation
and use cross-reference commands rather than fixed numbers.  The AMS guidance
and \texttt{amsmath} documentation provide the publication conventions used
here \cite{AMSAuthorHandbook,AMSMathGuide}.

\section{Auditing source and rendered output}
\label{sec:latex-style-source-render}

An audit of rendered output includes verifying that every rendered glyph is
the intended character; the normative catalogue of mathematical symbol
characters is \cite{Unicode2025}.

The following small example exhibits the whole path from semantic source to a
rendered claim.  The source deliberately contains a typed signature, quantified
semantic parameters, semantic macros, and a stable label:
\begin{lstlisting}[language={[LaTeX]TeX},caption={Source for a semantic implication}]
\begin{theorem}\label{thm:sample-semantic-implication}
Let \(\Lang\) contain unary relation symbols \(P,Q\).
For every \(\Lang\)-structure \(\mathcal M\) and assignment \(s\),
if \(\mathcal M,s\entails P(x)\) and
\(\mathcal M,s\entails\forall y(P(y)\to Q(y))\), then
\(\mathcal M,s\entails Q(x)\).
\end{theorem}
\begin{proof}
By the universal clause, if \(P^{\mathcal M}(s(x))\) then
\(Q^{\mathcal M}(s(x))\).  The first hypothesis gives
\(P^{\mathcal M}(s(x))\), so \(Q^{\mathcal M}(s(x))\), which is exactly
\(\mathcal M,s\entails Q(x)\).
\end{proof}
\end{lstlisting}

The rendered statement must visibly read as a claim about every
\(\Lang\)-structure and assignment; both hypotheses must remain inside the
conditional, and \(\entails\) must render as semantic satisfaction rather than
derivability.  The rendered proof must expose the instance \(s(x)\), because
that calculation is the bridge between the universal semantic clause and the
conclusion.

Audit this example in the following order.
\begin{enumerate}
  \item \emph{Source audit.}  Confirm that \(P,Q\) are relation symbols,
        \(\mathcal M\) is a structure, \(s\) is an assignment, and \(x,y\) are
        object-language variables.  Confirm that the label is unique and that
        a later reference uses \verb|\cref| rather than a typed theorem number.
  \item \emph{Formal-logical audit.}  Expand the satisfaction clauses.  The
        hypotheses become \(P^{\mathcal M}(s(x))\) and
        \(\forall a\in\abs{\mathcal M}
        (P^{\mathcal M}(a)\to Q^{\mathcal M}(a))\); instantiation at \(s(x)\)
        yields the conclusion.  This independently checks the prose proof.
  \item \emph{Compilation audit.}  Compile enough times to resolve the label
        and bibliography.  Reject undefined references, missing citations,
        duplicate labels, malformed environments, and any overfull box that
        obscures a symbol.
  \item \emph{Rendered-page audit.}  Inspect the theorem number, line breaks,
        subscripts, turnstile/model glyphs, proof-ending mark, and surrounding
        page break.  Verify that no line makes the quantifier appear to govern
        less text than the source says it governs.
  \item \emph{Extraction audit.}  Extract text from the final PDF and check that
        the theorem heading, label target, hypotheses, and conclusion remain
        searchable.  Text extraction supplements visual inspection; it does
        not replace it.
\end{enumerate}
Passing only the source or only the rendered page is insufficient: the audit
requires agreement among the formal claim, the LaTeX source, compilation
records, cross-reference target, and visible output.

\section{Rule displays}
\label{sec:latex-style-04}

Each rule display must show:
\begin{itemize}
  \item all premises;
  \item the conclusion;
  \item the rule name;
  \item discharged assumptions; and
  \item freshness, admissibility, or occurrence side conditions.
\end{itemize}
If contexts are sequences, multisets, or sets, state that choice before the
rule sheet.  Structural rules depend on it.

\section{Worked correction of an invalid rule application}
\label{sec:latex-style-rule-correction}

The following tempting use of universal introduction is invalid:
\[
  \Rule[invalid \(\forall I\)]
  {P(x)\derives P(x)}
  {P(x)\derives\forall xP(x)}.
\]
The variable \(x\) occurs free in the open assumption \(P(x)\), contrary to the
side condition \(x\notin\FV(\Gamma)\).  The failure is semantic as well: on a
two-element domain, let \(P\) hold only of the element assigned to \(x\).
Then \(P(x)\) is satisfied under that assignment while
\(\forall xP(x)\) is not, so \(\set{P(x)}\notentails\forall xP(x)\):
the semantic counterpart of the conclusion judgment fails.

One corrected derivation begins with a genuinely universal premise:
\[
  \Rule[\(\forall I\)]
  {\Rule[\(\forall E\)]
    {\forall yP(y)\derives\forall yP(y)}
    {\forall yP(y)\derives P(x)}}
  {\forall yP(y)\derives\forall xP(x)}
  \side{\(x\notin\FV(\set{\forall yP(y)})\)}.
\]
Universal elimination supplies the arbitrary instance \(P(x)\); now \(x\) is
absent from the only open assumption, so universal introduction is licensed.
The correction changes the premise rather than hiding the violated side
condition.

\section{Citation practice}
\label{sec:latex-style-05}

Separate formal-logical proof from historical attribution.  A citation can
identify a source or authorize an imported theorem, but it does not turn an
unstated theorem into a proved theorem.  When importing a result, reproduce the
hypotheses and conclusion needed.  When citing a continuously updated manual,
record the version or retrieval date.

Primary sources control historical priority and original theorem statements.
Canonical monographs control modern notation and standard proofs.  Official
manuals control current software and language behavior.  Secondary summaries
may orient research but should not control a disputed technical claim when the
owning source is available.

\section{Revision audit}
\label{sec:latex-style-06}

\begin{method}[Release audit for a formal document]
\begin{enumerate}
  \item Search for every technical term at first use and verify a definition.
  \item Search every symbol and verify that it retains one role within each
        declared scope.
  \item Check every theorem for quantified objects and exact hypotheses.
  \item Check every use of soundness, completeness, consistency, compactness,
        decidability, and categoricity for the intended meaning.
  \item Check every substitution for capture avoidance.
  \item Check every quantifier rule for its side condition.
  \item Check every imported theorem against the cited source.
  \item Check that proof sketches are labelled and that no later result treats a
        sketch as a formal derivation.
  \item Compile until there are no undefined references, missing citations,
        duplicate labels, or overfull boxes that damage readability.
  \item Inspect the rendered pages, including rule displays, tables, index, and
        bibliography.
  \item Report formal-logical review, machine checking, compilation, visual
        inspection, and human pedagogical review separately.
\end{enumerate}
\end{method}

\begin{exercise}
\label{ex:latex-style-01}
Rewrite a proof paragraph that uses ``clearly'' twice and omits the type of a
function.  Replace each occurrence with the exact definition, calculation, or
typed declaration required.
\end{exercise}

\begin{exercise}
\label{ex:latex-style-02}
Create a notation register for one chapter of a logic text.  Identify every
overloaded glyph and either remove the overloading or specify how argument
shape disambiguates it.
\end{exercise}

\chapter{Research Workflow and Reproducibility}
\label{ch:research-workflow}

\section{Table for theorem status}
\label{sec:research-workflow-01}

Maintain a table with one row per result:
\begin{longtable}{@{}p{0.16\textwidth}p{0.26\textwidth}p{0.24\textwidth}p{0.24\textwidth}@{}}
\toprule
Field & Content & Purpose & Failure detected \\
\midrule
\endhead
Identifier & Stable label & Cross-reference & Duplicate or drifting names \\
Statement & Exact quantified formula & Scope & Strength mismatch \\
Dependencies & Definitions and prior theorems & Order & Definition after use \\
Status & complete proof, proof sketch, or cited theorem & Evidence & Unsupported result \\
System & Language, calculus, semantics & Meaning & Logic shift \\
Verification & Checks actually performed & Reporting & Compilation reported as proof \\
\bottomrule
\end{longtable}

\section{Counterexample-first testing}
\label{sec:research-workflow-02}

Before investing in a proof, test the statement against boundary cases:
\begin{itemize}
  \item empty premise set;
  \item one-element and two-element structures;
  \item empty relation and universal relation;
  \item formulas with no free variable occurrence;
  \item vacuous quantification;
  \item inconsistent theories;
  \item finite versus infinite models;
  \item countable versus uncountable signatures; and
  \item constructive versus classical metalogic.
\end{itemize}
These tests do not prove the statement.  They expose missing hypotheses and
incorrect generality early.

\section{Reproducible computational supplements}
\label{sec:research-workflow-03}

A computational supplement should include exact input files, tool versions,
commands, expected results, and certificate-checking steps.  Generated examples
should have fixed seeds.  If a proof assistant is used, the source should build
from a clean environment and print or otherwise expose the axiom surface of the
principal results.

\section{Independent review}
\label{sec:research-workflow-04}

An author should ask a reviewer to begin from the theorem statements and seek
counterexamples before reading the proofs.  A second pass should trace every
definition through the proof.  A third pass should compare the final statement
with its formalization or computational encoding.  These passes test different
failure modes and should be recorded separately.

\begin{warning}
No finite audit establishes that a long reference contains no formal-logical
error.  A release report should state the checks completed and the review not
performed.  It should not replace those facts with a global assurance.
\end{warning}

\appendix

\chapter{Registered Notation}
\label{ch:registered-notation}
\label{app:notation}

This is a scoped register, not a list of global reservations.  Each row governs
only the scope named in its second column.  The same glyph may have another
standing in a disjoint scope only after that standing has been declared before
substantive use.  Within one row's scope, the registered standing and meaning
remain fixed.

\begin{longtable}{@{}p{0.14\textwidth}p{0.23\textwidth}p{0.20\textwidth}p{0.34\textwidth}@{}}
\toprule
Notation & Declared scope & Standing & Meaning \\
\midrule
\endhead
\(\Lang\) & after a language declaration & signature & Nonlogical symbols with
declared arities or sorts. \\
\(x,y,z\) & first-order syntax, semantics, and rule sheets & object-language
variables & Variables occurring in terms and formulas. \\
\(r,t\) & first-order term formation, substitution, and identity rules & term
metavariables & Members of \(\Term(\Lang)\). \\
\(\varphi,\psi,\chi\) & propositional and first-order chapters, unless a local
declaration says otherwise & formula metavariables & Members of the active
formula set. \\
\(\sigma,\tau\) & theory, model-theory, and arithmetic passages after the
sentence set is fixed & sentence metavariables & Members of
\(\Sent(\Lang)\). \\
\(\theta,\rho,\eta\) & first-order resolution chapter & simultaneous term
substitutions & Finite-support variable-to-term maps with the action and
composition declared there. \\
\(\theta_{\mathrm u},\rho_{\mathrm u},\eta_{\mathrm u}\) & automated-reasoning
chapter in this part & simultaneous term substitutions & Local visible aliases
for the preceding substitution role, used to avoid collision with a formula
metavariable in the Tseitin example. \\
\(\Gamma,\Delta,\Sigma,\Pi\) & a natural-deduction or sequent-calculus display
that declares its context representation & formula contexts & Finite sets,
multisets, or sequences exactly as declared by that active calculus. \\
\(\mathcal A_{\mathrm{in}},\mathcal A_{\mathrm{tape}}\) & Turing-machine
definitions & alphabets & Input and tape alphabets, respectively, with no
formula-context standing in that scope. \\
\(u,v,w\) & formal-language and computability passages that declare an
alphabet & strings & Finite words over the declared alphabet. \\
\(\Sigma^0_n,\Pi^0_n\) & arithmetical-hierarchy section & syntactic classes
of formulas & Classes fixed by the declared normalization and unbounded
quantifier alternation. \\
\(\boldsymbol{\Sigma}^0_n,\boldsymbol{\Pi}^0_n,
\boldsymbol{\Delta}^0_n\) & arithmetical-hierarchy section & definability
classes of relations & Classes fixed by definability in \(\mathcal N\) using
formulas from the corresponding syntactic class. \\
\(\alpha,\beta\) & cumulative-hierarchy passages & ordinal metavariables &
Ordinals in the ambient set theory. \\
\(\alpha,\beta\) & simple-type-theory passages after the local type grammar &
type metavariables & Simple types; this scope is disjoint from the preceding
ordinal scope. \\
\(T,U\) & first-order theories and metatheory after a signature is fixed &
theories & Sets of sentences in that signature. \\
\(\mathcal M,\mathcal N\) & semantic and model-theory passages & structures &
Structures for the active signature, with nonempty domains unless a local
convention says otherwise. \\
\(\abs{\mathcal M}\) & whenever \(\mathcal M\) is a structure & set & Domain
of \(\mathcal M\). \\
\(v\) & propositional and finite-matrix semantics & valuation & Map from
propositional variables to the active truth-value set; its unique
homomorphic extension \(\widehat v\) to formulas is fixed by the active
semantics. \\
\(s\) & first-order semantics & variable assignment & Map from first-order
variables to \(\abs{\mathcal M}\). \\
\(\mathcal M,s\entails\varphi\) & first-order semantics & satisfaction &
\(\varphi\) is satisfied by \(\mathcal M\) under \(s\). \\
\(\mathcal M\entails\sigma\) & sentence semantics & truth & Sentence \(\sigma\)
is true in \(\mathcal M\). \\
\(\Gamma\entails\varphi\) & after a semantic class is declared & semantic
consequence & Every interpretation and assignment in that class satisfying
\(\Gamma\) satisfies \(\varphi\). \\
\(\Gamma\derives_{\mathcal C}\varphi\) & after calculus \(\mathcal C\) is
declared & derivability & A derivation exists in \(\mathcal C\). \\
\(\varphi\equiv\psi\) & after a semantic class is declared & metalinguistic
relation & The formulas have the stated logical equivalence in that class. \\
\(\varphi\leftrightarrow\psi\) & a language containing biconditional & formula
constructor & Object-language biconditional. \\
\(a\defeq b\) & definitions in the metalanguage & metalinguistic abbreviation &
The expression on the left is defined by the expression on the right. \\
\(\varphi[t/x]\) & first-order syntax & formula operation & Capture-avoiding
substitution of \(t\) for free occurrences of \(x\) in \(\varphi\). \\
\(\FV(\varphi)\), \(\FV(\Gamma)\) & first-order syntax & set-valued operation &
Variables with a free occurrence in \(\varphi\), and their union over all
members of the formula context \(\Gamma\), respectively. \\
\(\Mod(T)\) & model theory & class & Structures satisfying every sentence of
\(T\). \\
\(\Th(\mathcal M)\) & model theory & theory & Sentences true in
\(\mathcal M\). \\
\(R^*\) & relation-closure and common-knowledge passages & relation & Reflexive
transitive closure of the locally declared relation \(R\). \\
\(\bottom,\topv\) & a formula language that declares these constants & formulas
& The bottom formula and the defined top-formula abbreviation. \\
\bottomrule
\end{longtable}

\section{Structure convention}
\label{sec:registered-notation-01}

Unless a chapter explicitly states otherwise, every first-order structure and
every sort domain is nonempty.  This convention affects quantifier laws,
Skolemization, and translations between many-sorted and one-sorted logic.

\section{Natural numbers convention}
\label{sec:registered-notation-02}

\(\N\) contains \(0\).  A finite ordinal \(n\) is
\(\set{0,\ldots,n-1}\).  The word ``positive'' therefore means greater
than \(0\), not merely a member of \(\N\).

\chapter{Binding and Precedence}
\label{ch:binding-precedence}
\label{app:precedence}

The official grammar is primary.  The following precedence convention permits
parentheses to be omitted in expository formulas after unique readability has
been established.

\begin{center}
\begin{tabularx}{\textwidth}{@{}c>{\raggedright\arraybackslash}p{0.38\textwidth}>{\raggedright\arraybackslash}X@{}}
\toprule
Rank & Symbols & Association \\
\midrule
1 & function application, predicate application & determined by arity \\
2 & \(=\) & nonassociative \\
3 & \(\neg,\Box,\Diamond,\forall x,\exists x\) & prefix; governs the immediately following formula \\
4 & \(\land\) & left associative \\
5 & \(\lor\) & left associative \\
6 & \(\to\) & right associative \\
7 & \(\leftrightarrow\) & nonassociative \\
\bottomrule
\end{tabularx}
\end{center}

Thus \(\neg P\land Q\to R\) abbreviates
\(((\neg P)\land Q)\to R\), while
\(P\to Q\to R\) abbreviates \(P\to(Q\to R)\).  Parentheses must be restored
whenever omission could make the intended parse uncertain.  In particular,
\(\forall x\,P(x)\to Q(x)\) abbreviates
\((\forall x\,P(x))\to Q(x)\); write
\(\forall x\,(P(x)\to Q(x))\) when the implication belongs to the
quantifier's scope.

\chapter{Rule Sheets}
\label{ch:rule-sheets}
\label{app:rules}

\section{Classical natural deduction}
\label{sec:rule-sheets-01}

Contexts in this sheet are finite sets of formulas.  A displayed rule is a rule
schema: uniformly substituting formulas for its formula metavariables produces
a rule instance.  In the propositional portion, \(\derives\) abbreviates
\(\derives_{\ND_{\CPL}}\); in the quantifier and identity portion it
abbreviates \(\derives_{\ND_{\FOL}}\).  The latter calculus contains every
rule of the former.  Assumptions enter derivations through
\[
\Rule[assumption]{}{\Gamma\cup\set{\varphi}\derives\varphi}.
\]

\subsection*{Conjunction}
\[
\Rule[\(\land I\)]{\Gamma\derives\varphi\qquad\Delta\derives\psi}
{\Gamma\cup\Delta\derives\varphi\land\psi}
\qquad
\Rule[\(\land E_1\)]{\Gamma\derives\varphi\land\psi}
{\Gamma\derives\varphi}
\qquad
\Rule[\(\land E_2\)]{\Gamma\derives\varphi\land\psi}
{\Gamma\derives\psi}.
\]

\subsection*{Disjunction}
\[
\Rule[\(\lor I_1\)]{\Gamma\derives\varphi}
{\Gamma\derives\varphi\lor\psi}
\qquad
\Rule[\(\lor I_2\)]{\Gamma\derives\psi}
{\Gamma\derives\varphi\lor\psi}
\]
\[
\Rule[\(\lor E\)]{
\Gamma\derives\varphi\lor\psi\qquad
\Delta\cup\set\varphi\derives\chi\qquad
\Sigma\cup\set\psi\derives\chi}
{\Gamma\cup\Delta\cup\Sigma\derives\chi}.
\]
The displayed occurrences of \(\varphi\) and \(\psi\) in the minor
derivations are discharged.

\subsection*{Implication}
\[
\Rule[\(\to I\)]{\Gamma\cup\set\varphi\derives\psi}
{\Gamma\derives\varphi\to\psi}
\qquad
\Rule[\(\to E\)]{\Gamma\derives\varphi\to\psi\qquad
\Delta\derives\varphi}
{\Gamma\cup\Delta\derives\psi}.
\]
The displayed occurrence of \(\varphi\) in implication introduction is
discharged.

\subsection*{Primitive negation, the bottom formula, and RAA}
\[
\Rule[\(\neg I\)]{\Gamma\cup\set{\varphi}\derives\bottom}
{\Gamma\derives\neg\varphi}
\qquad
\Rule[\(\neg E\)]{\Gamma\derives\varphi\qquad
\Delta\derives\neg\varphi}
{\Gamma\cup\Delta\derives\bottom}
\]
\[
\Rule[\(\bottom E\)]{\Gamma\derives\bottom}
{\Gamma\derives\varphi}
\qquad
\Rule[RAA]{\Gamma\cup\set{\neg\varphi}\derives\bottom}
{\Gamma\derives\varphi}.
\]
The rules \(\to I\), \(\neg I\), and RAA discharge the displayed assumption.
Negation is primitive in the active classical language.  Removing RAA gives
the primitive-negation presentation of its intuitionistic fragment.

\subsection*{Quantifiers}
\[
\Rule[\(\forall I\)]{\Gamma\derives\varphi}
{\Gamma\derives\forall x\,\varphi}
\side{\(x\notin\FV(\Gamma)\)}
\]
\[
\Rule[\(\forall E\)]{\Gamma\derives\forall x\,\varphi}
{\Gamma\derives\varphi[t/x]}
\]
\[
\Rule[\(\exists I\)]{\Gamma\derives\varphi[t/x]}
{\Gamma\derives\exists x\,\varphi}
\]
\[
\Rule[\(\exists E\)]{
\Gamma\derives\exists x\,\varphi\qquad
\Delta\cup\set{\varphi[y/x]}\derives\psi}
{\Gamma\cup\Delta\derives\psi}.
\]
\noindent\emph{Side condition.}  \(y\) is absent from \(\psi\), from every
other open assumption, and from \(\exists x\,\varphi\).
The displayed assumption \(\varphi[y/x]\) is discharged.  In the contextual
display, ``every other open assumption'' means every formula in
\(\Gamma\cup\Delta\); the side condition is therefore literally the one used
by the active natural-deduction calculus.

\subsection*{Alpha-equivalence}
\[
\Rule[\(\alpha\)]{\Gamma\derives\varphi}{\Gamma\derives\psi}
\side{\(\varphi\equiv_\alpha\psi\)}.
\]
This is the explicit bound-variable conversion rule of the active calculus;
it does not permit replacement by an arbitrary logically equivalent
formula.

\subsection*{Logical identity}
\[
\Rule[\(=I\)]{}{\derives t=t}
\qquad
\Rule[\(=E\)]{\Gamma\derives r=t\qquad\Delta\derives\varphi[r/x]}
{\Gamma\cup\Delta\derives\varphi[t/x]}.
\]
Identity elimination is symmetric because symmetry is derivable from these
rules.  Every substitution displayed in this sheet is the total
capture-avoiding substitution fixed in the syntax chapter, so no separate
substitution side condition is imposed.

\section{Classical sequent calculus}
\label{sec:rule-sheets-02}

A sequent is \(\Gamma\Rightarrow\Delta\), with finite multisets on both sides.
It is interpreted as
\(\bigwedge\Gamma\to\bigvee\Delta\), with the empty conjunction interpreted as
\(\topv\) and empty disjunction as \(\bottom\).

The propositional rules below are exactly the active calculus
\(\LK_{\CPL}\).  The subsequent quantifier and identity rules extend it to
\(\LK_{\FOL}\).

\subsection*{Initial sequents and cut}
\[
\Rule[id]{}{\Gamma,\varphi\Rightarrow\varphi,\Delta}
\qquad
\Rule[\(\bottom L\)]{}{\Gamma,\bottom\Rightarrow\Delta}
\]
\[
\Rule[cut]{\Gamma\Rightarrow\Delta,\varphi\qquad
\Sigma,\varphi\Rightarrow\Pi}
{\Gamma,\Sigma\Rightarrow\Delta,\Pi}.
\]

\subsection*{Structural rules}
\[
\Rule[WL]{\Gamma\Rightarrow\Delta}{\varphi,\Gamma\Rightarrow\Delta}
\quad
\Rule[WR]{\Gamma\Rightarrow\Delta}{\Gamma\Rightarrow\Delta,\varphi}
\]
\[
\Rule[CL]{\varphi,\varphi,\Gamma\Rightarrow\Delta}
{\varphi,\Gamma\Rightarrow\Delta}
\quad
\Rule[CR]{\Gamma\Rightarrow\Delta,\varphi,\varphi}
{\Gamma\Rightarrow\Delta,\varphi}.
\]
Exchange is built into the multiset convention.

\subsection*{Logical rules}
\[
\Rule[\(\neg L\)]{\Gamma\Rightarrow\Delta,\varphi}
{\Gamma,\neg\varphi\Rightarrow\Delta}
\qquad
\Rule[\(\neg R\)]{\Gamma,\varphi\Rightarrow\Delta}
{\Gamma\Rightarrow\Delta,\neg\varphi}
\]
\[
\Rule[\(\land L\)]{\varphi,\psi,\Gamma\Rightarrow\Delta}
{\varphi\land\psi,\Gamma\Rightarrow\Delta}
\qquad
\Rule[\(\land R\)]{\Gamma\Rightarrow\Delta,\varphi\qquad
\Gamma\Rightarrow\Delta,\psi}
{\Gamma\Rightarrow\Delta,\varphi\land\psi}
\]
\[
\Rule[\(\lor L\)]{\varphi,\Gamma\Rightarrow\Delta\qquad
\psi,\Gamma\Rightarrow\Delta}
{\varphi\lor\psi,\Gamma\Rightarrow\Delta}
\qquad
\Rule[\(\lor R\)]{\Gamma\Rightarrow\Delta,\varphi,\psi}
{\Gamma\Rightarrow\Delta,\varphi\lor\psi}
\]
\[
\Rule[\(\to L\)]{\Gamma\Rightarrow\Delta,\varphi\qquad
\Gamma,\psi\Rightarrow\Delta}
{\Gamma,\varphi\to\psi\Rightarrow\Delta}
\qquad
\Rule[\(\to R\)]{\varphi,\Gamma\Rightarrow\Delta,\psi}
{\Gamma\Rightarrow\Delta,\varphi\to\psi}.
\]

\subsection*{Quantifiers}
\[
\Rule[\(\forall L\)]{\Gamma,\varphi[t/x]\Rightarrow\Delta}
{\Gamma,\forall x\,\varphi\Rightarrow\Delta}
\qquad
\Rule[\(\exists R\)]{\Gamma\Rightarrow\Delta,\varphi[t/x]}
{\Gamma\Rightarrow\Delta,\exists x\,\varphi}
\]
\[
\Rule[\(\forall R\)]{\Gamma\Rightarrow\Delta,\varphi[y/x]}
{\Gamma\Rightarrow\Delta,\forall x\,\varphi}
\qquad
\Rule[\(\exists L\)]{\Gamma,\varphi[y/x]\Rightarrow\Delta}
{\Gamma,\exists x\,\varphi\Rightarrow\Delta}.
\]
In \(\forall R\), the eigenvariable \(y\) is free in no formula of
\(\Gamma,\Delta,\forall x\,\varphi\).  In \(\exists L\), it is free in no
formula of \(\Gamma,\Delta,\exists x\,\varphi\).  Every displayed
substitution is total and capture-avoiding.

\subsection*{Logical identity}
\[
\Rule[\(=I\)]{}{\Rightarrow t=t}
\qquad
\Rule[\(=L\)]{\Gamma,\varphi[r/x]\Rightarrow\Delta}
{\Gamma,r=t,\varphi[t/x]\Rightarrow\Delta}.
\]
The symmetric replacement instance is derivable.  These displays, together
with the propositional rules, are the complete active first-order sequent
calculus used in this reference.

\section[Modal calculus for K]{Modal calculus for \(\K\)}
\label{sec:rule-sheets-03}

The Hilbert calculus \(\K\) contains all substitution instances of classical
classically valid propositional formulas, the axiom schema
\[
  \Box(\varphi\to\psi)\to(\Box\varphi\to\Box\psi),
\]
modus ponens, and necessitation:
\[
  \Rule[nec]{\derives_{\K}\varphi}{\derives_{\K}\Box\varphi}.
\]
Necessitation applies to theorems, not to an arbitrary formula depending on an
undischarged assumption.  Adding \(\Box\varphi\to\varphi\),
\(\Box\varphi\to\Box\Box\varphi\), and
\(\Diamond\varphi\to\Box\Diamond\varphi\) yields the standard systems
\(\mathsf T\), \(\Sfour\), and \(\Sfive\) in the combinations defined in the
modal chapter.

\chapter{Qualifying-Examination Problems and Selected Solutions}
\label{ch:qualifying-problems}
\label{app:qualifying}

\section{Problem 1: substitution and satisfaction}
\label{sec:qualifying-problems-01}
\label{prob:substitution}

Let \(\varphi\) be a first-order formula, let \(x\) be a variable, let
\(t\) be a term, let
\(\Meta\) be a structure, and let \(s\) be a variable assignment.  Prove
\[
  \Meta,s\entails\varphi[t/x]
  \quad\Longleftrightarrow\quad
  \Meta,s[x\mapsto\sem{t}^{\Meta}_s]\entails\varphi.
\]

\begin{solution}
First prove the term-substitution lemma by structural induction on terms.  Then
proceed by structural induction on \(\varphi\).  Atomic formulas follow from
the term lemma.  For each connective, unfold its semantic clause and apply the
induction hypotheses to its immediate subformulas.  Let
\(Q\in\set{\forall,\exists}\) and consider \(Qy\,\psi\).  If \(y=x\),
substitution stops at the binder.  On the right side, the quantifier update of
\(x\) overrides the preceding update of \(x\), so the truth value of
\(Qx\,\psi\) is independent of that preceding value.

Now suppose \(y\neq x\).  If \(y\) occurs in \(t\), the deterministic
capture-avoiding definition first renames the binder to a variable fresh for
\(t\) and \(\psi\).  Thus it remains to consider
\(y\notin\operatorname{Var}(t)\).  For each \(a\in\abs{\Meta}\), the assignments
\[
  s[y\mapsto a][x\mapsto\sem{t}^{\Meta}_s]
  \quad\text{and}\quad
  s[x\mapsto\sem{t}^{\Meta}_s][y\mapsto a]
\]
agree because \(x\neq y\), and term evaluation of \(t\) is unchanged by the
update at \(y\).  Apply the induction hypothesis to \(\psi\) under each updated
assignment and then apply the semantic clause for \(Q\).  Alpha-invariance
justifies the preliminary binder renaming.  Substitution is total and
capture-avoiding, so the displayed definition requires no additional
definedness hypothesis.
\end{solution}

\section{Problem 2: compactness}
\label{sec:qualifying-problems-02}
\label{prob:compactness}

Let \(\Lang\) contain a binary relation symbol \(<\) and constants
\(c_n\) for \(n\in\N\).  Let \(T\) say that \(<\) is a strict linear order and
include \(c_{n+1}<c_n\) for every \(n\).  Prove that \(T\) has a model, and
explain why this does not prove that every strict linear order has an infinite
descending chain.

\begin{solution}
Every finite subset \(T_0\subseteq T\) mentions only finitely many of the
descending-chain axioms.  Interpret the relevant constants as a finite
descending segment of \(\Z\) with its usual order, and interpret unused
constants arbitrarily.  Then \(T_0\) has a model.  Compactness gives a model of
all of \(T\), in which the interpretations of the \(c_n\) form an infinite
descending chain.  The conclusion is existential: at least one strict linear
order with such a chain exists.  It makes no universal assertion about all
strict linear orders; \(\N\) with its usual order is a counterexample to that
universal assertion.
\end{solution}

\section{Problem 3: incompleteness hypotheses}
\label{sec:qualifying-problems-03}
\label{prob:incompleteness}

State a version of the first incompleteness theorem for a computably enumerable
theory \(T\), and explain why each hypothesis is needed.

\begin{solution}
One standard version is: if \(T\) is a consistent, computably enumerable
theory in the language of arithmetic extending a fixed weak theory sufficient
to represent the required computable syntactic relations, then there is a
sentence \(G_T\) such that \(T\notderives G_T\); with the stronger
\(\omega\)-consistency hypothesis in G\"odel's original form,
\(T\notderives\neg G_T\) as well.  Rosser's refinement replaces
\(\omega\)-consistency by consistency for a differently constructed sentence.
The computably enumerable axiom set makes provability arithmetically expressible;
arithmetic strength supports coding and diagonalization; consistency prevents
the theory from deriving every sentence; and the strengthened hypothesis in
the original proof excludes derivation of the negation.  No version licenses a
claim about an arbitrary collection of axioms lacking these properties.
\end{solution}

\section{Problem 4: modal frame condition}
\label{sec:qualifying-problems-04}
\label{prob:modal-frame}

Prove that \(\Box P\to P\) is valid on a frame iff the accessibility relation is
reflexive.

\begin{solution}
If the relation is reflexive and
\(\mathcal M,w\entails\Box P\), then every accessible world satisfies \(P\);
since \(w\) is accessible from itself, \(\mathcal M,w\entails P\).  Conversely,
suppose the frame is not reflexive at \(w\).  Let \(\mathcal M\) be the
structure on that frame whose valuation makes \(P\) true at every world
accessible from \(w\) and false at \(w\).  Then
\(\mathcal M,w\entails\Box P\) but \(\mathcal M,w\notentails P\), so the
formula is not valid on the frame.
\end{solution}

\section{Problem 5: cut and subformula property}
\label{sec:qualifying-problems-05}
\label{prob:cut-subformula}

Explain why cut elimination implies a subformula property for the cut-free
classical propositional sequent calculus, and state one consequence.

\begin{solution}
Inspect every propositional rule except cut.  Each formula in a premise is either copied from
the conclusion or is an immediate subformula of the principal formula in the
conclusion.  Induction upward through a cut-free derivation therefore shows
that every formula occurrence is a subformula of a formula in the end sequent,
apart from the declared treatment of structural copies.  Cut is the rule that
can introduce an arbitrary intermediate formula and then remove it.  A
consequence is consistency: a cut-free derivation of the empty sequent
would have a last rule, but no rule of the cut-free calculus---initial,
structural, or logical---has the empty sequent as its conclusion, since
every weakening or contraction conclusion contains a displayed formula
occurrence.
For the first-order calculus, quantifier premises contain substitution
instances.  Its analytic property must therefore be stated modulo term
substitution and eigenvariable renaming rather than as the literal subformula
property just proved.
\end{solution}

\section{Additional problems}
\label{sec:qualifying-problems-06}

\begin{enumerate}
  \item Prove \L{}o\'s's theorem for existential formulas from its induction
        hypothesis, explicitly identifying the use of the Axiom of Choice in
        selecting coordinatewise witnesses.
  \item Give a complete Henkin construction for a countable consistent theory
        and prove the truth lemma for its term structure.
  \item Prove that the full second-order Peano axioms are categorical, and state
        why the proof uses full semantics.
  \item Construct a three-valued matrix in which explosion fails.  State the
        designated values and calculate one countervaluation.
  \item Formalize one theorem in a proof assistant and record the exact theorem
        declaration, printed axiom surface, build instructions, and a paragraph
        comparing the declaration with the informal theorem.
\end{enumerate}

\chapter{Exercise Feedback}
\label{ch:exercise-feedback}
\label{app:exercise-feedback}

This appendix follows the exercise order in the reference.  Every heading gives
the exact source label and links back to the exercise.  ``Solution'' marks a
complete answer at the level requested.  ``Solution sketch'' gives the decisive
construction or invariant and an explicit criterion by which a completed proof
can be checked.

\section{Language levels, formal objects, and proof writing}

\paragraph{\hyperref[ex:levels-01]{\texttt{ex:levels-01}}.}
\textit{Solution.}  Relative to classical propositional logic, \(P\) and
\(P\land Q\) are object-language formulas, while \(\varphi\), the phrase ``is
valid,'' and \(\entails\) belong to the metalanguage.  A left-parenthesis token
used in the fully parenthesized formula \((P\land Q)\) is object-language
punctuation; the phrase ``the left parenthesis'' that mentions that token is
metalinguistic.

\paragraph{\hyperref[ex:levels-02]{\texttt{ex:levels-02}}.}
\textit{Solution.}  Put \(v(P)=\False\) and \(v(Q)=\True\).  Then
\(v(P\to Q)=\True\) and \(v(Q)=\True\), while \(v(P)=\False\).  Thus this
single row satisfies both premises and falsifies the conclusion.

\paragraph{\hyperref[ex:strings-01]{\texttt{ex:strings-01}}.}
\textit{Solution.}  Let \(B\) be the least set of strings such that
\(\epsilon\in B\), \(w\in B\) implies \((w)\in B\), and \(u,v\in B\) implies
\(uv\in B\).  Write \(L(w),R(w)\) for the two parenthesis counts.  Structural
induction gives \(L(\epsilon)=R(\epsilon)=0\); if \(L(w)=R(w)\), then both
counts increase by one in \((w)\); and if the equality holds for \(u,v\), then
\(L(uv)=L(u)+L(v)=R(u)+R(v)=R(uv)\).

\paragraph{\hyperref[ex:strings-02]{\texttt{ex:strings-02}}.}
\textit{Solution.}  Give every initial derivation height \(0\), and define the
height of a derivation ending with immediate subderivations
\(D_1,\ldots,D_k\) to be \(1+\max_i h(D_i)\).  Each immediate subderivation
has smaller height.  Repeating that strict inequality down the finite
constructor path shows that every proper subderivation has strictly smaller
height.

\paragraph{\hyperref[ex:writing-foundations-01]{\texttt{ex:writing-foundations-01}}.}
\textit{Solution.}  A correct completeness statement is: for every set
\(\Gamma\) of first-order sentences and sentence \(\sigma\), if
\(\Gamma\entails\sigma\) over all nonempty first-order structures, then
\(\Gamma\derives_{\ND_{\FOL}}\sigma\).  A false correction is: for every
arithmetic sentence \(\sigma\), if \(\mathcal N\entails\sigma\), then
\(\PA\derives\sigma\).  The first fixes a calculus and semantic class; the
second fails because true arithmetic is not the theorem set of \(\PA\).

\paragraph{\hyperref[ex:writing-foundations-02]{\texttt{ex:writing-foundations-02}}.}
\textit{Solution.}  The original leaves ``theory,'' ``true,'' the intended
structures, and the proof calculus unspecified.  Deductive completeness is:
for every sentence \(\sigma\), either \(T\derives\sigma\) or
\(T\derives\neg\sigma\).  The separate semantic property is that all models of
\(T\) agree on every sentence.  For a satisfiable \(T\), soundness and
completeness of the fixed first-order calculus make these two properties
equivalent.

\paragraph{\hyperref[ex:set-prerequisites-01]{\texttt{ex:set-prerequisites-01}}.}
\textit{Solution.}  Let \(i:A\to\N\) be an injection.  For each
\(b\in f[A]\), define
\[
  j(b)\defeq\min\set{i(a):a\in A\text{ and }f(a)=b}.
\]
The set minimized over is a nonempty subset of \(\N\).  If
\(j(b)=j(b')\), injectivity of \(i\) gives the same selected preimage, so
\(b=b'\).  Thus \(j:f[A]\to\N\) is an injection.  This also covers
\(A=\varnothing\), when \(j\) is the empty function.

\paragraph{\hyperref[ex:set-prerequisites-02]{\texttt{ex:set-prerequisites-02}}.}
\textit{Solution.}  Define \([a]_E+[b]_E=[a+b]_E\).  If
\(aE a'\) and \(bE b'\), then \(5\mid(a-a')\) and \(5\mid(b-b')\), so
\(5\mid((a+b)-(a'+b'))\).  Therefore \([a+b]_E=[a'+b']_E\), which proves
independence from both representatives.

\paragraph{\hyperref[ex:set-prerequisites-03]{\texttt{ex:set-prerequisites-03}}.}
\textit{Solution.}  For a family \((A_n)_{n\in\N}\), the usual proof first
selects, for every \(n\), one enumeration or surjection
\(e_n:\N\to A_n\) (with a separate convention for empty members), and only
then maps \((n,k)\) to \(e_n(k)\).  The simultaneous selection of the sequence
\((e_n)_n\) is the simultaneous selection step; pairing \(\N^2\) with \(\N\) is not.  If the
enumerations are part of the given data, that step is unnecessary.

\section{Classical propositional logic}

\paragraph{\hyperref[ex:prop-metatheory-01]{\texttt{ex:prop-metatheory-01}}.}
\textit{Solution.}  The full parse is
\(((\neg P\land Q)\to(R\lor S))\).  Its subformulas are
\(P,Q,R,S,\neg P,(\neg P\land Q),(R\lor S)\), and the whole implication.
Every binary node shown is fixed by the declared precedence, and implication
associativity is not needed because there is only one implication token.

\paragraph{\hyperref[ex:prop-metatheory-02]{\texttt{ex:prop-metatheory-02}}.}
\textit{Solution.}  Let \(l(\varphi)\) count syntax-tree leaves---occurrences
of propositional variables or \(\bottom\)---and let \(b(\varphi)\) count
binary-connective occurrences.  For a variable or \(\bottom\),
\(l=1=0+1=b+1\).  A negation has the same leaves and binary nodes as its sole
subformula, so the induction hypothesis is preserved.  If
\(\varphi\circ\psi\) is formed with a binary connective, then
\[
l(\varphi\circ\psi)=l(\varphi)+l(\psi)
=(b(\varphi)+1)+(b(\psi)+1)
=b(\varphi\circ\psi)+1.
\]
These are all constructors, so structural induction proves the invariant.

\paragraph{\hyperref[ex:prop-metatheory-03]{\texttt{ex:prop-metatheory-03}}.}
\textit{Solution.}  Set \(P=\False,Q=\True\).  Both \(P\to Q\) and \(Q\) are
true, but \(P\) is false, so semantic consequence fails.  This valuation
inspects truth values only; nonexistence of a derivation follows separately
from soundness of the specified calculus.

\paragraph{\hyperref[ex:prop-metatheory-04]{\texttt{ex:prop-metatheory-04}}.}
\textit{Solution.}  Let \(P\) mean that the process terminates, \(Q\)
that the input is finite, and \(R\) that the search is bounded.  Taking
``or'' as inclusive before applying the stated restriction, the phrase
following ``only if'' requires exactly one of \(Q\) and \(R\) whenever
\(P\) holds.  A faithful translation is
\[
  P\to\bigl((Q\lor R)\land\neg(Q\land R)\bigr).
\]
The stronger formula
\((P\to(Q\lor R))\land\neg(Q\land R)\) would forbid \(Q\land R\) even when
the process does not terminate; that global restriction is not stated by the
sentence.

\paragraph{\hyperref[ex:prop-metatheory-05]{\texttt{ex:prop-metatheory-05}}.}
\textit{Solution.}  A conjunctive normal form is
\((\neg P\lor Q)\land(\neg Q\lor R)\).  A disjunctive normal form is
\((\neg P\land\neg Q)\lor(Q\land R)\).  In the four rows with
\(Q=\False\), the formula is true iff \(P=\False\); in the four rows
with \(Q=\True\), it is true iff \(R=\True\).  Those two cases are
exactly the two displayed
disjuncts.

\paragraph{\hyperref[ex:prop-metatheory-06]{\texttt{ex:prop-metatheory-06}}.}
\textit{Solution.}  Assume \(P\to R\), then \(Q\to R\), then \(P\lor Q\).
For the left case assume \(P\) and infer \(R\); for the right case assume
\(Q\) and infer \(R\).  Disjunction elimination discharges the two case
assumptions.  Three successive applications of \(\to I\) discharge
\(P\lor Q\), \(Q\to R\), and \(P\to R\), in that order, yielding the required
nested implication.

\paragraph{\hyperref[ex:prop-metatheory-07]{\texttt{ex:prop-metatheory-07}}.}
\textit{Solution.}  Assume \(P\land\neg P\).  Conjunction elimination gives
\(P\) and \(\neg P\); negation elimination gives \(\bottom\); negation
introduction discharges the assumption and yields \(\neg(P\land\neg P)\).
The decisive final rule is \(\neg I\), which is available in both displayed
classical and intuitionistic natural deduction; no reductio rule is used.

\paragraph{\hyperref[ex:prop-metatheory-08]{\texttt{ex:prop-metatheory-08}}.}
\textit{Solution sketch.}  Apply \(\to R\) to reduce the target to
\(P\to Q,Q\to R,P\Rightarrow R\).  Apply \(\to L\) to \(Q\to R\); its right
premise is the identity sequent \(R\Rightarrow R\).  Its left premise derives
\(Q\) by applying \(\to L\) to \(P\to Q\), with identities
\(P\Rightarrow P\) and \(Q\Rightarrow Q\).  Context weakening supplies the
displayed companions.  A checked tree contains no cut, and every branch ends in
identity.

\paragraph{\hyperref[ex:prop-metatheory-09]{\texttt{ex:prop-metatheory-09}}.}
\textit{Solution sketch.}  Reflexivity, symmetry, and transitivity follow from
equality of truth values under every valuation.  For congruence, assume
corresponding inputs have equal truth values; each primitive connective is a
fixed truth function, so applying it to equal input tuples gives equal outputs.
Check the unary negation case and each binary truth table separately; these are
all primitive constructors.

\paragraph{\hyperref[ex:prop-metatheory-10]{\texttt{ex:prop-metatheory-10}}.}
\textit{Solution sketch.}  Choose the underivable formula \(P\).  Then
\(\set{\neg P}\) is consistent; extend it to a maximal consistent set
\(\Delta\).  Define \(v(Q)=\True\) iff \(Q\in\Delta\).  The truth
lemma, proved by induction on formulas, gives
\(v(\varphi)=\True\) iff \(\varphi\in\Delta\).  Since \(\neg P\in\Delta\),
maximal consistency gives \(P\notin\Delta\), hence \(v(P)=\False\).  A complete
answer must include the truth-lemma cases for every primitive connective.

\paragraph{\hyperref[ex:prop-metatheory-11]{\texttt{ex:prop-metatheory-11}}.}
\textit{Solution.}  The clause set is
\(\set{\set{P,Q},\set{\neg P,R},\set{\neg Q,R},\set{\neg R}}\).
Resolve the second and fourth clauses to obtain \(\set{\neg P}\), and the
third and fourth to obtain \(\set{\neg Q}\).  Resolve \(\set{P,Q}\) with
\(\set{\neg P}\) to get \(\set{Q}\), then with \(\set{\neg Q}\) to obtain
the empty clause.

\paragraph{\hyperref[ex:prop-metatheory-12]{\texttt{ex:prop-metatheory-12}}.}
\textit{Solution.}  A tautological clause is true under every valuation, so
removing it changes neither the satisfying valuations nor satisfiability.  For
the second claim take \(S=\set{\set{P},\set{Q}}\).  The valuation
\(P=Q=\True\) satisfies the clause \(\set{P}\) that is deleted, but the reduced
set \(\set{\set{Q}}\) is also satisfied by \(P=\False,Q=\True\), whereas
\(S\) is not.  Thus satisfaction by the current valuation does not establish
equivalence under all valuations.

\section{Classical first-order logic}

\paragraph{\hyperref[ex:fol-deduction-01]{\texttt{ex:fol-deduction-01}}.}
\textit{Solution.}  The free-variable set is \(\set{x,y,z}\).  The occurrence
of \(x\) in \(P(x,z)\) is free; the occurrences of \(x\) in \(Q(x,y)\) and
\(R(x,y,z)\) are bound by \(\forall x\).  The occurrence of \(y\) in
\(Q(x,y)\) is free because it lies outside the scope of \(\exists y\); the
occurrence of \(y\) in \(R(x,y,z)\) is bound by that existential quantifier.
Both occurrences of \(z\) are free.

\paragraph{\hyperref[ex:fol-deduction-02]{\texttt{ex:fol-deduction-02}}.}
\textit{Solution.}  In the first formula, rename the bound \(y\) to fresh
\(z\) before substitution, obtaining
\((\exists yR(x,y))[y/x]=\exists zR(y,z)\).  In the second, rename bound
\(x\) to fresh \(z\), because the replacing term contains \(x\); the result is
\((\forall x(P(x)\to Q(y)))[f(x)/y]
=\forall z(P(z)\to Q(f(x)))\).  In each result the formerly free variables in
the replacing term remain free.

\paragraph{\hyperref[ex:fol-deduction-03]{\texttt{ex:fol-deduction-03}}.}
\textit{Solution sketch.}  Prove simultaneously: term values agree whenever
two assignments agree on the term's variables, and formula truth agrees
whenever they agree on the formula's free variables.  The variable case is the
assignment hypothesis; the binary-function case applies both term induction
hypotheses; the binary-relation case applies term coincidence to both
arguments.  Boolean cases use formula induction.  In a quantifier case, update
both assignments at the bound variable by the same arbitrary element, then use
the induction hypothesis.  The proof for an arbitrary signature differs only
by replacing two arguments with the declared finite arity of the active symbol.

\paragraph{\hyperref[ex:fol-deduction-04]{\texttt{ex:fol-deduction-04}}.}
\textit{Solution.}  Choose \(y\) absent from both premises and the conclusion.
Under the assumption \(P(y)\), to be discharged by \(\exists E\),
universal elimination gives
\(P(y)\to Q(y)\); implication elimination gives \(Q(y)\); existential
introduction gives \(\exists xQ(x)\).  Apply \(\exists E\) to the premise
\(\exists xP(x)\) and discharge \(P(y)\).  The eigenvariable \(y\) is absent
from the conclusion, every other open assumption, and the existential premise.

\paragraph{\hyperref[ex:fol-deduction-05]{\texttt{ex:fol-deduction-05}}.}
\textit{Solution.}  From \(x=y\) and \(y=z\), apply equality elimination to
the latter equality with the formula \(\varphi(w)\defeq x=w\).  Its required
minor premise is \(\varphi(y)\), namely \(x=y\), and its conclusion is
\(\varphi(z)\), namely \(x=z\).  Thus \(x=y,y=z\derives x=z\).

\paragraph{\hyperref[ex:fol-deduction-06]{\texttt{ex:fol-deduction-06}}.}
\textit{Solution.}  Use one domain containing people and papers, unary
predicates \(\mathsf{Rev},\mathsf{Paper}\), and binary relations
\(\mathsf{Read},\mathsf{Cites}\).  The translations are
\[
\begin{aligned}
&\forall x\bigl(\mathsf{Rev}(x)\to\exists y(\mathsf{Paper}(y)\land
 \mathsf{Read}(x,y)\land
 \forall z((\mathsf{Paper}(z)\land\mathsf{Read}(x,z))\to z=y))\bigr),\\
&\exists y\bigl(\mathsf{Paper}(y)\land
 \forall x(\mathsf{Rev}(x)\to\mathsf{Read}(x,y))\bigr),\\
&\forall x\bigl(\mathsf{Rev}(x)\to\exists y(\mathsf{Paper}(y)\land
 \mathsf{Read}(x,y)\land\neg\mathsf{Cites}(y,y))\bigr).
\end{aligned}
\]
The last formula makes the nonvacuous reading explicit: every reviewer read at
least one paper that did not cite itself.

\paragraph{\hyperref[ex:fol-deduction-07]{\texttt{ex:fol-deduction-07}}.}
\textit{Solution.}  If a \(\Lang_{\mathrm{grp}}\)-structure satisfies the three
displayed sentences, the interpretation of \(\cdot\) is associative,
\(e\) is a two-sided identity, and \(i(a)\) is a two-sided inverse of each
\(a\); hence the reduct is a group.  Conversely, interpret \(\cdot,e,i\) in
any group by its multiplication, identity, and inverse operation.  The group
laws verify each universal sentence under every assignment.  Equality is
logical identity in both directions, so the model class is exactly the class of
groups in that signature.

\paragraph{\hyperref[ex:fol-deduction-08]{\texttt{ex:fol-deduction-08}}.}
\textit{Solution.}  In the signature with one binary relation \(E\), take
\[
T_{\mathrm{graph}}=\set{
\forall x\,\neg E(x,x),\quad
\forall x\forall y(E(x,y)\leftrightarrow E(y,x))}.
\]
Its models are precisely nonempty simple undirected graphs.  The first axiom
forbids loops, and logical identity needs no additional equivalence axioms.

\paragraph{\hyperref[ex:fol-deduction-09]{\texttt{ex:fol-deduction-09}}.}
\textit{Solution.}  Begin with the identity sequent \(P(y)\Rightarrow P(y)\).
Apply \(\forall L\) to obtain \(\forall xP(x)\Rightarrow P(y)\), then
\(\exists R\) to obtain
\(\forall xP(x)\Rightarrow\exists xP(x)\).  No cut occurs.  The semantic
soundness of choosing the instance denoted by \(y\) uses that assignments into
the domain exist, hence that the domain is nonempty.

\paragraph{\hyperref[ex:fol-cardinality-normalforms-01]{\texttt{ex:fol-cardinality-normalforms-01}}.}
\textit{Solution sketch.}  For \(\mathcal M\subseteq\mathcal N\), prove by
structural induction that formulas with parameters from \(M\) have the same
truth value in both structures.  Atomic and Boolean cases use the substructure
hypothesis and induction.  For \(\exists x\varphi\), the forward direction
keeps a witness in \(M\); the reverse direction uses the Tarski--Vaught witness
condition to replace an \(N\)-witness by one in \(M\), then applies induction.
For \(\forall x\varphi\), one direction restricts the quantified domain, while
the other argues that a counterexample in \(N\) would make
\(\exists x\neg\varphi\) true and hence have an \(M\)-witness.  A complete
proof must state both quantifier directions without invoking duality as a
substitute for these arguments.

\paragraph{\hyperref[ex:fol-cardinality-normalforms-02]{\texttt{ex:fol-cardinality-normalforms-02}}.}
\textit{Solution sketch.}  Let \(A_0\) be \(A\) together with \(m_0\)
and the interpretations of the constant symbols.  Given countable
\(A_n\), let \(A_{n+1}\) add the witness value \(F_\varphi(\bar a)\) for
every formula \(\exists y\,\varphi(y,\bar x)\) and tuple \(\bar a\) from
\(A_n\); the witness functions for the formulas
\(\exists y\,(y=f(\bar x))\) put the value of every signature function on
every tuple from \(A_n\) into \(A_{n+1}\).  These prescriptions explicitly
give \(A_0,A_1,A_2\), and countability follows from the countable signature.
If a finite tuple lies in \(\bigcup_nA_n\), it lies in some common stage; the
value of every signature function on that tuple lies in the next stage.  This
is the required closure check.

\paragraph{\hyperref[ex:fol-cardinality-normalforms-03]{\texttt{ex:fol-cardinality-normalforms-03}}.}
\textit{Solution.}  Add constants \(c_0,c_1,\ldots\) and all inequalities
\(c_i\ne c_j\) for \(i\ne j\) to \(T\).  Every finite subset has a model:
choose a sufficiently large finite model of the mentioned part of \(T\) and
interpret the finitely many constants distinctly.  Compactness gives a model
of the whole expanded theory, whose reduct is an infinite model of \(T\).  No
map from any particular finite model was constructed, so no embedding claim
follows.

\paragraph{\hyperref[ex:fol-cardinality-normalforms-04]{\texttt{ex:fol-cardinality-normalforms-04}}.}
\textit{Solution.}  Suppose an empty-signature sentence \(\sigma\) held exactly
in finite nonempty structures.  Add, for each \(n\geq1\), a sentence saying
that at least \(n\) distinct elements exist.  Every finite subset together with
\(\sigma\) has a sufficiently large finite model.  Compactness therefore gives
an infinite model of \(\sigma\), contradicting the assumed characterization.

\paragraph{\hyperref[ex:fol-cardinality-normalforms-05]{\texttt{ex:fol-cardinality-normalforms-05}}.}
\textit{Solution.}  Successively applying quantifier negation, implication
elimination, De Morgan's law, and vacuous-variable movement gives
\[
\begin{aligned}
&\neg\forall x(P(x)\to\exists y(Q(y)\land R(x,y)))\\
&\equiv\exists x\bigl(P(x)\land\neg\exists y(Q(y)\land R(x,y))\bigr)\\
&\equiv\exists x\forall y
  \bigl(P(x)\land(\neg Q(y)\lor\neg R(x,y))\bigr).
\end{aligned}
\]
The variables are distinct, and moving \(\forall y\) across the conjunction is
licensed because \(y\notin\FV(P(x))\).  The final formula is rectified and in
prenex normal form.

\paragraph{\hyperref[ex:fol-cardinality-normalforms-06]{\texttt{ex:fol-cardinality-normalforms-06}}.}
\textit{Solution.}  Replace \(u\) by a new constant \(c\); replace \(y\) by a
new unary function \(f(x)\); and replace \(w\) by a new binary function
\(g(x,z)\).  The universal closure is
\[
\forall x\forall z\,R(c,x,f(x),z,g(x,z)).
\]
The arities \(0,1,2\) record exactly the preceding universal variables on which
the corresponding existential witness may depend.

\paragraph{\hyperref[ex:fol-cardinality-normalforms-07]{\texttt{ex:fol-cardinality-normalforms-07}}.}
\textit{Solution.}  Let the original sentence be
\(\forall x\exists yR(x,y)\), and take the two-element structure with
\(R=\set{(0,0),(1,1)}\).  Its reduct satisfies the sentence.  Expanding by the
identity function makes \(\forall xR(x,f(x))\) true, while expanding by the
constant-zero function makes it false at \(1\).  Thus Skolemization is a claim
about existence of a suitable expansion and its reduct, not logical
equivalence of formulas from different signatures.

\paragraph{\hyperref[ex:fol-cardinality-normalforms-08]{\texttt{ex:fol-cardinality-normalforms-08}}.}
\textit{Solution.}  The first pair produces \(y=f(x)\) and
\(x=g(y)\), hence \(x=g(f(x))\); the occurs check proves there is no unifier.
The second pair has most general unifier \(\set{x\mapsto y}\) (equivalently,
with variables reversed, \(\set{y\mapsto x}\)).  The third produces
\(x=y\) and \(x=f(y)\), hence \(y=f(y)\); the occurs check again proves there
is no finite-term unifier.

\paragraph{\hyperref[ex:fol-cardinality-normalforms-09]{\texttt{ex:fol-cardinality-normalforms-09}}.}
\textit{Solution.}  After Skolemizing the existential premise with new constant
\(c\), the clauses are
\(\set{\neg P(x),Q(x)}\), \(\set{P(c)}\), and
\(\set{\neg Q(z)}\), with parent variables renamed apart.  Resolve the first
two using \(\set{x\mapsto c}\) to obtain \(\set{Q(c)}\).  Resolve that clause
with the third using \(\set{z\mapsto c}\) to obtain the empty clause.

\paragraph{\hyperref[ex:fol-cardinality-normalforms-10]{\texttt{ex:fol-cardinality-normalforms-10}}.}
\textit{Solution.}  A time bound truncates the search, so the result reports
only that no refutation was found among the inferences explored within that
resource limit.  Refutational completeness guarantees eventual discovery for
an unsatisfiable set only under an unbounded fair search.  It does not convert
finite unsuccessful search into a model or a satisfiability certificate.

\section{Model theory}

\paragraph{\hyperref[ex:maps-elementarity-01]{\texttt{ex:maps-elementarity-01}};
alias \hyperref[ex:diagram-test]{\texttt{ex:diagram-test}}.}
\textit{Solution.}  Interpret each new constant \(c_a\) in \(\mathcal N\) by
the same element \(a\in M\subseteq N\).  Substructure preservation of
constants and functions, and agreement of relations on tuples from \(M\), make
every true or negated atomic \(\Lang(M)\)-sentence from \(\mathcal M\) true in
the expansion of \(\mathcal N\).  The expansion satisfies the elementary
diagram iff, for every \(\Lang\)-formula \(\varphi(\bar x)\) and
\(\bar a\in M\),
\(\mathcal M\entails\varphi(\bar a)\) iff
\(\mathcal N\entails\varphi(\bar a)\), which is precisely
\(\mathcal M\preccurlyeq\mathcal N\).

\paragraph{\hyperref[ex:maps-elementarity-02]{\texttt{ex:maps-elementarity-02}}.}
\textit{Solution.}  Every formula in the language expanded by constants for
\(A\) becomes a formula of the original language with the corresponding
members of \(A\) as parameters.  Since \(\mathcal M\preccurlyeq\mathcal N\),
the two structures agree on that parameter instance.  Translating back gives
elementarity of the two constant expansions.

\paragraph{\hyperref[ex:maps-elementarity-03]{\texttt{ex:maps-elementarity-03}}.}\mbox{}\par
\noindent
\textit{Solution.}  Work in the equality-only signature.  Let
\(M_n=\set{0,1,\ldots,n}\), with inclusions as substructures.  Their union is
\(\N\), but \(M_0\) is not elementary in the union: with parameter \(0\), the
formula \(\exists x\,(x\ne0)\) is false in \(M_0\) and true in \(\N\).

\paragraph{\hyperref[ex:types-ultraproducts-01]{\texttt{ex:types-ultraproducts-01}};
alias \hyperref[ex:type-space-basis]{\texttt{ex:type-space-basis}}.}
\textit{Solution.}  A complete type contains both \(\varphi\) and \(\psi\)
iff it contains \(\varphi\land\psi\), giving the intersection
identity; it does not contain \(\varphi\) iff completeness puts \(\neg\varphi\)
in it, giving the complement identity.  For the equality-only theory of a
two-element structure, \(S_1(T)\) has one point, isolated by \([x=x]\).
The space \(S_2(T)\) has the equality type and inequality type, whose singleton
basic opens are \([x_1=x_2]\) and \([x_1\ne x_2]\), respectively.

\paragraph{\hyperref[ex:types-ultraproducts-02]{\texttt{ex:types-ultraproducts-02}};
alias \hyperref[ex:homogeneity-back-step]{\texttt{ex:homogeneity-back-step}}.}
\textit{Solution.}  The inverse \(h^{-1}:h[A]\to M\) is elementary.  Apply
\(\kappa\)-homogeneity to extend it to an elementary map
\(g:h[A]\cup\set{d}\to M\).  Elementary maps are injective, so
\(g^{-1}\) is elementary on its range; it extends \(h\), because
\(g(h(a))=a\), and its range contains \(d\), because
\(g^{-1}(g(d))=d\).  The relevant domains have size below \(\kappa\).

\paragraph{\hyperref[ex:types-ultraproducts-03]{\texttt{ex:types-ultraproducts-03}};
alias \hyperref[ex:omit-constant-type]{\texttt{ex:omit-constant-type}}.}
\textit{Solution.}  Suppose that \(\theta(x)\) were principal for
\(p(x)=\set{x\ne c_n:n\in\N}\).  Choose a model \(\mathcal M\models T\) and
\(a\in M\) with \(\mathcal M\models\theta(a)\).  If
\(a=c_k^{\mathcal M}\) for some \(k\), then this same witness shows that
\(T\not\entails\forall x(\theta(x)\to x\ne c_k)\), contrary to principality.
Otherwise \(a\ne c_n^{\mathcal M}\) for every \(n\).  The formula
\(\theta\) mentions only finitely many constants, so choose \(c_m\) that it
does not mention and form \(\mathcal M'\) by changing only the interpretation
of \(c_m\) to \(a\).  The interpretations of all constants remain pairwise
distinct, because \(a\) differed from every old constant interpretation;
hence \(\mathcal M'\models T\).  The truth of \(\theta(a)\) is unchanged,
but \(\mathcal M'\not\models a\ne c_m\), again contradicting principality.
Thus \(p\) is nonprincipal.  In the displayed model every element equals some
\(c_n\), so no element realizes \(p\).

\paragraph{\hyperref[ex:types-ultraproducts-04]{\texttt{ex:types-ultraproducts-04}}.}
\textit{Solution.}  Define
\(F([a]_{\mathcal U})=a(i_0)\).  Equality modulo an ultrafilter principal at
\(i_0\) means exactly equality at \(i_0\), so \(F\) is well-defined and
injective.  Every \(b\in M_{i_0}\) is the image of a choice function with
value \(b\) at \(i_0\), so \(F\) is surjective.  Functions and relations are
interpreted coordinatewise, hence \(F\) is an isomorphism.

\paragraph{\hyperref[ex:types-ultraproducts-05]{\texttt{ex:types-ultraproducts-05}}.}
\textit{Solution.}  In the ultrapower over \(\N\) by a nonprincipal
ultrafilter, let \(a=[i\mapsto i]_{\mathcal U}\).  For each standard \(n\),
the set \(\set{i:n<i}\) is cofinite and therefore belongs to
\(\mathcal U\).  By \L{}o\'s's theorem, \(d(n)<a\).  Thus \(a\) is greater
than every diagonal image while remaining one element of the ultrapower.

\paragraph{\hyperref[ex:qe-categoricity-01]{\texttt{ex:qe-categoricity-01}}.}
\textit{Solution sketch.}  Enumerate the domains of a countable dense order
without endpoints \(M\) and of \(\Q\).  Starting from the empty partial
isomorphism, alternately map the next unused element of \(M\) into the unique
open interval determined by its already mapped neighbors, and map the next
unused rational back into the corresponding interval of \(M\).  Density and
the absence of endpoints supply a fresh point in every step.  The union is
total, surjective, and order preserving and reflecting.  Check all four interval
cases: two bounds, lower bound only, upper bound only, and no bounds.

\paragraph{\hyperref[ex:qe-categoricity-02]{\texttt{ex:qe-categoricity-02}}.}
\textit{Solution.}  Given \(\varphi(\bar x)\), choose quantifier-free
\(\psi(\bar x)\) with
\(T\entails\forall\bar x(\varphi\leftrightarrow\psi)\).  Since both
structures model \(T\), each agrees with \(\psi\) on \(\varphi\).  A
substructure and its superstructure agree on quantifier-free formulas with
parameters from the substructure, so they agree on \(\varphi\).  If either
structure fails to model \(T\), the relative equivalence to \(\psi\) is not
available there, and the argument stops.

\paragraph{\hyperref[ex:qe-categoricity-03]{\texttt{ex:qe-categoricity-03}};
alias \hyperref[ex:back-and-forth-equivalence]{\texttt{ex:back-and-forth-equivalence}}.}
\textit{Solution sketch.}  Enumerate two countable models and maintain a finite
partial bijection preserving the equivalence relation.  To add an element
equivalent to an old domain element, choose a fresh member of the corresponding
infinite target class.  To add an element in a new class, choose a fresh target
class, possible because infinitely many classes remain.  Apply the symmetric
construction on alternate steps.  The union is a total isomorphism; checking
the two extension cases proves preservation and reflection of equivalence.

\section{Computability and arithmetical hierarchy}

\paragraph{\hyperref[ex:turing-machines-01]{\texttt{ex:turing-machines-01}}.}
\textit{Solution.}  Complement a decider by exchanging its two halting answers.
For finite union or intersection, run the finitely many deciders and combine
their answers by disjunction or conjunction.  For computably enumerable sets,
dovetail the recognizers: accept a union when one accepts, and accept an
intersection after every recognizer has accepted its common input.  No
complement construction is used, because a recognizer's failure to accept need
not halt.

\paragraph{\hyperref[ex:universal-reductions-01]{\texttt{ex:universal-reductions-01}}.}
\textit{Solution.}  Each property depends only on the computed partial
function, not on its program text, and is nontrivial.  The everywhere-undefined
function is not total while the zero function is; the former's range omits
\(0\) while the latter's contains \(0\).  Under the exercise's stated
nonempty-domain convention for a constant partial function, the zero function
has the third property and the everywhere-undefined function does not.  Rice's
theorem therefore makes all three index sets undecidable.

\paragraph{\hyperref[ex:decision-hierarchy-01]{\texttt{ex:decision-hierarchy-01}}.}
\textit{Solution sketch.}  With the fixed bounded arithmetic computation
formula \(\mathsf{Comp}(e,x,s)\), nonempty domain is defined by
\(\exists x\exists s\,\mathsf{Comp}(e,x,s)\), a \(\Sigma^0_1\) formula.  Given \(e\),
effectively produce a program that is undefined except that on input \(0\) it
simulates the normalized wrapper for program \(e\) on input \(e\), accepts
with a valid numerical output iff \(\varphi_e(e)\) is defined, and diverges
otherwise.
Its domain is nonempty iff \(e\in K\).  The uniform program-index map
is the required many-one reduction; writing it via the parameter theorem
completes the formal check.  Since \(K\) is
\(\boldsymbol{\Sigma}^0_1\)-complete under many-one reductions, this proves
the same completeness classification for the nonempty-domain index set.

\paragraph{\hyperref[ex:decision-hierarchy-02]{\texttt{ex:decision-hierarchy-02}}.}
\textit{Solution.}  Using the same bounded arithmetic formula
\(\mathsf{Comp}\), finite domain has
the immediate form
\(\exists b\forall x>b\forall s\,\neg\mathsf{Comp}(e,x,s)\), hence belongs to
\(\boldsymbol{\Sigma}^0_2\).  Infinite domain has form
\(\forall b\exists x>b\exists s\,\mathsf{Comp}(e,x,s)\), hence belongs to
\(\boldsymbol{\Pi}^0_2\).
Cofinite domain has form
\(\exists b\forall x>b\exists s\,\mathsf{Comp}(e,x,s)\), hence belongs to
\(\boldsymbol{\Sigma}^0_3\).
These are upper bounds only; no completeness claim follows without reductions.

\paragraph{\hyperref[ex:decision-hierarchy-03]{\texttt{ex:decision-hierarchy-03}}.}
\textit{Solution.}  Enumerate all finite derivations in a fixed effective,
sound, and complete first-order calculus and output their conclusions.  Every
output is valid by soundness.  Conversely, completeness is used exactly to
infer that every semantically valid sentence has some finite derivation and
therefore eventually appears.  The enumeration never ranges directly over
structures.

\section{Arithmetic and incompleteness}

\paragraph{\hyperref[ex:arithmetic-coding-01]{\texttt{ex:arithmetic-coding-01}};
alias \hyperref[ex:nonstandard-true-arithmetic]{\texttt{ex:nonstandard-true-arithmetic}}.}
\textit{Solution.}  Taking \(T=\Th(\mathcal N)\), the compactness construction
produces a countable \(\mathcal M\models T\) with an element unequal to every
numeral.  Since \(T\) contains exactly the sentences true in \(\mathcal N\),
\(\mathcal M\) and \(\mathcal N\) agree on every arithmetic sentence.  They
are not isomorphic, because every element of \(\mathcal N\) is denoted by a
numeral and the displayed element of \(\mathcal M\) is not.  Thus elementary
equivalence is strictly weaker than isomorphism.

\paragraph{\hyperref[ex:arithmetic-coding-02]{\texttt{ex:arithmetic-coding-02}}.}
\textit{Solution sketch.}  Fix primitive recursive coding of finite sequences
with length and component functions.  Let \(A(c,a,b)\) say: \(c\) has length
\(b+1\), its zeroth component is \(a\), and for every \(i<b\) its
\((i+1)\)-st component is the successor of its \(i\)-th component.  Bounded
quantification over components keeps \(A\) primitive recursive.  The last
component is \(a+b\); representability of \(A\), together with existence and
uniqueness of the coded recursion, yields a representing formula for addition.
A completed account must verify the sequence-code operations and the unique
last value.

\paragraph{\hyperref[ex:arithmetic-coding-03]{\texttt{ex:arithmetic-coding-03}}.}
\textit{Solution.}  The first assertion is a metalinguistic claim that a
\(T\)-derivation of \(\sigma\) exists.  The second is a semantic assertion in
the standard arithmetic structure about the numeral coding \(\sigma\).  The
third is again a metalinguistic derivability claim, now with an arithmetized
provability sentence as conclusion.  Correct representation of the proof
predicate connects an actual proof with standard truth of the second assertion;
if \(T\supseteq\mathsf Q\), numeralwise representability gives the first
derivability condition and connects an actual proof to the third.  The reverse
move from \(T\)-provability to standard truth requires the corresponding
soundness hypothesis.

\paragraph{\hyperref[ex:first-incompleteness-01]{\texttt{ex:first-incompleteness-01}}.}
\textit{Solution sketch.}  Effectiveness makes finite \(T\)-proof checking
primitive recursive and \(\mathsf{Prov}_T\) arithmetically expressible.
Representability turns a concrete proof into a \(T\)-proof of its proof-code
instance and verifies each standard nonproof code.  Consistency rules out
\(T\derives G_T\), since the fixed point would then yield both provability and
its negation.  \(\Sigma^0_1\)-soundness rules out
\(T\derives\neg G_T\), because \(\neg G_T\) is a false
\(\Sigma^0_1\) sentence once consistency shows \(G_T\) true in
\(\mathcal N\).  These four uses should appear at those exact transitions.

\paragraph{\hyperref[ex:first-incompleteness-02]{\texttt{ex:first-incompleteness-02}}.}
\textit{Solution.}  Consistency gives \(T\notderives G_T\).  Hence no standard
number codes a \(T\)-proof of \(G_T\), so
\(\mathcal N\entails\neg\mathsf{Prov}_T(\ulcorner G_T\urcorner)\).  The fixed
point is already provable in \(\mathsf Q\), which is true in \(\mathcal N\),
so \(\mathcal N\entails G_T\).  This semantic conclusion alone does not show
that \(T\cup\set{G_T}\) is consistent: that additionally requires
\(T\notderives\neg G_T\).  Under, for example, the stated
\(\Sigma^0_1\)-soundness hypothesis, that condition holds; then the one-sentence
extension is effective and consistent, so an incompleteness theorem supplies
another sentence it does not decide.  Thus adjoining \(G_T\) does not produce
all arithmetical truths under the strengthened hypothesis, while mere
consistency licenses no such claim about the extension.

\paragraph{\hyperref[ex:first-incompleteness-03]{\texttt{ex:first-incompleteness-03}}.}
\textit{Solution sketch.}  Once a proof with least standard code \(m\) is
assumed, every comparison with a code below \(m\) is a bounded finite check
whose true and false instances are numeralwise representable in \(\mathsf Q\);
\cref{lem:q-bounded-cases} turns the bounded existential over codes at most
\(m\) into a finite disjunction, and \cref{lem:q-dichotomy} splits an
arbitrary proof code against \(\bar m\) inside \(\mathsf Q\).
The Rosser sentence uses that bounded comparison to force an opposite proof
with no larger code, contradicting consistency.  Omega-consistency is not
needed because the argument never infers a universal standard claim from all
of its numeral instances; the least code converts the relevant search into a
single bounded verification.

\paragraph{\hyperref[ex:second-incompleteness-01]{\texttt{ex:second-incompleteness-01}}.}
\textit{Solution sketch.}  From the fixed point
\(G_T\leftrightarrow\neg\mathsf{Prov}_T(\ulcorner G_T\urcorner)\), internalize
the implication \(G_T\to\neg\mathsf{Prov}_T(\ulcorner G_T\urcorner)\) using
derivability conditions 1 and 2.  Condition 3 gives
\(\mathsf{Prov}_T(\ulcorner G_T\urcorner)\to
\mathsf{Prov}_T(\ulcorner
  \mathsf{Prov}_T(\ulcorner G_T\urcorner)\urcorner)\); internalized
modus ponens then yields
\(\mathsf{Prov}_T(\ulcorner G_T\urcorner)\to
\mathsf{Prov}_T(\ulcorner\bottom_A\urcorner)\).  Therefore
\(\mathsf{Con}(T)\) implies
\(\neg\mathsf{Prov}_T(\ulcorner G_T\urcorner)\), and the fixed point
gives \(\mathsf{Con}(T)\to G_T\).  If \(T\) proved its consistency, it would
prove \(G_T\), contradicting consistency.  A full derivation must display the
instances of all three conditions just identified.

\paragraph{\hyperref[ex:second-incompleteness-02]{\texttt{ex:second-incompleteness-02}}.}\mbox{}\par
\noindent
\textit{Solution.}  The assumed biconditional includes
\(T\derives\mathsf{Prov}_T(\ulcorner\varphi\urcorner)\to\varphi\).  This is
exactly the hypothesis of L\"ob's theorem, so
\(T\derives\varphi\).  The reverse half of the biconditional is not needed.

\paragraph{\hyperref[ex:second-incompleteness-03]{\texttt{ex:second-incompleteness-03}}.}
\textit{Solution.}  Suppose one formula \(\mathsf{Tr}(x)\) had all the stated
biconditionals provable in consistent \(T\supseteq\mathsf Q\).  By the diagonal
lemma choose \(L\) with
\(T\derives L\leftrightarrow
\neg\mathsf{Tr}(\overline{\ulcorner L\urcorner})\).  The
assumed instance for \(L\) gives
\(T\derives\mathsf{Tr}(\overline{\ulcorner L\urcorner})\leftrightarrow L\).
Propositional reasoning combines them to derive \(L\leftrightarrow\neg L\),
and hence a contradiction.  Consistency is used only at the final step to rule
out that derivation.

\paragraph{\hyperref[ex:second-incompleteness-04]{\texttt{ex:second-incompleteness-04}}.}
\textit{Solution.}  Three correct replacements are: (i) if \(T\) is a
consistent computably enumerable extension of \(\mathsf Q\), then its fixed
G\"odel sentence \(G_T\) is not provable in \(T\), and it is true in
\(\mathcal N\); (ii) if in addition the fixed provability predicate satisfies
the derivability conditions, then \(T\notderives\mathsf{Con}(T)\); and
(iii) \(\Th(\mathcal N)\) contains every true arithmetic sentence by
definition, is deductively complete, and is not computably enumerable.  These
claims concern different theories, predicates, and evidence.

\section{Set theory and forcing}

\paragraph{\hyperref[ex:zf-zfc-01]{\texttt{ex:zf-zfc-01}}.}
\textit{Solution sketch.}  Form \(a\cup b\) using Pairing and Union.  For
\(x\in a,y\in b\), the Kuratowski pair
\(\langle x,y\rangle=\set{\set{x},\set{x,y}}\) belongs to
\(\Pow(\Pow(a\cup b))\).  This double power set is the required bounding set.
Apply Separation to it with the formula
\(\exists x\in a\exists y\in b\,z=\langle x,y\rangle\).  Extensionality
identifies the separated set with \(a\times b\).  A complete derivation should
cite Pairing, Union, Power Set twice, Separation, and Extensionality at these
exact points.

\paragraph{\hyperref[ex:zf-zfc-02]{\texttt{ex:zf-zfc-02}}.}
\textit{Solution.}  By Pairing form the nonempty set \(c=\set{a,b}\).
Foundation supplies \(x\in c\) with \(x\cap c=\varnothing\).  If \(x=a\),
then \(b\in a\cap c\); if \(x=b\), then \(a\in b\cap c\).  Both cases
contradict the supplied disjointness, so no two-cycle \(a\in b\in a\) exists.

\paragraph{\hyperref[ex:cumulative-hierarchy-01]{\texttt{ex:cumulative-hierarchy-01}}.}
\textit{Solution.}  Fix \(\beta\) and argue directly from transitivity;
no transfinite induction is required.  If
\(\alpha=\beta\), inclusion is immediate.  If \(\alpha<\beta\), then
\(\alpha\in\beta\) by the ordinal order definition; since every ordinal
\(\beta\) is transitive, every member of \(\alpha\) is also a member of
\(\beta\), so \(\alpha\subseteq\beta\).  The use of transitivity is exactly the
step from \(u\in\alpha\in\beta\) to \(u\in\beta\).  The argument
invokes no induction hypothesis, which is why no transfinite induction is
required.

\paragraph{\hyperref[ex:cumulative-hierarchy-02]{\texttt{ex:cumulative-hierarchy-02}}.}
\textit{Solution.}  Put \(X=\set{\varnothing,\set\varnothing}\).  Since
\(\rank(\varnothing)=0\) and \(\rank(\set\varnothing)=1\),
\[
\rank(X)=\sup\set{0+1,1+1}=2.
\]
Also \(X\subseteq V_2\), so \(X\in V_3\), while
\(\set\varnothing\notin V_1\), so \(X\notin V_2\).  This confirms the rank
calculation from the hierarchy stages.

\paragraph{\hyperref[ex:cumulative-hierarchy-03]{\texttt{ex:cumulative-hierarchy-03}}.}
\textit{Solution.}  Let \(X=\set{a,b}\) with \(a\ne b\) and let
\(R=\varnothing\).  The relation is well-founded but not extensional, because
\(a\) and \(b\) have the same predecessor set.  The collapse recursion sends
both to \(\set{}=\varnothing\), so the collapse map is not injective.

\paragraph{\hyperref[ex:constructibility-forcing-01]{\texttt{ex:constructibility-forcing-01}}.}
\textit{Solution.}  For a finite partial function \(p\), if
\(n\in\dom(p)\) then \(p\in D_n\); otherwise
\(p\cup\set{(n,0)}\) extends \(p\) and belongs to \(D_n\).  Thus \(D_n\) is
dense.  If \((n,i)\in p\in G\) and \((n,j)\in q\in G\), directedness of the
filter gives \(r\in G\) extending both \(p,q\).  Since \(r\) is a function,
\(i=j\).  Hence \(\bigcup G\) is a function; meeting every \(D_n\) further
makes its domain all of \(\omega\).

\paragraph{\hyperref[ex:constructibility-forcing-02]{\texttt{ex:constructibility-forcing-02}}.}
\textit{Solution sketch.}  \(L\entails\mathsf{CH}\) is an external semantic
claim about the constructible universe, proved from the fine structure and
definable well-order of \(L\).  The formula
\(\ZFC\derives V=L\to\mathsf{CH}\) is a syntactic derivability claim obtained
by formalizing that result in \(\ZFC\).  The implication
\(\mathsf{Con}(\ZFC)\to\mathsf{Con}(\ZFC+\mathsf{CH})\) is a
metatheoretic relative-consistency claim: from a model or syntactic consistency
construction for \(\ZFC\), pass to its constructible inner model, which
satisfies \(V=L\) and hence \(\mathsf{CH}\).  A completed proof must state the
soundness/completeness or syntactic formalization used when moving between
model existence and consistency.

\section{Sequent calculus, normalization, and typed proof terms}

\paragraph{\hyperref[ex:lk-lj-01]{\texttt{ex:lk-lj-01}}.}
\textit{Solution sketch.}  Apply \(\land L\) to the antecedent, making
\(A\) and \(B\lor C\) available; then apply \(\lor L\) to \(B\lor C\).  In
the \(B\)-branch, identities for \(A,B\) and \(\land R\) yield
\(A\land B\), followed by the left \(\lor R\) injection.  In the
\(C\)-branch, the analogous steps yield \(A\land C\), followed by the right
injection.  The principal formulas, from bottom upward, are
\(A\land(B\lor C)\), \(B\lor C\), the concluding disjunction, and the
relevant conjunction.  The finished tree contains no cut and has one succedent at every
node.

\paragraph{\hyperref[ex:lk-lj-02]{\texttt{ex:lk-lj-02}}.}
\textit{Solution.}  Without the eigenvariable condition, identity
\(P(y)\Rightarrow P(y)\) would permit \(\forall R\) to infer
\(P(y)\Rightarrow\forall xP(x)\).  Take domain \(\set{0,1}\), interpret
\(P\) as \(\set{0}\), and assign \(y\) to \(0\).  The antecedent is true under
that assignment and the succedent is false.  Thus the relaxed rule derives an
invalid sequent.

\paragraph{\hyperref[ex:cut-elimination-01]{\texttt{ex:cut-elimination-01}}.}
\textit{Solution sketch.}  Suppose the left cut premise obtains
\(A\land B\) by \(\land R\) from derivations of \(A\) and \(B\), while the
right cut premise obtains its sequent by \(\land L\) from a derivation with
both \(A,B\) in the antecedent.  First cut the derivation of \(B\) against the
right premise to remove \(B\), retaining \(A\); then cut the derivation of
\(A\) against that result.  Context contractions or exchanges are applied as
licensed by the active calculus.  Both new cut formulas, \(A\) and \(B\), are
proper subformulas of \(A\land B\), so each has strictly smaller cut rank.

\paragraph{\hyperref[ex:cut-elimination-02]{\texttt{ex:cut-elimination-02}}.}
\textit{Solution.}  In propositional \(\LJ\), a cut-free backwards search uses
only finitely many subformulas of the endsequent; with the declared context
representation, only finitely many relevant sequents can arise, so exhaustive
search with repetition control decides derivability.  In first-order
\(\LJ\), quantifier rules introduce arbitrarily many term instances, and the
set of possible terms is generally infinite.  Literal subformula finiteness
therefore does not by itself bound proof search.

\paragraph{\hyperref[ex:normalization-curry-howard-01]{\texttt{ex:normalization-curry-howard-01}}.}
\textit{Solution.}  The term has type \(A\to(B\to A)\).  In the typing
calculus, the hypothesis \(y:B\) enters the context and is never consumed
by the \textsc{Var} rule; that unused context entry is the structural
weakening implicit in \(\lambda x{:}A.\lambda y{:}B.x\).  In the tree
natural-deduction calculus, a discharging application must discharge at
least one open occurrence, so the corresponding derivation of
\(A\to(B\to A)\) routes the assumption \(B\) through a conjunction
detour: assume \(A\) and \(B\), form \(B\land A\) by \(\land I\), recover
\(A\) by \(\land E_2\), discharge \(B\) to obtain \(B\to A\), and
discharge \(A\).  Both discharged occurrences are genuinely used by
\(\land I\).

\paragraph{\hyperref[ex:normalization-curry-howard-02]{\texttt{ex:normalization-curry-howard-02}}.}
\textit{Solution sketch.}  The required typing substitution lemma is: if
\(\Theta_{\mathrm{ty}},x{:}A\derives M:B\) and
\(\Theta_{\mathrm{ty}}\derives N:A\), then
\(\Theta_{\mathrm{ty}}\derives M[N/x]:B\), after capture-avoiding
renaming of bound variables.  Typing inversion on
\(\Theta_{\mathrm{ty}}\derives(\lambda x{:}A.M)N:B\) gives its two
premises; the lemma then types the reduct \(M[N/x]\) by \(B\).  Capture
avoidance is used in the lambda-abstraction case of the lemma when the
binder occurs free in \(N\).
Checking that case and the variable case completes the outer-redex proof.

\paragraph{\hyperref[ex:normalization-curry-howard-03]{\texttt{ex:normalization-curry-howard-03}}.}
\textit{Solution.}  If every reduction sequence from a term terminates, then in
particular at least one terminating sequence exists, so strong normalization
implies weak normalization.  For the converse failure, take a relation with
\(a\to a\) and \(a\to b\), where \(b\) has no successor.  The object \(a\) has
the terminating path \(a\to b\), but also the infinite path
\(a\to a\to a\to\cdots\).

\section{Second-order, typed, intuitionistic, and modal systems}

\paragraph{\hyperref[ex:second-order-logic-01]{\texttt{ex:second-order-logic-01}}.}
\textit{Solution.}  Let \(X\subseteq\N\) be arbitrary, not assumed definable,
and suppose \(0\in X\) and \(n\in X\Rightarrow n+1\in X\).  External induction
on \(n\) gives \(n\in X\) for every \(n\in\N\), hence \(X=\N\).  Since full
second-order semantics ranges its unary-relation variables over every subset of
\(\N\), this verifies the full induction axiom in the standard structure.

\paragraph{\hyperref[ex:second-order-logic-02]{\texttt{ex:second-order-logic-02}}.}
\textit{Solution.}  The categoricity proof forms the range of the recursively
defined map \(f:\N\to M\) and instantiates second-order induction with that
range to show it is all of \(M\).  In a Henkin structure, the unary-relation
domain may omit \(\operatorname{ran}(f)\).  The induction axiom then cannot be
instantiated with that subset, so the surjectivity step fails.

\paragraph{\hyperref[ex:second-order-logic-03]{\texttt{ex:second-order-logic-03}}.}
\textit{Solution sketch.}  Any finitary strongly complete calculus would make
full second-order consequence compact: from \(\Gamma\entails\varphi\), strong
completeness gives a finite derivation, which uses only a finite
\(\Gamma_0\subseteq\Gamma\), so \(\Gamma_0\entails\varphi\) by soundness.  The
displayed full-semantics counterexample to compactness contradicts this.  This
argument excludes finitary sound strong completeness; the stronger claim that
full second-order validities are not computably enumerable requires a separate
effectivity theorem.

\paragraph{\hyperref[ex:simple-type-theory-01]{\texttt{ex:simple-type-theory-01}}.}
\textit{Solution.}  Here
\(R:\iota\to\iota\to o\), \(x:\iota\),
\(Rx:\iota\to o\), and \(Rxx:o\).  Therefore
\(\lambda x{:}\iota.Rxx:\iota\to o\), and the whole term has type
\((\iota\to\iota\to o)\to(\iota\to o)\).

\paragraph{\hyperref[ex:simple-type-theory-02]{\texttt{ex:simple-type-theory-02}}.}
\textit{Solution.}  Under standard semantics,
\(\forall f{:}(\alpha\to\beta).A\) is true under assignment \(s\) iff
\(A\) is true under \(s[f\mapsto g]\) for every function
\(g:D_\alpha\to D_\beta\).  Under Henkin semantics the same clause quantifies
only over the designated domain \(D_{\alpha\to\beta}\), whose members need not
be all such functions.  The quantifier range is the only difference in the two
clauses.

\paragraph{\hyperref[ex:simple-type-theory-03]{\texttt{ex:simple-type-theory-03}}.}
\textit{Solution.}  A precise true statement is: the standard typed deductive
calculus is sound and complete for its declared class of Henkin structures.  A
precise false statement is: some finitary effective sound calculus is strongly
complete for standard full higher-order semantics.  Thus the bare word
``complete'' omits both the calculus and the semantic class and can conceal a
change from Henkin to full semantics.

\paragraph{\hyperref[ex:intuitionistic-logic-01]{\texttt{ex:intuitionistic-logic-01}}.}
\textit{Solution.}  Let \(w<v\), and let \(\mathcal M\) make \(A\) true only
at \(v\) and \(B\) true nowhere.  At both \(w\) and \(v\), \(A\to B\) fails
because \(\mathcal M,v\entails A\) and \(\mathcal M,v\notentails B\).
Hence \(\mathcal M,w\entails(A\to B)\to A\) vacuously, while
\(\mathcal M,w\notentails A\).  Peirce's formula therefore fails at \(w\).

\paragraph{\hyperref[ex:intuitionistic-logic-02]{\texttt{ex:intuitionistic-logic-02}}.}
\textit{Solution.}  If \(\mathcal M,w\entails\neg(A\lor B)\), then no
extension of \(w\) satisfies \(A\lor B\); consequently no extension satisfies
\(A\), and none satisfies \(B\).  Thus
\(\mathcal M,w\entails\neg A\land\neg B\).  Conversely, if
\(\mathcal M,w\entails\neg A\land\neg B\), an extension satisfying
\(A\lor B\) would satisfy one disjunct and contradict the corresponding
negation.  Hence the converse also holds intuitionistically.

\paragraph{\hyperref[ex:modal-temporal-epistemic-01]{\texttt{ex:modal-temporal-epistemic-01}}.}
\textit{Solution.}  Let \(W=\set{w,u}\) and
\(R=\set{(w,w),(w,u),(u,u)}\).  This frame is reflexive and is not Euclidean,
since \(wRw\) and \(wRu\) but not \(uRw\).  Let \(\mathcal M\) be the
structure on this frame that makes \(p\) true only at \(w\).
Then \(\mathcal M,w\entails\Diamond p\), but
\(\mathcal M,u\notentails\Diamond p\); because \(wRu\),
\(\mathcal M,w\notentails\Box\Diamond p\).

\paragraph{\hyperref[ex:modal-temporal-epistemic-02]{\texttt{ex:modal-temporal-epistemic-02}}.}
\textit{Solution.}  At positions \(0,1,2\), choose position \(3\) as the
witness for \(q\); \(p\) holds at every earlier position, so
\(p\mathbin{\mathsf U}q\) is true.  It is also true at position \(3\), with
the current position as witness and an empty intervening interval.  Its value
at position \(4\) is undetermined from the data: one must specify the truth of
\(p,q\) there and on the infinite continuation, because a past occurrence of
\(q\) cannot witness a future-until formula.

\paragraph{\hyperref[ex:modal-temporal-epistemic-03]{\texttt{ex:modal-temporal-epistemic-03}}.}
\textit{Solution.}  Take agents \(a,b\) and worlds \(w,u,v\).  Let the
\(a\)-equivalence classes be \(\set{w,u}\), \(\set v\), and the
\(b\)-classes be \(\set w\), \(\set{u,v}\).  Make \(p\) true at \(w,u\) and
false at \(v\).  At \(w\), both agents know \(p\), but the two-step path
\(w\sim_a u\sim_b v\) reaches a world where \(p\) is false.  Therefore
\(C_{\set{a,b}}p\) fails at \(w\).

\section{Many-valued, relevance, and linear logic}

\paragraph{\hyperref[ex:many-valued-paraconsistent-01]{\texttt{ex:many-valued-paraconsistent-01}}.}
\textit{Solution sketch.}  For \(\mathbf K_3\), the two roots are
\((p\land q){:}1,p{:}0\) and
\((p\land q){:}1,p{:}\tfrac12\).  Expanding the conjunction labelled \(1\)
forces \(p{:}1,q{:}1\), so each root closes on conflicting labels for \(p\).
For \(\mathsf{LP}\), test designated premise labels \(1,\tfrac12\) against
the sole nondesignated conclusion label \(0\).  Every branch of the
\(\min\)-expansion of a designated conjunction gives \(p\) a designated
value, so every branch again conflicts with \(p{:}0\).  Listing all
\(\min(a,b)=\tfrac12\) input pairs completes the LP tree.

\paragraph{\hyperref[ex:many-valued-paraconsistent-02]{\texttt{ex:many-valued-paraconsistent-02}}.}
\textit{Solution.}  Use \(p\lor\neg p\).  At values \(0,\tfrac12,1\), it has
values \(1,\tfrac12,1\), respectively.  All three are designated in
\(\mathsf{LP}\), whereas \(\tfrac12\) is not designated in
\(\mathbf K_3\).  Hence the formula is LP-valid but not K3-valid.

\paragraph{\hyperref[ex:relevance-linear-01]{\texttt{ex:relevance-linear-01}}.}
\textit{Solution.}  Write \(g\) for the tensor of \(\Gamma\).  The premise of
\(\multimap R\) is \(g\otimes a\leq b\), which is equivalent by residuation to
\(g\leq a\multimap b\), its conclusion.  For \(\multimap L\), suppose
\(g\leq a\) and \(d\otimes b\leq c\).  Since
\(a\otimes(a\multimap b)\leq b\), monotonicity gives
\(d\otimes g\otimes(a\multimap b)\leq d\otimes b\leq c\), exactly the
conclusion sequent's inequality.

\paragraph{\hyperref[ex:relevance-linear-02]{\texttt{ex:relevance-linear-02}}.}
\textit{Solution.}  Use identities \(A\Rightarrow A\) and
\(B\Rightarrow B\) as the two premises of \(\multimap L\).  It yields
\(A,A\multimap B\Rightarrow B\), and exchange built into multiset contexts
gives \(A\multimap B,A\Rightarrow B\).  Every antecedent occurrence is consumed
by the rule, so no weakening is required.

\paragraph{\hyperref[ex:relevance-linear-03]{\texttt{ex:relevance-linear-03}}.}
\textit{Solution.}  Sequence contexts build in none of exchange, weakening, or
contraction.  Multiset contexts identify permutations and thus build in
exchange, but retain multiplicity and do not build in weakening or contraction.
Set contexts ignore order and multiplicity, so exchange and contraction are
built in; weakening is still a substantive monotonicity principle because a
set does not identify \(\Gamma\) with \(\Gamma\cup\set A\).

\paragraph{\hyperref[ex:relevance-linear-04]{\texttt{ex:relevance-linear-04}}.}
\textit{Solution.}  The decisive clause says that \(A\to B\) at \(w\) tests
every pair \(u,v\) with \(R(w,u,v)\), requiring truth of \(A\) at \(u\) to
yield truth of \(B\) at the potentially different world \(v\).  In the
displayed structure, \(R(a,u,v)\), \(q\) holds at \(u\), and \(p\) fails at
\(v\), so \(q\to p\) fails at \(a\), which then defeats the outer
conditional at the normal world.  An ordinary binary Kripke clause evaluates
antecedent and consequent at one accessible world and therefore does not encode
this ternary split between the two tested worlds.

\section{Formal practice, automation, mechanization, and typesetting}

\paragraph{\hyperref[ex:formal-practice-01]{\texttt{ex:formal-practice-01}}.}
\textit{Solution.}  On a one-element domain \(\set{a}\), the premise makes
\(P(a)\lor Q(a)\) true, so one disjunct is true; that disjunct then holds of
every domain element, making one universal true.  For a two-element
countermodel, take domain \(\set{a,b}\), let \(P=\set{a}\) and
\(Q=\set{b}\).  Every element satisfies \(P\lor Q\), but neither predicate
holds of every element, so the conclusion is false.

\paragraph{\hyperref[ex:formal-practice-02]{\texttt{ex:formal-practice-02}}.}
\textit{Solution sketch.}  Induct structurally on the formula, after proving
the term-substitution lemma by induction on terms.  Strengthen the claim to all
structures, assignments, variables, and replacing terms at once:
\(\mathcal M,s\entails\varphi[t/x]\) iff
\(\mathcal M,s[x\mapsto\sem{t}^{\mathcal M}_s]\entails\varphi\).  Atomic and
Boolean cases use the corresponding induction hypotheses.  In
\(\forall y\psi\) and \(\exists y\psi\), if \(y=x\) substitution stops below
the binder; if \(y\ne x\) but \(y\in\operatorname{Var}(t)\), first rename
\(y\) fresh.  A
completed proof must show the two assignment updates commute after that
renaming.

\paragraph{\hyperref[ex:automated-reasoning-01]{\texttt{ex:automated-reasoning-01}}.}
\textit{Solution.}  The unit clause \(\set{\neg R}\) sets \(R=\False\).
Then \(\set{\neg P,R}\) and \(\set{\neg Q,R}\) reduce to units
\(\set{\neg P}\) and \(\set{\neg Q}\), setting \(P=Q=\False\).  The first
clause \(\set{P,Q}\) then becomes empty.  DPLL returns \(\UNSAT\).

\paragraph{\hyperref[ex:automated-reasoning-02]{\texttt{ex:automated-reasoning-02}}.}
\textit{Solution.}  Decomposing the two atoms gives
\(x\doteq g(y)\) and \(x\doteq y\).  Substituting the second equation into the
first gives \(y\doteq g(y)\).  Since \(y\) occurs properly in \(g(y)\), the
occurs check fails, and no finite-term unifier exists.

\paragraph{\hyperref[ex:automated-reasoning-03]{\texttt{ex:automated-reasoning-03}}.}
\textit{Solution.}  Before Skolemization, equivalence-preserving steps give
\[
\neg\forall x(P(x)\to\exists yR(x,y))
\equiv\exists x\bigl(P(x)\land\forall y\neg R(x,y)\bigr).
\]
Replace the existential witness by a new constant \(c\).  This step preserves
only equisatisfiability between the original signature and suitable expansions.
The remaining universal formula is equivalent in the expanded signature to the
universal closures of the clauses \(\set{P(c)}\) and
\(\set{\neg R(c,y)}\).  Thus only the Skolem step weakens equivalence to
equisatisfiability.

\paragraph{\hyperref[ex:mechanization-01]{\texttt{ex:mechanization-01}}.}
\textit{Solution.}  Suitable terms are
\[
\lambda x{:}A.x,\qquad
\lambda p{:}(A\times B).\pi_1p,\qquad
\lambda f{:}(A\to B).\lambda g{:}(B\to C).\lambda x{:}A.g(fx).
\]
They have types \(A\to A\), \((A\times B)\to A\), and
\((A\to B)\to(B\to C)\to(A\to C)\), respectively, by variable,
product-elimination, and function-elimination typing rules.

\paragraph{\hyperref[ex:mechanization-02]{\texttt{ex:mechanization-02}}.}
\textit{Solution sketch.}  In a deep embedding, let
\(\mathsf{Form}_{d}\) be an inductive datatype with constructor
\(\mathsf{And}:\mathsf{Form}_{d}\to\mathsf{Form}_{d}\to\mathsf{Form}_{d}\),
and define \(\mathsf{eval}_{d}(A,v)\) by recursion.  The required adequacy
theorem is
\[
  \mathsf{eval}_{d}(A,v)=\mathsf{true}
  \quad\Longleftrightarrow\quad
  v\entails_{\mathrm{CPL}} A
\]
for every object-language formula \(A\) and valuation \(v\); its conjunction
case uses
\(\mathsf{eval}_{d}(\mathsf{And}(A,B),v)=
 \mathsf{eval}_{d}(A,v)\land_{\mathrm{Bool}}\mathsf{eval}_{d}(B,v)\).

For a shallow embedding, fix for each valuation \(v\) a translation
\(\llbracket-\rrbracket^{s}_{v}\) from object-language formulas to host
propositions and stipulate
\[
  \llbracket A\land B\rrbracket^{s}_{v}
  \defeq
  \llbracket A\rrbracket^{s}_{v}\land_{\mathrm{host}}
  \llbracket B\rrbracket^{s}_{v}.
\]
Its required adequacy theorem is
\[
  \llbracket A\rrbracket^{s}_{v}
  \quad\Longleftrightarrow\quad
  v\entails_{\mathrm{CPL}}A
\]
for every \(A\) and \(v\).  The deep representation retains syntax for
induction, substitution, and quotation; the shallow representation supplies
those operations only if additional encoded syntax is introduced.

\paragraph{\hyperref[ex:latex-style-01]{\texttt{ex:latex-style-01}}.}
\textit{Solution.}  Replace ``Let \(f\) be the inclusion.  Clearly atomic
formulas are preserved.  Clearly the result follows by induction'' with:
``Let \(f:M\hookrightarrow N\) be the inclusion map of the
\(\Lang\)-substructure \(\mathcal M\subseteq\mathcal N\).  For every function
symbol \(g\), substructure closure gives
\(f(g^{\mathcal M}(\bar a))=g^{\mathcal N}(f\bar a)\); for every relation
symbol \(R\), the substructure definition gives
\(R^{\mathcal M}(\bar a)\) iff \(R^{\mathcal N}(f\bar a)\).  These are the
atomic induction cases.  The Boolean constructor cases follow from their truth
clauses and the induction hypotheses.''  The typed map and the two omitted
justifications are now explicit.

\paragraph{\hyperref[ex:latex-style-02]{\texttt{ex:latex-style-02}}.}
\textit{Solution.}  For a first-order resolution chapter, register
\(x,y,z\) as object variables, \(r,t\) as terms,
\(A,B\) as atoms, \(C,D\) as finite literal clauses, and
\(\theta_{\mathrm u},\rho_{\mathrm u},\eta_{\mathrm u}\) as finite-support
term substitutions.  Register \(\operatorname{ucl}_{\Lang}(C)\) as the formula
interpretation of a clause.  Do not also use \(C\) for a calculus or
\(\theta_{\mathrm u}\) for a formula in that chapter.  If \(\entails\) occurs,
its argument shapes must separately declare structure satisfaction and clause
set consequence, as in the registered-notation appendix.

\section{Keys for the qualifying problems}

The five labelled qualifying problems retain their existing solutions.  This
list makes their keys part of the same deterministic lookup system.

\paragraph{\hyperref[prob:substitution]{\texttt{prob:substitution}}.}
\textit{Solution.}  The solution immediately following the problem
identifies the term-substitution lemma, sets up the structural induction,
and displays the quantified assignment-commutation calculation; the
term-lemma induction and the atomic and connective cases are left to the
reader.

\paragraph{\hyperref[prob:compactness]{\texttt{prob:compactness}}.}
\textit{Solution.}  The complete solution immediately following the problem
builds models for finite subsets, applies compactness, and separates the
existential model claim from the rejected assertion about all orders.

\paragraph{\hyperref[prob:incompleteness]{\texttt{prob:incompleteness}}.}
\textit{Solution.}  The solution immediately following the problem states the
 computably enumerable axiom-set, arithmetic-strength, and consistency hypotheses and
distinguishes the original G\"odel conclusion from Rosser's refinement.

\paragraph{\hyperref[prob:modal-frame]{\texttt{prob:modal-frame}}.}
\textit{Solution.}  The complete solution immediately following the problem
proves the reflexive direction by the satisfaction clause and constructs a
valuation at a nonreflexive world for the converse.

\paragraph{\hyperref[prob:cut-subformula]{\texttt{prob:cut-subformula}}.}
\textit{Solution.}  The solution immediately following the problem traces the
subformula invariant through cut-free propositional rules, derives consistency
as a consequence, and states the required first-order qualification modulo term
substitution and eigenvariable renaming.

\chapter{Editorial and Formal-Logical Release Checklist}
\label{ch:release-checklist}
\label{app:release-checklist}

\section{Definitions and notation}
\label{sec:release-checklist-01}
\begin{itemize}
  \item Every technical term is defined before use.
  \item Every new technical name has necessity, an exact eliminable definition
        or explicit axiom, standing, nonredundancy, and later use at its first
        introduction.
  \item Every symbol has one registered role and type within each declared
        scope.
  \item The signature, grammar, binding, and precedence rules are explicit.
  \item Every recursive definition has a well-founded basis.
  \item Every quotient definition is proved independent of representatives.
\end{itemize}

\section{Theorems and proofs}
\label{sec:release-checklist-02}
\begin{itemize}
  \item Every theorem states all quantified objects and hypotheses.
  \item Every derivability statement identifies a calculus.
  \item Every semantic statement identifies a structure class.
  \item Every completeness, consistency, compactness, categoricity, and
        decidability claim identifies the object having that property.
  \item Every induction identifies the induction object and all cases.
  \item Every quantifier-rule side condition is checked.
  \item Every substitution is capture-avoiding.
  \item Every proof sketch is labelled and is not cited later as a complete
        proof.
  \item Every relative-consistency theorem remains explicitly relative.
\end{itemize}

\section{Sources and computation}
\label{sec:release-checklist-03}
\begin{itemize}
  \item Historical claims use primary sources when available.
  \item Current tool behavior uses exact-version official manuals.
  \item Automated transformations state equivalence or equisatisfiability.
  \item Solver results include input, status, resource limit, and certificate or
        trust assumption.
  \item Mechanized proofs report formal statement, correspondence, imports,
        axioms, tool version, and kernel result.
\end{itemize}

\section{Production}
\label{sec:release-checklist-04}
\begin{itemize}
  \item The source compiles without undefined references or missing citations.
  \item Duplicate labels and multiply defined macros are absent.
  \item The index and bibliography are generated.
  \item No display is clipped or overlaps surrounding text.
  \item Headers, footers, page numbers, tables, and rule sheets are inspected.
  \item Formal-logical proof review, static checks, compilation, visual
        inspection, review by an automated software agent (reported with
        the agent identity and configuration), and human review are
        reported separately.
\end{itemize}

\backmatter

\cleardoublepage
\phantomsection
\addcontentsline{toc}{chapter}{Bibliography}
\begin{thebibliography}{99}

\bibitem{AMSAuthorHandbook}
American Mathematical Society.
\emph{AMS Author Handbook: Journal Classes}.
American Mathematical Society, current online edition.
\url{https://www.ams.org/arc/journals/packages/proc/proc_amslatex/Author_Handbook_Journals.pdf}.

\bibitem{AMSMathGuide}
American Mathematical Society and the \LaTeX{} Project.
\emph{User's Guide for the amsmath Package}.
\url{https://ctan.org/pkg/latex-amsmath}.

\bibitem{Avigad2018}
Jeremy Avigad.
The mechanization of mathematics.
\emph{Notices of the American Mathematical Society}, 65(6):681--690, 2018.
\url{https://www.ams.org/journals/notices/201806/rnoti-p681.pdf}.

\bibitem{BaazEglyLeitsch2001}
Matthias Baaz, Uwe Egly, and Alexander Leitsch.
Normal form transformations.
In J. Alan Robinson and Andrei Voronkov, editors,
\emph{Handbook of Automated Reasoning}, volume I, pages 273--333.
Elsevier, 2001.
\href{https://doi.org/10.1016/B978-044450813-3/50007-2}{doi:10.1016/B978-044450813-3/50007-2}.

\bibitem{BachmairGanzinger2001}
Leo Bachmair and Harald Ganzinger.
Resolution theorem proving.
In J. Alan Robinson and Andrei Voronkov, editors,
\emph{Handbook of Automated Reasoning}, volume I, pages 19--99.
Elsevier, 2001.
\href{https://doi.org/10.1016/B978-044450813-3/50004-7}{doi:10.1016/B978-044450813-3/50004-7}.

\bibitem{Barwise1977}
Jon Barwise, editor.
\emph{Handbook of Mathematical Logic}.
North-Holland, 1977.

\bibitem{BarwiseFeferman1985}
Jon Barwise and Solomon Feferman, editors.
\emph{Model-Theoretic Logics}.
North-Holland, 1985; Cambridge University Press digital edition, 2017.
\href{https://doi.org/10.1017/9781316717158}{doi:10.1017/9781316717158}.

\bibitem{BlackburnDeRijkeVenema2001}
Patrick Blackburn, Maarten de Rijke, and Yde Venema.
\emph{Modal Logic}.
Cambridge University Press, 2001.
\href{https://doi.org/10.1017/CBO9781107050884}{doi:10.1017/CBO9781107050884}.

\bibitem{BoolosBurgessJeffrey2007}
George S. Boolos, John P. Burgess, and Richard C. Jeffrey.
\emph{Computability and Logic}. 5th edition.
Cambridge University Press, 2007.
\href{https://doi.org/10.1017/CBO9780511804076}{doi:10.1017/CBO9780511804076}.

\bibitem{Buss1998}
Samuel R. Buss, editor.
\emph{Handbook of Proof Theory}.
Elsevier, 1998.

\bibitem{ChangKeisler1990}
C. C. Chang and H. J. Keisler.
\emph{Model Theory}. 3rd edition.
North-Holland, 1990.

\bibitem{ChiswellHodges2007}
Ian Chiswell and Wilfrid Hodges.
\emph{Mathematical Logic}.
Oxford University Press, 2007.
\href{https://doi.org/10.1093/oso/9780198571001.001.0001}{doi:10.1093/oso/9780198571001.001.0001}.

\bibitem{Church1936}
Alonzo Church.
An unsolvable problem of elementary number theory.
\emph{American Journal of Mathematics}, 58(2):345--363, 1936.
\href{https://doi.org/10.2307/2371045}{doi:10.2307/2371045}.

\bibitem{Cook1971}
Stephen A. Cook.
The complexity of theorem-proving procedures.
In \emph{Proceedings of the Third Annual ACM Symposium on Theory of Computing},
pages 151--158, 1971.
\href{https://doi.org/10.1145/800157.805047}{doi:10.1145/800157.805047}.

\bibitem{DavisLogemannLoveland1962}
Martin Davis, George Logemann, and Donald W. Loveland.
A machine program for theorem-proving.
\emph{Communications of the ACM}, 5(7):394--397, 1962.
\href{https://doi.org/10.1145/368273.368557}{doi:10.1145/368273.368557}.

\bibitem{DavisPutnam1960}
Martin Davis and Hilary Putnam.
A computing procedure for quantification theory.
\emph{Journal of the ACM}, 7(3):201--215, 1960.
\href{https://doi.org/10.1145/321033.321034}{doi:10.1145/321033.321034}.

\bibitem{DeMouraUllrich2021}
Leonardo de Moura and Sebastian Ullrich.
The Lean 4 theorem prover and programming language.
In \emph{Automated Deduction -- CADE 28}, pages 625--635. Springer, 2021.
\href{https://doi.org/10.1007/978-3-030-79876-5_37}{doi:10.1007/978-3-030-79876-5\_37}.

\bibitem{EbbinghausFlumThomas1994}
H.-D. Ebbinghaus, J. Flum, and W. Thomas.
\emph{Mathematical Logic}. 2nd edition.
Springer, 1994.
\href{https://doi.org/10.1007/978-1-4757-2355-7}{doi:10.1007/978-1-4757-2355-7}.

\bibitem{Enderton2001}
Herbert B. Enderton.
\emph{A Mathematical Introduction to Logic}. 2nd edition.
Academic Press, 2001.

\bibitem{GabbayPnueliShelahStavi1980}
Dov Gabbay, Amir Pnueli, Saharon Shelah, and Jonathan Stavi.
On the temporal analysis of fairness.
In \emph{Proceedings of the Seventh ACM Symposium on Principles of
Programming Languages}, pages 163--173, 1980.
\href{https://doi.org/10.1145/567446.567462}{doi:10.1145/567446.567462}.

\bibitem{Gentzen1934a}
Gerhard Gentzen.
Untersuchungen \"uber das logische Schlie\ss en I.
\emph{Mathematische Zeitschrift}, 39:176--210, 1934.
\href{https://doi.org/10.1007/BF01201353}{doi:10.1007/BF01201353}.

\bibitem{Gentzen1935b}
Gerhard Gentzen.
Untersuchungen \"uber das logische Schlie\ss en II.
\emph{Mathematische Zeitschrift}, 39:405--431, 1935.
\href{https://doi.org/10.1007/BF01201363}{doi:10.1007/BF01201363}.

\bibitem{Girard1987}
Jean-Yves Girard.
Linear logic.
\emph{Theoretical Computer Science}, 50:1--101, 1987.
\href{https://doi.org/10.1016/0304-3975(87)90045-4}{doi:10.1016/0304-3975(87)90045-4}.

\bibitem{Godel1930}
Kurt G\"odel.
Die Vollst\"andigkeit der Axiome des logischen Funktionenkalk\"uls.
\emph{Monatshefte f\"ur Mathematik und Physik}, 37:349--360, 1930.
\href{https://doi.org/10.1007/BF01696781}{doi:10.1007/BF01696781}.

\bibitem{Godel1931}
Kurt G\"odel.
\"Uber formal unentscheidbare S\"atze der \emph{Principia Mathematica} und
verwandter Systeme I.
\emph{Monatshefte f\"ur Mathematik und Physik}, 38:173--198, 1931.
\href{https://doi.org/10.1007/BF01700692}{doi:10.1007/BF01700692}.

\bibitem{Halmos1970}
Paul R. Halmos.
How to write mathematics.
\emph{L'Enseignement Math\'ematique}, 16(1):123--152, 1970.
\href{https://doi.org/10.5169/seals-43857}{doi:10.5169/seals-43857}.

\bibitem{Henkin1949}
Leon Henkin.
The completeness of the first-order functional calculus.
\emph{The Journal of Symbolic Logic}, 14(3):159--166, 1949.
\href{https://doi.org/10.2307/2267044}{doi:10.2307/2267044}.

\bibitem{Henkin1950}
Leon Henkin.
Completeness in the theory of types.
\emph{The Journal of Symbolic Logic}, 15(2):81--91, 1950.
\href{https://doi.org/10.2307/2266967}{doi:10.2307/2266967}.

\bibitem{HilbertBernays1939}
David Hilbert and Paul Bernays.
\emph{Grundlagen der Mathematik}, volume 2.
Julius Springer, Berlin, first edition, 1939.
\href{https://search.worldcat.org/title/Grundlagen-der-Mathematik.-Bd.-2/oclc/1424451730}{WorldCat catalog record}.

\bibitem{Hodges1993}
Wilfrid Hodges.
\emph{Model Theory}.
Cambridge University Press, 1993.
\href{https://doi.org/10.1017/CBO9780511551574}{doi:10.1017/CBO9780511551574}.

\bibitem{Isabelle2025}
The Isabelle Development Team.
\emph{Isabelle2025-2 Documentation}.
\url{https://isabelle.in.tum.de/documentation.html}.

\bibitem{ISO80000}
International Organization for Standardization.
\emph{ISO 80000-2:2019, Quantities and Units -- Part 2: Mathematics}.
2nd edition, corrected 2021.
\url{https://www.iso.org/standard/64973.html}.

\bibitem{Jech2003}
Thomas Jech.
\emph{Set Theory: The Third Millennium Edition, Revised and Expanded}.
3rd edition. Springer, 2003.
\href{https://doi.org/10.1007/3-540-44761-X}{doi:10.1007/3-540-44761-X}.

\bibitem{Kleene1952}
Stephen Cole Kleene.
\emph{Introduction to Metamathematics}.
North-Holland, 1952.

\bibitem{KnuthLarrabeeRoberts1989}
Donald E. Knuth, Tracy L. Larrabee, and Paul M. Roberts.
\emph{Mathematical Writing}.
Mathematical Association of America, 1989.
\url{https://cs-faculty.stanford.edu/~knuth/klr.html}.

\bibitem{Kripke1963}
Saul A. Kripke.
Semantical analysis of modal logic I: Normal modal propositional calculi.
\emph{Zeitschrift f\"ur mathematische Logik und Grundlagen der Mathematik},
9(5--6):67--96, 1963.
\href{https://doi.org/10.1002/malq.19630090502}{doi:10.1002/malq.19630090502}.

\bibitem{Kripke1965}
Saul A. Kripke.
Semantical analysis of intuitionistic logic I.
In J. N. Crossley and M. A. E. Dummett, editors,
\emph{Formal Systems and Recursive Functions}, pages 92--130.
North-Holland, 1965.
\href{https://doi.org/10.1016/S0049-237X(08)71685-9}{doi:10.1016/S0049-237X(08)71685-9}.

\bibitem{LeanManual2026}
Lean Project.
\emph{The Lean Language Reference}, version 4.32.1.
Retrieved 9 August 2026.
\url{https://lean-lang.org/doc/reference/4.32.1/}.

\bibitem{LeanTPIL}
Jeremy Avigad, Leonardo de Moura, Soonho Kong, and Sebastian Ullrich.
\emph{Theorem Proving in Lean 4}.
Official online tutorial, retrieved 9 August 2026.
\url{https://docs.lean-lang.org/theorem_proving_in_lean4/}.

\bibitem{Lindstrom1969}
Per Lindstr\"om.
On extensions of elementary logic.
\emph{Theoria}, 35(1):1--11, 1969.
\href{https://doi.org/10.1111/j.1755-2567.1969.tb00356.x}{doi:10.1111/j.1755-2567.1969.tb00356.x}.

\bibitem{Lob1955}
M. H. L\"ob.
Solution of a problem of Leon Henkin.
\emph{The Journal of Symbolic Logic}, 20(2):115--118, 1955.
\href{https://doi.org/10.2307/2266895}{doi:10.2307/2266895}.

\bibitem{Los1955}
Jerzy \L{}o\'s.
On the extending of models I.
\emph{Fundamenta Mathematicae}, 42:38--54, 1955.
\href{https://doi.org/10.4064/fm-42-1-38-54}{doi:10.4064/fm-42-1-38-54}.

\bibitem{Manzano1996}
Mar\'ia Manzano.
\emph{Extensions of First Order Logic}.
Cambridge University Press, 1996.

\bibitem{Marker2002}
David Marker.
\emph{Model Theory: An Introduction}.
Springer, 2002.
\href{https://doi.org/10.1007/b98860}{doi:10.1007/b98860}.

\bibitem{Marker2024}
David Marker.
\emph{An Invitation to Mathematical Logic}.
Springer, 2024.
\href{https://doi.org/10.1007/978-3-031-55368-4}{doi:10.1007/978-3-031-55368-4}.

\bibitem{NegriVonPlato2001}
Sara Negri and Jan von Plato.
\emph{Structural Proof Theory}.
Cambridge University Press, 2001.
\href{https://doi.org/10.1017/CBO9780511527340}{doi:10.1017/CBO9780511527340}.

\bibitem{OpenLogic2026}
The Open Logic Project.
\emph{The Open Logic Text}.
Release 9f12419, 25 May 2026.
\url{https://builds.openlogicproject.org/open-logic-complete.pdf}.

\bibitem{Prawitz1965}
Dag Prawitz.
\emph{Natural Deduction: A Proof-Theoretical Study}.
Almqvist \& Wiksell, 1965; Dover reprint, 2006.

\bibitem{Robinson1965}
J. A. Robinson.
A machine-oriented logic based on the resolution principle.
\emph{Journal of the ACM}, 12(1):23--41, 1965.
\href{https://doi.org/10.1145/321250.321253}{doi:10.1145/321250.321253}.

\bibitem{RocqManual2026}
The Rocq Development Team.
\emph{The Rocq Prover Reference Manual}, version 9.1.1.
\url{https://rocq-prover.org/doc/V9.1.1/refman/index.html}.

\bibitem{RobinsonVoronkov2001}
J. Alan Robinson and Andrei Voronkov, editors.
\emph{Handbook of Automated Reasoning}.
Elsevier and MIT Press, 2001.

\bibitem{Rosser1936}
J. Barkley Rosser.
Extensions of some theorems of G\"odel and Church.
\emph{The Journal of Symbolic Logic}, 1(3):87--91, 1936.
\href{https://doi.org/10.2307/2269028}{doi:10.2307/2269028}.

\bibitem{RoutleyMeyer1972}
Richard Routley and Robert K. Meyer.
The semantics of entailment---II.
\emph{Journal of Philosophical Logic}, 1(1):53--73, 1972.
\href{https://doi.org/10.1007/BF00649991}{doi:10.1007/BF00649991}.

\bibitem{Shapiro1991}
Stewart Shapiro.
\emph{Foundations without Foundationalism: A Case for Second-Order Logic}.
Oxford University Press, 1991.
\href{https://doi.org/10.1093/0198250290.001.0001}{doi:10.1093/0198250290.001.0001}.

\bibitem{Shoenfield1967}
Joseph R. Shoenfield.
\emph{Mathematical Logic}.
Addison-Wesley, 1967.

\bibitem{SMTLIB2025}
Clark Barrett, Pascal Fontaine, and Cesare Tinelli.
\emph{The SMT-LIB Standard}, version 2.7, release 7 July 2025.
\url{https://smt-lib.org/papers/smt-lib-reference-v2.7-r2025-07-07.pdf}.

\bibitem{Tarski1983}
Alfred Tarski.
\emph{Logic, Semantics, Metamathematics: Papers from 1923 to 1938}.
2nd edition, edited by John Corcoran, translated by J. H. Woodger.
Hackett, 1983.

\bibitem{TPTP}
Geoff Sutcliffe and contributors.
\emph{The TPTP Language: Hyperlinked BNF}, BNF revision 9.3.0.3.
Retrieved 9 August 2026.
SHA-256 of the retrieved UTF-8 page content:
\texttt{6bcbfc86f2e606e3}\allowbreak
\texttt{b0db8c5b1de29409}\allowbreak
\texttt{3b2753c868296b50}\allowbreak
\texttt{78ad028b9f863243}.
\url{https://tptp.org/UserDocs/TPTPLanguage/SyntaxBNF.html}.

\bibitem{Trakhtenbrot1950}
Boris A. Trakhtenbrot.
The impossibility of an algorithm for the decision problem on finite classes.
\emph{Doklady Akademii Nauk SSSR}, 70:569--572, 1950.

\bibitem{TroelstraSchwichtenberg2000}
A. S. Troelstra and H. Schwichtenberg.
\emph{Basic Proof Theory}. 2nd edition.
Cambridge University Press, 2000.
\href{https://doi.org/10.1017/CBO9781139168717}{doi:10.1017/CBO9781139168717}.

\bibitem{TroelstraVanDalen1988}
A. S. Troelstra and D. van Dalen.
\emph{Constructivism in Mathematics: An Introduction}.
North-Holland, 1988.

\bibitem{Tseitin1968}
G. S. Tseitin.
On the complexity of derivation in propositional calculus.
In J\"org H. Siekmann and Graham Wrightson, editors,
\emph{Automation of Reasoning 2: Classical Papers on Computational Logic
1967--1970}, pages 466--483.  Springer, 1983.  Originally published in 1968.
\href{https://doi.org/10.1007/978-3-642-81955-1_28}{doi:10.1007/978-3-642-81955-1\_28}.

\bibitem{Turing1937}
A. M. Turing.
On computable numbers, with an application to the Entscheidungsproblem.
\emph{Proceedings of the London Mathematical Society}, series 2, 42:230--265,
1937.
\href{https://doi.org/10.1112/plms/s2-42.1.230}{doi:10.1112/plms/s2-42.1.230}.

\bibitem{Unicode2025}
The Unicode Consortium.
\emph{The Unicode Standard}, version 17.0, Chapter 22: Symbols.
\url{https://www.unicode.org/versions/Unicode17.0.0/core-spec/chapter-22/}.

\end{thebibliography}

\cleardoublepage
\phantomsection
\addcontentsline{toc}{chapter}{Index}
\begingroup
\sloppy
\printindex
\endgroup

\end{document}
